\chapter{Sprach- und Beweisstruktur der Metatheorie}

\section{Syntax und Wahrheit in Mengenstrukturen}

\FormulaDefDeltaK[Sprache und Beweis]
{\begin{aligned}
 \mathcal L_\in&\coloneqq\{=,\in\},\\
 \operatorname{Var}&\coloneq\mathbb N,\qquad v_i\coloneq i,\\
 T\vdash\varphi&\coloneqq
       \exists p\,\operatorname{Prf}_{1}(p,T,\varphi)
 \end{aligned}}
{B48Syntax}{
  \DeltaRow{Theorie}{T\subseteq\operatorname{Sent}(\mathcal L_\in)}
  \DeltaRow{Endliche Herleitung nach Band 1}{p}
}
\noindent
\(\operatorname{Prf}_{1}(p,T,\varphi)\) bezeichnet die in Band~1
festgelegte syntaktische Bedingung: \(p\) ist eine gueltige endliche
elementare Herleitung mit Schlussformel \(\varphi\), deren noch offene
Voraussetzungen zu \(T\) gehoeren. Dies ist eine Abkuerzung fuer
das dortige Herleitungsformat, kein neues Axiom und insbesondere
keine semantische Definition der Beweisbarkeit. Eine spaetere
Umformung von Beweisobjekten braucht jeweils einen eigenen Nachweis.
Abkuerzungen fuer mehrere Regelschritte werden in
\(\operatorname{Prf}_1\) als die ausgeschriebene endliche Folge
gelesen. Die Gueltigkeit eines blossen Bibliotheksverweises ohne
zugehoerige Herleitung wird damit nicht vorausgesetzt.

Es gibt abzählbar viele Variablen. Aus atomaren Formeln werden mit
\(\neg,\land,\exists\) die Formeln aufgebaut; die übrigen Junktoren und
\(\forall\) sind die üblichen Abkürzungen.
Die folgenden Definitionen realisieren die Formelcodes als endliche
Baumwoerter auf einem mengentheoretisch festgelegten Alphabet.
Eine zusaetzliche Nummerierung aller Codes durch natuerliche Zahlen
wird fuer den Existenzsatz der Erfuellungsrelation nicht vorausgesetzt.
Wir verwenden die klassische Prädikatenlogik mit Gleichheit und den
Einführungs- und Beseitigungsregeln aus Band~1. Die Eigenvariablenbedingungen
bei Quantorenregeln gehören ausdrücklich zu diesen Regeln.

% Mengeninterne Syntax und Semantik. Keine neuen Mengenaxiome.
% Alle Hilfssaetze stehen als Proofparts beim Existenzsatz.
\newcommand{\BXLVIIIIf}[3]{\mathsf{if}(#1,#2,#3)}
\subsection*{Formelcodes und axiomatische Charakterisierung}
Wie bei den induktiven Mengen in
\FormulaRefAuto{InductiveSetDef}[def] stehen die kennzeichnenden Axiome
eines Strukturbegriffs einzeln in dessen Definition. Die Definition
dient als Einfuehrungsregel: Ihre Axiome werden durch getrennte
Beweiszeilen nachgewiesen. Die anschliessenden Zugriffsaxiome liefern
aus dem Strukturpraedikat jeweils genau eine dieser Bedingungen.
Keine dieser Regeln behauptet die Existenz einer solchen Struktur.
In verbleibenden logischen Formeln werden mehrgliedrige Konjunktionen
rechts geklammert.
Die Mengen \(D,E\) sind feste Parameter dieses Abschnitts; \(M=(D,E)\).
\FormulaDefDeltaK[Kurzschreibweise des Fallterms]
{\BXLVIIIIf{P}{U}{V}\coloneq\IfThenElse{P}{U}{V}}
{B48IfNotation}{\DeltaRow{Aussage}{P}\DeltaRow{Mengen}{U,V}}


Die folgenden Bedingungen charakterisieren eine gegebene Relation.
Sie behaupten noch nicht deren Existenz. Als Codebereich dienen die
bereits konstruierten Klammerungsbaeume aus Band~27.
Gebundene Tupel sind die uebliche Abkuerzung fuer die eindeutige
Produktzerlegung; es werden keine neuen Arten von Variablen eingefuehrt.
Die Strukturbegriffe sind Praedikate ueber Mengen; die uebrigen
Kurzzeichen bezeichnen Mengen und Mengenoperationen.

\FormulaDefDeltaK[Marken und Codebereich]
{\begin{gathered}
 t_{\rm e}\coloneq0,\quad t_{\rm m}\coloneq\Succ(t_{\rm e}),\quad
 t_{\rm n}\coloneq\Succ(t_{\rm m}),\\
 t_{\rm c}\coloneq\Succ(t_{\rm n}),\quad
 t_{\rm q}\coloneq\Succ(t_{\rm c}),\quad
 t_{\rm b}\coloneq\Succ(t_{\rm q}),\\
 \Sigma\coloneq\mathbb N\times(\mathbb N\times\mathbb N),\qquad
 \mathcal T\coloneq\BracketTreeSet{\Sigma},\\
 \mathsf l(c)\coloneq\LeafTree{\Sigma}{c},\qquad
 \mathsf n(p,q)\coloneq\NodeTree{\Sigma}{p}{q}
\end{gathered}}
{B48CodeCarrier}{\DeltaRow{Mengen}{c\in\Sigma,\quad p,q\in\mathcal T}}

\FormulaDefDeltaK[Formelkonstruktoren]
{\begin{gathered}
 \begin{aligned}
 \mathsf e(i,j)&\coloneq\mathsf l((t_{\rm e},(i,j))),&
 \mathsf m(i,j)&\coloneq\mathsf l((t_{\rm m},(i,j))),\\
 \mathsf b&\coloneq\mathsf l((t_{\rm b},(0,0))),&
 \mathsf n_0&\coloneq\mathsf l((t_{\rm n},(0,0))),\\
 \mathsf c_0&\coloneq\mathsf l((t_{\rm c},(0,0))),&
 \mathsf q_i&\coloneq\mathsf l((t_{\rm q},(i,0))),\\
 \mathsf{neg}(p)&\coloneq\mathsf n(\mathsf n_0,p),&
 \mathsf{and}(p,q)&\coloneq\mathsf n(\mathsf n(\mathsf c_0,p),q),\\
 \mathsf{ex}_i(p)&\coloneq\mathsf n(\mathsf q_i,p)
 \end{aligned}\\
 \begin{aligned}
 \Sigma_{\rm at}\coloneq\{(k,(i,j))\in\Sigma\mid{}&
 k=t_{\rm e}\lor k=t_{\rm m}\\
 &{}\lor(k=t_{\rm b}\land i=0\land j=0)\}.
 \end{aligned}
\end{gathered}}
{B48FormulaConstructors}{\DeltaRow{Indizes}{i,j\in\mathbb N}
 \DeltaRow{Baumcodes}{p,q\in\mathcal T}}

\FormulaDefDeltaK[Formelabgeschlossene Menge]
{\mathsf{Cl}(X)}{B48FormulaClosure}{
 \DeltaRow{Mengen}{X,p,c,q,i}
 \DeltaRow{\textbf{Axiome}}{}
 \DeltaRow{Baumcodes}{X\subseteq\mathcal T}
 \DeltaRow{Atome}{\forall c\in\Sigma_{\rm at}\;(\mathsf l(c)\in X)}
 \DeltaRow{Negation}{\forall p\in X\;(\mathsf{neg}(p)\in X)}
 \DeltaRow{Konjunktion}{\begin{aligned}[t]&\forall p\in X\,\forall q\in X\\&\quad(\mathsf{and}(p,q)\in X)\end{aligned}}
 \DeltaRow{Existenzquantor}{\begin{aligned}[t]&\forall i\in\mathbb N\,\forall p\in X\\&\quad(\mathsf{ex}_i(p)\in X)\end{aligned}}
 \DeltaRow{\textbf{Neues Symbol}}{}
 \DeltaPrem{Formelabschluss}{\mathsf{Cl}(X)}
}

\FormulaAxiomDeltaK[Traeger einer formelabgeschlossenen Menge]
{\mathsf{Cl}(X)\vdash X\subseteq\mathcal T}
{B48ClDomainAxiom}{\DeltaRow{Mengen}{X,p,c,q,i}}

\FormulaAxiomDeltaK[Atomabschluss]
{\mathsf{Cl}(X)\vdash \forall c\in\Sigma_{\rm at}\;(\mathsf l(c)\in X)}
{B48ClAtomsAxiom}{\DeltaRow{Mengen}{X,p,c,q,i}}

\FormulaAxiomDeltaK[Abschluss unter Negation]
{\mathsf{Cl}(X)\vdash \forall p\in X\;(\mathsf{neg}(p)\in X)}
{B48ClNegAxiom}{\DeltaRow{Mengen}{X,p,c,q,i}}

\FormulaAxiomDeltaK[Abschluss unter Konjunktion]
{\begin{gathered}\mathsf{Cl}(X)\vdash{}\\
 \forall p\in X\,\forall q\in X\;(\mathsf{and}(p,q)\in X)\end{gathered}}
{B48ClAndAxiom}{\DeltaRow{Mengen}{X,p,c,q,i}}

\FormulaAxiomDeltaK[Abschluss unter Existenzquantifikation]
{\begin{gathered}\mathsf{Cl}(X)\vdash{}\\
 \forall i\in\mathbb N\,\forall p\in X\;(\mathsf{ex}_i(p)\in X)\end{gathered}}
{B48ClExistsAxiom}{\DeltaRow{Mengen}{X,p,c,q,i}}

\FormulaDefDeltaK[Formelmenge]
{\mathcal F\coloneq
 \{p\in\mathcal T\mid\forall X(\mathsf{Cl}(X)\to p\in X)\}}
{B48FormulaSet}{\DeltaRow{Mengen}{p,X}}

Dies ist eine beschraenkte Aussonderung aus einer bereits vorhandenen
Menge, keine rekursive Existenzbehauptung.
Die Codes \(\mathsf e(i,j)\), \(\mathsf m(i,j)\), \(\mathsf b\),
\(\mathsf{neg}(p)\), \(\mathsf{and}(p,q)\) und \(\mathsf{ex}_i(p)\)
codieren jeweils \(v_i=v_j\), \(v_i\in v_j\), \(\bot\), \(\neg p\),
\(p\land q\) und \(\exists v_i\,p\).
Innerhalb von \(\models\) werden spaeter die lesbaren Formeln statt
dieser Codes geschrieben. \(\mathsf n_0,\mathsf c_0,\mathsf q_i\)
sind Hilfsbaeume, nicht zusaetzliche atomare Formeln.

\FormulaDefDeltaK[Mengenstruktur]
{\mathsf{Str}(M)}{B48SetStructureDef}{
 \DeltaRow{Mengen}{M,D,E}
 \DeltaRow{\textbf{Axiome}}{}
 \DeltaRow{Strukturpaar}{M=(D,E)}
 \DeltaRow{Grundmenge}{D\neq\varnothing}
 \DeltaRow{Relation}{E\subseteq D\times D}
 \DeltaRow{\textbf{Neues Symbol}}{}
 \DeltaPrem{Mengenstruktur}{\mathsf{Str}(M)}
}

\FormulaAxiomDeltaK[Paar einer Mengenstruktur]
{\mathsf{Str}(M)\vdash M=(D,E)}
{B48StrPairAxiom}{\DeltaRow{Mengen}{M,D,E}}

\FormulaAxiomDeltaK[Nichtleere Grundmenge]
{\mathsf{Str}(M)\vdash D\neq\varnothing}
{B48StrDomainAxiom}{\DeltaRow{Mengen}{M,D,E}}

\FormulaAxiomDeltaK[Interpretation der Elementrelation]
{\mathsf{Str}(M)\vdash E\subseteq D\times D}
{B48StrRelationAxiom}{\DeltaRow{Mengen}{M,D,E}}

\FormulaDefDeltaK[Belegungen und Belegungswechsel]
{\begin{aligned}
 A_D&\coloneq\{s\in\powerset(\mathbb N\times D)\mid s\colon\mathbb N\to D\},\\
 s[i\mapsto a]&\coloneq
 \{(j,b)\in\mathbb N\times D\mid
       b=\BXLVIIIIf{j=i}{a}{s(j)}\}.
\end{aligned}}
{B48Assignments}{\DeltaRow{Mengen}{D,s,i,a,j,b}}

\FormulaDefDeltaK[Atomare Wertemengen und Quantorenoperator]
{\begin{aligned}
 B_M((k,(i,j)))\coloneq{}&
 \mathsf{if}\bigl(k=t_{\rm e},\{s\in A_D\mid s(i)=s(j)\},\\
 &\qquad\BXLVIIIIf{k=t_{\rm m}}
     {\{s\in A_D\mid(s(i),s(j))\in E\}}{\varnothing}\bigr),\\
 Q_i(U)\coloneq{}&\{s\in A_D\mid\exists a\in D\;(s[i\mapsto a]\in U)\},\\
 Q(t,U)\coloneq{}&\{s\in A_D\mid
       \exists i\in\mathbb N\;(t=\mathsf q_i\land
                              \exists a\in D\;(s[i\mapsto a]\in U))\},\\
 R_S(p)\coloneq{}&\{s\in A_D\mid(p,s)\in S\}.
\end{aligned}}
{B48SemanticTerms}{\DeltaRow{Mengen}{M=(D,E),\ (k,(i,j))\in\Sigma,\ t,U,S,p}}

\FormulaDefDeltaK[Atomklausel der Erfuellung]
{C_{\rm at}(M,S)\coloneq \forall c\in\Sigma_{\rm at}\;(R_S(\mathsf l(c))=B_M(c))}
{B48SatisfactionAtomsDef}{\DeltaRow{Mengen}{M=(D,E),S,p,q,c,i}}

\FormulaDefDeltaK[Negationsklausel der Erfuellung]
{C_{\rm n}(M,S)\coloneq \forall p\in\mathcal F\;(R_S(\mathsf{neg}(p))=A_D\setminus R_S(p))}
{B48SatisfactionNegDef}{\DeltaRow{Mengen}{M=(D,E),S,p,q,c,i}}

\FormulaDefDeltaK[Konjunktionsklausel der Erfuellung]
{C_{\rm c}(M,S)\coloneq \begin{aligned}[t]&\forall p\in\mathcal F\,\forall q\in\mathcal F\\&\quad(R_S(\mathsf{and}(p,q))=R_S(p)\cap R_S(q))\end{aligned}}
{B48SatisfactionAndDef}{\DeltaRow{Mengen}{M=(D,E),S,p,q,c,i}}

\FormulaDefDeltaK[Quantorenklausel der Erfuellung]
{C_{\rm q}(M,S)\coloneq \begin{aligned}[t]&\forall i\in\mathbb N\,\forall p\in\mathcal F\\&\quad(R_S(\mathsf{ex}_i(p))=Q_i(R_S(p)))\end{aligned}}
{B48SatisfactionExistsDef}{\DeltaRow{Mengen}{M=(D,E),S,p,q,c,i}}

\FormulaDefDeltaK[Erfuellungsrelation]
{\mathsf{Erf}(M,S)}{B48SatisfactionStructure}{
 \DeltaRow{Mengen}{M=(D,E),S}
 \DeltaRow{\textbf{Axiome}}{}
 \DeltaRow{Struktur}{\mathsf{Str}(M)}
 \DeltaRow{Relationstraeger}{S\subseteq\mathcal F\times A_D}
 \DeltaRow{Atome}{C_{\rm at}(M,S)}[\FormulaRefAuto{B48SatisfactionAtomsDef}[def]]
 \DeltaRow{Negation}{C_{\rm n}(M,S)}[\FormulaRefAuto{B48SatisfactionNegDef}[def]]
 \DeltaRow{Konjunktion}{C_{\rm c}(M,S)}[\FormulaRefAuto{B48SatisfactionAndDef}[def]]
 \DeltaRow{Existenzquantor}{C_{\rm q}(M,S)}[\FormulaRefAuto{B48SatisfactionExistsDef}[def]]
 \DeltaRow{\textbf{Neues Symbol}}{}
 \DeltaPrem{Erfuellung}{\mathsf{Erf}(M,S)}
}

\FormulaAxiomDeltaK[Struktur einer Erfuellungsrelation]
{\mathsf{Erf}(M,S)\vdash \mathsf{Str}(M)}
{B48ErfStructureAxiom}{\DeltaRow{Mengen}{M=(D,E),S}}

\FormulaAxiomDeltaK[Traeger einer Erfuellungsrelation]
{\mathsf{Erf}(M,S)\vdash S\subseteq\mathcal F\times A_D}
{B48ErfSubsetAxiom}{\DeltaRow{Mengen}{M=(D,E),S}}

\FormulaAxiomDeltaK[Atomklausel der Erfuellung]
{\begin{gathered}\mathsf{Erf}(M,S)\vdash{}\\
 \forall c\in\Sigma_{\rm at}\;(R_S(\mathsf l(c))=B_M(c))\end{gathered}}
{B48ErfAtomsAxiom}{\DeltaRow{Mengen}{M=(D,E),S}}

\FormulaAxiomDeltaK[Negationsklausel der Erfuellung]
{\begin{gathered}\mathsf{Erf}(M,S)\vdash{}\\
 \forall p\in\mathcal F\;(R_S(\mathsf{neg}(p))=A_D\setminus R_S(p))\end{gathered}}
{B48ErfNegAxiom}{\DeltaRow{Mengen}{M=(D,E),S}}

\FormulaAxiomDeltaK[Konjunktionsklausel der Erfuellung]
{\begin{gathered}\mathsf{Erf}(M,S)\vdash{}\\
 \forall p\in\mathcal F\,\forall q\in\mathcal F\;(R_S(\mathsf{and}(p,q))=R_S(p)\cap R_S(q))\end{gathered}}
{B48ErfAndAxiom}{\DeltaRow{Mengen}{M=(D,E),S}}

\FormulaAxiomDeltaK[Quantorenklausel der Erfuellung]
{\begin{gathered}\mathsf{Erf}(M,S)\vdash{}\\
 \forall i\in\mathbb N\,\forall p\in\mathcal F\;(R_S(\mathsf{ex}_i(p))=Q_i(R_S(p)))\end{gathered}}
{B48ErfExistsAxiom}{\DeltaRow{Mengen}{M=(D,E),S}}

Die vier Klauselpraedikate wurden jeweils durch eine eigene Formel
definiert. Zusammen mit den Bedingungen an Struktur und Relation bilden
sie die sechs Axiome des Begriffs \(\mathsf{Erf}(M,S)\).
Aus dem Nachweis dieser sechs Bedingungen wird der Begriff eingefuehrt;
aus dem Begriff wird jede Bedingung durch ihr eigenes Zugriffsaxiom
gewonnen. Die Existenz einer solchen Relation folgt erst aus dem
nachstehenden Konstruktionsbeweis.

\FormulaDefDeltaK[Totaler Auswertungsoperator auf Hilfsbaeumen]
{\begin{aligned}
 J(t)&\coloneq\exists p\in\mathcal T\;(t=\mathsf n(\mathsf c_0,p)),\\
 H_{2,M}(t,U,V)&\coloneq\BXLVIIIIf{J(t)}{U\cap V}{Q(t,V)},\\
 H_{1,M}(t,U,V)&\coloneq\BXLVIIIIf{t=\mathsf c_0}{V}{H_{2,M}(t,U,V)},\\
 H_M(t,U,V)&\coloneq\BXLVIIIIf{t=\mathsf n_0}{A_D\setminus V}{H_{1,M}(t,U,V)},\\
 Y_D&\coloneq\mathcal T\times\powerset(A_D),\\
 f_M&\coloneq\{(c,z)\in\Sigma\times Y_D\mid z=(\mathsf l(c),B_M(c))\},\\
 g_M&\coloneq
 \{(((p,U),(q,V)),z)\in(Y_D\times Y_D)\times Y_D\mid\\
 &\hspace{42mm}z=(\mathsf n(p,q),H_M(p,U,V))\}.
\end{aligned}}
{B48EvaluationOperators}{\DeltaRow{Mengen}{M=(D,E),t,U,V,c,z,p,q}}

Hier wird auch auf Hilfsbaeumen ausgewertet. Der erste Bestandteil eines
Auswertungswerts speichert den Baumcode selbst. Dadurch kann der zweite
Bestandteil die drei logischen Konstruktoren unterscheiden.
Diese Hilfsauswertung erweitert nicht die Objektlogik.

\FormulaDefDeltaK[Rekursive Auswertung]
{\mathsf{Rec}(e,M)}{B48RecursiveEvaluationDef}{
 \DeltaRow{Mengen}{M=(D,E),e,p,q,c}
 \DeltaRow{\textbf{Axiome}}{}
 \DeltaRow{Auswertung}{e\colon\mathcal T\to Y_D}
 \DeltaRow{Blaetter}{\forall c\in\Sigma\;(e(\mathsf l(c))=f_M(c))}
 \DeltaRow{Knoten}{\begin{aligned}[t]&\forall p\in\mathcal T\,\forall q\in\mathcal T\\&\quad(e(\mathsf n(p,q))=g_M(e(p),e(q)))\end{aligned}}
 \DeltaRow{\textbf{Neues Symbol}}{}
 \DeltaPrem{Auswertung}{\mathsf{Rec}(e,M)}
}

\FormulaAxiomDeltaK[Typ einer rekursiven Auswertung]
{\mathsf{Rec}(e,M)\vdash e\colon\mathcal T\to Y_D}
{B48RecTypeAxiom}{\DeltaRow{Mengen}{M=(D,E),e,p,q,c}}

\FormulaAxiomDeltaK[Blattklausel der rekursiven Auswertung]
{\mathsf{Rec}(e,M)\vdash \forall c\in\Sigma\;(e(\mathsf l(c))=f_M(c))}
{B48RecLeafAxiom}{\DeltaRow{Mengen}{M=(D,E),e,p,q,c}}

\FormulaAxiomDeltaK[Knotenklausel der rekursiven Auswertung]
{\begin{gathered}\mathsf{Rec}(e,M)\vdash{}\\
 \forall p\in\mathcal T\,\forall q\in\mathcal T\;(e(\mathsf n(p,q))=g_M(e(p),e(q)))\end{gathered}}
{B48RecNodeAxiom}{\DeltaRow{Mengen}{M=(D,E),e,p,q,c}}

\FormulaDefDeltaK[Auslesen und Kandidat]
{\begin{aligned}
 V_A(z,p)&\coloneq\{s\in A\mid
           \exists U\in\powerset(A)\;(z=(p,U)\land s\in U)\},\\
 \beta_e(p)&\coloneq V_{A_D}(e(p),p),\\
 S_e&\coloneq\{(p,s)\in\mathcal F\times A_D\mid s\in\beta_e(p)\}.
\end{aligned}}
{B48EvaluationCandidate}{\DeltaRow{Mengen}{M=(D,E),e,p,z,A,s,U}}

\FormulaDefDeltaK[Gleichheitstraeger zweier Relationen]
{K(S,T)\coloneq\{p\in\mathcal F\mid R_S(p)=R_T(p)\}}
{B48EqualityCarrier}{\DeltaRow{Mengen}{S,T,p}}

\FormulaThmDeltaK[Existenz und Eindeutigkeit der Erfuellungsrelation]
{\mathsf{Str}(M)\vdash\exists!S\;\mathsf{Erf}(M,S)}
{B48SatisfactionExists}{\DeltaRow{Struktur}{M=(D,E)}
 \DeltaRow{Mengen}{S,T,e,p,q,c,i,j,U,V,X,z,s,a}
 \DeltaRow{Praedikate}{P}}
\noindent
Die folgenden Proofparts bilden den Beweis. Jeder Hilfssatz wird vor
seiner Verwendung hergeleitet. In Zeugenbeweisen sind die gewaelten
Mengen frisch fuer die ausserhalb des jeweiligen Teilbeweises offenen
Voraussetzungen.
\begin{tabproofsplitwide}
\par\medskip\noindent\textbf{A. Konkrete Marken und Baumcodes}\par
\proofpartwideindR[Typisierung der Marke e]{t_{\rm e}\in\mathbb N}{t_{\rm e}\in\mathbb N}[B48TagType1]
 \proofstepwidestar[]{0\in\mathbb N}{\FormulaRefAuto{PeanoZero}[ax]}
 \proofstepwidestar[]{t_{\rm e}\in\mathbb N}{\FormulaRefAuto{B48CodeCarrier}[def]{1}}
\closeproofpartwideindR

\proofpartwideindR[Typisierung der Marke m]{t_{\rm m}\in\mathbb N}{t_{\rm m}\in\mathbb N}[B48TagType2]
 \proofstepwidestar[]{t_{\rm e}\in\mathbb N}{\FormulaRefAuto{B48TagType1}[thm]}
 \proofstepwidestar[]{\Succ(t_{\rm e})\in\mathbb N}{\FormulaRefAuto{PeanoSuccClosure}[ax]{1}}
 \proofstepwidestar[]{t_{\rm m}\in\mathbb N}{\FormulaRefAuto{B48CodeCarrier}[def]{2}}
\closeproofpartwideindR

\proofpartwideindR[Typisierung der Marke n]{t_{\rm n}\in\mathbb N}{t_{\rm n}\in\mathbb N}[B48TagType3]
 \proofstepwidestar[]{t_{\rm m}\in\mathbb N}{\FormulaRefAuto{B48TagType2}[thm]}
 \proofstepwidestar[]{\Succ(t_{\rm m})\in\mathbb N}{\FormulaRefAuto{PeanoSuccClosure}[ax]{1}}
 \proofstepwidestar[]{t_{\rm n}\in\mathbb N}{\FormulaRefAuto{B48CodeCarrier}[def]{2}}
\closeproofpartwideindR

\proofpartwideindR[Typisierung der Marke c]{t_{\rm c}\in\mathbb N}{t_{\rm c}\in\mathbb N}[B48TagType4]
 \proofstepwidestar[]{t_{\rm n}\in\mathbb N}{\FormulaRefAuto{B48TagType3}[thm]}
 \proofstepwidestar[]{\Succ(t_{\rm n})\in\mathbb N}{\FormulaRefAuto{PeanoSuccClosure}[ax]{1}}
 \proofstepwidestar[]{t_{\rm c}\in\mathbb N}{\FormulaRefAuto{B48CodeCarrier}[def]{2}}
\closeproofpartwideindR

\proofpartwideindR[Typisierung der Marke q]{t_{\rm q}\in\mathbb N}{t_{\rm q}\in\mathbb N}[B48TagType5]
 \proofstepwidestar[]{t_{\rm c}\in\mathbb N}{\FormulaRefAuto{B48TagType4}[thm]}
 \proofstepwidestar[]{\Succ(t_{\rm c})\in\mathbb N}{\FormulaRefAuto{PeanoSuccClosure}[ax]{1}}
 \proofstepwidestar[]{t_{\rm q}\in\mathbb N}{\FormulaRefAuto{B48CodeCarrier}[def]{2}}
\closeproofpartwideindR

\proofpartwideindR[Typisierung der Marke b]{t_{\rm b}\in\mathbb N}{t_{\rm b}\in\mathbb N}[B48TagType6]
 \proofstepwidestar[]{t_{\rm q}\in\mathbb N}{\FormulaRefAuto{B48TagType5}[thm]}
 \proofstepwidestar[]{\Succ(t_{\rm q})\in\mathbb N}{\FormulaRefAuto{PeanoSuccClosure}[ax]{1}}
 \proofstepwidestar[]{t_{\rm b}\in\mathbb N}{\FormulaRefAuto{B48CodeCarrier}[def]{2}}
\closeproofpartwideindR

\proofpartwideindR[Trennung der Marken e und m]{t_{\rm e}\neq t_{\rm m}}{t_{\rm e}\neq t_{\rm m}}[B48TagNe01]
 \proofstepwidestar[]{t_{\rm e}\in\mathbb N}{\FormulaRefAuto{B48TagType1}[thm]}
 \proofstepwidestar[]{0\neq\Succ(t_{\rm e})}{\FormulaRefAuto{PeanoZeroNeSuccessor}[thm]{1}}
 \proofstepwidestar[]{t_{\rm e}\neq t_{\rm m}}{\FormulaRefAuto{B48CodeCarrier}[def]{2}}
\closeproofpartwideindR

\proofpartwideindR[Trennung der Marken e und n]{t_{\rm e}\neq t_{\rm n}}{t_{\rm e}\neq t_{\rm n}}[B48TagNe02]
 \proofstepwidestar[]{t_{\rm m}\in\mathbb N}{\FormulaRefAuto{B48TagType2}[thm]}
 \proofstepwidestar[]{0\neq\Succ(t_{\rm m})}{\FormulaRefAuto{PeanoZeroNeSuccessor}[thm]{1}}
 \proofstepwidestar[]{t_{\rm e}\neq t_{\rm n}}{\FormulaRefAuto{B48CodeCarrier}[def]{2}}
\closeproofpartwideindR

\proofpartwideindR[Trennung der Marken e und c]{t_{\rm e}\neq t_{\rm c}}{t_{\rm e}\neq t_{\rm c}}[B48TagNe03]
 \proofstepwidestar[]{t_{\rm n}\in\mathbb N}{\FormulaRefAuto{B48TagType3}[thm]}
 \proofstepwidestar[]{0\neq\Succ(t_{\rm n})}{\FormulaRefAuto{PeanoZeroNeSuccessor}[thm]{1}}
 \proofstepwidestar[]{t_{\rm e}\neq t_{\rm c}}{\FormulaRefAuto{B48CodeCarrier}[def]{2}}
\closeproofpartwideindR

\proofpartwideindR[Trennung der Marken e und q]{t_{\rm e}\neq t_{\rm q}}{t_{\rm e}\neq t_{\rm q}}[B48TagNe04]
 \proofstepwidestar[]{t_{\rm c}\in\mathbb N}{\FormulaRefAuto{B48TagType4}[thm]}
 \proofstepwidestar[]{0\neq\Succ(t_{\rm c})}{\FormulaRefAuto{PeanoZeroNeSuccessor}[thm]{1}}
 \proofstepwidestar[]{t_{\rm e}\neq t_{\rm q}}{\FormulaRefAuto{B48CodeCarrier}[def]{2}}
\closeproofpartwideindR

\proofpartwideindR[Trennung der Marken e und b]{t_{\rm e}\neq t_{\rm b}}{t_{\rm e}\neq t_{\rm b}}[B48TagNe05]
 \proofstepwidestar[]{t_{\rm q}\in\mathbb N}{\FormulaRefAuto{B48TagType5}[thm]}
 \proofstepwidestar[]{0\neq\Succ(t_{\rm q})}{\FormulaRefAuto{PeanoZeroNeSuccessor}[thm]{1}}
 \proofstepwidestar[]{t_{\rm e}\neq t_{\rm b}}{\FormulaRefAuto{B48CodeCarrier}[def]{2}}
\closeproofpartwideindR

\proofpartwideindR[Trennung der Marken m und n]{t_{\rm m}\neq t_{\rm n}}{t_{\rm m}\neq t_{\rm n}}[B48TagNe12]
 \proofstepwidestar[]{t_{\rm e}\in\mathbb N}{\FormulaRefAuto{B48TagType1}[thm]}
 \proofstepwidestar[]{t_{\rm m}\in\mathbb N}{\FormulaRefAuto{B48TagType2}[thm]}
 \proofstepwidestar[]{t_{\rm e}\neq t_{\rm m}}{\FormulaRefAuto{B48TagNe01}[thm]}
 \proofstepwidestar[]{\Succ(t_{\rm e})\neq\Succ(t_{\rm m})}{\FormulaRefAuto{PeanoSuccPreservesNe}[thm]{1,2,3}}
 \proofstepwidestar[]{t_{\rm m}\neq t_{\rm n}}{\FormulaRefAuto{B48CodeCarrier}[def]{4}}
\closeproofpartwideindR

\proofpartwideindR[Trennung der Marken m und c]{t_{\rm m}\neq t_{\rm c}}{t_{\rm m}\neq t_{\rm c}}[B48TagNe13]
 \proofstepwidestar[]{t_{\rm e}\in\mathbb N}{\FormulaRefAuto{B48TagType1}[thm]}
 \proofstepwidestar[]{t_{\rm n}\in\mathbb N}{\FormulaRefAuto{B48TagType3}[thm]}
 \proofstepwidestar[]{t_{\rm e}\neq t_{\rm n}}{\FormulaRefAuto{B48TagNe02}[thm]}
 \proofstepwidestar[]{\Succ(t_{\rm e})\neq\Succ(t_{\rm n})}{\FormulaRefAuto{PeanoSuccPreservesNe}[thm]{1,2,3}}
 \proofstepwidestar[]{t_{\rm m}\neq t_{\rm c}}{\FormulaRefAuto{B48CodeCarrier}[def]{4}}
\closeproofpartwideindR

\proofpartwideindR[Trennung der Marken m und q]{t_{\rm m}\neq t_{\rm q}}{t_{\rm m}\neq t_{\rm q}}[B48TagNe14]
 \proofstepwidestar[]{t_{\rm e}\in\mathbb N}{\FormulaRefAuto{B48TagType1}[thm]}
 \proofstepwidestar[]{t_{\rm c}\in\mathbb N}{\FormulaRefAuto{B48TagType4}[thm]}
 \proofstepwidestar[]{t_{\rm e}\neq t_{\rm c}}{\FormulaRefAuto{B48TagNe03}[thm]}
 \proofstepwidestar[]{\Succ(t_{\rm e})\neq\Succ(t_{\rm c})}{\FormulaRefAuto{PeanoSuccPreservesNe}[thm]{1,2,3}}
 \proofstepwidestar[]{t_{\rm m}\neq t_{\rm q}}{\FormulaRefAuto{B48CodeCarrier}[def]{4}}
\closeproofpartwideindR

\proofpartwideindR[Trennung der Marken m und b]{t_{\rm m}\neq t_{\rm b}}{t_{\rm m}\neq t_{\rm b}}[B48TagNe15]
 \proofstepwidestar[]{t_{\rm e}\in\mathbb N}{\FormulaRefAuto{B48TagType1}[thm]}
 \proofstepwidestar[]{t_{\rm q}\in\mathbb N}{\FormulaRefAuto{B48TagType5}[thm]}
 \proofstepwidestar[]{t_{\rm e}\neq t_{\rm q}}{\FormulaRefAuto{B48TagNe04}[thm]}
 \proofstepwidestar[]{\Succ(t_{\rm e})\neq\Succ(t_{\rm q})}{\FormulaRefAuto{PeanoSuccPreservesNe}[thm]{1,2,3}}
 \proofstepwidestar[]{t_{\rm m}\neq t_{\rm b}}{\FormulaRefAuto{B48CodeCarrier}[def]{4}}
\closeproofpartwideindR

\proofpartwideindR[Trennung der Marken n und c]{t_{\rm n}\neq t_{\rm c}}{t_{\rm n}\neq t_{\rm c}}[B48TagNe23]
 \proofstepwidestar[]{t_{\rm m}\in\mathbb N}{\FormulaRefAuto{B48TagType2}[thm]}
 \proofstepwidestar[]{t_{\rm n}\in\mathbb N}{\FormulaRefAuto{B48TagType3}[thm]}
 \proofstepwidestar[]{t_{\rm m}\neq t_{\rm n}}{\FormulaRefAuto{B48TagNe12}[thm]}
 \proofstepwidestar[]{\Succ(t_{\rm m})\neq\Succ(t_{\rm n})}{\FormulaRefAuto{PeanoSuccPreservesNe}[thm]{1,2,3}}
 \proofstepwidestar[]{t_{\rm n}\neq t_{\rm c}}{\FormulaRefAuto{B48CodeCarrier}[def]{4}}
\closeproofpartwideindR

\proofpartwideindR[Trennung der Marken n und q]{t_{\rm n}\neq t_{\rm q}}{t_{\rm n}\neq t_{\rm q}}[B48TagNe24]
 \proofstepwidestar[]{t_{\rm m}\in\mathbb N}{\FormulaRefAuto{B48TagType2}[thm]}
 \proofstepwidestar[]{t_{\rm c}\in\mathbb N}{\FormulaRefAuto{B48TagType4}[thm]}
 \proofstepwidestar[]{t_{\rm m}\neq t_{\rm c}}{\FormulaRefAuto{B48TagNe13}[thm]}
 \proofstepwidestar[]{\Succ(t_{\rm m})\neq\Succ(t_{\rm c})}{\FormulaRefAuto{PeanoSuccPreservesNe}[thm]{1,2,3}}
 \proofstepwidestar[]{t_{\rm n}\neq t_{\rm q}}{\FormulaRefAuto{B48CodeCarrier}[def]{4}}
\closeproofpartwideindR

\proofpartwideindR[Trennung der Marken n und b]{t_{\rm n}\neq t_{\rm b}}{t_{\rm n}\neq t_{\rm b}}[B48TagNe25]
 \proofstepwidestar[]{t_{\rm m}\in\mathbb N}{\FormulaRefAuto{B48TagType2}[thm]}
 \proofstepwidestar[]{t_{\rm q}\in\mathbb N}{\FormulaRefAuto{B48TagType5}[thm]}
 \proofstepwidestar[]{t_{\rm m}\neq t_{\rm q}}{\FormulaRefAuto{B48TagNe14}[thm]}
 \proofstepwidestar[]{\Succ(t_{\rm m})\neq\Succ(t_{\rm q})}{\FormulaRefAuto{PeanoSuccPreservesNe}[thm]{1,2,3}}
 \proofstepwidestar[]{t_{\rm n}\neq t_{\rm b}}{\FormulaRefAuto{B48CodeCarrier}[def]{4}}
\closeproofpartwideindR

\proofpartwideindR[Trennung der Marken c und q]{t_{\rm c}\neq t_{\rm q}}{t_{\rm c}\neq t_{\rm q}}[B48TagNe34]
 \proofstepwidestar[]{t_{\rm n}\in\mathbb N}{\FormulaRefAuto{B48TagType3}[thm]}
 \proofstepwidestar[]{t_{\rm c}\in\mathbb N}{\FormulaRefAuto{B48TagType4}[thm]}
 \proofstepwidestar[]{t_{\rm n}\neq t_{\rm c}}{\FormulaRefAuto{B48TagNe23}[thm]}
 \proofstepwidestar[]{\Succ(t_{\rm n})\neq\Succ(t_{\rm c})}{\FormulaRefAuto{PeanoSuccPreservesNe}[thm]{1,2,3}}
 \proofstepwidestar[]{t_{\rm c}\neq t_{\rm q}}{\FormulaRefAuto{B48CodeCarrier}[def]{4}}
\closeproofpartwideindR

\proofpartwideindR[Trennung der Marken c und b]{t_{\rm c}\neq t_{\rm b}}{t_{\rm c}\neq t_{\rm b}}[B48TagNe35]
 \proofstepwidestar[]{t_{\rm n}\in\mathbb N}{\FormulaRefAuto{B48TagType3}[thm]}
 \proofstepwidestar[]{t_{\rm q}\in\mathbb N}{\FormulaRefAuto{B48TagType5}[thm]}
 \proofstepwidestar[]{t_{\rm n}\neq t_{\rm q}}{\FormulaRefAuto{B48TagNe24}[thm]}
 \proofstepwidestar[]{\Succ(t_{\rm n})\neq\Succ(t_{\rm q})}{\FormulaRefAuto{PeanoSuccPreservesNe}[thm]{1,2,3}}
 \proofstepwidestar[]{t_{\rm c}\neq t_{\rm b}}{\FormulaRefAuto{B48CodeCarrier}[def]{4}}
\closeproofpartwideindR

\proofpartwideindR[Trennung der Marken q und b]{t_{\rm q}\neq t_{\rm b}}{t_{\rm q}\neq t_{\rm b}}[B48TagNe45]
 \proofstepwidestar[]{t_{\rm c}\in\mathbb N}{\FormulaRefAuto{B48TagType4}[thm]}
 \proofstepwidestar[]{t_{\rm q}\in\mathbb N}{\FormulaRefAuto{B48TagType5}[thm]}
 \proofstepwidestar[]{t_{\rm c}\neq t_{\rm q}}{\FormulaRefAuto{B48TagNe34}[thm]}
 \proofstepwidestar[]{\Succ(t_{\rm c})\neq\Succ(t_{\rm q})}{\FormulaRefAuto{PeanoSuccPreservesNe}[thm]{1,2,3}}
 \proofstepwidestar[]{t_{\rm q}\neq t_{\rm b}}{\FormulaRefAuto{B48CodeCarrier}[def]{4}}
\closeproofpartwideindR

\proofpartwideindR[Typisierung eines Blattzeichens]{k\in\mathbb N,\ i\in\mathbb N,\ j\in\mathbb N\vdash(k,(i,j))\in\Sigma}{k\in\mathbb N,\ i\in\mathbb N,\ j\in\mathbb N\vdash(k,(i,j))\in\Sigma}[B48LabelType]
 \proofstepwidestar[1]{k\in\mathbb N}{\rA}
 \proofstepwidestar[2]{i\in\mathbb N}{\rA}
 \proofstepwidestar[3]{j\in\mathbb N}{\rA}
 \proofstepwidestar[2,3]{(i,j)\in\mathbb N\times\mathbb N}{\FormulaRefAuto{(a,b)\in A\times B\eqvdash a\in A\land b\in B}[thm]{\rAI{2,3}}}
 \proofstepwidestar[1,2,3]{(k,(i,j))\in\mathbb N\times(\mathbb N\times\mathbb N)}{\FormulaRefAuto{(a,b)\in A\times B\eqvdash a\in A\land b\in B}[thm]{\rAI{1,4}}}
 \proofstepwidestar[1,2,3]{(k,(i,j))\in\Sigma}{\FormulaRefAuto{B48CodeCarrier}[def]{5}}
\closeproofpartwideindR

\proofpartwideindR[Blattcodes sind Baumcodes]{c\in\Sigma\vdash\mathsf l(c)\in\mathcal T}{c\in\Sigma\vdash\mathsf l(c)\in\mathcal T}[B48LeafType]
 \proofstepwidestar[1]{c\in\Sigma}{\rA}
 \proofstepwidestar[1]{\LeafTree{\Sigma}{c}\in\BracketTreeSet{\Sigma}}{\FormulaRefAuto{TreeLeafClosure}[thm]{1}}
 \proofstepwidestar[1]{\mathsf l(c)\in\mathcal T}{\FormulaRefAuto{B48CodeCarrier}[def]{2}}
\closeproofpartwideindR

\proofpartwideindR[Knotencodes sind Baumcodes]{p\in\mathcal T,\ q\in\mathcal T\vdash\mathsf n(p,q)\in\mathcal T}{p\in\mathcal T,\ q\in\mathcal T\vdash\mathsf n(p,q)\in\mathcal T}[B48NodeType]
 \proofstepwidestar[1]{p\in\mathcal T}{\rA}
 \proofstepwidestar[2]{q\in\mathcal T}{\rA}
 \proofstepwidestar[1,2]{\NodeTree{\Sigma}{p}{q}\in\BracketTreeSet{\Sigma}}{\FormulaRefAuto{TreeNodeClosure}[thm]{1,2}}
 \proofstepwidestar[1,2]{\mathsf n(p,q)\in\mathcal T}{\FormulaRefAuto{B48CodeCarrier}[def]{3}}
\closeproofpartwideindR

\proofpartwideindR[Baumcode der Negationsmarke]{\mathsf n_0\in\mathcal T}{\mathsf n_0\in\mathcal T}[B48NegRootType]
 \proofstepwidestar[]{t_{\rm n}\in\mathbb N}{\FormulaRefAuto{B48TagType3}[thm]}
 \proofstepwidestar[]{0\in\mathbb N}{\FormulaRefAuto{PeanoZero}[ax]}
 \proofstepwidestar[]{(t_{\rm n},(0,0))\in\Sigma}{\FormulaRefAuto{B48LabelType}[thm]{1,2,2}}
 \proofstepwidestar[]{\mathsf l((t_{\rm n},(0,0)))\in\mathcal T}{\FormulaRefAuto{B48LeafType}[thm]{3}}
 \proofstepwidestar[]{\mathsf n_0\in\mathcal T}{\FormulaRefAuto{B48FormulaConstructors}[def]{4}}
\closeproofpartwideindR

\proofpartwideindR[Baumcode der Konjunktionsmarke]{\mathsf c_0\in\mathcal T}{\mathsf c_0\in\mathcal T}[B48AndRootType]
 \proofstepwidestar[]{t_{\rm c}\in\mathbb N}{\FormulaRefAuto{B48TagType4}[thm]}
 \proofstepwidestar[]{0\in\mathbb N}{\FormulaRefAuto{PeanoZero}[ax]}
 \proofstepwidestar[]{(t_{\rm c},(0,0))\in\Sigma}{\FormulaRefAuto{B48LabelType}[thm]{1,2,2}}
 \proofstepwidestar[]{\mathsf l((t_{\rm c},(0,0)))\in\mathcal T}{\FormulaRefAuto{B48LeafType}[thm]{3}}
 \proofstepwidestar[]{\mathsf c_0\in\mathcal T}{\FormulaRefAuto{B48FormulaConstructors}[def]{4}}
\closeproofpartwideindR

\proofpartwideindR[Baumcode einer Quantorenmarke]{i\in\mathbb N\vdash\mathsf q_i\in\mathcal T}{i\in\mathbb N\vdash\mathsf q_i\in\mathcal T}[B48ExistsRootType]
 \proofstepwidestar[1]{i\in\mathbb N}{\rA}
 \proofstepwidestar[]{t_{\rm q}\in\mathbb N}{\FormulaRefAuto{B48TagType5}[thm]}
 \proofstepwidestar[]{0\in\mathbb N}{\FormulaRefAuto{PeanoZero}[ax]}
 \proofstepwidestar[1]{(t_{\rm q},(i,0))\in\Sigma}{\FormulaRefAuto{B48LabelType}[thm]{2,1,3}}
 \proofstepwidestar[1]{\mathsf l((t_{\rm q},(i,0)))\in\mathcal T}{\FormulaRefAuto{B48LeafType}[thm]{4}}
 \proofstepwidestar[1]{\mathsf q_i\in\mathcal T}{\FormulaRefAuto{B48FormulaConstructors}[def]{5}}
\closeproofpartwideindR

\proofpartwideindR[Typisierung des Negationscodes]{p\in\mathcal T\vdash\mathsf{neg}(p)\in\mathcal T}{p\in\mathcal T\vdash\mathsf{neg}(p)\in\mathcal T}[B48NegCodeType]
 \proofstepwidestar[1]{p\in\mathcal T}{\rA}
 \proofstepwidestar[]{\mathsf n_0\in\mathcal T}{\FormulaRefAuto{B48NegRootType}[thm]}
 \proofstepwidestar[1]{\mathsf n(\mathsf n_0,p)\in\mathcal T}{\FormulaRefAuto{B48NodeType}[thm]{2,1}}
 \proofstepwidestar[1]{\mathsf{neg}(p)\in\mathcal T}{\FormulaRefAuto{B48FormulaConstructors}[def]{3}}
\closeproofpartwideindR

\proofpartwideindR[Typisierung des Konjunktionscodes]{p\in\mathcal T,\ q\in\mathcal T\vdash\mathsf{and}(p,q)\in\mathcal T}{p\in\mathcal T,\ q\in\mathcal T\vdash\mathsf{and}(p,q)\in\mathcal T}[B48AndCodeType]
 \proofstepwidestar[1]{p\in\mathcal T}{\rA}
 \proofstepwidestar[2]{q\in\mathcal T}{\rA}
 \proofstepwidestar[]{\mathsf c_0\in\mathcal T}{\FormulaRefAuto{B48AndRootType}[thm]}
 \proofstepwidestar[1]{\mathsf n(\mathsf c_0,p)\in\mathcal T}{\FormulaRefAuto{B48NodeType}[thm]{3,1}}
 \proofstepwidestar[1,2]{\mathsf n(\mathsf n(\mathsf c_0,p),q)\in\mathcal T}{\FormulaRefAuto{B48NodeType}[thm]{4,2}}
 \proofstepwidestar[1,2]{\mathsf{and}(p,q)\in\mathcal T}{\FormulaRefAuto{B48FormulaConstructors}[def]{5}}
\closeproofpartwideindR

\proofpartwideindR[Typisierung des Existenzcodes]{i\in\mathbb N,\ p\in\mathcal T\vdash\mathsf{ex}_i(p)\in\mathcal T}{i\in\mathbb N,\ p\in\mathcal T\vdash\mathsf{ex}_i(p)\in\mathcal T}[B48ExistsCodeType]
 \proofstepwidestar[1]{i\in\mathbb N}{\rA}
 \proofstepwidestar[2]{p\in\mathcal T}{\rA}
 \proofstepwidestar[1]{\mathsf q_i\in\mathcal T}{\FormulaRefAuto{B48ExistsRootType}[thm]{1}}
 \proofstepwidestar[1,2]{\mathsf n(\mathsf q_i,p)\in\mathcal T}{\FormulaRefAuto{B48NodeType}[thm]{3,2}}
 \proofstepwidestar[1,2]{\mathsf{ex}_i(p)\in\mathcal T}{\FormulaRefAuto{B48FormulaConstructors}[def]{4}}
\closeproofpartwideindR

\par\medskip\noindent\textbf{B. Belegungsmenge und Formelabschluss}\par
\proofpartwideindR[Wahrer Zweig des Fallterms]{P\vdash\BXLVIIIIf{P}{U}{V}=U}{P\vdash\BXLVIIIIf{P}{U}{V}=U}[B48IfTrue]
 \proofstepwidestar[1]{P}{\rA}
 \proofstepwidestar[1]{\IfThenElse{P}{U}{V}=U}{\FormulaRefAuto{P\vdash\IfThenElse{P}{a}{b}=a}[thm]{1}}
 \proofstepwidestar[1]{\BXLVIIIIf{P}{U}{V}=U}{\FormulaRefAuto{B48IfNotation}[def]{2}}
\closeproofpartwideindR

\proofpartwideindR[Falscher Zweig des Fallterms]{\neg P\vdash\BXLVIIIIf{P}{U}{V}=V}{\neg P\vdash\BXLVIIIIf{P}{U}{V}=V}[B48IfFalse]
 \proofstepwidestar[1]{\neg P}{\rA}
 \proofstepwidestar[1]{\IfThenElse{P}{U}{V}=V}{\FormulaRefAuto{\neg P\vdash\IfThenElse{P}{a}{b}=b}[thm]{1}}
 \proofstepwidestar[1]{\BXLVIIIIf{P}{U}{V}=V}{\FormulaRefAuto{B48IfNotation}[def]{2}}
\closeproofpartwideindR

\proofpartwideindR[Typisierung beider Fallzweige]{U\in X,\ V\in X\vdash\BXLVIIIIf{P}{U}{V}\in X}{U\in X,\ V\in X\vdash\BXLVIIIIf{P}{U}{V}\in X}[B48IfTyped]
 \proofstepwidestar[1]{U\in X}{\rA}
 \proofstepwidestar[2]{V\in X}{\rA}
 \proofstepwidestar[]{P\lor\neg P}{\FormulaRefAuto{P\lor\neg P}[thm]}
 \proofstepwidestar[4]{P}{\rA}
 \proofstepwidestar[4]{\BXLVIIIIf{P}{U}{V}=U}{\FormulaRefAuto{B48IfTrue}[thm]{4}}
 \proofstepwidestar[1,4]{\BXLVIIIIf{P}{U}{V}\in X}{\FormulaRefAuto{a=b,\ b\in C\vdash a\in C}[thm]{5,1}}
 \proofstepwidestar[7]{\neg P}{\rA}
 \proofstepwidestar[7]{\BXLVIIIIf{P}{U}{V}=V}{\FormulaRefAuto{B48IfFalse}[thm]{7}}
 \proofstepwidestar[2,7]{\BXLVIIIIf{P}{U}{V}\in X}{\FormulaRefAuto{a=b,\ b\in C\vdash a\in C}[thm]{8,2}}
 \proofstepwidestar[1,2]{\BXLVIIIIf{P}{U}{V}\in X}{\rOE{3,4:6,7:9}}
\closeproofpartwideindR

\proofpartwideindR[Beschraenkte Praedikate liefern Wertemengen]{\{x\in A\mid P(x)\}\in\powerset(A)}{\{x\in A\mid P(x)\}\in\powerset(A)}[B48SeparationPower]
 \proofstepwidestar[]{\{x\in A\mid P(x)\}\subseteq A}{\FormulaRefAuto{\{x\in A\mid P(x)\}\subseteq A}[thm]}
 \proofstepwidestar[]{\{x\in A\mid P(x)\}\in\powerset(A)}{\FormulaRefAuto{x\in\powerset(A)\eqvdash x\subseteq A}[thm]{1}}
\closeproofpartwideindR

\proofpartwideindR[Elementkriterium der Belegungsmenge]{s\in A_D\eqvdash s\colon\mathbb N\to D}{s\in A_D\eqvdash s\colon\mathbb N\to D}[B48AssignmentCriterion]
 \proofstepwidestar[1]{s\in A_D}{\rA}
 \proofstepwidestar[1]{s\in\powerset(\mathbb N\times D)\land s\colon\mathbb N\to D}{\FormulaRefAuto{B48Assignments}[def]{1}}
 \proofstepwidestar[1]{s\colon\mathbb N\to D}{\rAEb{2}}
 \proofstepwidestar[]{s\in A_D\to s\colon\mathbb N\to D}{\rRI{1,3}}
 \proofstepwidestar[5]{s\colon\mathbb N\to D}{\rA}
 \proofstepwidestar[5]{s\subseteq\mathbb N\times D}{\FormulaRefAuto{FunctionGraphSubsetProduct}[thm]{5}}
 \proofstepwidestar[5]{s\in\powerset(\mathbb N\times D)}{\FormulaRefAuto{x\in\powerset(A)\eqvdash x\subseteq A}[thm]{6}}
 \proofstepwidestar[5]{s\in A_D}{\FormulaRefAuto{B48Assignments}[def]{\rAI{7,5}}}
 \proofstepwidestar[]{(s\colon\mathbb N\to D)\to s\in A_D}{\rRI{5,8}}
 \proofstepwidestar[]{s\in A_D\leftrightarrow s\colon\mathbb N\to D}{\rLRI{4,9}}
\closeproofpartwideindR

\proofpartwideindR[Ein Belegungswechsel bleibt eine Belegung]{s\in A_D,\ i\in\mathbb N,\ a\in D\vdash s[i\mapsto a]\in A_D}{s\in A_D,\ i\in\mathbb N,\ a\in D\vdash s[i\mapsto a]\in A_D}[B48AssignmentUpdate]
 \proofstepwidestar[1]{s\in A_D}{\rA}
 \proofstepwidestar[2]{i\in\mathbb N}{\rA}
 \proofstepwidestar[3]{a\in D}{\rA}
 \proofstepwidestar[1]{s\colon\mathbb N\to D}{\FormulaRefAuto{B48AssignmentCriterion}[thm]{1}}
 \proofstepwidestar[5]{j\in\mathbb N}{\rA}
 \proofstepwidestar[1,5]{s(j)\in D}{\FormulaRefAuto{F\colon A\to B,\ x\in A\vdash F(x)\in B}[thm]{4,5}}
 \proofstepwidestar[1,3,5]{\BXLVIIIIf{j=i}{a}{s(j)}\in D}{\FormulaRefAuto{B48IfTyped}[thm]{3,6}}
 \proofstepwidestar[1,3]{\forall j\in\mathbb N\;(\BXLVIIIIf{j=i}{a}{s(j)}\in D)}{\rUI{\rRI{5,7}}}
 \proofstepwidestar[1,3]{\{(j,b)\in\mathbb N\times D\mid b=\BXLVIIIIf{j=i}{a}{s(j)}\}\colon\mathbb N\to D}{\FormulaRefAuto{TypedTermGraphFunction}[thm]{8}}
 \proofstepwidestar[1,3]{s[i\mapsto a]\colon\mathbb N\to D}{\FormulaRefAuto{B48Assignments}[def]{9}}
 \proofstepwidestar[1,3]{s[i\mapsto a]\in A_D}{\FormulaRefAuto{B48AssignmentCriterion}[thm]{10}}
\closeproofpartwideindR

\proofpartwideindR[Abschlussklausel 1]{\mathsf{Cl}(X)\vdash X\subseteq\mathcal T}{\mathsf{Cl}(X)\vdash X\subseteq\mathcal T}[B48ClDomain]
 \proofstepwidestar[1]{\mathsf{Cl}(X)}{\rA}
 \proofstepwidestar[1]{X\subseteq\mathcal T}{\FormulaRefAuto{B48ClDomainAxiom}[ax]{1}}
\closeproofpartwideindR

\proofpartwideindR[Abschlussklausel 2]{\mathsf{Cl}(X)\vdash \forall c\in\Sigma_{\rm at}\;(\mathsf l(c)\in X)}{\mathsf{Cl}(X)\vdash \forall c\in\Sigma_{\rm at}\;(\mathsf l(c)\in X)}[B48ClAtoms]
 \proofstepwidestar[1]{\mathsf{Cl}(X)}{\rA}
 \proofstepwidestar[1]{\forall c\in\Sigma_{\rm at}\;(\mathsf l(c)\in X)}{\FormulaRefAuto{B48ClAtomsAxiom}[ax]{1}}
\closeproofpartwideindR

\proofpartwideindR[Abschlussklausel 3]{\mathsf{Cl}(X)\vdash \forall p\in X\;(\mathsf{neg}(p)\in X)}{\mathsf{Cl}(X)\vdash \forall p\in X\;(\mathsf{neg}(p)\in X)}[B48ClNeg]
 \proofstepwidestar[1]{\mathsf{Cl}(X)}{\rA}
 \proofstepwidestar[1]{\forall p\in X\;(\mathsf{neg}(p)\in X)}{\FormulaRefAuto{B48ClNegAxiom}[ax]{1}}
\closeproofpartwideindR

\proofpartwideindR[Abschlussklausel 4]{\mathsf{Cl}(X)\vdash \forall p\in X\,\forall q\in X\;(\mathsf{and}(p,q)\in X)}{\mathsf{Cl}(X)\vdash \forall p\in X\,\forall q\in X\;(\mathsf{and}(p,q)\in X)}[B48ClAnd]
 \proofstepwidestar[1]{\mathsf{Cl}(X)}{\rA}
 \proofstepwidestar[1]{\forall p\in X\,\forall q\in X\;(\mathsf{and}(p,q)\in X)}{\FormulaRefAuto{B48ClAndAxiom}[ax]{1}}
\closeproofpartwideindR

\proofpartwideindR[Abschlussklausel 5]{\mathsf{Cl}(X)\vdash \forall i\in\mathbb N\,\forall p\in X\;(\mathsf{ex}_i(p)\in X)}{\mathsf{Cl}(X)\vdash \forall i\in\mathbb N\,\forall p\in X\;(\mathsf{ex}_i(p)\in X)}[B48ClExists]
 \proofstepwidestar[1]{\mathsf{Cl}(X)}{\rA}
 \proofstepwidestar[1]{\forall i\in\mathbb N\,\forall p\in X\;(\mathsf{ex}_i(p)\in X)}{\FormulaRefAuto{B48ClExistsAxiom}[ax]{1}}
\closeproofpartwideindR

\proofpartwideindR[Elementkriterium fuer Formelcodes]{p\in\mathcal F\eqvdash p\in\mathcal T\land\forall X(\mathsf{Cl}(X)\to p\in X)}{p\in\mathcal F\eqvdash p\in\mathcal T\land\forall X(\mathsf{Cl}(X)\to p\in X)}[B48FormulaMember]
 \proofstepwidestar[]{p\in\mathcal F\leftrightarrow p\in\mathcal T\land\forall X(\mathsf{Cl}(X)\to p\in X)}{\FormulaRefAuto{B48FormulaSet}[def]}
\closeproofpartwideindR

\proofpartwideindR[Jeder Formelcode ist ein Baumcode]{p\in\mathcal F\vdash p\in\mathcal T}{p\in\mathcal F\vdash p\in\mathcal T}[B48FormulaInTree]
 \proofstepwidestar[1]{p\in\mathcal F}{\rA}
 \proofstepwidestar[1]{p\in\mathcal T\land\forall X(\mathsf{Cl}(X)\to p\in X)}{\FormulaRefAuto{B48FormulaMember}[thm]{1}}
 \proofstepwidestar[1]{p\in\mathcal T}{\rAEa{2}}
\closeproofpartwideindR

\proofpartwideindR[Minimalitaet der Formelcodes]{\mathsf{Cl}(X),\ p\in\mathcal F\vdash p\in X}{\mathsf{Cl}(X),\ p\in\mathcal F\vdash p\in X}[B48FormulaMinimalMember]
 \proofstepwidestar[1]{\mathsf{Cl}(X)}{\rA}
 \proofstepwidestar[2]{p\in\mathcal F}{\rA}
 \proofstepwidestar[2]{p\in\mathcal T\land\forall Y(\mathsf{Cl}(Y)\to p\in Y)}{\FormulaRefAuto{B48FormulaMember}[thm]{2}}
 \proofstepwidestar[2]{\forall Y(\mathsf{Cl}(Y)\to p\in Y)}{\rAEb{3}}
 \proofstepwidestar[2]{\mathsf{Cl}(X)\to p\in X}{\rUE{4}}
 \proofstepwidestar[1,2]{p\in X}{\rRE{5,1}}
\closeproofpartwideindR

\proofpartwideindR[Atomare Zeichen liegen im Alphabet]{c\in\Sigma_{\rm at}\vdash c\in\Sigma}{c\in\Sigma_{\rm at}\vdash c\in\Sigma}[B48AtomicLabelType]
 \proofstepwidestar[1]{c\in\Sigma_{\rm at}}{\rA}
 \proofstepwidestar[]{\Sigma_{\rm at}\subseteq\Sigma}{\FormulaRefAuto{\{x\in A\mid P(x)\}\subseteq A}[thm]}
 \proofstepwidestar[1]{c\in\Sigma}{\FormulaRefAuto{A\subseteq B,\ x\in A\vdash x\in B}[thm]{2,1}}
\closeproofpartwideindR

\proofpartwideindR[Atomare Formelcodes]{c\in\Sigma_{\rm at}\vdash\mathsf l(c)\in\mathcal F}{c\in\Sigma_{\rm at}\vdash\mathsf l(c)\in\mathcal F}[B48FormulaAtom]
 \proofstepwidestar[1]{c\in\Sigma_{\rm at}}{\rA}
 \proofstepwidestar[1]{c\in\Sigma}{\FormulaRefAuto{B48AtomicLabelType}[thm]{1}}
 \proofstepwidestar[1]{\mathsf l(c)\in\mathcal T}{\FormulaRefAuto{B48LeafType}[thm]{2}}
 \proofstepwidestar[4]{\mathsf{Cl}(X)}{\rA}
 \proofstepwidestar[4]{\forall d\in\Sigma_{\rm at}\;(\mathsf l(d)\in X)}{\FormulaRefAuto{B48ClAtoms}[thm]{4}}
 \proofstepwidestar[1,4]{\mathsf l(c)\in X}{\rRE{\rUE{5},1}}
 \proofstepwidestar[1]{\forall X(\mathsf{Cl}(X)\to\mathsf l(c)\in X)}{\rUI{\rRI{4,6}}}
 \proofstepwidestar[1]{\mathsf l(c)\in\mathcal F}{\FormulaRefAuto{B48FormulaMember}[thm]{\rAI{3,7}}}
\closeproofpartwideindR

\proofpartwideindR[Abschluss unter Negationscodes]{p\in\mathcal F\vdash\mathsf{neg}(p)\in\mathcal F}{p\in\mathcal F\vdash\mathsf{neg}(p)\in\mathcal F}[B48FormulaNeg]
 \proofstepwidestar[1]{p\in\mathcal F}{\rA}
 \proofstepwidestar[1]{p\in\mathcal T}{\FormulaRefAuto{B48FormulaInTree}[thm]{1}}
 \proofstepwidestar[1]{\mathsf{neg}(p)\in\mathcal T}{\FormulaRefAuto{B48NegCodeType}[thm]{2}}
 \proofstepwidestar[4]{\mathsf{Cl}(X)}{\rA}
 \proofstepwidestar[1,4]{p\in X}{\FormulaRefAuto{B48FormulaMinimalMember}[thm]{4,1}}
 \proofstepwidestar[4]{\forall q\in X\;(\mathsf{neg}(q)\in X)}{\FormulaRefAuto{B48ClNeg}[thm]{4}}
 \proofstepwidestar[1,4]{\mathsf{neg}(p)\in X}{\rRE{\rUE{6},5}}
 \proofstepwidestar[1]{\forall X(\mathsf{Cl}(X)\to\mathsf{neg}(p)\in X)}{\rUI{\rRI{4,7}}}
 \proofstepwidestar[1]{\mathsf{neg}(p)\in\mathcal F}{\FormulaRefAuto{B48FormulaMember}[thm]{\rAI{3,8}}}
\closeproofpartwideindR

\proofpartwideindR[Abschluss unter Konjunktionscodes]{p\in\mathcal F,\ q\in\mathcal F\vdash\mathsf{and}(p,q)\in\mathcal F}{p\in\mathcal F,\ q\in\mathcal F\vdash\mathsf{and}(p,q)\in\mathcal F}[B48FormulaAnd]
 \proofstepwidestar[1]{p\in\mathcal F}{\rA}
 \proofstepwidestar[2]{q\in\mathcal F}{\rA}
 \proofstepwidestar[1]{p\in\mathcal T}{\FormulaRefAuto{B48FormulaInTree}[thm]{1}}
 \proofstepwidestar[2]{q\in\mathcal T}{\FormulaRefAuto{B48FormulaInTree}[thm]{2}}
 \proofstepwidestar[1,2]{\mathsf{and}(p,q)\in\mathcal T}{\FormulaRefAuto{B48AndCodeType}[thm]{3,4}}
 \proofstepwidestar[6]{\mathsf{Cl}(X)}{\rA}
 \proofstepwidestar[1,6]{p\in X}{\FormulaRefAuto{B48FormulaMinimalMember}[thm]{6,1}}
 \proofstepwidestar[2,6]{q\in X}{\FormulaRefAuto{B48FormulaMinimalMember}[thm]{6,2}}
 \proofstepwidestar[6]{\forall u\in X\,\forall v\in X\;(\mathsf{and}(u,v)\in X)}{\FormulaRefAuto{B48ClAnd}[thm]{6}}
 \proofstepwidestar[1,2,6]{\mathsf{and}(p,q)\in X}{\rRE{\rUE{\rRE{\rUE{9},7}},8}}
 \proofstepwidestar[1,2]{\forall X(\mathsf{Cl}(X)\to\mathsf{and}(p,q)\in X)}{\rUI{\rRI{6,10}}}
 \proofstepwidestar[1,2]{\mathsf{and}(p,q)\in\mathcal F}{\FormulaRefAuto{B48FormulaMember}[thm]{\rAI{5,11}}}
\closeproofpartwideindR

\proofpartwideindR[Abschluss unter Existenzcodes]{i\in\mathbb N,\ p\in\mathcal F\vdash\mathsf{ex}_i(p)\in\mathcal F}{i\in\mathbb N,\ p\in\mathcal F\vdash\mathsf{ex}_i(p)\in\mathcal F}[B48FormulaExists]
 \proofstepwidestar[1]{i\in\mathbb N}{\rA}
 \proofstepwidestar[2]{p\in\mathcal F}{\rA}
 \proofstepwidestar[2]{p\in\mathcal T}{\FormulaRefAuto{B48FormulaInTree}[thm]{2}}
 \proofstepwidestar[1,2]{\mathsf{ex}_i(p)\in\mathcal T}{\FormulaRefAuto{B48ExistsCodeType}[thm]{1,3}}
 \proofstepwidestar[5]{\mathsf{Cl}(X)}{\rA}
 \proofstepwidestar[2,5]{p\in X}{\FormulaRefAuto{B48FormulaMinimalMember}[thm]{5,2}}
 \proofstepwidestar[5]{\forall j\in\mathbb N\,\forall q\in X\;(\mathsf{ex}_j(q)\in X)}{\FormulaRefAuto{B48ClExists}[thm]{5}}
 \proofstepwidestar[1,2,5]{\mathsf{ex}_i(p)\in X}{\rRE{\rUE{\rRE{\rUE{7},1}},6}}
 \proofstepwidestar[1,2]{\forall X(\mathsf{Cl}(X)\to\mathsf{ex}_i(p)\in X)}{\rUI{\rRI{5,8}}}
 \proofstepwidestar[1,2]{\mathsf{ex}_i(p)\in\mathcal F}{\FormulaRefAuto{B48FormulaMember}[thm]{\rAI{4,9}}}
\closeproofpartwideindR

\proofpartwideindR[Induktion auf der definierten Formelmenge]{\mathsf{Cl}(\{p\in\mathcal F\mid P(p)\})\vdash\forall p\in\mathcal F\;P(p)}{\mathsf{Cl}(\{p\in\mathcal F\mid P(p)\})\vdash\forall p\in\mathcal F\;P(p)}[B48FormulaInduction]
 \proofstepwidestar[1]{\mathsf{Cl}(\{p\in\mathcal F\mid P(p)\})}{\rA}
 \proofstepwidestar[2]{q\in\mathcal F}{\rA}
 \proofstepwidestar[1,2]{q\in\{p\in\mathcal F\mid P(p)\}}{\FormulaRefAuto{B48FormulaMinimalMember}[thm]{1,2}}
 \proofstepwidestar[1,2]{q\in\mathcal F\land P(q)}{\FormulaRefAuto{x\in\{u\in A\mid P(u)\}\eqvdash x\in A\land P(x)}[thm]{3}}
 \proofstepwidestar[1,2]{P(q)}{\rAEb{4}}
 \proofstepwidestar[1]{\forall q\in\mathcal F\;P(q)}{\rUI{\rRI{2,5}}}
\closeproofpartwideindR

\par\medskip\noindent\textbf{C. Unterscheidung der Auswertungsfaelle}\par
\proofpartwideindR[Baumcodes sind nichtleere Woerter]{p\in\mathcal T\vdash p\in\NonemptyWords{\TreeAlphabet{\Sigma}}}{p\in\mathcal T\vdash p\in\NonemptyWords{\TreeAlphabet{\Sigma}}}[B48TreeWord]
 \proofstepwidestar[1]{p\in\mathcal T}{\rA}
 \proofstepwidestar[1]{p\in\NonemptyWords{\TreeAlphabet{\Sigma}}\land\exists n\in\mathbb N\;(p\in\TreeStage{\Sigma}{n})}{\FormulaRefAuto{FullPlanarBinaryTreeSetDef}[def]{1}}
 \proofstepwidestar[1]{p\in\NonemptyWords{\TreeAlphabet{\Sigma}}}{\rAEa{2}}
\closeproofpartwideindR

\proofpartwideindR[Injektivitaet der Blattcodes]{c\in\Sigma,\ d\in\Sigma,\ \mathsf l(c)=\mathsf l(d)\vdash c=d}{c\in\Sigma,\ d\in\Sigma,\ \mathsf l(c)=\mathsf l(d)\vdash c=d}[B48LeafInjective]
 \proofstepwidestar[1]{c\in\Sigma}{\rA}
 \proofstepwidestar[2]{d\in\Sigma}{\rA}
 \proofstepwidestar[3]{\mathsf l(c)=\mathsf l(d)}{\rA}
 \proofstepwidestar[1,2]{\mathsf l(c)=\mathsf l(d)\leftrightarrow c=d}{\FormulaRefAuto{TreeLeafConstructorInjective}[thm]{1,2}}
 \proofstepwidestar[1,2,3]{c=d}{\rLREa{4,3}}
\closeproofpartwideindR

\proofpartwideindR[Verschiedene Marken liefern verschiedene Blattcodes]{(k,a)\in\Sigma,\ (l,b)\in\Sigma,\ k\neq l\vdash\mathsf l((k,a))\neq\mathsf l((l,b))}{(k,a)\in\Sigma,\ (l,b)\in\Sigma,\ k\neq l\vdash\mathsf l((k,a))\neq\mathsf l((l,b))}[B48TaggedLeafNe]
 \proofstepwidestar[1]{(k,a)\in\Sigma}{\rA}
 \proofstepwidestar[2]{(l,b)\in\Sigma}{\rA}
 \proofstepwidestar[3]{k\neq l}{\rA}
 \proofstepwidestar[4]{\mathsf l((k,a))=\mathsf l((l,b))}{\rA}
 \proofstepwidestar[1,2,4]{(k,a)=(l,b)}{\FormulaRefAuto{B48LeafInjective}[thm]{1,2,4}}
 \proofstepwidestar[1,2,4]{k=l\land a=b}{\FormulaRefAuto{(a,b)=(c,d)\eqvdash a=c\land b=d}[thm]{5}}
 \proofstepwidestar[1,2,4]{k=l}{\rAEa{6}}
 \proofstepwidestar[1,2,3,4]{\bot}{\rBI{7,3}}
 \proofstepwidestar[1,2,3]{\mathsf l((k,a))\neq\mathsf l((l,b))}{\rCI{4,8}}
\closeproofpartwideindR

\proofpartwideindR[Blatt und Knoten sind verschieden]{c\in\Sigma,\ p\in\mathcal T,\ q\in\mathcal T\vdash\mathsf l(c)\neq\mathsf n(p,q)}{c\in\Sigma,\ p\in\mathcal T,\ q\in\mathcal T\vdash\mathsf l(c)\neq\mathsf n(p,q)}[B48LeafNotNode]
 \proofstepwidestar[1]{c\in\Sigma}{\rA}
 \proofstepwidestar[2]{p\in\mathcal T}{\rA}
 \proofstepwidestar[3]{q\in\mathcal T}{\rA}
 \proofstepwidestar[2]{p\in\NonemptyWords{\TreeAlphabet{\Sigma}}}{\FormulaRefAuto{B48TreeWord}[thm]{2}}
 \proofstepwidestar[3]{q\in\NonemptyWords{\TreeAlphabet{\Sigma}}}{\FormulaRefAuto{B48TreeWord}[thm]{3}}
 \proofstepwidestar[1,2,3]{\mathsf l(c)\neq\mathsf n(p,q)}{\FormulaRefAuto{TreeConstructorsDisjoint}[thm]{1,4,5}}
\closeproofpartwideindR

\proofpartwideindR[Trennung der Konjunktionsmarke]{\mathsf c_0\neq\mathsf n_0}{\mathsf c_0\neq\mathsf n_0}[B48AndNeNeg]
 \proofstepwidestar[]{0\in\mathbb N}{\FormulaRefAuto{PeanoZero}[ax]}
 \proofstepwidestar[]{t_{\rm c}\in\mathbb N}{\FormulaRefAuto{B48TagType4}[thm]}
 \proofstepwidestar[]{t_{\rm n}\in\mathbb N}{\FormulaRefAuto{B48TagType3}[thm]}
 \proofstepwidestar[]{(t_{\rm c},(0,0))\in\Sigma}{\FormulaRefAuto{B48LabelType}[thm]{2,1,1}}
 \proofstepwidestar[]{(t_{\rm n},(0,0))\in\Sigma}{\FormulaRefAuto{B48LabelType}[thm]{3,1,1}}
 \proofstepwidestar[]{t_{\rm n}\neq t_{\rm c}}{\FormulaRefAuto{B48TagNe23}[thm]}
 \proofstepwidestar[]{t_{\rm c}\neq t_{\rm n}}{\FormulaRefAuto{a\neq b\vdash b\neq a}[thm]{6}}
 \proofstepwidestar[]{\mathsf c_0\neq\mathsf n_0}{\FormulaRefAuto{B48TaggedLeafNe}[thm]{4,5,7}}
\closeproofpartwideindR

\proofpartwideindR[Trennung der Quantorenmarke]{i\in\mathbb N\vdash \mathsf q_i\neq\mathsf n_0}{i\in\mathbb N\vdash \mathsf q_i\neq\mathsf n_0}[B48QNeNeg]
 \proofstepwidestar[1]{i\in\mathbb N}{\rA}
 \proofstepwidestar[]{0\in\mathbb N}{\FormulaRefAuto{PeanoZero}[ax]}
 \proofstepwidestar[]{t_{\rm q}\in\mathbb N}{\FormulaRefAuto{B48TagType5}[thm]}
 \proofstepwidestar[]{t_{\rm n}\in\mathbb N}{\FormulaRefAuto{B48TagType3}[thm]}
 \proofstepwidestar[1]{(t_{\rm q},(i,0))\in\Sigma}{\FormulaRefAuto{B48LabelType}[thm]{3,1,2}}
 \proofstepwidestar[]{(t_{\rm n},(0,0))\in\Sigma}{\FormulaRefAuto{B48LabelType}[thm]{4,2,2}}
 \proofstepwidestar[]{t_{\rm n}\neq t_{\rm q}}{\FormulaRefAuto{B48TagNe24}[thm]}
 \proofstepwidestar[]{t_{\rm q}\neq t_{\rm n}}{\FormulaRefAuto{a\neq b\vdash b\neq a}[thm]{7}}
 \proofstepwidestar[1]{\mathsf q_i\neq\mathsf n_0}{\FormulaRefAuto{B48TaggedLeafNe}[thm]{5,6,8}}
\closeproofpartwideindR

\proofpartwideindR[Trennung der Quantorenmarke]{i\in\mathbb N\vdash \mathsf q_i\neq\mathsf c_0}{i\in\mathbb N\vdash \mathsf q_i\neq\mathsf c_0}[B48QNeAnd]
 \proofstepwidestar[1]{i\in\mathbb N}{\rA}
 \proofstepwidestar[]{0\in\mathbb N}{\FormulaRefAuto{PeanoZero}[ax]}
 \proofstepwidestar[]{t_{\rm q}\in\mathbb N}{\FormulaRefAuto{B48TagType5}[thm]}
 \proofstepwidestar[]{t_{\rm c}\in\mathbb N}{\FormulaRefAuto{B48TagType4}[thm]}
 \proofstepwidestar[1]{(t_{\rm q},(i,0))\in\Sigma}{\FormulaRefAuto{B48LabelType}[thm]{3,1,2}}
 \proofstepwidestar[]{(t_{\rm c},(0,0))\in\Sigma}{\FormulaRefAuto{B48LabelType}[thm]{4,2,2}}
 \proofstepwidestar[]{t_{\rm c}\neq t_{\rm q}}{\FormulaRefAuto{B48TagNe34}[thm]}
 \proofstepwidestar[]{t_{\rm q}\neq t_{\rm c}}{\FormulaRefAuto{a\neq b\vdash b\neq a}[thm]{7}}
 \proofstepwidestar[1]{\mathsf q_i\neq\mathsf c_0}{\FormulaRefAuto{B48TaggedLeafNe}[thm]{5,6,8}}
\closeproofpartwideindR

\proofpartwideindR[Quantorenmarken bestimmen den Variablenindex]{i\in\mathbb N,\ j\in\mathbb N,\ \mathsf q_i=\mathsf q_j\vdash i=j}{i\in\mathbb N,\ j\in\mathbb N,\ \mathsf q_i=\mathsf q_j\vdash i=j}[B48QuantifierRootInjective]
 \proofstepwidestar[1]{i\in\mathbb N}{\rA}
 \proofstepwidestar[2]{j\in\mathbb N}{\rA}
 \proofstepwidestar[3]{\mathsf q_i=\mathsf q_j}{\rA}
 \proofstepwidestar[]{t_{\rm q}\in\mathbb N}{\FormulaRefAuto{B48TagType5}[thm]}
 \proofstepwidestar[]{0\in\mathbb N}{\FormulaRefAuto{PeanoZero}[ax]}
 \proofstepwidestar[1]{(t_{\rm q},(i,0))\in\Sigma}{\FormulaRefAuto{B48LabelType}[thm]{4,1,5}}
 \proofstepwidestar[2]{(t_{\rm q},(j,0))\in\Sigma}{\FormulaRefAuto{B48LabelType}[thm]{4,2,5}}
 \proofstepwidestar[1,2,3]{(t_{\rm q},(i,0))=(t_{\rm q},(j,0))}{\FormulaRefAuto{B48LeafInjective}[thm]{6,7,3}}
 \proofstepwidestar[1,2,3]{i=j\land0=0}{\FormulaRefAuto{TaggedNestedPairInjective}[thm]{8}}
 \proofstepwidestar[1,2,3]{i=j}{\rAEa{9}}
\closeproofpartwideindR

\proofpartwideindR[Ein Blatt ist kein Konjunktions-Zwischencode]{c\in\Sigma\vdash\neg J(\mathsf l(c))}{c\in\Sigma\vdash\neg J(\mathsf l(c))}[B48LeafNotAndSlot]
 \proofstepwidestar[1]{c\in\Sigma}{\rA}
 \proofstepwidestar[2]{J(\mathsf l(c))}{\rA}
 \proofstepwidestar[2]{\exists p\in\mathcal T\;(\mathsf l(c)=\mathsf n(\mathsf c_0,p))}{\FormulaRefAuto{B48EvaluationOperators}[def]{2}}
 \proofstepwidestar[4]{p\in\mathcal T\land\mathsf l(c)=\mathsf n(\mathsf c_0,p)}{\rA}
 \proofstepwidestar[4]{p\in\mathcal T}{\rAEa{4}}
 \proofstepwidestar[]{\mathsf c_0\in\mathcal T}{\FormulaRefAuto{B48AndRootType}[thm]}
 \proofstepwidestar[1,4]{\mathsf l(c)\neq\mathsf n(\mathsf c_0,p)}{\FormulaRefAuto{B48LeafNotNode}[thm]{1,6,5}}
 \proofstepwidestar[4]{\mathsf l(c)=\mathsf n(\mathsf c_0,p)}{\rAEb{4}}
 \proofstepwidestar[1,4]{\bot}{\rBI{8,7}}
 \proofstepwidestar[1,2]{\bot}{\rEE{3,4:9}}
 \proofstepwidestar[1]{\neg J(\mathsf l(c))}{\rCI{2,10}}
\closeproofpartwideindR

\proofpartwideindR[Konjunktions-Zwischencodes werden erkannt]{p\in\mathcal T\vdash J(\mathsf n(\mathsf c_0,p))}{p\in\mathcal T\vdash J(\mathsf n(\mathsf c_0,p))}[B48AndSlotSelector]
 \proofstepwidestar[1]{p\in\mathcal T}{\rA}
 \proofstepwidestar[]{\mathsf n(\mathsf c_0,p)=\mathsf n(\mathsf c_0,p)}{\rII}
 \proofstepwidestar[1]{\exists q\in\mathcal T\;(\mathsf n(\mathsf c_0,p)=\mathsf n(\mathsf c_0,q))}{\rEI{\rAI{1,2}}}
 \proofstepwidestar[1]{J(\mathsf n(\mathsf c_0,p))}{\FormulaRefAuto{B48EvaluationOperators}[def]{3}}
\closeproofpartwideindR

\proofpartwideindR[Ein Konjunktions-Zwischencode ist kein Blatt]{p\in\mathcal T\vdash\mathsf n(\mathsf c_0,p)\neq\mathsf n_0}{p\in\mathcal T\vdash\mathsf n(\mathsf c_0,p)\neq\mathsf n_0}[B48AndSlotNeNeg]
 \proofstepwidestar[1]{p\in\mathcal T}{\rA}
 \proofstepwidestar[]{0\in\mathbb N}{\FormulaRefAuto{PeanoZero}[ax]}
 \proofstepwidestar[]{t_{\rm n}\in\mathbb N}{\FormulaRefAuto{B48TagType3}[thm]}
 \proofstepwidestar[]{(t_{\rm n},(0,0))\in\Sigma}{\FormulaRefAuto{B48LabelType}[thm]{3,2,2}}
 \proofstepwidestar[]{\mathsf c_0\in\mathcal T}{\FormulaRefAuto{B48AndRootType}[thm]}
 \proofstepwidestar[1]{\mathsf n_0\neq\mathsf n(\mathsf c_0,p)}{\FormulaRefAuto{B48LeafNotNode}[thm]{4,5,1}}
 \proofstepwidestar[1]{\mathsf n(\mathsf c_0,p)\neq\mathsf n_0}{\FormulaRefAuto{a\neq b\vdash b\neq a}[thm]{6}}
\closeproofpartwideindR

\proofpartwideindR[Ein Konjunktions-Zwischencode ist kein Blatt]{p\in\mathcal T\vdash\mathsf n(\mathsf c_0,p)\neq\mathsf c_0}{p\in\mathcal T\vdash\mathsf n(\mathsf c_0,p)\neq\mathsf c_0}[B48AndSlotNeAnd]
 \proofstepwidestar[1]{p\in\mathcal T}{\rA}
 \proofstepwidestar[]{0\in\mathbb N}{\FormulaRefAuto{PeanoZero}[ax]}
 \proofstepwidestar[]{t_{\rm c}\in\mathbb N}{\FormulaRefAuto{B48TagType4}[thm]}
 \proofstepwidestar[]{(t_{\rm c},(0,0))\in\Sigma}{\FormulaRefAuto{B48LabelType}[thm]{3,2,2}}
 \proofstepwidestar[]{\mathsf c_0\in\mathcal T}{\FormulaRefAuto{B48AndRootType}[thm]}
 \proofstepwidestar[1]{\mathsf c_0\neq\mathsf n(\mathsf c_0,p)}{\FormulaRefAuto{B48LeafNotNode}[thm]{4,5,1}}
 \proofstepwidestar[1]{\mathsf n(\mathsf c_0,p)\neq\mathsf c_0}{\FormulaRefAuto{a\neq b\vdash b\neq a}[thm]{6}}
\closeproofpartwideindR

\proofpartwideindR[Einzeln referenzierter Auswertungszweig]{t=\mathsf n_0\vdash H_M(t,U,V)=A_D\setminus V}{t=\mathsf n_0\vdash H_M(t,U,V)=A_D\setminus V}[B48HNegSwitch]
 \proofstepwidestar[1]{t=\mathsf n_0}{\rA}
 \proofstepwidestar[]{H_M(t,U,V)=\BXLVIIIIf{t=\mathsf n_0}{A_D\setminus V}{H_{1,M}(t,U,V)}}{\FormulaRefAuto{B48EvaluationOperators}[def]}
 \proofstepwidestar[1]{\BXLVIIIIf{t=\mathsf n_0}{A_D\setminus V}{H_{1,M}(t,U,V)}=A_D\setminus V}{\FormulaRefAuto{B48IfTrue}[thm]{1}}
 \proofstepwidestar[1]{H_M(t,U,V)=A_D\setminus V}{\rIE{3,2}}
\closeproofpartwideindR

\proofpartwideindR[Einzeln referenzierter Auswertungszweig]{\neg(t=\mathsf n_0)\vdash H_M(t,U,V)=H_{1,M}(t,U,V)}{\neg(t=\mathsf n_0)\vdash H_M(t,U,V)=H_{1,M}(t,U,V)}[B48HSkipNeg]
 \proofstepwidestar[1]{\neg(t=\mathsf n_0)}{\rA}
 \proofstepwidestar[]{H_M(t,U,V)=\BXLVIIIIf{t=\mathsf n_0}{A_D\setminus V}{H_{1,M}(t,U,V)}}{\FormulaRefAuto{B48EvaluationOperators}[def]}
 \proofstepwidestar[1]{\BXLVIIIIf{t=\mathsf n_0}{A_D\setminus V}{H_{1,M}(t,U,V)}=H_{1,M}(t,U,V)}{\FormulaRefAuto{B48IfFalse}[thm]{1}}
 \proofstepwidestar[1]{H_M(t,U,V)=H_{1,M}(t,U,V)}{\rIE{3,2}}
\closeproofpartwideindR

\proofpartwideindR[Einzeln referenzierter Auswertungszweig]{t=\mathsf c_0\vdash H_{1,M}(t,U,V)=V}{t=\mathsf c_0\vdash H_{1,M}(t,U,V)=V}[B48HAndSwitch]
 \proofstepwidestar[1]{t=\mathsf c_0}{\rA}
 \proofstepwidestar[]{H_{1,M}(t,U,V)=\BXLVIIIIf{t=\mathsf c_0}{V}{H_{2,M}(t,U,V)}}{\FormulaRefAuto{B48EvaluationOperators}[def]}
 \proofstepwidestar[1]{\BXLVIIIIf{t=\mathsf c_0}{V}{H_{2,M}(t,U,V)}=V}{\FormulaRefAuto{B48IfTrue}[thm]{1}}
 \proofstepwidestar[1]{H_{1,M}(t,U,V)=V}{\rIE{3,2}}
\closeproofpartwideindR

\proofpartwideindR[Einzeln referenzierter Auswertungszweig]{\neg(t=\mathsf c_0)\vdash H_{1,M}(t,U,V)=H_{2,M}(t,U,V)}{\neg(t=\mathsf c_0)\vdash H_{1,M}(t,U,V)=H_{2,M}(t,U,V)}[B48HSkipAnd]
 \proofstepwidestar[1]{\neg(t=\mathsf c_0)}{\rA}
 \proofstepwidestar[]{H_{1,M}(t,U,V)=\BXLVIIIIf{t=\mathsf c_0}{V}{H_{2,M}(t,U,V)}}{\FormulaRefAuto{B48EvaluationOperators}[def]}
 \proofstepwidestar[1]{\BXLVIIIIf{t=\mathsf c_0}{V}{H_{2,M}(t,U,V)}=H_{2,M}(t,U,V)}{\FormulaRefAuto{B48IfFalse}[thm]{1}}
 \proofstepwidestar[1]{H_{1,M}(t,U,V)=H_{2,M}(t,U,V)}{\rIE{3,2}}
\closeproofpartwideindR

\proofpartwideindR[Einzeln referenzierter Auswertungszweig]{J(t)\vdash H_{2,M}(t,U,V)=U\cap V}{J(t)\vdash H_{2,M}(t,U,V)=U\cap V}[B48HSlotSwitch]
 \proofstepwidestar[1]{J(t)}{\rA}
 \proofstepwidestar[]{H_{2,M}(t,U,V)=\BXLVIIIIf{J(t)}{U\cap V}{Q(t,V)}}{\FormulaRefAuto{B48EvaluationOperators}[def]}
 \proofstepwidestar[1]{\BXLVIIIIf{J(t)}{U\cap V}{Q(t,V)}=U\cap V}{\FormulaRefAuto{B48IfTrue}[thm]{1}}
 \proofstepwidestar[1]{H_{2,M}(t,U,V)=U\cap V}{\rIE{3,2}}
\closeproofpartwideindR

\proofpartwideindR[Einzeln referenzierter Auswertungszweig]{\neg(J(t))\vdash H_{2,M}(t,U,V)=Q(t,V)}{\neg(J(t))\vdash H_{2,M}(t,U,V)=Q(t,V)}[B48HSkipSlot]
 \proofstepwidestar[1]{\neg(J(t))}{\rA}
 \proofstepwidestar[]{H_{2,M}(t,U,V)=\BXLVIIIIf{J(t)}{U\cap V}{Q(t,V)}}{\FormulaRefAuto{B48EvaluationOperators}[def]}
 \proofstepwidestar[1]{\BXLVIIIIf{J(t)}{U\cap V}{Q(t,V)}=Q(t,V)}{\FormulaRefAuto{B48IfFalse}[thm]{1}}
 \proofstepwidestar[1]{H_{2,M}(t,U,V)=Q(t,V)}{\rIE{3,2}}
\closeproofpartwideindR

\proofpartwideindR[Auswertung der Negationsmarke]{H_M(\mathsf n_0,U,V)=A_D\setminus V}{H_M(\mathsf n_0,U,V)=A_D\setminus V}[B48HNeg]
 \proofstepwidestar[]{\mathsf n_0=\mathsf n_0}{\rII}
 \proofstepwidestar[]{H_M(\mathsf n_0,U,V)=A_D\setminus V}{\FormulaRefAuto{B48HNegSwitch}[thm]{1}}
\closeproofpartwideindR

\proofpartwideindR[Auswertung des ersten Konjunktionsarguments]{H_M(\mathsf c_0,U,V)=V}{H_M(\mathsf c_0,U,V)=V}[B48HAndPrefix]
 \proofstepwidestar[]{\mathsf c_0\neq\mathsf n_0}{\FormulaRefAuto{B48AndNeNeg}[thm]}
 \proofstepwidestar[]{H_M(\mathsf c_0,U,V)=H_{1,M}(\mathsf c_0,U,V)}{\FormulaRefAuto{B48HSkipNeg}[thm]{1}}
 \proofstepwidestar[]{\mathsf c_0=\mathsf c_0}{\rII}
 \proofstepwidestar[]{H_{1,M}(\mathsf c_0,U,V)=V}{\FormulaRefAuto{B48HAndSwitch}[thm]{3}}
 \proofstepwidestar[]{H_M(\mathsf c_0,U,V)=V}{\rIE{4,2}}
\closeproofpartwideindR

\proofpartwideindR[Auswertung des zweiten Konjunktionsarguments]{p\in\mathcal T\vdash H_M(\mathsf n(\mathsf c_0,p),U,V)=U\cap V}{p\in\mathcal T\vdash H_M(\mathsf n(\mathsf c_0,p),U,V)=U\cap V}[B48HAnd]
 \proofstepwidestar[1]{p\in\mathcal T}{\rA}
 \proofstepwidestar[1]{\mathsf n(\mathsf c_0,p)\neq\mathsf n_0}{\FormulaRefAuto{B48AndSlotNeNeg}[thm]{1}}
 \proofstepwidestar[1]{\mathsf n(\mathsf c_0,p)\neq\mathsf c_0}{\FormulaRefAuto{B48AndSlotNeAnd}[thm]{1}}
 \proofstepwidestar[1]{J(\mathsf n(\mathsf c_0,p))}{\FormulaRefAuto{B48AndSlotSelector}[thm]{1}}
 \proofstepwidestar[1]{H_M(\mathsf n(\mathsf c_0,p),U,V)=H_{1,M}(\mathsf n(\mathsf c_0,p),U,V)}{\FormulaRefAuto{B48HSkipNeg}[thm]{2}}
 \proofstepwidestar[1]{H_{1,M}(\mathsf n(\mathsf c_0,p),U,V)=H_{2,M}(\mathsf n(\mathsf c_0,p),U,V)}{\FormulaRefAuto{B48HSkipAnd}[thm]{3}}
 \proofstepwidestar[1]{H_{2,M}(\mathsf n(\mathsf c_0,p),U,V)=U\cap V}{\FormulaRefAuto{B48HSlotSwitch}[thm]{4}}
 \proofstepwidestar[1]{H_M(\mathsf n(\mathsf c_0,p),U,V)=U\cap V}{\rIE{7,\rIE{6,5}}}
\closeproofpartwideindR

\proofpartwideindR[Elementkriterium bei bekanntem Traeger]{s\in A\vdash\bigl(s\in\{x\in A\mid P(x)\}\leftrightarrow P(s)\bigr)}{s\in A\vdash\bigl(s\in\{x\in A\mid P(x)\}\leftrightarrow P(s)\bigr)}[B48SeparationWithin]
 \proofstepwidestar[1]{s\in A}{\rA}
 \proofstepwidestar[2]{s\in\{x\in A\mid P(x)\}}{\rA}
 \proofstepwidestar[2]{s\in A\land P(s)}{\FormulaRefAuto{x\in\{u\in A\mid P(u)\}\eqvdash x\in A\land P(x)}[thm]{2}}
 \proofstepwidestar[2]{P(s)}{\rAEb{3}}
 \proofstepwidestar[]{s\in\{x\in A\mid P(x)\}\to P(s)}{\rRI{2,4}}
 \proofstepwidestar[6]{P(s)}{\rA}
 \proofstepwidestar[1,6]{s\in A\land P(s)}{\rAI{1,6}}
 \proofstepwidestar[1,6]{s\in\{x\in A\mid P(x)\}}{\FormulaRefAuto{x\in\{u\in A\mid P(u)\}\eqvdash x\in A\land P(x)}[thm]{7}}
 \proofstepwidestar[1]{P(s)\to s\in\{x\in A\mid P(x)\}}{\rRI{6,8}}
 \proofstepwidestar[1]{s\in\{x\in A\mid P(x)\}\leftrightarrow P(s)}{\rLRI{5,9}}
\closeproofpartwideindR

\proofpartwideindR[Eine Quantorenmarke waehlt genau ihren Index]{i\in\mathbb N\vdash Q(\mathsf q_i,U)=Q_i(U)}{i\in\mathbb N\vdash Q(\mathsf q_i,U)=Q_i(U)}[B48QValue]
 \proofstepwidestar[1]{i\in\mathbb N}{\rA}
 \proofstepwidestar[2]{s\in Q(\mathsf q_i,U)}{\rA}
 \proofstepwidestar[2]{s\in A_D\land\exists j\in\mathbb N\;(\mathsf q_i=\mathsf q_j\land\exists a\in D\;(s[j\mapsto a]\in U))}{\FormulaRefAuto{B48SemanticTerms}[def]{2}}
 \proofstepwidestar[2]{s\in A_D}{\rAEa{3}}
 \proofstepwidestar[2]{\exists j\in\mathbb N\;(\mathsf q_i=\mathsf q_j\land\exists a\in D\;(s[j\mapsto a]\in U))}{\rAEb{3}}
 \proofstepwidestar[6]{j\in\mathbb N\land(\mathsf q_i=\mathsf q_j\land\exists a\in D\;(s[j\mapsto a]\in U))}{\rA}
 \proofstepwidestar[6]{j\in\mathbb N}{\rAEa{6}}
 \proofstepwidestar[6]{\mathsf q_i=\mathsf q_j}{\rAEa{\rAEb{6}}}
 \proofstepwidestar[1,6]{i=j}{\FormulaRefAuto{B48QuantifierRootInjective}[thm]{1,7,8}}
 \proofstepwidestar[1,6]{j=i}{\FormulaRefAuto{a=b\vdash b=a}[thm]{9}}
 \proofstepwidestar[6]{\exists a\in D\;(s[j\mapsto a]\in U)}{\rAEb{\rAEb{6}}}
 \proofstepwidestar[1,6]{\exists a\in D\;(s[i\mapsto a]\in U)}{\rIE{10,11}}
 \proofstepwidestar[1,2]{\exists a\in D\;(s[i\mapsto a]\in U)}{\rEE{5,6:12}}
 \proofstepwidestar[1,2]{s\in Q_i(U)}{\FormulaRefAuto{B48SemanticTerms}[def]{\rAI{4,13}}}
 \proofstepwidestar[1]{s\in Q(\mathsf q_i,U)\to s\in Q_i(U)}{\rRI{2,14}}
 \proofstepwidestar[16]{s\in Q_i(U)}{\rA}
 \proofstepwidestar[16]{s\in A_D\land\exists a\in D\;(s[i\mapsto a]\in U)}{\FormulaRefAuto{B48SemanticTerms}[def]{16}}
 \proofstepwidestar[]{\mathsf q_i=\mathsf q_i}{\rII}
 \proofstepwidestar[1,16]{\exists j\in\mathbb N\;(\mathsf q_i=\mathsf q_j\land\exists a\in D\;(s[j\mapsto a]\in U))}{\rEI{\rAI{1,\rAI{18,\rAEb{17}}}}}
 \proofstepwidestar[1,16]{s\in Q(\mathsf q_i,U)}{\FormulaRefAuto{B48SemanticTerms}[def]{\rAI{\rAEa{17},19}}}
 \proofstepwidestar[1]{s\in Q_i(U)\to s\in Q(\mathsf q_i,U)}{\rRI{16,20}}
 \proofstepwidestar[1]{s\in Q(\mathsf q_i,U)\leftrightarrow s\in Q_i(U)}{\rLRI{15,21}}
 \proofstepwidestar[1]{Q(\mathsf q_i,U)=Q_i(U)}{\FormulaRefAuto{\forall x(x\in A\leftrightarrow x\in B)\eqvdash A=B}[thm]{\rUI{22}}}
\closeproofpartwideindR

\proofpartwideindR[Auswertung einer Quantorenmarke]{i\in\mathbb N\vdash H_M(\mathsf q_i,U,V)=Q_i(V)}{i\in\mathbb N\vdash H_M(\mathsf q_i,U,V)=Q_i(V)}[B48HExists]
 \proofstepwidestar[1]{i\in\mathbb N}{\rA}
 \proofstepwidestar[1]{\mathsf q_i\neq\mathsf n_0}{\FormulaRefAuto{B48QNeNeg}[thm]{1}}
 \proofstepwidestar[1]{\mathsf q_i\neq\mathsf c_0}{\FormulaRefAuto{B48QNeAnd}[thm]{1}}
 \proofstepwidestar[]{t_{\rm q}\in\mathbb N}{\FormulaRefAuto{B48TagType5}[thm]}
 \proofstepwidestar[]{0\in\mathbb N}{\FormulaRefAuto{PeanoZero}[ax]}
 \proofstepwidestar[1]{(t_{\rm q},(i,0))\in\Sigma}{\FormulaRefAuto{B48LabelType}[thm]{4,1,5}}
 \proofstepwidestar[1]{\neg J(\mathsf q_i)}{\FormulaRefAuto{B48LeafNotAndSlot}[thm]{6}}
 \proofstepwidestar[1]{H_M(\mathsf q_i,U,V)=H_{1,M}(\mathsf q_i,U,V)}{\FormulaRefAuto{B48HSkipNeg}[thm]{2}}
 \proofstepwidestar[1]{H_{1,M}(\mathsf q_i,U,V)=H_{2,M}(\mathsf q_i,U,V)}{\FormulaRefAuto{B48HSkipAnd}[thm]{3}}
 \proofstepwidestar[1]{H_{2,M}(\mathsf q_i,U,V)=Q(\mathsf q_i,V)}{\FormulaRefAuto{B48HSkipSlot}[thm]{7}}
 \proofstepwidestar[1]{Q(\mathsf q_i,V)=Q_i(V)}{\FormulaRefAuto{B48QValue}[thm]{1}}
 \proofstepwidestar[1]{H_M(\mathsf q_i,U,V)=Q_i(V)}{\rIE{11,\rIE{10,\rIE{9,8}}}}
\closeproofpartwideindR

\par\medskip\noindent\textbf{D. Mengenbildung und Strukturrekursion}\par
\proofpartwideindR[Die leere Wertemenge]{\varnothing\in\powerset(A)}{\varnothing\in\powerset(A)}[B48EmptyPower]
 \proofstepwidestar[]{\varnothing\subseteq A}{\FormulaRefAuto{\varnothing\subseteq A}[thm]}
 \proofstepwidestar[]{\varnothing\in\powerset(A)}{\FormulaRefAuto{x\in\powerset(A)\eqvdash x\subseteq A}[thm]{1}}
\closeproofpartwideindR

\proofpartwideindR[Schnitt zweier Wertemengen]{U\in\powerset(A),\ V\in\powerset(A)\vdash U\cap V\in\powerset(A)}{U\in\powerset(A),\ V\in\powerset(A)\vdash U\cap V\in\powerset(A)}[B48IntersectionPower]
 \proofstepwidestar[1]{U\in\powerset(A)}{\rA}
 \proofstepwidestar[2]{V\in\powerset(A)}{\rA}
 \proofstepwidestar[1]{U\subseteq A}{\FormulaRefAuto{x\in\powerset(A)\eqvdash x\subseteq A}[thm]{1}}
 \proofstepwidestar[]{U\cap V\subseteq U}{\FormulaRefAuto{A\cap B\subseteq A}[thm]}
 \proofstepwidestar[1]{U\cap V\subseteq A}{\FormulaRefAuto{A\subseteq B,\ B\subseteq C\vdash A\subseteq C}[thm]{4,3}}
 \proofstepwidestar[1]{U\cap V\in\powerset(A)}{\FormulaRefAuto{x\in\powerset(A)\eqvdash x\subseteq A}[thm]{5}}
\closeproofpartwideindR

\proofpartwideindR[Eindeutig zerlegbare Alphabetdaten]{c\in\Sigma\vdash\exists k\in\mathbb N\,\exists i\in\mathbb N\,\exists j\in\mathbb N\;(c=(k,(i,j)))}{c\in\Sigma\vdash\exists k\in\mathbb N\,\exists i\in\mathbb N\,\exists j\in\mathbb N\;(c=(k,(i,j)))}[B48LabelDecomposition]
 \proofstepwidestar[1]{c\in\Sigma}{\rA}
 \proofstepwidestar[1]{\exists k\in\mathbb N\,\exists u\in\mathbb N\times\mathbb N\;(c=(k,u))}{\FormulaRefAuto{x\in A\times B\eqvdash\exists a\in A\,\exists b\in B\,x=(a,b)}[thm]{1}}
 \proofstepwidestar[3]{k\in\mathbb N\land(u\in\mathbb N\times\mathbb N\land c=(k,u))}{\rA}
 \proofstepwidestar[3]{u\in\mathbb N\times\mathbb N}{\rAEa{\rAEb{3}}}
 \proofstepwidestar[3]{\exists i\in\mathbb N\,\exists j\in\mathbb N\;(u=(i,j))}{\FormulaRefAuto{x\in A\times B\eqvdash\exists a\in A\,\exists b\in B\,x=(a,b)}[thm]{4}}
 \proofstepwidestar[6]{i\in\mathbb N\land(j\in\mathbb N\land u=(i,j))}{\rA}
 \proofstepwidestar[3,6]{c=(k,(i,j))}{\rIE{\rAEb{\rAEb{6}},\rAEb{\rAEb{3}}}}
 \proofstepwidestar[3,6]{\exists a\in\mathbb N\,\exists b\in\mathbb N\,\exists d\in\mathbb N\;(c=(a,(b,d)))}{\rEI{\rAI{\rAEa{3},\rEI{\rAI{\rAEa{6},\rEI{\rAI{\rAEa{\rAEb{6}},7}}}}}}}
 \proofstepwidestar[3]{\exists a\in\mathbb N\,\exists b\in\mathbb N\,\exists d\in\mathbb N\;(c=(a,(b,d)))}{\rEEStar{5,6:8}}
 \proofstepwidestar[1]{\exists a\in\mathbb N\,\exists b\in\mathbb N\,\exists d\in\mathbb N\;(c=(a,(b,d)))}{\rEEStar{2,3:9}}
\closeproofpartwideindR

\proofpartwideindR[Atomare Wertemengen sind Mengen von Belegungen]{c\in\Sigma\vdash B_M(c)\in\powerset(A_D)}{c\in\Sigma\vdash B_M(c)\in\powerset(A_D)}[B48BType]
 \proofstepwidestar[1]{c\in\Sigma}{\rA}
 \proofstepwidestar[1]{\exists k\in\mathbb N\,\exists i\in\mathbb N\,\exists j\in\mathbb N\;(c=(k,(i,j)))}{\FormulaRefAuto{B48LabelDecomposition}[thm]{1}}
 \proofstepwidestar[3]{k\in\mathbb N\land(i\in\mathbb N\land(j\in\mathbb N\land c=(k,(i,j))))}{\rA}
 \proofstepwidestar[]{\{s\in A_D\mid s(i)=s(j)\}\in\powerset(A_D)}{\FormulaRefAuto{B48SeparationPower}[thm]}
 \proofstepwidestar[]{\{s\in A_D\mid(s(i),s(j))\in E\}\in\powerset(A_D)}{\FormulaRefAuto{B48SeparationPower}[thm]}
 \proofstepwidestar[]{\varnothing\in\powerset(A_D)}{\FormulaRefAuto{B48EmptyPower}[thm]}
 \proofstepwidestar[]{\BXLVIIIIf{k=t_{\rm m}}{\{s\in A_D\mid(s(i),s(j))\in E\}}{\varnothing}\in\powerset(A_D)}{\FormulaRefAuto{B48IfTyped}[thm]{5,6}}
 \proofstepwidestar[]{\BXLVIIIIf{k=t_{\rm e}}{\{s\in A_D\mid s(i)=s(j)\}}{\BXLVIIIIf{k=t_{\rm m}}{\{s\in A_D\mid(s(i),s(j))\in E\}}{\varnothing}}\in\powerset(A_D)}{\FormulaRefAuto{B48IfTyped}[thm]{4,7}}
 \proofstepwidestar[3]{B_M((k,(i,j)))\in\powerset(A_D)}{\FormulaRefAuto{B48SemanticTerms}[def]{8}}
 \proofstepwidestar[3]{c=(k,(i,j))}{\rAEb{\rAEb{\rAEb{3}}}}
 \proofstepwidestar[3]{(k,(i,j))=c}{\FormulaRefAuto{a=b\vdash b=a}[thm]{10}}
 \proofstepwidestar[3]{B_M(c)\in\powerset(A_D)}{\rIE{11,9}}
 \proofstepwidestar[1]{B_M(c)\in\powerset(A_D)}{\rEEStar{2,3:12}}
\closeproofpartwideindR

\proofpartwideindR[Der Knotenoperator bleibt im Wertemengenbereich]{U\in\powerset(A_D),\ V\in\powerset(A_D)\vdash H_M(t,U,V)\in\powerset(A_D)}{U\in\powerset(A_D),\ V\in\powerset(A_D)\vdash H_M(t,U,V)\in\powerset(A_D)}[B48HType]
 \proofstepwidestar[1]{U\in\powerset(A_D)}{\rA}
 \proofstepwidestar[2]{V\in\powerset(A_D)}{\rA}
 \proofstepwidestar[]{A_D\setminus V\in\powerset(A_D)}{\FormulaRefAuto{B48SeparationPower}[thm]}
 \proofstepwidestar[1,2]{U\cap V\in\powerset(A_D)}{\FormulaRefAuto{B48IntersectionPower}[thm]{1,2}}
 \proofstepwidestar[]{Q(t,V)\in\powerset(A_D)}{\FormulaRefAuto{B48SeparationPower}[thm]}
 \proofstepwidestar[1,2]{\BXLVIIIIf{J(t)}{U\cap V}{Q(t,V)}\in\powerset(A_D)}{\FormulaRefAuto{B48IfTyped}[thm]{4,5}}
 \proofstepwidestar[1,2]{H_{2,M}(t,U,V)\in\powerset(A_D)}{\FormulaRefAuto{B48EvaluationOperators}[def]{6}}
 \proofstepwidestar[1,2]{\BXLVIIIIf{t=\mathsf c_0}{V}{H_{2,M}(t,U,V)}\in\powerset(A_D)}{\FormulaRefAuto{B48IfTyped}[thm]{2,7}}
 \proofstepwidestar[1,2]{H_{1,M}(t,U,V)\in\powerset(A_D)}{\FormulaRefAuto{B48EvaluationOperators}[def]{8}}
 \proofstepwidestar[1,2]{\BXLVIIIIf{t=\mathsf n_0}{A_D\setminus V}{H_{1,M}(t,U,V)}\in\powerset(A_D)}{\FormulaRefAuto{B48IfTyped}[thm]{3,9}}
 \proofstepwidestar[1,2]{H_M(t,U,V)\in\powerset(A_D)}{\FormulaRefAuto{B48EvaluationOperators}[def]{10}}
\closeproofpartwideindR

\proofpartwideindR[Typisierung des Blattauswertungswerts]{c\in\Sigma\vdash(\mathsf l(c),B_M(c))\in Y_D}{c\in\Sigma\vdash(\mathsf l(c),B_M(c))\in Y_D}[B48LeafGraphTerm]
 \proofstepwidestar[1]{c\in\Sigma}{\rA}
 \proofstepwidestar[1]{\mathsf l(c)\in\mathcal T}{\FormulaRefAuto{B48LeafType}[thm]{1}}
 \proofstepwidestar[1]{B_M(c)\in\powerset(A_D)}{\FormulaRefAuto{B48BType}[thm]{1}}
 \proofstepwidestar[1]{(\mathsf l(c),B_M(c))\in Y_D}{\FormulaRefAuto{(a,b)\in A\times B\eqvdash a\in A\land b\in B}[thm]{\rAI{2,3}}}
\closeproofpartwideindR

\proofpartwideindR[Der Blattgraph ist eine Funktion]{f_M\colon\Sigma\to Y_D}{f_M\colon\Sigma\to Y_D}[B48LeafGraphType]
 \proofstepwidestar[1]{c\in\Sigma}{\rA}
 \proofstepwidestar[1]{(\mathsf l(c),B_M(c))\in Y_D}{\FormulaRefAuto{B48LeafGraphTerm}[thm]{1}}
 \proofstepwidestar[]{\forall c\in\Sigma\;((\mathsf l(c),B_M(c))\in Y_D)}{\rUI{\rRI{1,2}}}
 \proofstepwidestar[]{f_M\colon\Sigma\to Y_D}{\FormulaRefAuto{TypedTermGraphFunction}[thm]{3}}
\closeproofpartwideindR

\proofpartwideindR[Werte der Blattauswertungsfunktion]{c\in\Sigma\vdash f_M(c)=(\mathsf l(c),B_M(c))}{c\in\Sigma\vdash f_M(c)=(\mathsf l(c),B_M(c))}[B48LeafGraphValue]
 \proofstepwidestar[1]{c\in\Sigma}{\rA}
 \proofstepwidestar[]{f_M\colon\Sigma\to Y_D}{\FormulaRefAuto{B48LeafGraphType}[thm]}
 \proofstepwidestar[1]{(c,f_M(c))\in f_M}{\FormulaRefAuto{F\colon A\to B,\ x\in A\vdash(x,F(x))\in F}[thm]{2,1}}
 \proofstepwidestar[1]{(c,f_M(c))\in\Sigma\times Y_D\land f_M(c)=(\mathsf l(c),B_M(c))}{\FormulaRefAuto{B48EvaluationOperators}[def]{3}}
 \proofstepwidestar[1]{f_M(c)=(\mathsf l(c),B_M(c))}{\rAEb{4}}
\closeproofpartwideindR

\proofpartwideindR[Typisierung des Knotenauswertungswerts]{(p,U)\in Y_D,\ (q,V)\in Y_D\vdash(\mathsf n(p,q),H_M(p,U,V))\in Y_D}{(p,U)\in Y_D,\ (q,V)\in Y_D\vdash(\mathsf n(p,q),H_M(p,U,V))\in Y_D}[B48NodeGraphTerm]
 \proofstepwidestar[1]{(p,U)\in Y_D}{\rA}
 \proofstepwidestar[2]{(q,V)\in Y_D}{\rA}
 \proofstepwidestar[1]{p\in\mathcal T\land U\in\powerset(A_D)}{\FormulaRefAuto{(a,b)\in A\times B\eqvdash a\in A\land b\in B}[thm]{1}}
 \proofstepwidestar[2]{q\in\mathcal T\land V\in\powerset(A_D)}{\FormulaRefAuto{(a,b)\in A\times B\eqvdash a\in A\land b\in B}[thm]{2}}
 \proofstepwidestar[1,2]{\mathsf n(p,q)\in\mathcal T}{\FormulaRefAuto{B48NodeType}[thm]{\rAEa{3},\rAEa{4}}}
 \proofstepwidestar[1,2]{H_M(p,U,V)\in\powerset(A_D)}{\FormulaRefAuto{B48HType}[thm]{\rAEb{3},\rAEb{4}}}
 \proofstepwidestar[1,2]{(\mathsf n(p,q),H_M(p,U,V))\in Y_D}{\FormulaRefAuto{(a,b)\in A\times B\eqvdash a\in A\land b\in B}[thm]{\rAI{5,6}}}
\closeproofpartwideindR

\proofpartwideindR[Der Knotengraph ist eine Funktion]{g_M\colon Y_D\times Y_D\to Y_D}{g_M\colon Y_D\times Y_D\to Y_D}[B48NodeGraphType]
 \proofstepwidestar[1]{(p,U)\in Y_D}{\rA}
 \proofstepwidestar[2]{(q,V)\in Y_D}{\rA}
 \proofstepwidestar[1,2]{(\mathsf n(p,q),H_M(p,U,V))\in Y_D}{\FormulaRefAuto{B48NodeGraphTerm}[thm]{1,2}}
 \proofstepwidestar[]{\forall(p,U)\in Y_D\,\forall(q,V)\in Y_D\;((\mathsf n(p,q),H_M(p,U,V))\in Y_D)}{\rUI{\rRI{1,\rUI{\rRI{2,3}}}}}
 \proofstepwidestar[]{g_M\colon Y_D\times Y_D\to Y_D}{\FormulaRefAuto{TypedTermGraphFunction}[thm]{4}}
\closeproofpartwideindR

\proofpartwideindR[Werte der Knotenauswertungsfunktion]{(p,U)\in Y_D,\ (q,V)\in Y_D\vdash g_M((p,U),(q,V))=(\mathsf n(p,q),H_M(p,U,V))}{(p,U)\in Y_D,\ (q,V)\in Y_D\vdash g_M((p,U),(q,V))=(\mathsf n(p,q),H_M(p,U,V))}[B48NodeGraphValue]
 \proofstepwidestar[1]{(p,U)\in Y_D}{\rA}
 \proofstepwidestar[2]{(q,V)\in Y_D}{\rA}
 \proofstepwidestar[]{g_M\colon Y_D\times Y_D\to Y_D}{\FormulaRefAuto{B48NodeGraphType}[thm]}
 \proofstepwidestar[1,2]{((p,U),(q,V))\in Y_D\times Y_D}{\FormulaRefAuto{(a,b)\in A\times B\eqvdash a\in A\land b\in B}[thm]{\rAI{1,2}}}
 \proofstepwidestar[1,2]{(((p,U),(q,V)),g_M((p,U),(q,V)))\in g_M}{\FormulaRefAuto{F\colon A\to B,\ x\in A\vdash(x,F(x))\in F}[thm]{3,4}}
 \proofstepwidestar[1,2]{\begin{gathered}(((p,U),(q,V)),g_M((p,U),(q,V)))\\{}\in(Y_D\times Y_D)\times Y_D\\{}\land g_M((p,U),(q,V))=(\mathsf n(p,q),H_M(p,U,V))\end{gathered}}{\FormulaRefAuto{B48EvaluationOperators}[def]{5}}
 \proofstepwidestar[1,2]{g_M((p,U),(q,V))=(\mathsf n(p,q),H_M(p,U,V))}{\rAEb{6}}
\closeproofpartwideindR

\proofpartwideindR[Anwendung des bereits bewiesenen Rekursionssatzes]{\exists e\;\mathsf{Rec}(e,M)}{\exists e\;\mathsf{Rec}(e,M)}[B48TreeEvaluationExists]
 \proofstepwidestar[]{f_M\colon\Sigma\to Y_D}{\FormulaRefAuto{B48LeafGraphType}[thm]}
 \proofstepwidestar[]{g_M\colon Y_D\times Y_D\to Y_D}{\FormulaRefAuto{B48NodeGraphType}[thm]}
 \proofstepwidestar[]{\begin{gathered}\exists!e\colon\mathcal T\to Y_D\;\bigl((\forall c\in\Sigma\;e(\mathsf l(c))=f_M(c))\\{}\land(\forall p\in\mathcal T\,\forall q\in\mathcal T\;e(\mathsf n(p,q))=g_M(e(p),e(q)))\bigr)\end{gathered}}{\FormulaRefAuto{FullPlanarBinaryTreeRecursion}[thm]{1,2}}
 \proofstepwidestar[]{\begin{gathered}\exists e\colon\mathcal T\to Y_D\;\bigl((\forall c\in\Sigma\;e(\mathsf l(c))=f_M(c))\\{}\land(\forall p\in\mathcal T\,\forall q\in\mathcal T\;e(\mathsf n(p,q))=g_M(e(p),e(q)))\bigr)\end{gathered}}{\FormulaRefAuto{\exists! x\,P(x)\vdash \exists x\,P(x)}[thm]{3}}
 \proofstepwidestar[5]{\begin{gathered}e\colon\mathcal T\to Y_D\\{}\land(\forall c\in\Sigma\;e(\mathsf l(c))=f_M(c))\\{}\land(\forall p\in\mathcal T\,\forall q\in\mathcal T\;e(\mathsf n(p,q))=g_M(e(p),e(q)))\end{gathered}}{\rA}
 \proofstepwidestar[5]{e\colon\mathcal T\to Y_D}{\rAEa{5}}
 \proofstepwidestar[5]{\forall c\in\Sigma\;(e(\mathsf l(c))=f_M(c))}{\rAEa{\rAEb{5}}}
 \proofstepwidestar[5]{\forall p\in\mathcal T\,\forall q\in\mathcal T\;(e(\mathsf n(p,q))=g_M(e(p),e(q)))}{\rAEb{\rAEb{5}}}
 \proofstepwidestar[5]{\mathsf{Rec}(e,M)}{\FormulaRefAuto{B48RecursiveEvaluationDef}[def]{6,7,8}}
 \proofstepwidestar[5]{\exists e\;\mathsf{Rec}(e,M)}{\rEI{9}}
 \proofstepwidestar[]{\exists e\;\mathsf{Rec}(e,M)}{\rEE{4,5,10}}
\closeproofpartwideindR

\proofpartwideindR[Rekursionsklausel 1]{\mathsf{Rec}(e,M)\vdash e\colon\mathcal T\to Y_D}{\mathsf{Rec}(e,M)\vdash e\colon\mathcal T\to Y_D}[B48RecType]
 \proofstepwidestar[1]{\mathsf{Rec}(e,M)}{\rA}
 \proofstepwidestar[1]{e\colon\mathcal T\to Y_D}{\FormulaRefAuto{B48RecTypeAxiom}[ax]{1}}
\closeproofpartwideindR

\proofpartwideindR[Rekursionsklausel 2]{\mathsf{Rec}(e,M)\vdash \forall c\in\Sigma\;(e(\mathsf l(c))=f_M(c))}{\mathsf{Rec}(e,M)\vdash \forall c\in\Sigma\;(e(\mathsf l(c))=f_M(c))}[B48RecLeaf]
 \proofstepwidestar[1]{\mathsf{Rec}(e,M)}{\rA}
 \proofstepwidestar[1]{\forall c\in\Sigma\;(e(\mathsf l(c))=f_M(c))}{\FormulaRefAuto{B48RecLeafAxiom}[ax]{1}}
\closeproofpartwideindR

\proofpartwideindR[Rekursionsklausel 3]{\mathsf{Rec}(e,M)\vdash \forall p\in\mathcal T\,\forall q\in\mathcal T\;(e(\mathsf n(p,q))=g_M(e(p),e(q)))}{\mathsf{Rec}(e,M)\vdash \forall p\in\mathcal T\,\forall q\in\mathcal T\;(e(\mathsf n(p,q))=g_M(e(p),e(q)))}[B48RecNode]
 \proofstepwidestar[1]{\mathsf{Rec}(e,M)}{\rA}
 \proofstepwidestar[1]{\forall p\in\mathcal T\,\forall q\in\mathcal T\;(e(\mathsf n(p,q))=g_M(e(p),e(q)))}{\FormulaRefAuto{B48RecNodeAxiom}[ax]{1}}
\closeproofpartwideindR

\proofpartwideindR[Auslesen einer eindeutig bestimmten Wertemenge]{U\in\powerset(A),\ z=(p,U)\vdash U=V_A(z,p)}{U\in\powerset(A),\ z=(p,U)\vdash U=V_A(z,p)}[B48ReadValue]
 \proofstepwidestar[1]{U\in\powerset(A)}{\rA}
 \proofstepwidestar[2]{z=(p,U)}{\rA}
 \proofstepwidestar[1]{U\subseteq A}{\FormulaRefAuto{x\in\powerset(A)\eqvdash x\subseteq A}[thm]{1}}
 \proofstepwidestar[4]{s\in U}{\rA}
 \proofstepwidestar[1,4]{s\in A}{\FormulaRefAuto{A\subseteq B,\ x\in A\vdash x\in B}[thm]{3,4}}
 \proofstepwidestar[1,2,4]{\exists W\in\powerset(A)\;(z=(p,W)\land s\in W)}{\rEI{\rAI{1,\rAI{2,4}}}}
 \proofstepwidestar[1,2,4]{s\in V_A(z,p)}{\FormulaRefAuto{B48EvaluationCandidate}[def]{\rAI{5,6}}}
 \proofstepwidestar[1,2]{s\in U\to s\in V_A(z,p)}{\rRI{4,7}}
 \proofstepwidestar[9]{s\in V_A(z,p)}{\rA}
 \proofstepwidestar[9]{s\in A\land\exists W\in\powerset(A)\;(z=(p,W)\land s\in W)}{\FormulaRefAuto{B48EvaluationCandidate}[def]{9}}
 \proofstepwidestar[9]{\exists W\in\powerset(A)\;(z=(p,W)\land s\in W)}{\rAEb{10}}
 \proofstepwidestar[12]{W\in\powerset(A)\land(z=(p,W)\land s\in W)}{\rA}
 \proofstepwidestar[12]{z=(p,W)}{\rAEa{\rAEb{12}}}
 \proofstepwidestar[2,12]{(p,U)=(p,W)}{\rIE{2,13}}
 \proofstepwidestar[2,12]{U=W}{\FormulaRefAuto{TaggedPairRightInjective}[thm]{14}}
 \proofstepwidestar[2,12]{W=U}{\FormulaRefAuto{a=b\vdash b=a}[thm]{15}}
 \proofstepwidestar[12]{s\in W}{\rAEb{\rAEb{12}}}
 \proofstepwidestar[2,12]{s\in U}{\rIE{16,17}}
 \proofstepwidestar[2,9]{s\in U}{\rEE{11,12:18}}
 \proofstepwidestar[2]{s\in V_A(z,p)\to s\in U}{\rRI{9,19}}
 \proofstepwidestar[1,2]{s\in U\leftrightarrow s\in V_A(z,p)}{\rLRI{8,20}}
 \proofstepwidestar[1,2]{U=V_A(z,p)}{\FormulaRefAuto{\forall x(x\in A\leftrightarrow x\in B)\eqvdash A=B}[thm]{\rUI{21}}}
\closeproofpartwideindR

\proofpartwideindR[Typisierung jeder ausgelesenen Wertemenge]{\beta_e(p)\in\powerset(A_D)}{\beta_e(p)\in\powerset(A_D)}[B48BetaType]
 \proofstepwidestar[]{\beta_e(p)\in\powerset(A_D)}{\FormulaRefAuto{B48SeparationPower}[thm]}
\closeproofpartwideindR

\proofpartwideindR[Blattauswertung mit sichtbarem Baumcode]{\mathsf{Rec}(e,M),\ c\in\Sigma\vdash e(\mathsf l(c))=(\mathsf l(c),B_M(c))}{\mathsf{Rec}(e,M),\ c\in\Sigma\vdash e(\mathsf l(c))=(\mathsf l(c),B_M(c))}[B48EvalLeafEquation]
 \proofstepwidestar[1]{\mathsf{Rec}(e,M)}{\rA}
 \proofstepwidestar[2]{c\in\Sigma}{\rA}
 \proofstepwidestar[1]{\forall d\in\Sigma\;(e(\mathsf l(d))=f_M(d))}{\FormulaRefAuto{B48RecLeaf}[thm]{1}}
 \proofstepwidestar[1,2]{e(\mathsf l(c))=f_M(c)}{\rRE{\rUE{3},2}}
 \proofstepwidestar[2]{f_M(c)=(\mathsf l(c),B_M(c))}{\FormulaRefAuto{B48LeafGraphValue}[thm]{2}}
 \proofstepwidestar[1,2]{e(\mathsf l(c))=(\mathsf l(c),B_M(c))}{\rIE{5,4}}
\closeproofpartwideindR

\proofpartwideindR[Wertemenge eines Blattes]{\mathsf{Rec}(e,M),\ c\in\Sigma\vdash\beta_e(\mathsf l(c))=B_M(c)}{\mathsf{Rec}(e,M),\ c\in\Sigma\vdash\beta_e(\mathsf l(c))=B_M(c)}[B48EvalLeafValue]
 \proofstepwidestar[1]{\mathsf{Rec}(e,M)}{\rA}
 \proofstepwidestar[2]{c\in\Sigma}{\rA}
 \proofstepwidestar[1,2]{e(\mathsf l(c))=(\mathsf l(c),B_M(c))}{\FormulaRefAuto{B48EvalLeafEquation}[thm]{1,2}}
 \proofstepwidestar[2]{B_M(c)\in\powerset(A_D)}{\FormulaRefAuto{B48BType}[thm]{2}}
 \proofstepwidestar[1,2]{B_M(c)=\beta_e(\mathsf l(c))}{\FormulaRefAuto{B48ReadValue}[thm]{4,3}}
 \proofstepwidestar[1,2]{\beta_e(\mathsf l(c))=B_M(c)}{\FormulaRefAuto{a=b\vdash b=a}[thm]{5}}
\closeproofpartwideindR

\proofpartwideindR[Erhaltung des Blattcodes]{\mathsf{Rec}(e,M),\ c\in\Sigma\vdash e(\mathsf l(c))=(\mathsf l(c),\beta_e(\mathsf l(c)))}{\mathsf{Rec}(e,M),\ c\in\Sigma\vdash e(\mathsf l(c))=(\mathsf l(c),\beta_e(\mathsf l(c)))}[B48MetadataLeaf]
 \proofstepwidestar[1]{\mathsf{Rec}(e,M)}{\rA}
 \proofstepwidestar[2]{c\in\Sigma}{\rA}
 \proofstepwidestar[1,2]{e(\mathsf l(c))=(\mathsf l(c),B_M(c))}{\FormulaRefAuto{B48EvalLeafEquation}[thm]{1,2}}
 \proofstepwidestar[1,2]{\beta_e(\mathsf l(c))=B_M(c)}{\FormulaRefAuto{B48EvalLeafValue}[thm]{1,2}}
 \proofstepwidestar[1,2]{B_M(c)=\beta_e(\mathsf l(c))}{\FormulaRefAuto{a=b\vdash b=a}[thm]{4}}
 \proofstepwidestar[1,2]{e(\mathsf l(c))=(\mathsf l(c),\beta_e(\mathsf l(c)))}{\rIE{5,3}}
\closeproofpartwideindR

\proofpartwideindR[Erhaltung des Knotencodes]{\begin{gathered}\mathsf{Rec}(e,M),\ p\in\mathcal T,\ q\in\mathcal T,\\ e(p)=(p,\beta_e(p))\land e(q)=(q,\beta_e(q))\\ \vdash e(\mathsf n(p,q))=(\mathsf n(p,q),\beta_e(\mathsf n(p,q)))\end{gathered}}{\begin{gathered}\mathsf{Rec}(e,M),\ p\in\mathcal T,\ q\in\mathcal T,\\ e(p)=(p,\beta_e(p))\land e(q)=(q,\beta_e(q))\\ \vdash e(\mathsf n(p,q))=(\mathsf n(p,q),\beta_e(\mathsf n(p,q)))\end{gathered}}[B48MetadataNode]
 \proofstepwidestar[1]{\mathsf{Rec}(e,M)}{\rA}
 \proofstepwidestar[2]{p\in\mathcal T}{\rA}
 \proofstepwidestar[3]{q\in\mathcal T}{\rA}
 \proofstepwidestar[4]{e(p)=(p,\beta_e(p))\land e(q)=(q,\beta_e(q))}{\rA}
 \proofstepwidestar[1]{\forall u\in\mathcal T\,\forall v\in\mathcal T\;(e(\mathsf n(u,v))=g_M(e(u),e(v)))}{\FormulaRefAuto{B48RecNode}[thm]{1}}
 \proofstepwidestar[1,2,3]{e(\mathsf n(p,q))=g_M(e(p),e(q))}{\rRE{\rUE{\rRE{\rUE{5},2}},3}}
 \proofstepwidestar[]{\beta_e(p)\in\powerset(A_D)}{\FormulaRefAuto{B48BetaType}[thm]}
 \proofstepwidestar[]{\beta_e(q)\in\powerset(A_D)}{\FormulaRefAuto{B48BetaType}[thm]}
 \proofstepwidestar[2]{(p,\beta_e(p))\in Y_D}{\FormulaRefAuto{(a,b)\in A\times B\eqvdash a\in A\land b\in B}[thm]{\rAI{2,7}}}
 \proofstepwidestar[3]{(q,\beta_e(q))\in Y_D}{\FormulaRefAuto{(a,b)\in A\times B\eqvdash a\in A\land b\in B}[thm]{\rAI{3,8}}}
 \proofstepwidestar[2,3]{g_M((p,\beta_e(p)),(q,\beta_e(q)))=(\mathsf n(p,q),H_M(p,\beta_e(p),\beta_e(q)))}{\FormulaRefAuto{B48NodeGraphValue}[thm]{9,10}}
 \proofstepwidestar[1,2,3,4]{e(\mathsf n(p,q))=g_M((p,\beta_e(p)),(q,\beta_e(q)))}{\rIE{\rAEb{4},\rIE{\rAEa{4},6}}}
 \proofstepwidestar[1,2,3,4]{e(\mathsf n(p,q))=(\mathsf n(p,q),H_M(p,\beta_e(p),\beta_e(q)))}{\rIE{11,12}}
 \proofstepwidestar[]{H_M(p,\beta_e(p),\beta_e(q))\in\powerset(A_D)}{\FormulaRefAuto{B48HType}[thm]{7,8}}
 \proofstepwidestar[1,2,3,4]{H_M(p,\beta_e(p),\beta_e(q))=\beta_e(\mathsf n(p,q))}{\FormulaRefAuto{B48ReadValue}[thm]{14,13}}
 \proofstepwidestar[1,2,3,4]{e(\mathsf n(p,q))=(\mathsf n(p,q),\beta_e(\mathsf n(p,q)))}{\rIE{15,13}}
\closeproofpartwideindR

\proofpartwideindR[Strukturinduktion ueber alle Hilfsbaeume]{\mathsf{Rec}(e,M)\vdash\forall p\in\mathcal T\;(e(p)=(p,\beta_e(p)))}{\mathsf{Rec}(e,M)\vdash\forall p\in\mathcal T\;(e(p)=(p,\beta_e(p)))}[B48MetadataAll]
 \proofstepwidestar[1]{\mathsf{Rec}(e,M)}{\rA}
 \proofstepwidestar[2]{c\in\Sigma}{\rA}
 \proofstepwidestar[1,2]{e(\mathsf l(c))=(\mathsf l(c),\beta_e(\mathsf l(c)))}{\FormulaRefAuto{B48MetadataLeaf}[thm]{1,2}}
 \proofstepwidestar[1]{\forall c\in\Sigma\;(e(\mathsf l(c))=(\mathsf l(c),\beta_e(\mathsf l(c))))}{\rUI{\rRI{2,3}}}
 \proofstepwidestar[5]{p\in\mathcal T}{\rA}
 \proofstepwidestar[6]{q\in\mathcal T}{\rA}
 \proofstepwidestar[7]{e(p)=(p,\beta_e(p))\land e(q)=(q,\beta_e(q))}{\rA}
 \proofstepwidestar[1,5,6,7]{e(\mathsf n(p,q))=(\mathsf n(p,q),\beta_e(\mathsf n(p,q)))}{\FormulaRefAuto{B48MetadataNode}[thm]{1,5,6,7}}
 \proofstepwidestar[1]{\begin{gathered}\forall p\in\mathcal T\,\forall q\in\mathcal T\;\\
 ((e(p)=(p,\beta_e(p))\land e(q)=(q,\beta_e(q)))\\
 \to e(\mathsf n(p,q))=(\mathsf n(p,q),\beta_e(\mathsf n(p,q))))\end{gathered}}{\rUI{\rRI{5,\rUI{\rRI{6,\rRI{7,8}}}}}}
 \proofstepwidestar[1]{\forall p\in\mathcal T\;(e(p)=(p,\beta_e(p)))}{\FormulaRefAuto{FullPlanarBinaryTreeStructuralInduction}[thm]{4,9}}
\closeproofpartwideindR

\proofpartwideindR[Wertemenge eines Knotens]{\mathsf{Rec}(e,M),\ p\in\mathcal T,\ q\in\mathcal T\vdash\beta_e(\mathsf n(p,q))=H_M(p,\beta_e(p),\beta_e(q))}{\mathsf{Rec}(e,M),\ p\in\mathcal T,\ q\in\mathcal T\vdash\beta_e(\mathsf n(p,q))=H_M(p,\beta_e(p),\beta_e(q))}[B48EvalNodeValue]
 \proofstepwidestar[1]{\mathsf{Rec}(e,M)}{\rA}
 \proofstepwidestar[2]{p\in\mathcal T}{\rA}
 \proofstepwidestar[3]{q\in\mathcal T}{\rA}
 \proofstepwidestar[1]{\forall t\in\mathcal T\;(e(t)=(t,\beta_e(t)))}{\FormulaRefAuto{B48MetadataAll}[thm]{1}}
 \proofstepwidestar[1,2]{e(p)=(p,\beta_e(p))}{\rRE{\rUE{4},2}}
 \proofstepwidestar[1,3]{e(q)=(q,\beta_e(q))}{\rRE{\rUE{4},3}}
 \proofstepwidestar[1]{\forall u\in\mathcal T\,\forall v\in\mathcal T\;(e(\mathsf n(u,v))=g_M(e(u),e(v)))}{\FormulaRefAuto{B48RecNode}[thm]{1}}
 \proofstepwidestar[1,2,3]{e(\mathsf n(p,q))=g_M(e(p),e(q))}{\rRE{\rUE{\rRE{\rUE{7},2}},3}}
 \proofstepwidestar[]{\beta_e(p)\in\powerset(A_D)}{\FormulaRefAuto{B48BetaType}[thm]}
 \proofstepwidestar[]{\beta_e(q)\in\powerset(A_D)}{\FormulaRefAuto{B48BetaType}[thm]}
 \proofstepwidestar[2]{(p,\beta_e(p))\in Y_D}{\FormulaRefAuto{(a,b)\in A\times B\eqvdash a\in A\land b\in B}[thm]{\rAI{2,9}}}
 \proofstepwidestar[3]{(q,\beta_e(q))\in Y_D}{\FormulaRefAuto{(a,b)\in A\times B\eqvdash a\in A\land b\in B}[thm]{\rAI{3,10}}}
 \proofstepwidestar[2,3]{g_M((p,\beta_e(p)),(q,\beta_e(q)))=(\mathsf n(p,q),H_M(p,\beta_e(p),\beta_e(q)))}{\FormulaRefAuto{B48NodeGraphValue}[thm]{11,12}}
 \proofstepwidestar[1,2,3]{e(\mathsf n(p,q))=(\mathsf n(p,q),H_M(p,\beta_e(p),\beta_e(q)))}{\rIE{13,\rIE{6,\rIE{5,8}}}}
 \proofstepwidestar[]{H_M(p,\beta_e(p),\beta_e(q))\in\powerset(A_D)}{\FormulaRefAuto{B48HType}[thm]{9,10}}
 \proofstepwidestar[1,2,3]{H_M(p,\beta_e(p),\beta_e(q))=\beta_e(\mathsf n(p,q))}{\FormulaRefAuto{B48ReadValue}[thm]{15,14}}
 \proofstepwidestar[1,2,3]{\beta_e(\mathsf n(p,q))=H_M(p,\beta_e(p),\beta_e(q))}{\FormulaRefAuto{a=b\vdash b=a}[thm]{16}}
\closeproofpartwideindR

\par\medskip\noindent\textbf{E. Nachweis der Erfuellungsbedingungen}\par
\proofpartwideindR[Negationsgleichung der konstruierten Wertemengen]{\mathsf{Rec}(e,M),\ p\in\mathcal T\vdash\beta_e(\mathsf{neg}(p))=A_D\setminus\beta_e(p)}{\mathsf{Rec}(e,M),\ p\in\mathcal T\vdash\beta_e(\mathsf{neg}(p))=A_D\setminus\beta_e(p)}[B48BetaNeg]
 \proofstepwidestar[1]{\mathsf{Rec}(e,M)}{\rA}
 \proofstepwidestar[2]{p\in\mathcal T}{\rA}
 \proofstepwidestar[]{\mathsf n_0\in\mathcal T}{\FormulaRefAuto{B48NegRootType}[thm]}
 \proofstepwidestar[1,2]{\beta_e(\mathsf n(\mathsf n_0,p))=H_M(\mathsf n_0,\beta_e(\mathsf n_0),\beta_e(p))}{\FormulaRefAuto{B48EvalNodeValue}[thm]{1,3,2}}
 \proofstepwidestar[]{H_M(\mathsf n_0,\beta_e(\mathsf n_0),\beta_e(p))=A_D\setminus\beta_e(p)}{\FormulaRefAuto{B48HNeg}[thm]}
 \proofstepwidestar[1,2]{\beta_e(\mathsf{neg}(p))=A_D\setminus\beta_e(p)}{\rIE{5,4}}
\closeproofpartwideindR

\proofpartwideindR[Konjunktionsgleichung der konstruierten Wertemengen]{\mathsf{Rec}(e,M),\ p\in\mathcal T,\ q\in\mathcal T\vdash\beta_e(\mathsf{and}(p,q))=\beta_e(p)\cap\beta_e(q)}{\mathsf{Rec}(e,M),\ p\in\mathcal T,\ q\in\mathcal T\vdash\beta_e(\mathsf{and}(p,q))=\beta_e(p)\cap\beta_e(q)}[B48BetaAnd]
 \proofstepwidestar[1]{\mathsf{Rec}(e,M)}{\rA}
 \proofstepwidestar[2]{p\in\mathcal T}{\rA}
 \proofstepwidestar[3]{q\in\mathcal T}{\rA}
 \proofstepwidestar[]{\mathsf c_0\in\mathcal T}{\FormulaRefAuto{B48AndRootType}[thm]}
 \proofstepwidestar[2]{\mathsf n(\mathsf c_0,p)\in\mathcal T}{\FormulaRefAuto{B48NodeType}[thm]{4,2}}
 \proofstepwidestar[1,2]{\beta_e(\mathsf n(\mathsf c_0,p))=H_M(\mathsf c_0,\beta_e(\mathsf c_0),\beta_e(p))}{\FormulaRefAuto{B48EvalNodeValue}[thm]{1,4,2}}
 \proofstepwidestar[]{H_M(\mathsf c_0,\beta_e(\mathsf c_0),\beta_e(p))=\beta_e(p)}{\FormulaRefAuto{B48HAndPrefix}[thm]}
 \proofstepwidestar[1,2]{\beta_e(\mathsf n(\mathsf c_0,p))=\beta_e(p)}{\rIE{7,6}}
 \proofstepwidestar[1,2,3]{\beta_e(\mathsf{and}(p,q))=H_M(\mathsf n(\mathsf c_0,p),\beta_e(\mathsf n(\mathsf c_0,p)),\beta_e(q))}{\FormulaRefAuto{B48EvalNodeValue}[thm]{1,5,3}}
 \proofstepwidestar[2]{H_M(\mathsf n(\mathsf c_0,p),\beta_e(\mathsf n(\mathsf c_0,p)),\beta_e(q))=\beta_e(\mathsf n(\mathsf c_0,p))\cap\beta_e(q)}{\FormulaRefAuto{B48HAnd}[thm]{2}}
 \proofstepwidestar[1,2,3]{\beta_e(\mathsf{and}(p,q))=\beta_e(\mathsf n(\mathsf c_0,p))\cap\beta_e(q)}{\rIE{10,9}}
 \proofstepwidestar[1,2,3]{\beta_e(\mathsf{and}(p,q))=\beta_e(p)\cap\beta_e(q)}{\rIE{8,11}}
\closeproofpartwideindR

\proofpartwideindR[Existenzgleichung der konstruierten Wertemengen]{\mathsf{Rec}(e,M),\ i\in\mathbb N,\ p\in\mathcal T\vdash\beta_e(\mathsf{ex}_i(p))=Q_i(\beta_e(p))}{\mathsf{Rec}(e,M),\ i\in\mathbb N,\ p\in\mathcal T\vdash\beta_e(\mathsf{ex}_i(p))=Q_i(\beta_e(p))}[B48BetaExists]
 \proofstepwidestar[1]{\mathsf{Rec}(e,M)}{\rA}
 \proofstepwidestar[2]{i\in\mathbb N}{\rA}
 \proofstepwidestar[3]{p\in\mathcal T}{\rA}
 \proofstepwidestar[2]{\mathsf q_i\in\mathcal T}{\FormulaRefAuto{B48ExistsRootType}[thm]{2}}
 \proofstepwidestar[1,2,3]{\beta_e(\mathsf{ex}_i(p))=H_M(\mathsf q_i,\beta_e(\mathsf q_i),\beta_e(p))}{\FormulaRefAuto{B48EvalNodeValue}[thm]{1,4,3}}
 \proofstepwidestar[2]{H_M(\mathsf q_i,\beta_e(\mathsf q_i),\beta_e(p))=Q_i(\beta_e(p))}{\FormulaRefAuto{B48HExists}[thm]{2}}
 \proofstepwidestar[1,2,3]{\beta_e(\mathsf{ex}_i(p))=Q_i(\beta_e(p))}{\rIE{6,5}}
\closeproofpartwideindR

\proofpartwideindR[Der Kandidat als Formel-Belegungsrelation]{S_e\subseteq\mathcal F\times A_D}{S_e\subseteq\mathcal F\times A_D}[B48CandidateSubset]
 \proofstepwidestar[]{S_e\subseteq\mathcal F\times A_D}{\FormulaRefAuto{\{x\in A\mid P(x)\}\subseteq A}[thm]}
\closeproofpartwideindR

\proofpartwideindR[Fasern des Kandidaten]{p\in\mathcal F\vdash R_{S_e}(p)=\beta_e(p)}{p\in\mathcal F\vdash R_{S_e}(p)=\beta_e(p)}[B48CandidateRow]
 \proofstepwidestar[1]{p\in\mathcal F}{\rA}
 \proofstepwidestar[2]{s\in R_{S_e}(p)}{\rA}
 \proofstepwidestar[2]{s\in A_D\land(p,s)\in S_e}{\FormulaRefAuto{B48SemanticTerms}[def]{2}}
 \proofstepwidestar[2]{(p,s)\in S_e}{\rAEb{3}}
 \proofstepwidestar[2]{(p,s)\in\mathcal F\times A_D\land s\in\beta_e(p)}{\FormulaRefAuto{B48EvaluationCandidate}[def]{4}}
 \proofstepwidestar[2]{s\in\beta_e(p)}{\rAEb{5}}
 \proofstepwidestar[]{s\in R_{S_e}(p)\to s\in\beta_e(p)}{\rRI{2,6}}
 \proofstepwidestar[8]{s\in\beta_e(p)}{\rA}
 \proofstepwidestar[]{\beta_e(p)\in\powerset(A_D)}{\FormulaRefAuto{B48BetaType}[thm]}
 \proofstepwidestar[]{\beta_e(p)\subseteq A_D}{\FormulaRefAuto{x\in\powerset(A)\eqvdash x\subseteq A}[thm]{9}}
 \proofstepwidestar[8]{s\in A_D}{\FormulaRefAuto{A\subseteq B,\ x\in A\vdash x\in B}[thm]{10,8}}
 \proofstepwidestar[1,8]{(p,s)\in\mathcal F\times A_D}{\FormulaRefAuto{(a,b)\in A\times B\eqvdash a\in A\land b\in B}[thm]{\rAI{1,11}}}
 \proofstepwidestar[1,8]{(p,s)\in S_e}{\FormulaRefAuto{B48EvaluationCandidate}[def]{\rAI{12,8}}}
 \proofstepwidestar[1,8]{s\in R_{S_e}(p)}{\FormulaRefAuto{B48SemanticTerms}[def]{\rAI{11,13}}}
 \proofstepwidestar[1]{s\in\beta_e(p)\to s\in R_{S_e}(p)}{\rRI{8,14}}
 \proofstepwidestar[1]{s\in R_{S_e}(p)\leftrightarrow s\in\beta_e(p)}{\rLRI{7,15}}
 \proofstepwidestar[1]{R_{S_e}(p)=\beta_e(p)}{\FormulaRefAuto{\forall x(x\in A\leftrightarrow x\in B)\eqvdash A=B}[thm]{\rUI{16}}}
\closeproofpartwideindR

\proofpartwideindR[Der Kandidat erfuellt die atomare Klausel]{\mathsf{Rec}(e,M)\vdash C_{\rm at}(M,S_e)}{\mathsf{Rec}(e,M)\vdash C_{\rm at}(M,S_e)}[B48CandidateAtoms]
 \proofstepwidestar[1]{\mathsf{Rec}(e,M)}{\rA}
 \proofstepwidestar[2]{c\in\Sigma_{\rm at}}{\rA}
 \proofstepwidestar[2]{c\in\Sigma}{\FormulaRefAuto{B48AtomicLabelType}[thm]{2}}
 \proofstepwidestar[2]{\mathsf l(c)\in\mathcal F}{\FormulaRefAuto{B48FormulaAtom}[thm]{2}}
 \proofstepwidestar[2]{R_{S_e}(\mathsf l(c))=\beta_e(\mathsf l(c))}{\FormulaRefAuto{B48CandidateRow}[thm]{4}}
 \proofstepwidestar[1,2]{\beta_e(\mathsf l(c))=B_M(c)}{\FormulaRefAuto{B48EvalLeafValue}[thm]{1,3}}
 \proofstepwidestar[1,2]{R_{S_e}(\mathsf l(c))=B_M(c)}{\rIE{6,5}}
 \proofstepwidestar[1]{\forall c\in\Sigma_{\rm at}\;(R_{S_e}(\mathsf l(c))=B_M(c))}{\rUI{\rRI{2,7}}}
 \proofstepwidestar[1]{C_{\rm at}(M,S_e)}{\FormulaRefAuto{B48SatisfactionAtomsDef}[def]{8}}
\closeproofpartwideindR

\proofpartwideindR[Der Kandidat erfuellt die Negationsklausel]{\mathsf{Rec}(e,M)\vdash C_{\rm n}(M,S_e)}{\mathsf{Rec}(e,M)\vdash C_{\rm n}(M,S_e)}[B48CandidateNeg]
 \proofstepwidestar[1]{\mathsf{Rec}(e,M)}{\rA}
 \proofstepwidestar[2]{p\in\mathcal F}{\rA}
 \proofstepwidestar[2]{p\in\mathcal T}{\FormulaRefAuto{B48FormulaInTree}[thm]{2}}
 \proofstepwidestar[2]{\mathsf{neg}(p)\in\mathcal F}{\FormulaRefAuto{B48FormulaNeg}[thm]{2}}
 \proofstepwidestar[2]{R_{S_e}(\mathsf{neg}(p))=\beta_e(\mathsf{neg}(p))}{\FormulaRefAuto{B48CandidateRow}[thm]{4}}
 \proofstepwidestar[1,2]{\beta_e(\mathsf{neg}(p))=A_D\setminus\beta_e(p)}{\FormulaRefAuto{B48BetaNeg}[thm]{1,3}}
 \proofstepwidestar[2]{R_{S_e}(p)=\beta_e(p)}{\FormulaRefAuto{B48CandidateRow}[thm]{2}}
 \proofstepwidestar[2]{\beta_e(p)=R_{S_e}(p)}{\FormulaRefAuto{a=b\vdash b=a}[thm]{7}}
 \proofstepwidestar[1,2]{R_{S_e}(\mathsf{neg}(p))=A_D\setminus R_{S_e}(p)}{\rIE{8,\rIE{6,5}}}
 \proofstepwidestar[1]{\forall p\in\mathcal F\;(R_{S_e}(\mathsf{neg}(p))=A_D\setminus R_{S_e}(p))}{\rUI{\rRI{2,9}}}
 \proofstepwidestar[1]{C_{\rm n}(M,S_e)}{\FormulaRefAuto{B48SatisfactionNegDef}[def]{10}}
\closeproofpartwideindR

\proofpartwideindR[Der Kandidat erfuellt die Konjunktionsklausel]{\mathsf{Rec}(e,M)\vdash C_{\rm c}(M,S_e)}{\mathsf{Rec}(e,M)\vdash C_{\rm c}(M,S_e)}[B48CandidateAnd]
 \proofstepwidestar[1]{\mathsf{Rec}(e,M)}{\rA}
 \proofstepwidestar[2]{p\in\mathcal F}{\rA}
 \proofstepwidestar[3]{q\in\mathcal F}{\rA}
 \proofstepwidestar[2]{p\in\mathcal T}{\FormulaRefAuto{B48FormulaInTree}[thm]{2}}
 \proofstepwidestar[3]{q\in\mathcal T}{\FormulaRefAuto{B48FormulaInTree}[thm]{3}}
 \proofstepwidestar[2,3]{\mathsf{and}(p,q)\in\mathcal F}{\FormulaRefAuto{B48FormulaAnd}[thm]{2,3}}
 \proofstepwidestar[2,3]{R_{S_e}(\mathsf{and}(p,q))=\beta_e(\mathsf{and}(p,q))}{\FormulaRefAuto{B48CandidateRow}[thm]{6}}
 \proofstepwidestar[1,2,3]{\beta_e(\mathsf{and}(p,q))=\beta_e(p)\cap\beta_e(q)}{\FormulaRefAuto{B48BetaAnd}[thm]{1,4,5}}
 \proofstepwidestar[2]{R_{S_e}(p)=\beta_e(p)}{\FormulaRefAuto{B48CandidateRow}[thm]{2}}
 \proofstepwidestar[3]{R_{S_e}(q)=\beta_e(q)}{\FormulaRefAuto{B48CandidateRow}[thm]{3}}
 \proofstepwidestar[2]{\beta_e(p)=R_{S_e}(p)}{\FormulaRefAuto{a=b\vdash b=a}[thm]{9}}
 \proofstepwidestar[3]{\beta_e(q)=R_{S_e}(q)}{\FormulaRefAuto{a=b\vdash b=a}[thm]{10}}
 \proofstepwidestar[1,2,3]{R_{S_e}(\mathsf{and}(p,q))=R_{S_e}(p)\cap R_{S_e}(q)}{\rIE{12,\rIE{11,\rIE{8,7}}}}
 \proofstepwidestar[1]{\forall p\in\mathcal F\,\forall q\in\mathcal F\;(R_{S_e}(\mathsf{and}(p,q))=R_{S_e}(p)\cap R_{S_e}(q))}{\rUI{\rRI{2,\rUI{\rRI{3,13}}}}}
 \proofstepwidestar[1]{C_{\rm c}(M,S_e)}{\FormulaRefAuto{B48SatisfactionAndDef}[def]{14}}
\closeproofpartwideindR

\proofpartwideindR[Der Kandidat erfuellt die Quantorenklausel]{\mathsf{Rec}(e,M)\vdash C_{\rm q}(M,S_e)}{\mathsf{Rec}(e,M)\vdash C_{\rm q}(M,S_e)}[B48CandidateExists]
 \proofstepwidestar[1]{\mathsf{Rec}(e,M)}{\rA}
 \proofstepwidestar[2]{i\in\mathbb N}{\rA}
 \proofstepwidestar[3]{p\in\mathcal F}{\rA}
 \proofstepwidestar[3]{p\in\mathcal T}{\FormulaRefAuto{B48FormulaInTree}[thm]{3}}
 \proofstepwidestar[2,3]{\mathsf{ex}_i(p)\in\mathcal F}{\FormulaRefAuto{B48FormulaExists}[thm]{2,3}}
 \proofstepwidestar[2,3]{R_{S_e}(\mathsf{ex}_i(p))=\beta_e(\mathsf{ex}_i(p))}{\FormulaRefAuto{B48CandidateRow}[thm]{5}}
 \proofstepwidestar[1,2,3]{\beta_e(\mathsf{ex}_i(p))=Q_i(\beta_e(p))}{\FormulaRefAuto{B48BetaExists}[thm]{1,2,4}}
 \proofstepwidestar[3]{R_{S_e}(p)=\beta_e(p)}{\FormulaRefAuto{B48CandidateRow}[thm]{3}}
 \proofstepwidestar[3]{\beta_e(p)=R_{S_e}(p)}{\FormulaRefAuto{a=b\vdash b=a}[thm]{8}}
 \proofstepwidestar[1,2,3]{R_{S_e}(\mathsf{ex}_i(p))=Q_i(R_{S_e}(p))}{\rIE{9,\rIE{7,6}}}
 \proofstepwidestar[1]{\forall i\in\mathbb N\,\forall p\in\mathcal F\;(R_{S_e}(\mathsf{ex}_i(p))=Q_i(R_{S_e}(p)))}{\rUI{\rRI{2,\rUI{\rRI{3,10}}}}}
 \proofstepwidestar[1]{C_{\rm q}(M,S_e)}{\FormulaRefAuto{B48SatisfactionExistsDef}[def]{11}}
\closeproofpartwideindR

\proofpartwideindR[Der konstruierte Kandidat ist eine Erfuellungsrelation]{\mathsf{Str}(M),\ \mathsf{Rec}(e,M)\vdash\mathsf{Erf}(M,S_e)}{\mathsf{Str}(M),\ \mathsf{Rec}(e,M)\vdash\mathsf{Erf}(M,S_e)}[B48CandidateSatisfies]
 \proofstepwidestar[1]{\mathsf{Str}(M)}{\rA}
 \proofstepwidestar[2]{\mathsf{Rec}(e,M)}{\rA}
 \proofstepwidestar[]{S_e\subseteq\mathcal F\times A_D}{\FormulaRefAuto{B48CandidateSubset}[thm]}
 \proofstepwidestar[2]{C_{\rm at}(M,S_e)}{\FormulaRefAuto{B48CandidateAtoms}[thm]{2}}
 \proofstepwidestar[2]{C_{\rm n}(M,S_e)}{\FormulaRefAuto{B48CandidateNeg}[thm]{2}}
 \proofstepwidestar[2]{C_{\rm c}(M,S_e)}{\FormulaRefAuto{B48CandidateAnd}[thm]{2}}
 \proofstepwidestar[2]{C_{\rm q}(M,S_e)}{\FormulaRefAuto{B48CandidateExists}[thm]{2}}
 \proofstepwidestar[1,2]{\mathsf{Erf}(M,S_e)}{\FormulaRefAuto{B48SatisfactionStructure}[def]{1,3,4,5,6,7}}
\closeproofpartwideindR

\par\medskip\noindent\textbf{F. Eindeutigkeit aus den Erfuellungsbedingungen}\par
\proofpartwideindR[Bedingung 1 einer Erfuellungsrelation]{\mathsf{Erf}(M,S)\vdash \mathsf{Str}(M)}{\mathsf{Erf}(M,S)\vdash \mathsf{Str}(M)}[B48ErfStructure]
 \proofstepwidestar[1]{\mathsf{Erf}(M,S)}{\rA}
 \proofstepwidestar[1]{\mathsf{Str}(M)}{\FormulaRefAuto{B48ErfStructureAxiom}[ax]{1}}
\closeproofpartwideindR

\proofpartwideindR[Bedingung 2 einer Erfuellungsrelation]{\mathsf{Erf}(M,S)\vdash S\subseteq\mathcal F\times A_D}{\mathsf{Erf}(M,S)\vdash S\subseteq\mathcal F\times A_D}[B48ErfSubset]
 \proofstepwidestar[1]{\mathsf{Erf}(M,S)}{\rA}
 \proofstepwidestar[1]{S\subseteq\mathcal F\times A_D}{\FormulaRefAuto{B48ErfSubsetAxiom}[ax]{1}}
\closeproofpartwideindR

\proofpartwideindR[Bedingung 3 einer Erfuellungsrelation]{\mathsf{Erf}(M,S)\vdash C_{\rm at}(M,S)}{\mathsf{Erf}(M,S)\vdash C_{\rm at}(M,S)}[B48ErfAtoms]
 \proofstepwidestar[1]{\mathsf{Erf}(M,S)}{\rA}
 \proofstepwidestar[1]{\forall c\in\Sigma_{\rm at}\;(R_S(\mathsf l(c))=B_M(c))}{\FormulaRefAuto{B48ErfAtomsAxiom}[ax]{1}}
 \proofstepwidestar[1]{C_{\rm at}(M,S)}{\FormulaRefAuto{B48SatisfactionAtomsDef}[def]{2}}
\closeproofpartwideindR

\proofpartwideindR[Bedingung 4 einer Erfuellungsrelation]{\mathsf{Erf}(M,S)\vdash C_{\rm n}(M,S)}{\mathsf{Erf}(M,S)\vdash C_{\rm n}(M,S)}[B48ErfNeg]
 \proofstepwidestar[1]{\mathsf{Erf}(M,S)}{\rA}
 \proofstepwidestar[1]{\forall p\in\mathcal F\;(R_S(\mathsf{neg}(p))=A_D\setminus R_S(p))}{\FormulaRefAuto{B48ErfNegAxiom}[ax]{1}}
 \proofstepwidestar[1]{C_{\rm n}(M,S)}{\FormulaRefAuto{B48SatisfactionNegDef}[def]{2}}
\closeproofpartwideindR

\proofpartwideindR[Bedingung 5 einer Erfuellungsrelation]{\mathsf{Erf}(M,S)\vdash C_{\rm c}(M,S)}{\mathsf{Erf}(M,S)\vdash C_{\rm c}(M,S)}[B48ErfAnd]
 \proofstepwidestar[1]{\mathsf{Erf}(M,S)}{\rA}
 \proofstepwidestar[1]{\forall p\in\mathcal F\,\forall q\in\mathcal F\;(R_S(\mathsf{and}(p,q))=R_S(p)\cap R_S(q))}{\FormulaRefAuto{B48ErfAndAxiom}[ax]{1}}
 \proofstepwidestar[1]{C_{\rm c}(M,S)}{\FormulaRefAuto{B48SatisfactionAndDef}[def]{2}}
\closeproofpartwideindR

\proofpartwideindR[Bedingung 6 einer Erfuellungsrelation]{\mathsf{Erf}(M,S)\vdash C_{\rm q}(M,S)}{\mathsf{Erf}(M,S)\vdash C_{\rm q}(M,S)}[B48ErfExists]
 \proofstepwidestar[1]{\mathsf{Erf}(M,S)}{\rA}
 \proofstepwidestar[1]{\forall i\in\mathbb N\,\forall p\in\mathcal F\;(R_S(\mathsf{ex}_i(p))=Q_i(R_S(p)))}{\FormulaRefAuto{B48ErfExistsAxiom}[ax]{1}}
 \proofstepwidestar[1]{C_{\rm q}(M,S)}{\FormulaRefAuto{B48SatisfactionExistsDef}[def]{2}}
\closeproofpartwideindR

\proofpartwideindR[Fasergleichung der Atome]{\mathsf{Erf}(M,S),\ c\in\Sigma_{\rm at}\vdash R_S(\mathsf l(c))=B_M(c)}{\mathsf{Erf}(M,S),\ c\in\Sigma_{\rm at}\vdash R_S(\mathsf l(c))=B_M(c)}[B48ErfAtomRow]
 \proofstepwidestar[1]{\mathsf{Erf}(M,S)}{\rA}
 \proofstepwidestar[2]{c\in\Sigma_{\rm at}}{\rA}
 \proofstepwidestar[1]{C_{\rm at}(M,S)}{\FormulaRefAuto{B48ErfAtoms}[thm]{1}}
 \proofstepwidestar[1]{\forall c\in\Sigma_{\rm at}\;(R_S(\mathsf l(c))=B_M(c))}{\FormulaRefAuto{B48SatisfactionAtomsDef}[def]{3}}
 \proofstepwidestar[1,2]{R_S(\mathsf l(c))=B_M(c)}{\rRE{\rUE{4},2}}
\closeproofpartwideindR

\proofpartwideindR[Fasergleichung der Negation]{\mathsf{Erf}(M,S),\ p\in\mathcal F\vdash R_S(\mathsf{neg}(p))=A_D\setminus R_S(p)}{\mathsf{Erf}(M,S),\ p\in\mathcal F\vdash R_S(\mathsf{neg}(p))=A_D\setminus R_S(p)}[B48ErfNegRow]
 \proofstepwidestar[1]{\mathsf{Erf}(M,S)}{\rA}
 \proofstepwidestar[2]{p\in\mathcal F}{\rA}
 \proofstepwidestar[1]{C_{\rm n}(M,S)}{\FormulaRefAuto{B48ErfNeg}[thm]{1}}
 \proofstepwidestar[1]{\forall p\in\mathcal F\;(R_S(\mathsf{neg}(p))=A_D\setminus R_S(p))}{\FormulaRefAuto{B48SatisfactionNegDef}[def]{3}}
 \proofstepwidestar[1,2]{R_S(\mathsf{neg}(p))=A_D\setminus R_S(p)}{\rRE{\rUE{4},2}}
\closeproofpartwideindR

\proofpartwideindR[Fasergleichung der Konjunktion]{\mathsf{Erf}(M,S),\ p\in\mathcal F,\ q\in\mathcal F\vdash R_S(\mathsf{and}(p,q))=R_S(p)\cap R_S(q)}{\mathsf{Erf}(M,S),\ p\in\mathcal F,\ q\in\mathcal F\vdash R_S(\mathsf{and}(p,q))=R_S(p)\cap R_S(q)}[B48ErfAndRow]
 \proofstepwidestar[1]{\mathsf{Erf}(M,S)}{\rA}
 \proofstepwidestar[2]{p\in\mathcal F}{\rA}
 \proofstepwidestar[3]{q\in\mathcal F}{\rA}
 \proofstepwidestar[1]{C_{\rm c}(M,S)}{\FormulaRefAuto{B48ErfAnd}[thm]{1}}
 \proofstepwidestar[1]{\forall u\in\mathcal F\,\forall v\in\mathcal F\;(R_S(\mathsf{and}(u,v))=R_S(u)\cap R_S(v))}{\FormulaRefAuto{B48SatisfactionAndDef}[def]{4}}
 \proofstepwidestar[1,2,3]{R_S(\mathsf{and}(p,q))=R_S(p)\cap R_S(q)}{\rRE{\rUE{\rRE{\rUE{5},2}},3}}
\closeproofpartwideindR

\proofpartwideindR[Fasergleichung des Existenzquantors]{\mathsf{Erf}(M,S),\ i\in\mathbb N,\ p\in\mathcal F\vdash R_S(\mathsf{ex}_i(p))=Q_i(R_S(p))}{\mathsf{Erf}(M,S),\ i\in\mathbb N,\ p\in\mathcal F\vdash R_S(\mathsf{ex}_i(p))=Q_i(R_S(p))}[B48ErfExistsRow]
 \proofstepwidestar[1]{\mathsf{Erf}(M,S)}{\rA}
 \proofstepwidestar[2]{i\in\mathbb N}{\rA}
 \proofstepwidestar[3]{p\in\mathcal F}{\rA}
 \proofstepwidestar[1]{C_{\rm q}(M,S)}{\FormulaRefAuto{B48ErfExists}[thm]{1}}
 \proofstepwidestar[1]{\forall j\in\mathbb N\,\forall q\in\mathcal F\;(R_S(\mathsf{ex}_j(q))=Q_j(R_S(q)))}{\FormulaRefAuto{B48SatisfactionExistsDef}[def]{4}}
 \proofstepwidestar[1,2,3]{R_S(\mathsf{ex}_i(p))=Q_i(R_S(p))}{\rRE{\rUE{\rRE{\rUE{5},2}},3}}
\closeproofpartwideindR

\proofpartwideindR[Elementkriterium des Gleichheitstraegers]{p\in K(S,T)\eqvdash p\in\mathcal F\land R_S(p)=R_T(p)}{p\in K(S,T)\eqvdash p\in\mathcal F\land R_S(p)=R_T(p)}[B48EqualityCarrierMember]
 \proofstepwidestar[]{p\in K(S,T)\leftrightarrow p\in\mathcal F\land R_S(p)=R_T(p)}{\FormulaRefAuto{B48EqualityCarrier}[def]}
\closeproofpartwideindR

\proofpartwideindR[Gleichheitsinduktion bei Atomen]{\mathsf{Erf}(M,S),\ \mathsf{Erf}(M,T),\ c\in\Sigma_{\rm at}\vdash\mathsf l(c)\in K(S,T)}{\mathsf{Erf}(M,S),\ \mathsf{Erf}(M,T),\ c\in\Sigma_{\rm at}\vdash\mathsf l(c)\in K(S,T)}[B48EqualityAtom]
 \proofstepwidestar[1]{\mathsf{Erf}(M,S)}{\rA}
 \proofstepwidestar[2]{\mathsf{Erf}(M,T)}{\rA}
 \proofstepwidestar[3]{c\in\Sigma_{\rm at}}{\rA}
 \proofstepwidestar[3]{\mathsf l(c)\in\mathcal F}{\FormulaRefAuto{B48FormulaAtom}[thm]{3}}
 \proofstepwidestar[1,3]{R_S(\mathsf l(c))=B_M(c)}{\FormulaRefAuto{B48ErfAtomRow}[thm]{1,3}}
 \proofstepwidestar[2,3]{R_T(\mathsf l(c))=B_M(c)}{\FormulaRefAuto{B48ErfAtomRow}[thm]{2,3}}
 \proofstepwidestar[1,2,3]{R_S(\mathsf l(c))=R_T(\mathsf l(c))}{\FormulaRefAuto{a=b,\ c=b\vdash a=c}[thm]{5,6}}
 \proofstepwidestar[1,2,3]{\mathsf l(c)\in K(S,T)}{\FormulaRefAuto{B48EqualityCarrierMember}[thm]{\rAI{4,7}}}
\closeproofpartwideindR

\proofpartwideindR[Gleichheitsinduktion bei Negation]{\mathsf{Erf}(M,S),\ \mathsf{Erf}(M,T),\ p\in K(S,T)\vdash\mathsf{neg}(p)\in K(S,T)}{\mathsf{Erf}(M,S),\ \mathsf{Erf}(M,T),\ p\in K(S,T)\vdash\mathsf{neg}(p)\in K(S,T)}[B48EqualityNeg]
 \proofstepwidestar[1]{\mathsf{Erf}(M,S)}{\rA}
 \proofstepwidestar[2]{\mathsf{Erf}(M,T)}{\rA}
 \proofstepwidestar[3]{p\in K(S,T)}{\rA}
 \proofstepwidestar[3]{p\in\mathcal F\land R_S(p)=R_T(p)}{\FormulaRefAuto{B48EqualityCarrierMember}[thm]{3}}
 \proofstepwidestar[3]{p\in\mathcal F}{\rAEa{4}}
 \proofstepwidestar[3]{R_S(p)=R_T(p)}{\rAEb{4}}
 \proofstepwidestar[3]{\mathsf{neg}(p)\in\mathcal F}{\FormulaRefAuto{B48FormulaNeg}[thm]{5}}
 \proofstepwidestar[1,3]{R_S(\mathsf{neg}(p))=A_D\setminus R_S(p)}{\FormulaRefAuto{B48ErfNegRow}[thm]{1,5}}
 \proofstepwidestar[2,3]{R_T(\mathsf{neg}(p))=A_D\setminus R_T(p)}{\FormulaRefAuto{B48ErfNegRow}[thm]{2,5}}
 \proofstepwidestar[1,3]{R_S(\mathsf{neg}(p))=A_D\setminus R_T(p)}{\rIE{6,8}}
 \proofstepwidestar[1,2,3]{R_S(\mathsf{neg}(p))=R_T(\mathsf{neg}(p))}{\FormulaRefAuto{a=b,\ c=b\vdash a=c}[thm]{10,9}}
 \proofstepwidestar[1,2,3]{\mathsf{neg}(p)\in K(S,T)}{\FormulaRefAuto{B48EqualityCarrierMember}[thm]{\rAI{7,11}}}
\closeproofpartwideindR

\proofpartwideindR[Gleichheitsinduktion bei Konjunktion]{\mathsf{Erf}(M,S),\ \mathsf{Erf}(M,T),\ p\in K(S,T),\ q\in K(S,T)\vdash\mathsf{and}(p,q)\in K(S,T)}{\mathsf{Erf}(M,S),\ \mathsf{Erf}(M,T),\ p\in K(S,T),\ q\in K(S,T)\vdash\mathsf{and}(p,q)\in K(S,T)}[B48EqualityAnd]
 \proofstepwidestar[1]{\mathsf{Erf}(M,S)}{\rA}
 \proofstepwidestar[2]{\mathsf{Erf}(M,T)}{\rA}
 \proofstepwidestar[3]{p\in K(S,T)}{\rA}
 \proofstepwidestar[4]{q\in K(S,T)}{\rA}
 \proofstepwidestar[3]{p\in\mathcal F\land R_S(p)=R_T(p)}{\FormulaRefAuto{B48EqualityCarrierMember}[thm]{3}}
 \proofstepwidestar[4]{q\in\mathcal F\land R_S(q)=R_T(q)}{\FormulaRefAuto{B48EqualityCarrierMember}[thm]{4}}
 \proofstepwidestar[3,4]{\mathsf{and}(p,q)\in\mathcal F}{\FormulaRefAuto{B48FormulaAnd}[thm]{\rAEa{5},\rAEa{6}}}
 \proofstepwidestar[1,3,4]{R_S(\mathsf{and}(p,q))=R_S(p)\cap R_S(q)}{\FormulaRefAuto{B48ErfAndRow}[thm]{1,\rAEa{5},\rAEa{6}}}
 \proofstepwidestar[2,3,4]{R_T(\mathsf{and}(p,q))=R_T(p)\cap R_T(q)}{\FormulaRefAuto{B48ErfAndRow}[thm]{2,\rAEa{5},\rAEa{6}}}
 \proofstepwidestar[1,3,4]{R_S(\mathsf{and}(p,q))=R_T(p)\cap R_T(q)}{\rIE{\rAEb{6},\rIE{\rAEb{5},8}}}
 \proofstepwidestar[1,2,3,4]{R_S(\mathsf{and}(p,q))=R_T(\mathsf{and}(p,q))}{\FormulaRefAuto{a=b,\ c=b\vdash a=c}[thm]{10,9}}
 \proofstepwidestar[1,2,3,4]{\mathsf{and}(p,q)\in K(S,T)}{\FormulaRefAuto{B48EqualityCarrierMember}[thm]{\rAI{7,11}}}
\closeproofpartwideindR

\proofpartwideindR[Gleichheitsinduktion beim Existenzquantor]{\mathsf{Erf}(M,S),\ \mathsf{Erf}(M,T),\ i\in\mathbb N,\ p\in K(S,T)\vdash\mathsf{ex}_i(p)\in K(S,T)}{\mathsf{Erf}(M,S),\ \mathsf{Erf}(M,T),\ i\in\mathbb N,\ p\in K(S,T)\vdash\mathsf{ex}_i(p)\in K(S,T)}[B48EqualityExists]
 \proofstepwidestar[1]{\mathsf{Erf}(M,S)}{\rA}
 \proofstepwidestar[2]{\mathsf{Erf}(M,T)}{\rA}
 \proofstepwidestar[3]{i\in\mathbb N}{\rA}
 \proofstepwidestar[4]{p\in K(S,T)}{\rA}
 \proofstepwidestar[4]{p\in\mathcal F\land R_S(p)=R_T(p)}{\FormulaRefAuto{B48EqualityCarrierMember}[thm]{4}}
 \proofstepwidestar[3,4]{\mathsf{ex}_i(p)\in\mathcal F}{\FormulaRefAuto{B48FormulaExists}[thm]{3,\rAEa{5}}}
 \proofstepwidestar[1,3,4]{R_S(\mathsf{ex}_i(p))=Q_i(R_S(p))}{\FormulaRefAuto{B48ErfExistsRow}[thm]{1,3,\rAEa{5}}}
 \proofstepwidestar[2,3,4]{R_T(\mathsf{ex}_i(p))=Q_i(R_T(p))}{\FormulaRefAuto{B48ErfExistsRow}[thm]{2,3,\rAEa{5}}}
 \proofstepwidestar[1,3,4]{R_S(\mathsf{ex}_i(p))=Q_i(R_T(p))}{\rIE{\rAEb{5},7}}
 \proofstepwidestar[1,2,3,4]{R_S(\mathsf{ex}_i(p))=R_T(\mathsf{ex}_i(p))}{\FormulaRefAuto{a=b,\ c=b\vdash a=c}[thm]{9,8}}
 \proofstepwidestar[1,2,3,4]{\mathsf{ex}_i(p)\in K(S,T)}{\FormulaRefAuto{B48EqualityCarrierMember}[thm]{\rAI{6,10}}}
\closeproofpartwideindR

\proofpartwideindR[Der Gleichheitstraeger ist abgeschlossen]{\mathsf{Erf}(M,S),\ \mathsf{Erf}(M,T)\vdash\mathsf{Cl}(K(S,T))}{\mathsf{Erf}(M,S),\ \mathsf{Erf}(M,T)\vdash\mathsf{Cl}(K(S,T))}[B48EqualityClosed]
 \proofstepwidestar[1]{\mathsf{Erf}(M,S)}{\rA}
 \proofstepwidestar[2]{\mathsf{Erf}(M,T)}{\rA}
 \proofstepwidestar[]{K(S,T)\subseteq\mathcal F}{\FormulaRefAuto{\{x\in A\mid P(x)\}\subseteq A}[thm]}
 \proofstepwidestar[]{\mathcal F\subseteq\mathcal T}{\FormulaRefAuto{\{x\in A\mid P(x)\}\subseteq A}[thm]}
 \proofstepwidestar[]{K(S,T)\subseteq\mathcal T}{\FormulaRefAuto{A\subseteq B,\ B\subseteq C\vdash A\subseteq C}[thm]{3,4}}
 \proofstepwidestar[6]{c\in\Sigma_{\rm at}}{\rA}
 \proofstepwidestar[1,2,6]{\mathsf l(c)\in K(S,T)}{\FormulaRefAuto{B48EqualityAtom}[thm]{1,2,6}}
 \proofstepwidestar[1,2]{\forall c\in\Sigma_{\rm at}\;(\mathsf l(c)\in K(S,T))}{\rUI{\rRI{6,7}}}
 \proofstepwidestar[9]{p\in K(S,T)}{\rA}
 \proofstepwidestar[1,2,9]{\mathsf{neg}(p)\in K(S,T)}{\FormulaRefAuto{B48EqualityNeg}[thm]{1,2,9}}
 \proofstepwidestar[1,2]{\forall p\in K(S,T)\;(\mathsf{neg}(p)\in K(S,T))}{\rUI{\rRI{9,10}}}
 \proofstepwidestar[12]{p\in K(S,T)}{\rA}
 \proofstepwidestar[13]{q\in K(S,T)}{\rA}
 \proofstepwidestar[1,2,12,13]{\mathsf{and}(p,q)\in K(S,T)}{\FormulaRefAuto{B48EqualityAnd}[thm]{1,2,12,13}}
 \proofstepwidestar[1,2]{\forall p\in K(S,T)\,\forall q\in K(S,T)\;(\mathsf{and}(p,q)\in K(S,T))}{\rUI{\rRI{12,\rUI{\rRI{13,14}}}}}
 \proofstepwidestar[16]{i\in\mathbb N}{\rA}
 \proofstepwidestar[17]{p\in K(S,T)}{\rA}
 \proofstepwidestar[1,2,16,17]{\mathsf{ex}_i(p)\in K(S,T)}{\FormulaRefAuto{B48EqualityExists}[thm]{1,2,16,17}}
 \proofstepwidestar[1,2]{\forall i\in\mathbb N\,\forall p\in K(S,T)\;(\mathsf{ex}_i(p)\in K(S,T))}{\rUI{\rRI{16,\rUI{\rRI{17,18}}}}}
 \proofstepwidestar[1,2]{\mathsf{Cl}(K(S,T))}{\FormulaRefAuto{B48FormulaClosure}[def]{5,8,11,15,19}}
\closeproofpartwideindR

\proofpartwideindR[Eindeutigkeit aller Formelfasern]{\mathsf{Erf}(M,S),\ \mathsf{Erf}(M,T)\vdash\forall p\in\mathcal F\;(R_S(p)=R_T(p))}{\mathsf{Erf}(M,S),\ \mathsf{Erf}(M,T)\vdash\forall p\in\mathcal F\;(R_S(p)=R_T(p))}[B48EqualityAllRows]
 \proofstepwidestar[1]{\mathsf{Erf}(M,S)}{\rA}
 \proofstepwidestar[2]{\mathsf{Erf}(M,T)}{\rA}
 \proofstepwidestar[1,2]{\mathsf{Cl}(K(S,T))}{\FormulaRefAuto{B48EqualityClosed}[thm]{1,2}}
 \proofstepwidestar[4]{p\in\mathcal F}{\rA}
 \proofstepwidestar[1,2,4]{p\in K(S,T)}{\FormulaRefAuto{B48FormulaMinimalMember}[thm]{3,4}}
 \proofstepwidestar[1,2,4]{p\in\mathcal F\land R_S(p)=R_T(p)}{\FormulaRefAuto{B48EqualityCarrierMember}[thm]{5}}
 \proofstepwidestar[1,2,4]{R_S(p)=R_T(p)}{\rAEb{6}}
 \proofstepwidestar[1,2]{\forall p\in\mathcal F\;(R_S(p)=R_T(p))}{\rUI{\rRI{4,7}}}
\closeproofpartwideindR

\proofpartwideindR[Gleiche Fasern liefern eine Inklusion]{S\subseteq\mathcal F\times A_D,\ \forall p\in\mathcal F\;(R_S(p)=R_T(p))\vdash S\subseteq T}{S\subseteq\mathcal F\times A_D,\ \forall p\in\mathcal F\;(R_S(p)=R_T(p))\vdash S\subseteq T}[B48RowsSubset]
 \proofstepwidestar[1]{S\subseteq\mathcal F\times A_D}{\rA}
 \proofstepwidestar[2]{\forall p\in\mathcal F\;(R_S(p)=R_T(p))}{\rA}
 \proofstepwidestar[3]{z\in S}{\rA}
 \proofstepwidestar[1,3]{z\in\mathcal F\times A_D}{\FormulaRefAuto{A\subseteq B,\ x\in A\vdash x\in B}[thm]{1,3}}
 \proofstepwidestar[1,3]{\exists p\in\mathcal F\,\exists s\in A_D\;(z=(p,s))}{\FormulaRefAuto{x\in A\times B\eqvdash\exists a\in A\,\exists b\in B\,x=(a,b)}[thm]{4}}
 \proofstepwidestar[6]{p\in\mathcal F\land(s\in A_D\land z=(p,s))}{\rA}
 \proofstepwidestar[6]{p\in\mathcal F}{\rAEa{6}}
 \proofstepwidestar[6]{s\in A_D}{\rAEa{\rAEb{6}}}
 \proofstepwidestar[6]{z=(p,s)}{\rAEb{\rAEb{6}}}
 \proofstepwidestar[3,6]{(p,s)\in S}{\rIE{9,3}}
 \proofstepwidestar[3,6]{s\in R_S(p)}{\FormulaRefAuto{B48SemanticTerms}[def]{\rAI{8,10}}}
 \proofstepwidestar[2,6]{R_S(p)=R_T(p)}{\rRE{\rUE{2},7}}
 \proofstepwidestar[2,3,6]{s\in R_T(p)}{\rIE{12,11}}
 \proofstepwidestar[2,3,6]{s\in A_D\land(p,s)\in T}{\FormulaRefAuto{B48SemanticTerms}[def]{13}}
 \proofstepwidestar[2,3,6]{(p,s)\in T}{\rAEb{14}}
 \proofstepwidestar[6]{(p,s)=z}{\FormulaRefAuto{a=b\vdash b=a}[thm]{9}}
 \proofstepwidestar[2,3,6]{z\in T}{\rIE{16,15}}
 \proofstepwidestar[1,2,3]{z\in T}{\rEEStar{5,6:17}}
 \proofstepwidestar[1,2]{S\subseteq T}{\FormulaRefAuto{SubsetIntroductionRule}[thm]{[3]\vdots 18}}
\closeproofpartwideindR

\proofpartwideindR[Eindeutigkeit der Erfuellungsrelation]{\mathsf{Erf}(M,S),\ \mathsf{Erf}(M,T)\vdash S=T}{\mathsf{Erf}(M,S),\ \mathsf{Erf}(M,T)\vdash S=T}[B48SatisfactionUnique]
 \proofstepwidestar[1]{\mathsf{Erf}(M,S)}{\rA}
 \proofstepwidestar[2]{\mathsf{Erf}(M,T)}{\rA}
 \proofstepwidestar[1]{S\subseteq\mathcal F\times A_D}{\FormulaRefAuto{B48ErfSubset}[thm]{1}}
 \proofstepwidestar[2]{T\subseteq\mathcal F\times A_D}{\FormulaRefAuto{B48ErfSubset}[thm]{2}}
 \proofstepwidestar[1,2]{\forall p\in\mathcal F\;(R_S(p)=R_T(p))}{\FormulaRefAuto{B48EqualityAllRows}[thm]{1,2}}
 \proofstepwidestar[1,2]{\forall p\in\mathcal F\;(R_T(p)=R_S(p))}{\FormulaRefAuto{B48EqualityAllRows}[thm]{2,1}}
 \proofstepwidestar[1,2]{S\subseteq T}{\FormulaRefAuto{B48RowsSubset}[thm]{3,5}}
 \proofstepwidestar[1,2]{T\subseteq S}{\FormulaRefAuto{B48RowsSubset}[thm]{4,6}}
 \proofstepwidestar[1,2]{S=T}{\FormulaRefAuto{A\subseteq B,\ B\subseteq A\vdash A=B}[thm]{7,8}}
\closeproofpartwideindR

\par\medskip\noindent\textbf{G. Punktweise Wahrheitsklauseln}\par
\proofpartwideindR[Umkehrung einer Aequivalenz]{P\leftrightarrow Q\vdash Q\leftrightarrow P}{P\leftrightarrow Q\vdash Q\leftrightarrow P}[B48IffSymmetry]
 \proofstepwidestar[1]{P\leftrightarrow Q}{\rA}
 \proofstepwidestar[2]{Q}{\rA}
 \proofstepwidestar[1,2]{P}{\rLREb{1,2}}
 \proofstepwidestar[1]{Q\to P}{\rRI{2,3}}
 \proofstepwidestar[5]{P}{\rA}
 \proofstepwidestar[1,5]{Q}{\rLREa{1,5}}
 \proofstepwidestar[1]{P\to Q}{\rRI{5,6}}
 \proofstepwidestar[1]{Q\leftrightarrow P}{\rLRI{4,7}}
\closeproofpartwideindR

\proofpartwideindR[Aequivalente beschraenkte Existenzbehauptungen]{\forall x\in A\;(P(x)\leftrightarrow Q(x))\vdash((\exists x\in A\;P(x))\leftrightarrow(\exists x\in A\;Q(x)))}{\forall x\in A\;(P(x)\leftrightarrow Q(x))\vdash((\exists x\in A\;P(x))\leftrightarrow(\exists x\in A\;Q(x)))}[B48ExistsIff]
 \proofstepwidestar[1]{\forall x\in A\;(P(x)\leftrightarrow Q(x))}{\rA}
 \proofstepwidestar[2]{\exists x\in A\;P(x)}{\rA}
 \proofstepwidestar[3]{a\in A\land P(a)}{\rA}
 \proofstepwidestar[1,3]{P(a)\leftrightarrow Q(a)}{\rRE{\rUE{1},\rAEa{3}}}
 \proofstepwidestar[1,3]{Q(a)}{\rLREa{4,\rAEb{3}}}
 \proofstepwidestar[1,3]{\exists x\in A\;Q(x)}{\rEI{\rAI{\rAEa{3},5}}}
 \proofstepwidestar[1,2]{\exists x\in A\;Q(x)}{\rEE{2,3:6}}
 \proofstepwidestar[1]{(\exists x\in A\;P(x))\to(\exists x\in A\;Q(x))}{\rRI{2,7}}
 \proofstepwidestar[9]{\exists x\in A\;Q(x)}{\rA}
 \proofstepwidestar[10]{b\in A\land Q(b)}{\rA}
 \proofstepwidestar[1,10]{P(b)\leftrightarrow Q(b)}{\rRE{\rUE{1},\rAEa{10}}}
 \proofstepwidestar[1,10]{P(b)}{\rLREb{11,\rAEb{10}}}
 \proofstepwidestar[1,10]{\exists x\in A\;P(x)}{\rEI{\rAI{\rAEa{10},12}}}
 \proofstepwidestar[1,9]{\exists x\in A\;P(x)}{\rEE{9,10:13}}
 \proofstepwidestar[1]{(\exists x\in A\;Q(x))\to(\exists x\in A\;P(x))}{\rRI{9,14}}
 \proofstepwidestar[1]{(\exists x\in A\;P(x))\leftrightarrow(\exists x\in A\;Q(x))}{\rLRI{8,15}}
\closeproofpartwideindR

\proofpartwideindR[Faserzugehoerigkeit und Relationszugehoerigkeit]{s\in A_D\vdash(s\in R_S(p)\leftrightarrow(p,s)\in S)}{s\in A_D\vdash(s\in R_S(p)\leftrightarrow(p,s)\in S)}[B48RowMember]
 \proofstepwidestar[1]{s\in A_D}{\rA}
 \proofstepwidestar[1]{s\in R_S(p)\leftrightarrow(p,s)\in S}{\FormulaRefAuto{B48SeparationWithin}[thm]{1}}
\closeproofpartwideindR

\proofpartwideindR[Punktweise atomare Erfuellungsklausel]{\mathsf{Erf}(M,S),\ c\in\Sigma_{\rm at},\ s\in A_D\vdash((\mathsf l(c),s)\in S\leftrightarrow s\in B_M(c))}{\mathsf{Erf}(M,S),\ c\in\Sigma_{\rm at},\ s\in A_D\vdash((\mathsf l(c),s)\in S\leftrightarrow s\in B_M(c))}[B48PointAtom]
 \proofstepwidestar[1]{\mathsf{Erf}(M,S)}{\rA}
 \proofstepwidestar[2]{c\in\Sigma_{\rm at}}{\rA}
 \proofstepwidestar[3]{s\in A_D}{\rA}
 \proofstepwidestar[1,2]{R_S(\mathsf l(c))=B_M(c)}{\FormulaRefAuto{B48ErfAtomRow}[thm]{1,2}}
 \proofstepwidestar[3]{s\in R_S(\mathsf l(c))\leftrightarrow(\mathsf l(c),s)\in S}{\FormulaRefAuto{B48RowMember}[thm]{3}}
 \proofstepwidestar[3]{(\mathsf l(c),s)\in S\leftrightarrow s\in R_S(\mathsf l(c))}{\FormulaRefAuto{B48IffSymmetry}[thm]{5}}
 \proofstepwidestar[1,2,3]{(\mathsf l(c),s)\in S\leftrightarrow s\in B_M(c)}{\rIE{4,6}}
\closeproofpartwideindR

\proofpartwideindR[Punktweise Negationsklausel]{\mathsf{Erf}(M,S),\ p\in\mathcal F,\ s\in A_D\vdash((\mathsf{neg}(p),s)\in S\leftrightarrow\neg((p,s)\in S))}{\mathsf{Erf}(M,S),\ p\in\mathcal F,\ s\in A_D\vdash((\mathsf{neg}(p),s)\in S\leftrightarrow\neg((p,s)\in S))}[B48PointNeg]
 \proofstepwidestar[1]{\mathsf{Erf}(M,S)}{\rA}
 \proofstepwidestar[2]{p\in\mathcal F}{\rA}
 \proofstepwidestar[3]{s\in A_D}{\rA}
 \proofstepwidestar[1,2]{R_S(\mathsf{neg}(p))=A_D\setminus R_S(p)}{\FormulaRefAuto{B48ErfNegRow}[thm]{1,2}}
 \proofstepwidestar[3]{s\in R_S(\mathsf{neg}(p))\leftrightarrow(\mathsf{neg}(p),s)\in S}{\FormulaRefAuto{B48RowMember}[thm]{3}}
 \proofstepwidestar[3]{(\mathsf{neg}(p),s)\in S\leftrightarrow s\in R_S(\mathsf{neg}(p))}{\FormulaRefAuto{B48IffSymmetry}[thm]{5}}
 \proofstepwidestar[1,2,3]{(\mathsf{neg}(p),s)\in S\leftrightarrow s\in A_D\setminus R_S(p)}{\rIE{4,6}}
 \proofstepwidestar[3]{s\in A_D\setminus R_S(p)\leftrightarrow\neg(s\in R_S(p))}{\FormulaRefAuto{B48SeparationWithin}[thm]{3}}
 \proofstepwidestar[1,2,3]{(\mathsf{neg}(p),s)\in S\leftrightarrow\neg(s\in R_S(p))}{\rChain{7,8}}
 \proofstepwidestar[3]{s\in R_S(p)\leftrightarrow(p,s)\in S}{\FormulaRefAuto{B48RowMember}[thm]{3}}
 \proofstepwidestar[1,2,3]{(\mathsf{neg}(p),s)\in S\leftrightarrow\neg((p,s)\in S)}{\rLRS{10,9}}
\closeproofpartwideindR

\proofpartwideindR[Punktweise Konjunktionsklausel]{\begin{gathered}\mathsf{Erf}(M,S),\ p\in\mathcal F,\ q\in\mathcal F,\ s\in A_D\\
 \vdash((\mathsf{and}(p,q),s)\in S\leftrightarrow((p,s)\in S\land(q,s)\in S))\end{gathered}}{\begin{gathered}\mathsf{Erf}(M,S),\ p\in\mathcal F,\ q\in\mathcal F,\ s\in A_D\\
 \vdash((\mathsf{and}(p,q),s)\in S\leftrightarrow((p,s)\in S\land(q,s)\in S))\end{gathered}}[B48PointAnd]
 \proofstepwidestar[1]{\mathsf{Erf}(M,S)}{\rA}
 \proofstepwidestar[2]{p\in\mathcal F}{\rA}
 \proofstepwidestar[3]{q\in\mathcal F}{\rA}
 \proofstepwidestar[4]{s\in A_D}{\rA}
 \proofstepwidestar[1,2,3]{R_S(\mathsf{and}(p,q))=R_S(p)\cap R_S(q)}{\FormulaRefAuto{B48ErfAndRow}[thm]{1,2,3}}
 \proofstepwidestar[4]{s\in R_S(\mathsf{and}(p,q))\leftrightarrow(\mathsf{and}(p,q),s)\in S}{\FormulaRefAuto{B48RowMember}[thm]{4}}
 \proofstepwidestar[4]{(\mathsf{and}(p,q),s)\in S\leftrightarrow s\in R_S(\mathsf{and}(p,q))}{\FormulaRefAuto{B48IffSymmetry}[thm]{6}}
 \proofstepwidestar[1,2,3,4]{(\mathsf{and}(p,q),s)\in S\leftrightarrow s\in R_S(p)\cap R_S(q)}{\rIE{5,7}}
 \proofstepwidestar[]{s\in R_S(p)\cap R_S(q)\leftrightarrow(s\in R_S(p)\land s\in R_S(q))}{\FormulaRefAuto{x\in A\cap B\eqvdash x\in A\land x\in B}[thm]}
 \proofstepwidestar[1,2,3,4]{(\mathsf{and}(p,q),s)\in S\leftrightarrow(s\in R_S(p)\land s\in R_S(q))}{\rChain{8,9}}
 \proofstepwidestar[4]{s\in R_S(p)\leftrightarrow(p,s)\in S}{\FormulaRefAuto{B48RowMember}[thm]{4}}
 \proofstepwidestar[4]{s\in R_S(q)\leftrightarrow(q,s)\in S}{\FormulaRefAuto{B48RowMember}[thm]{4}}
 \proofstepwidestar[1,2,3,4]{(\mathsf{and}(p,q),s)\in S\leftrightarrow((p,s)\in S\land(q,s)\in S)}{\rLRS{12,\rLRS{11,10}}}
\closeproofpartwideindR

\proofpartwideindR[Punktweise Quantorenklausel]{\begin{gathered}\mathsf{Erf}(M,S),\ i\in\mathbb N,\ p\in\mathcal F,\ s\in A_D\\
 \vdash((\mathsf{ex}_i(p),s)\in S\leftrightarrow
      \exists a\in D\;((p,s[i\mapsto a])\in S))\end{gathered}}{\begin{gathered}\mathsf{Erf}(M,S),\ i\in\mathbb N,\ p\in\mathcal F,\ s\in A_D\\
 \vdash((\mathsf{ex}_i(p),s)\in S\leftrightarrow
      \exists a\in D\;((p,s[i\mapsto a])\in S))\end{gathered}}[B48PointExists]
 \proofstepwidestar[1]{\mathsf{Erf}(M,S)}{\rA}
 \proofstepwidestar[2]{i\in\mathbb N}{\rA}
 \proofstepwidestar[3]{p\in\mathcal F}{\rA}
 \proofstepwidestar[4]{s\in A_D}{\rA}
 \proofstepwidestar[1,2,3]{R_S(\mathsf{ex}_i(p))=Q_i(R_S(p))}{\FormulaRefAuto{B48ErfExistsRow}[thm]{1,2,3}}
 \proofstepwidestar[4]{s\in R_S(\mathsf{ex}_i(p))\leftrightarrow(\mathsf{ex}_i(p),s)\in S}{\FormulaRefAuto{B48RowMember}[thm]{4}}
 \proofstepwidestar[4]{(\mathsf{ex}_i(p),s)\in S\leftrightarrow s\in R_S(\mathsf{ex}_i(p))}{\FormulaRefAuto{B48IffSymmetry}[thm]{6}}
 \proofstepwidestar[1,2,3,4]{(\mathsf{ex}_i(p),s)\in S\leftrightarrow s\in Q_i(R_S(p))}{\rIE{5,7}}
 \proofstepwidestar[4]{s\in Q_i(R_S(p))\leftrightarrow\exists a\in D\;(s[i\mapsto a]\in R_S(p))}{\FormulaRefAuto{B48SeparationWithin}[thm]{4}}
 \proofstepwidestar[10]{a\in D}{\rA}
 \proofstepwidestar[2,4,10]{s[i\mapsto a]\in A_D}{\FormulaRefAuto{B48AssignmentUpdate}[thm]{4,2,10}}
 \proofstepwidestar[2,4,10]{s[i\mapsto a]\in R_S(p)\leftrightarrow(p,s[i\mapsto a])\in S}{\FormulaRefAuto{B48RowMember}[thm]{11}}
 \proofstepwidestar[2,4]{\forall a\in D\;(s[i\mapsto a]\in R_S(p)\leftrightarrow(p,s[i\mapsto a])\in S)}{\rUI{\rRI{10,12}}}
 \proofstepwidestar[2,4]{(\exists a\in D\;s[i\mapsto a]\in R_S(p))\leftrightarrow(\exists a\in D\;(p,s[i\mapsto a])\in S)}{\FormulaRefAuto{B48ExistsIff}[thm]{13}}
 \proofstepwidestar[1,2,3,4]{(\mathsf{ex}_i(p),s)\in S\leftrightarrow\exists a\in D\;((p,s[i\mapsto a])\in S)}{\rChain{8,9,14}}
\closeproofpartwideindR

\proofpartwideindR[Atomarer Code fuer Gleichheit]{i\in\mathbb N,\ j\in\mathbb N\vdash (t_{\rm e},(i,j))\in\Sigma_{\rm at}}{i\in\mathbb N,\ j\in\mathbb N\vdash (t_{\rm e},(i,j))\in\Sigma_{\rm at}}[B48EqAtomic]
 \proofstepwidestar[1]{i\in\mathbb N}{\rA}
 \proofstepwidestar[2]{j\in\mathbb N}{\rA}
 \proofstepwidestar[]{t_{\rm e}\in\mathbb N}{\FormulaRefAuto{B48TagType1}[thm]}
 \proofstepwidestar[1,2]{(t_{\rm e},(i,j))\in\Sigma}{\FormulaRefAuto{B48LabelType}[thm]{3,1,2}}
 \proofstepwidestar[]{t_{\rm e}=t_{\rm e}}{\rII}
 \proofstepwidestar[]{t_{\rm e}=t_{\rm e}\lor(t_{\rm e}=t_{\rm m}\lor(t_{\rm e}=t_{\rm b}\land i=0\land j=0))}{\rOIa{5}}
 \proofstepwidestar[1,2]{(t_{\rm e},(i,j))\in\Sigma_{\rm at}}{\FormulaRefAuto{B48FormulaConstructors}[def]{\rAI{4,6}}}
\closeproofpartwideindR

\proofpartwideindR[Atomarer Code fuer Mitgliedschaft]{i\in\mathbb N,\ j\in\mathbb N\vdash (t_{\rm m},(i,j))\in\Sigma_{\rm at}}{i\in\mathbb N,\ j\in\mathbb N\vdash (t_{\rm m},(i,j))\in\Sigma_{\rm at}}[B48MemAtomic]
 \proofstepwidestar[1]{i\in\mathbb N}{\rA}
 \proofstepwidestar[2]{j\in\mathbb N}{\rA}
 \proofstepwidestar[]{t_{\rm m}\in\mathbb N}{\FormulaRefAuto{B48TagType2}[thm]}
 \proofstepwidestar[1,2]{(t_{\rm m},(i,j))\in\Sigma}{\FormulaRefAuto{B48LabelType}[thm]{3,1,2}}
 \proofstepwidestar[]{t_{\rm m}=t_{\rm m}}{\rII}
 \proofstepwidestar[]{t_{\rm m}=t_{\rm m}\lor(t_{\rm m}=t_{\rm b}\land i=0\land j=0)}{\rOIa{5}}
 \proofstepwidestar[]{t_{\rm m}=t_{\rm e}\lor(t_{\rm m}=t_{\rm m}\lor(t_{\rm m}=t_{\rm b}\land i=0\land j=0))}{\rOIb{6}}
 \proofstepwidestar[1,2]{(t_{\rm m},(i,j))\in\Sigma_{\rm at}}{\FormulaRefAuto{B48FormulaConstructors}[def]{\rAI{4,7}}}
\closeproofpartwideindR

\proofpartwideindR[Atomarer Widerspruchscode]{(t_{\rm b},(0,0))\in\Sigma_{\rm at}}{(t_{\rm b},(0,0))\in\Sigma_{\rm at}}[B48BottomAtomic]
 \proofstepwidestar[]{t_{\rm b}\in\mathbb N}{\FormulaRefAuto{B48TagType6}[thm]}
 \proofstepwidestar[]{0\in\mathbb N}{\FormulaRefAuto{PeanoZero}[ax]}
 \proofstepwidestar[]{(t_{\rm b},(0,0))\in\Sigma}{\FormulaRefAuto{B48LabelType}[thm]{1,2,2}}
 \proofstepwidestar[]{t_{\rm b}=t_{\rm b}}{\rII}
 \proofstepwidestar[]{0=0}{\rII}
 \proofstepwidestar[]{t_{\rm b}=t_{\rm b}\land(0=0\land0=0)}{\rAI{4,\rAI{5,5}}}
 \proofstepwidestar[]{t_{\rm b}=t_{\rm m}\lor(t_{\rm b}=t_{\rm b}\land0=0\land0=0)}{\rOIb{6}}
 \proofstepwidestar[]{t_{\rm b}=t_{\rm e}\lor(t_{\rm b}=t_{\rm m}\lor(t_{\rm b}=t_{\rm b}\land0=0\land0=0))}{\rOIb{7}}
 \proofstepwidestar[]{(t_{\rm b},(0,0))\in\Sigma_{\rm at}}{\FormulaRefAuto{B48FormulaConstructors}[def]{\rAI{3,8}}}
\closeproofpartwideindR

\proofpartwideindR[Wertemenge eines Gleichheitsatoms]{i\in\mathbb N,\ j\in\mathbb N\vdash B_M((t_{\rm e},(i,j)))=\{s\in A_D\mid s(i)=s(j)\}}{i\in\mathbb N,\ j\in\mathbb N\vdash B_M((t_{\rm e},(i,j)))=\{s\in A_D\mid s(i)=s(j)\}}[B48EqBValue]
 \proofstepwidestar[1]{i\in\mathbb N}{\rA}
 \proofstepwidestar[2]{j\in\mathbb N}{\rA}
 \proofstepwidestar[]{t_{\rm e}=t_{\rm e}}{\rII}
 \proofstepwidestar[]{\begin{gathered}\mathsf{if}\bigl(t_{\rm e}=t_{\rm e},\{s\in A_D\mid s(i)=s(j)\},\\
 \BXLVIIIIf{t_{\rm e}=t_{\rm m}}{\{s\in A_D\mid(s(i),s(j))\in E\}}{\varnothing}\bigr)\\
 =\{s\in A_D\mid s(i)=s(j)\}\end{gathered}}{\FormulaRefAuto{B48IfTrue}[thm]{3}}
 \proofstepwidestar[1,2]{B_M((t_{\rm e},(i,j)))=\{s\in A_D\mid s(i)=s(j)\}}{\FormulaRefAuto{B48SemanticTerms}[def]{4}}
\closeproofpartwideindR

\proofpartwideindR[Wertemenge eines Mitgliedschaftsatoms]{i\in\mathbb N,\ j\in\mathbb N\vdash B_M((t_{\rm m},(i,j)))=\{s\in A_D\mid(s(i),s(j))\in E\}}{i\in\mathbb N,\ j\in\mathbb N\vdash B_M((t_{\rm m},(i,j)))=\{s\in A_D\mid(s(i),s(j))\in E\}}[B48MemBValue]
 \proofstepwidestar[1]{i\in\mathbb N}{\rA}
 \proofstepwidestar[2]{j\in\mathbb N}{\rA}
 \proofstepwidestar[]{t_{\rm e}\neq t_{\rm m}}{\FormulaRefAuto{B48TagNe01}[thm]}
 \proofstepwidestar[]{t_{\rm m}\neq t_{\rm e}}{\FormulaRefAuto{a\neq b\vdash b\neq a}[thm]{3}}
 \proofstepwidestar[]{t_{\rm m}=t_{\rm m}}{\rII}
 \proofstepwidestar[]{\begin{gathered}\BXLVIIIIf{t_{\rm m}=t_{\rm m}}{\{s\in A_D\mid(s(i),s(j))\in E\}}{\varnothing}\\
 =\{s\in A_D\mid(s(i),s(j))\in E\}\end{gathered}}{\FormulaRefAuto{B48IfTrue}[thm]{5}}
 \proofstepwidestar[]{\begin{gathered}\mathsf{if}\bigl(t_{\rm m}=t_{\rm e},\{s\in A_D\mid s(i)=s(j)\},\\
 \BXLVIIIIf{t_{\rm m}=t_{\rm m}}{\{s\in A_D\mid(s(i),s(j))\in E\}}{\varnothing}\bigr)\\
 =\BXLVIIIIf{t_{\rm m}=t_{\rm m}}{\{s\in A_D\mid(s(i),s(j))\in E\}}{\varnothing}\end{gathered}}{\FormulaRefAuto{B48IfFalse}[thm]{4}}
 \proofstepwidestar[1,2]{\begin{gathered}B_M((t_{\rm m},(i,j)))\\
 =\BXLVIIIIf{t_{\rm m}=t_{\rm m}}{\{s\in A_D\mid(s(i),s(j))\in E\}}{\varnothing}\end{gathered}}{\FormulaRefAuto{B48SemanticTerms}[def]{7}}
 \proofstepwidestar[1,2]{B_M((t_{\rm m},(i,j)))=\{s\in A_D\mid(s(i),s(j))\in E\}}{\rIE{6,8}}
\closeproofpartwideindR

\proofpartwideindR[Wertemenge des Widerspruchscodes]{B_M((t_{\rm b},(0,0)))=\varnothing}{B_M((t_{\rm b},(0,0)))=\varnothing}[B48BottomBValue]
 \proofstepwidestar[]{t_{\rm e}\neq t_{\rm b}}{\FormulaRefAuto{B48TagNe05}[thm]}
 \proofstepwidestar[]{t_{\rm m}\neq t_{\rm b}}{\FormulaRefAuto{B48TagNe15}[thm]}
 \proofstepwidestar[]{t_{\rm b}\neq t_{\rm e}}{\FormulaRefAuto{a\neq b\vdash b\neq a}[thm]{1}}
 \proofstepwidestar[]{t_{\rm b}\neq t_{\rm m}}{\FormulaRefAuto{a\neq b\vdash b\neq a}[thm]{2}}
 \proofstepwidestar[]{\BXLVIIIIf{t_{\rm b}=t_{\rm m}}{\{s\in A_D\mid(s(0),s(0))\in E\}}{\varnothing}=\varnothing}{\FormulaRefAuto{B48IfFalse}[thm]{4}}
 \proofstepwidestar[]{\begin{gathered}\mathsf{if}\bigl(t_{\rm b}=t_{\rm e},\{s\in A_D\mid s(0)=s(0)\},\\
 \BXLVIIIIf{t_{\rm b}=t_{\rm m}}{\{s\in A_D\mid(s(0),s(0))\in E\}}{\varnothing}\bigr)\\
 =\BXLVIIIIf{t_{\rm b}=t_{\rm m}}{\{s\in A_D\mid(s(0),s(0))\in E\}}{\varnothing}\end{gathered}}{\FormulaRefAuto{B48IfFalse}[thm]{3}}
 \proofstepwidestar[]{B_M((t_{\rm b},(0,0)))=\BXLVIIIIf{t_{\rm b}=t_{\rm m}}{\{s\in A_D\mid(s(0),s(0))\in E\}}{\varnothing}}{\FormulaRefAuto{B48SemanticTerms}[def]{6}}
 \proofstepwidestar[]{B_M((t_{\rm b},(0,0)))=\varnothing}{\rIE{5,7}}
\closeproofpartwideindR

\proofpartwideindR[Punktweise Klausel fuer Gleichheit]{\begin{gathered}\mathsf{Erf}(M,S),\ i\in\mathbb N,\ j\in\mathbb N,\ s\in A_D\\\vdash((\mathsf e(i,j),s)\in S\leftrightarrow s(i)=s(j))\end{gathered}}{\begin{gathered}\mathsf{Erf}(M,S),\ i\in\mathbb N,\ j\in\mathbb N,\ s\in A_D\\\vdash((\mathsf e(i,j),s)\in S\leftrightarrow s(i)=s(j))\end{gathered}}[B48PointEq]
 \proofstepwidestar[1]{\mathsf{Erf}(M,S)}{\rA}
 \proofstepwidestar[2]{i\in\mathbb N}{\rA}
 \proofstepwidestar[3]{j\in\mathbb N}{\rA}
 \proofstepwidestar[4]{s\in A_D}{\rA}
 \proofstepwidestar[2,3]{(t_{\rm e},(i,j))\in\Sigma_{\rm at}}{\FormulaRefAuto{B48EqAtomic}[thm]{2,3}}
 \proofstepwidestar[1,2,3,4]{(\mathsf e(i,j),s)\in S\leftrightarrow s\in B_M((t_{\rm e},(i,j)))}{\FormulaRefAuto{B48PointAtom}[thm]{1,5,4}}
 \proofstepwidestar[2,3]{B_M((t_{\rm e},(i,j)))=\{s\in A_D\mid s(i)=s(j)\}}{\FormulaRefAuto{B48EqBValue}[thm]{2,3}}
 \proofstepwidestar[1,2,3,4]{(\mathsf e(i,j),s)\in S\leftrightarrow s\in\{u\in A_D\mid u(i)=u(j)\}}{\rIE{7,6}}
 \proofstepwidestar[4]{s\in\{u\in A_D\mid u(i)=u(j)\}\leftrightarrow s(i)=s(j)}{\FormulaRefAuto{B48SeparationWithin}[thm]{4}}
 \proofstepwidestar[1,2,3,4]{(\mathsf e(i,j),s)\in S\leftrightarrow s(i)=s(j)}{\rChain{8,9}}
\closeproofpartwideindR

\proofpartwideindR[Punktweise Klausel fuer Mitgliedschaft]{\begin{gathered}\mathsf{Erf}(M,S),\ i\in\mathbb N,\ j\in\mathbb N,\ s\in A_D\\\vdash((\mathsf m(i,j),s)\in S\leftrightarrow (s(i),s(j))\in E)\end{gathered}}{\begin{gathered}\mathsf{Erf}(M,S),\ i\in\mathbb N,\ j\in\mathbb N,\ s\in A_D\\\vdash((\mathsf m(i,j),s)\in S\leftrightarrow (s(i),s(j))\in E)\end{gathered}}[B48PointMem]
 \proofstepwidestar[1]{\mathsf{Erf}(M,S)}{\rA}
 \proofstepwidestar[2]{i\in\mathbb N}{\rA}
 \proofstepwidestar[3]{j\in\mathbb N}{\rA}
 \proofstepwidestar[4]{s\in A_D}{\rA}
 \proofstepwidestar[2,3]{(t_{\rm m},(i,j))\in\Sigma_{\rm at}}{\FormulaRefAuto{B48MemAtomic}[thm]{2,3}}
 \proofstepwidestar[1,2,3,4]{(\mathsf m(i,j),s)\in S\leftrightarrow s\in B_M((t_{\rm m},(i,j)))}{\FormulaRefAuto{B48PointAtom}[thm]{1,5,4}}
 \proofstepwidestar[2,3]{B_M((t_{\rm m},(i,j)))=\{s\in A_D\mid (s(i),s(j))\in E\}}{\FormulaRefAuto{B48MemBValue}[thm]{2,3}}
 \proofstepwidestar[1,2,3,4]{(\mathsf m(i,j),s)\in S\leftrightarrow s\in\{u\in A_D\mid (u(i),u(j))\in E\}}{\rIE{7,6}}
 \proofstepwidestar[4]{s\in\{u\in A_D\mid (u(i),u(j))\in E\}\leftrightarrow (s(i),s(j))\in E}{\FormulaRefAuto{B48SeparationWithin}[thm]{4}}
 \proofstepwidestar[1,2,3,4]{(\mathsf m(i,j),s)\in S\leftrightarrow (s(i),s(j))\in E}{\rChain{8,9}}
\closeproofpartwideindR

\proofpartwideindR[Keine Belegung erfuellt den Widerspruchscode]{\mathsf{Erf}(M,S),\ s\in A_D\vdash\neg((\mathsf b,s)\in S)}{\mathsf{Erf}(M,S),\ s\in A_D\vdash\neg((\mathsf b,s)\in S)}[B48PointBottom]
 \proofstepwidestar[1]{\mathsf{Erf}(M,S)}{\rA}
 \proofstepwidestar[2]{s\in A_D}{\rA}
 \proofstepwidestar[]{(t_{\rm b},(0,0))\in\Sigma_{\rm at}}{\FormulaRefAuto{B48BottomAtomic}[thm]}
 \proofstepwidestar[1,2]{(\mathsf b,s)\in S\leftrightarrow s\in B_M((t_{\rm b},(0,0)))}{\FormulaRefAuto{B48PointAtom}[thm]{1,3,2}}
 \proofstepwidestar[]{B_M((t_{\rm b},(0,0)))=\varnothing}{\FormulaRefAuto{B48BottomBValue}[thm]}
 \proofstepwidestar[1,2]{(\mathsf b,s)\in S\leftrightarrow s\in\varnothing}{\rIE{5,4}}
 \proofstepwidestar[7]{(\mathsf b,s)\in S}{\rA}
 \proofstepwidestar[1,2,7]{s\in\varnothing}{\rLREa{6,7}}
 \proofstepwidestar[]{s\notin\varnothing}{\FormulaRefAuto{x \not\in \varnothing}[thm]}
 \proofstepwidestar[1,2,7]{\bot}{\rBI{8,9}}
 \proofstepwidestar[1,2]{\neg((\mathsf b,s)\in S)}{\rCI{7,10}}
\closeproofpartwideindR


\proofpartwide{Schluss}
 \proofstepwidestar[1]{\mathsf{Str}(M)}{\rA}
 \proofstepwidestar[]{\exists e\;\mathsf{Rec}(e,M)}
   {\FormulaRefAuto{B48TreeEvaluationExists}[thm]}
 \proofstepwidestar[3]{\mathsf{Rec}(e,M)}{\rA}
 \proofstepwidestar[1,3]{\mathsf{Erf}(M,S_e)}
   {\FormulaRefAuto{B48CandidateSatisfies}[thm]{1,3}}
 \proofstepwidestar[1,3]{\exists S\;\mathsf{Erf}(M,S)}{\rEI{4}}
 \proofstepwidestar[1]{\exists S\;\mathsf{Erf}(M,S)}{\rEE{2,3:5}}
 \proofstepwidestar[7]{\mathsf{Erf}(M,S)}{\rA}
 \proofstepwidestar[8]{\mathsf{Erf}(M,T)}{\rA}
 \proofstepwidestar[7,8]{S=T}
   {\FormulaRefAuto{B48SatisfactionUnique}[thm]{7,8}}
 \proofstepwidestar[1]{\exists!S\;\mathsf{Erf}(M,S)}{\UEI{6,7,8,9}}
\closeproofpartwide
\end{tabproofsplitwide}



\subsection*{Die nun gerechtfertigte Schreibweise fuer Erfuellung}

Nach \FormulaRefAuto{B48SatisfactionExists}[thm] ist die folgende
Kennzeichnung wohldefiniert. Die rekursiven Klauseln werden hier als
gleichwertige definierende Gleichungen dieser eindeutigen Relation
mitgefuehrt, nicht als neue Existenzaxiome.
Die punktweisen Klauseln sind bereits hergeleitet.
Fuer Gleichheit und Elementrelation gelten\\
\FormulaRefAuto{B48PointEq}[thm] und \FormulaRefAuto{B48PointMem}[thm].

Fuer Negation und Konjunktion gelten\\
\FormulaRefAuto{B48PointNeg}[thm] und \FormulaRefAuto{B48PointAnd}[thm].

Fuer Existenzquantor und Widerspruch gelten\\
\FormulaRefAuto{B48PointExists}[thm] und \FormulaRefAuto{B48PointBottom}[thm].

\FormulaDefDeltaK[Erfuellung in einer Mengenstruktur]
{\begin{aligned}
 \operatorname{Sat}(M)&\coloneq\iota S\;\mathsf{Erf}(M,S),\\
 M,s\models p&\coloneq(p,s)\in\operatorname{Sat}(M),\\
 M,s\models v_i=v_j&\coloneq s(v_i)=s(v_j),\\
 M,s\models v_i\in v_j&\coloneq(s(v_i),s(v_j))\in E,\\
 M,s\models\neg p&\coloneq\neg(M,s\models p),\\
 M,s\models p\land q&\coloneq(M,s\models p)\land(M,s\models q),\\
 M,s\models\exists v_i\,p&\coloneq
          \exists a\in D\;(M,s[v_i\mapsto a]\models p),\\
 M,s\models\bot&\coloneq\bot.
\end{aligned}}
{B48Satisfaction}{
 \DeltaRow{Struktur}{M=(D,E),\quad D\neq\varnothing,\quad E\subseteq D\times D}
 \DeltaRow{Belegung}{s\colon\operatorname{Var}\to D}
 \DeltaRow{Formelcodes}{p,q\in\mathcal F}
 \DeltaRow{Variablencodes}{v_i=i,\quad v_j=j,\quad i,j\in\mathbb N}
}

Die punktweise Schreibweise ist damit auf eine konstruierte Menge
zurueckgefuehrt. Dies beweist noch nicht die Korrektheit beliebiger
syntaktischer Herleitungen: Dafuer werden spaeter zusaetzlich
Substitution, Koinzidenz und Herleitungsinduktion benoetigt.


% Einbindung nach 02-erfuellungsrelation.tex: B48FormulaInTree steht bereit.
% Die Positionsnotation ist in Band 27 lokal; providecommand ermoeglicht
% auch die eigenstaendige Uebersetzung von Band 48.
\providecommand{\TreePositions}[2]{\operatorname{Pos}_{#1}(#2)}
\providecommand{\AddressEdges}[1]{E_{#1}^{\mathrm{adr}}}
\providecommand{\AddressLeft}[1]{\lambda_{#1}^{\mathrm{adr}}}
\providecommand{\AddressRight}[1]{\rho_{#1}^{\mathrm{adr}}}

\subsection{Die endliche Baumstruktur der Formelcodes}

Die Formelmenge ist eine durch ihre Konstruktoren abgeschlossene
Teilmenge des Baumtraegers aus Band~27. Ihr Induktionsprinzip
\FormulaRefAuto{B48FormulaInduction}[thm] beruht auf dieser kleinsten
abgeschlossenen Menge. Fuer jeden einzelnen Code steht zusaetzlich
die endliche Graphstruktur seiner Vorkommen zur Verfuegung.

\FormulaDefDeltaK[Positionsdaten eines Formelcodes]
{\begin{aligned}
 P_p&\coloneq\TreePositions{\Sigma}{p},&
 E_p&\coloneq\AddressEdges{P_p},\\
 \lambda_p&\coloneq\AddressLeft{P_p},&
 \rho_p&\coloneq\AddressRight{P_p},\\
 \varepsilon_2&\coloneq\EmptyWord{\{0,1\}}.
\end{aligned}}
{B48FormulaPositionData}{\DeltaRow{Baumcode}{p\in\mathcal T}}

Die rechte Seite verwendet die bereits definierten Positionsmengen
und Adressgraphen aus \FormulaRefAuto{TreePositionSetDef}[def] und
\FormulaRefAuto{BinaryAddressGraphDef}[def]. Die neuen Kurzzeichen
erleichtern nur den Zugriff auf diese Daten.

\FormulaThmDeltaK[Formelcodes tragen endliche geordnete volle Binaerbaeume]
{\begin{aligned}
 p\in\mathcal F\vdash{}&\Finite{P_p}\\[-2pt]
 &\land\OrderedFullBinaryTreeStruct{
     P_p,E_p,\varepsilon_2,\lambda_p,\rho_p}
\end{aligned}}
{B48FormulaPositionTree}{\DeltaRow{Formelcode}{p}}
\begin{tabproofwide}
 \proofstepwidestar[1]{p\in\mathcal F}{\rA}
 \proofstepwidestar[1]{p\in\mathcal T}{%
  \FormulaRefAuto{B48FormulaInTree}[thm]{1}}
 \proofstepwidestar[1]{p\in\BracketTreeSet{\Sigma}}{%
  \FormulaRefAuto{B48CodeCarrier}[def]{2}}
 \proofstepwidestar[1]{\Finite{\TreePositions{\Sigma}{p}}}{%
  \FormulaRefAuto{BracketCodePositionSetFinite}[thm]{3}}
 \proofstepwidestar[1]{\begin{aligned}[t]
  &\mathsf{GVollBinBaum}\!\bigl(\TreePositions{\Sigma}{p},\\[-2pt]
  &\quad\AddressEdges{\TreePositions{\Sigma}{p}},
    \EmptyWord{\{0,1\}},\\[-2pt]
  &\quad\AddressLeft{\TreePositions{\Sigma}{p}},
    \AddressRight{\TreePositions{\Sigma}{p}}\bigr)
 \end{aligned}}{%
  \FormulaRefAuto{BracketCodePositionGraphOrderedFullBinaryTree}[thm]{3}}
 \proofstepwidestar[1]{\begin{aligned}[t]
   &\Finite{P_p}\\[-2pt]
   &\land\OrderedFullBinaryTreeStruct{
      P_p,E_p,\varepsilon_2,\lambda_p,\rho_p}
 \end{aligned}}{%
  \FormulaRefAuto{B48FormulaPositionData}[def]{\rAI{4,5}}}
\end{tabproofwide}

Die Knoten sind Vorkommen im Code. Derselbe Teilcode kann daher an
mehreren verschiedenen Positionen erscheinen. Der volle Binaerbaum
enthaelt auch die bereits eingefuehrten Hilfsmarken; ein einzelner
logischer Junktor muss deshalb nicht einem einzelnen Binaerknoten
entsprechen. Die Aussagen ueber Eltern, Kinder und Wurzelwege aus
Band~26 gelten fuer diesen Codegraphen. Auf seinen Teilmengen sind
die Saetze ueber endliche Mengen aus Band~20 anwendbar. Damit stehen
sowohl die Formelinduktion ueber Konstruktoren als auch die endliche
Menge ihrer jeweiligen Vorkommen auf einer bereits bewiesenen Grundlage.


% Lokale semantische Schlussformen. Keine neuen Mengenaxiome.
\subsection*{Einzeln referenzierte semantische Schlussformen}

In diesem Abschnitt wird die durch
\FormulaRefAuto{B48SatisfactionExists}[thm] konstruierte und eindeutig
bestimmte Erfuellungsrelation verwendet. Die definierenden Klauseln
aus \FormulaRefAuto{B48Satisfaction}[def] sind durch die dortigen
Proofparts gerechtfertigt. Die folgenden Saetze zeigen, welche
lokalen Schluesse aus diesen Klauseln folgen.
Alle Zeilen sind Aussagen der untersuchenden Mengenlehre.
Insbesondere sind \(M,s\models\neg\varphi\) und
\(\neg(M,s\models\varphi)\) verschiedene Schreibweisen, deren
Zusammenhang jeweils durch die angegebene Definition begruendet wird.

\FormulaThmDeltaK[Keine Belegung erfuellt die Widerspruchsformel]
{\neg(M,s\models\bot)}
{B48SatNotBottom}{
 \DeltaRow{Struktur und Belegung}{M=(D,E),\quad s:\operatorname{Var}\to D}
}
\begin{tabproof}
 \proofstep{1}{M,s\models\bot}{\rA}
 \proofstep{1}{\bot}{\FormulaRefAuto{B48Satisfaction}{1}}
 \proofstep{}{\neg(M,s\models\bot)}{\rCI{1,2}}
\end{tabproof}

\FormulaThmDeltaK[Semantische Konjunktionseinfuehrung]
{M,s\models\varphi,\ M,s\models\psi
 \ \vdash\ M,s\models(\varphi\land\psi)}
{B48SatAndIntro}{
 \DeltaRow{Formeln}{\varphi,\psi}
 \DeltaRow{Belegung}{s:\operatorname{Var}\to D}
}
\begin{tabproof}
 \proofstep{1}{M,s\models\varphi}{\rA}
 \proofstep{2}{M,s\models\psi}{\rA}
 \proofstep{1,2}{(M,s\models\varphi)\land(M,s\models\psi)}
 {\rAI{1,2}}
 \proofstep{1,2}{M,s\models(\varphi\land\psi)}
 {\FormulaRefAuto{B48Satisfaction}{3}}
\end{tabproof}

\FormulaThmDeltaK[Erste semantische Konjunktionsbeseitigung]
{M,s\models(\varphi\land\psi)\ \vdash\ M,s\models\varphi}
{B48SatAndLeft}{
 \DeltaRow{Formeln}{\varphi,\psi}
 \DeltaRow{Belegung}{s:\operatorname{Var}\to D}
}
\begin{tabproof}
 \proofstep{1}{M,s\models(\varphi\land\psi)}{\rA}
 \proofstep{1}{(M,s\models\varphi)\land(M,s\models\psi)}
 {\FormulaRefAuto{B48Satisfaction}{1}}
 \proofstep{1}{M,s\models\varphi}{\rAEa{2}}
\end{tabproof}

\FormulaThmDeltaK[Zweite semantische Konjunktionsbeseitigung]
{M,s\models(\varphi\land\psi)\ \vdash\ M,s\models\psi}
{B48SatAndRight}{
 \DeltaRow{Formeln}{\varphi,\psi}
 \DeltaRow{Belegung}{s:\operatorname{Var}\to D}
}
\begin{tabproof}
 \proofstep{1}{M,s\models(\varphi\land\psi)}{\rA}
 \proofstep{1}{(M,s\models\varphi)\land(M,s\models\psi)}
 {\FormulaRefAuto{B48Satisfaction}{1}}
 \proofstep{1}{M,s\models\psi}{\rAEb{2}}
\end{tabproof}

\FormulaThmDeltaK[Unvertraeglichkeit einer Formel mit ihrer Negation]
{M,s\models\varphi,\ M,s\models\neg\varphi\ \vdash\ \bot}
{B48SatContradiction}{
 \DeltaRow{Formel}{\varphi}
 \DeltaRow{Belegung}{s:\operatorname{Var}\to D}
}
\begin{tabproof}
 \proofstep{1}{M,s\models\varphi}{\rA}
 \proofstep{2}{M,s\models\neg\varphi}{\rA}
 \proofstep{2}{\neg(M,s\models\varphi)}
 {\FormulaRefAuto{B48Satisfaction}{2}}
 \proofstep{1,2}{\bot}{\rBI{1,3}}
\end{tabproof}

\FormulaThmDeltaK[Semantischer Negationsschluss]
{\bigl((M,s\models\varphi)\to\bot\bigr)
 \ \vdash\ M,s\models\neg\varphi}
{B48SatNegationIntro}{
 \DeltaRow{Formel}{\varphi}
 \DeltaRow{Belegung}{s:\operatorname{Var}\to D}
}
\begin{tabproof}
 \proofstep{1}{(M,s\models\varphi)\to\bot}{\rA}
 \proofstep{2}{M,s\models\varphi}{\rA}
 \proofstep{1,2}{\bot}{\rRE{1,2}}
 \proofstep{1}{\neg(M,s\models\varphi)}{\rCI{2,3}}
 \proofstep{1}{M,s\models\neg\varphi}
 {\FormulaRefAuto{B48Satisfaction}{4}}
\end{tabproof}

\FormulaThmDeltaK[Semantischer klassischer Widerspruchsschluss]
{\bigl((M,s\models\neg\varphi)\to\bot\bigr)
 \ \vdash\ M,s\models\varphi}
{B48SatClassical}{
 \DeltaRow{Formel}{\varphi}
 \DeltaRow{Belegung}{s:\operatorname{Var}\to D}
}
\begin{tabproof}
 \proofstep{1}{(M,s\models\neg\varphi)\to\bot}{\rA}
 \proofstep{2}{\neg(M,s\models\varphi)}{\rA}
 \proofstep{2}{M,s\models\neg\varphi}
 {\FormulaRefAuto{B48Satisfaction}{2}}
 \proofstep{1,2}{\bot}{\rRE{1,3}}
 \proofstep{1}{M,s\models\varphi}{\rCE{2,4}}
\end{tabproof}

\FormulaThmDeltaK[Semantische Existenzeinfuehrung mit bestimmtem Zeugen]
{a\in D,\ M,s[x\mapsto a]\models\varphi
 \ \vdash\ M,s\models\exists x\,\varphi}
{B48SatExistsIntro}{
 \DeltaRow{Formel und Variable}{\varphi,\quad x\in\operatorname{Var}}
 \DeltaRow{Belegung}{s:\operatorname{Var}\to D}
}
\begin{tabproof}
 \proofstep{1}{a\in D}{\rA}
 \proofstep{2}{M,s[x\mapsto a]\models\varphi}{\rA}
 \proofstep{1,2}{a\in D\land(M,s[x\mapsto a]\models\varphi)}
 {\rAI{1,2}}
 \proofstep{1,2}{\exists b\,
   (b\in D\land(M,s[x\mapsto b]\models\varphi))}
 {\rEI{3}}
 \proofstep{1,2}{M,s\models\exists x\,\varphi}
 {\FormulaRefAuto{B48Satisfaction}{4}}
\end{tabproof}

\FormulaThmDeltaK[Semantische Existenzbeseitigung mit Zeugenentladung]
{\begin{gathered}
 M,s\models\exists x\,\varphi,\\
 \forall a\,(a\in D\to
       ((M,s[x\mapsto a]\models\varphi)\to Q))
 \ \vdash\ Q
 \end{gathered}}
{B48SatExistsElim}{
 \DeltaRow{Formel und Variable}{\varphi,\quad x\in\operatorname{Var}}
 \DeltaRow{Belegung}{s:\operatorname{Var}\to D}
 \DeltaRow{Aussage der untersuchenden Theorie}{Q}
 \DeltaRow{Frischer Zeuge}{b}
}
Der Zeuge \(b\) kommt in \(Q,M,D,s,\varphi\) und in den ausserhalb
des Zeugenbeweises offenen Voraussetzungen nicht frei vor.
\begin{tabproof}
 \proofstep{1}{M,s\models\exists x\,\varphi}{\rA}
 \proofstep{2}{\forall a\,(a\in D\to
      ((M,s[x\mapsto a]\models\varphi)\to Q))}{\rA}
 \proofstep{1}{\exists a\,
   (a\in D\land(M,s[x\mapsto a]\models\varphi))}
 {\FormulaRefAuto{B48Satisfaction}{1}}
 \proofstep{4}{b\in D\land(M,s[x\mapsto b]\models\varphi)}{\rA}
 \proofstep{4}{b\in D}{\rAEa{4}}
 \proofstep{4}{M,s[x\mapsto b]\models\varphi}{\rAEb{4}}
 \proofstep{2}{b\in D\to
      ((M,s[x\mapsto b]\models\varphi)\to Q)}{\rUE{2}}
 \proofstep{2,4}{(M,s[x\mapsto b]\models\varphi)\to Q}
 {\rRE{7,5}}
 \proofstep{2,4}{Q}{\rRE{8,6}}
 \proofstep{1,2}{Q}{\rEE{3,4:9}}
\end{tabproof}

\noindent
Die letzten beiden Saetze benutzen ausdruecklich die veraenderte
Belegung \(s[x\mapsto a]\). Daraus folgt noch nicht ohne Weiteres
die Korrektheit der syntaktischen Einsetzung \(\varphi[t/x]\).
Die folgenden Abschnitte beweisen dafuer das Koinzidenz- und
Substitutionslemma. Erst danach werden diese lokalen Schlussformen
auf syntaktische Herleitungen mit offenen Annahmen uebertragen.

% Freie Variablen werden aus der vorhandenen Baumrekursion gewonnen.
\subsection{Freie Variablen und Koinzidenz}

In diesem Abschnitt sind \(M=(D,E)\) eine Mengenstruktur und
\(s,t\in A_D\). Variablen sind ihre Codes in \(\mathbb N\).
Alle Aussagen ueber Erfuellung gehoeren zur untersuchenden Theorie.
Insbesondere wird keine Eigenschaft der syntaktischen Beweisbarkeit
vorausgesetzt.

\FormulaDefDeltaK[Operatoren fuer freie Variablen]
{\begin{aligned}
 I(p)&\coloneq\{i\in\mathbb N\mid p=\mathsf q_i\},\\
 L((k,(i,j)))&\coloneq
   \BXLVIIIIf{k=t_{\rm e}\lor k=t_{\rm m}}{\{i,j\}}{\varnothing},\\
 Y&\coloneq\mathcal T\times\powerset(\mathbb N),\\
 f(c)&\coloneq(\mathsf l(c),L(c)),\\
 g((p,U),(q,V))&\coloneq
   (\mathsf n(p,q),U\cup(V\setminus I(p))).
\end{aligned}}
{B48FVOperators}{
 \DeltaRow{Bereiche}{c\in\Sigma,\quad p,q\in\mathcal T}
 \DeltaRow{Wertemengen}{U,V\subseteq\mathbb N}
}
Hier und bei den folgenden rekursiven Konstruktionen bezeichnet eine
Funktionsgleichung \(f(c)\coloneq u(c)\) den beschraenkten Graphen
\(\{(c,z)\in\Sigma\times Y\mid z=u(c)\}\).
Seine Funktionseigenschaft folgt aus
\FormulaRefAuto{TypedTermGraphFunction}[thm], sobald die Werte typisiert
sind; seine Werte liefert \FormulaRefAuto{TypedTermGraphValues}[thm].
Die Hilfsbaeume \(\mathsf n_0,\mathsf c_0,\mathsf q_i\) tragen keine
freien Variablen.

\FormulaThmDeltaK[Konstruktion der Freivariablenfunktion]
{\begin{gathered}
 \exists!V\colon\mathcal T\to\powerset(\mathbb N)\quad
 V(\mathsf l(c))=L(c),\\
 V(\mathsf n(p,q))=V(p)\cup(V(q)\setminus I(p))
\end{gathered}}
{B48FVExists}{
 \DeltaRow{Universell in den Gleichungen}{c\in\Sigma,\quad p,q\in\mathcal T}
}
\begin{tabproofsplitwide}
 \proofpartwide{Typisierung und Rekursion}
 \proofstepwidestar[]{L(c)\subseteq\mathbb N}
  {\FormulaRefAuto{B48FVOperators}[def]}
 \proofstepwidestar[]{U\cup(V\setminus I(p))\subseteq\mathbb N}
  {\FormulaRefAuto{B48FVOperators}[def]}
 \proofstepwidestar[]{f\colon\Sigma\to Y}
  {\FormulaRefAuto{TypedTermGraphFunction}[thm]{1},
   \FormulaRefAuto{B48LeafType}[thm]}
 \proofstepwidestar[]{g\colon Y\times Y\to Y}
  {\FormulaRefAuto{TypedTermGraphFunction}[thm]{2},
   \FormulaRefAuto{B48NodeType}[thm]}
 \proofstepwidestar[]{\begin{gathered}
  \exists!e\colon\mathcal T\to Y\quad
  e(\mathsf l(c))=f(c),\\
  e(\mathsf n(p,q))=g(e(p),e(q))
 \end{gathered}}
  {\FormulaRefAuto{FullPlanarBinaryTreeRecursion}[thm]{3,4}}
 \closeproofpartwide
\end{tabproofsplitwide}
Fixiere die eindeutig bestimmte Auswertung \(e\).
Die erste Komponente ist der ausgewertete Baum selbst. Dies wird
ausdruecklich durch Strukturinduktion gezeigt. Im folgenden Teilbeweis
steht \(H(p)\) fuer \(\exists U\subseteq\mathbb N\;(e(p)=(p,U))\).
\begin{tabproofsplitwide}
 \proofpartwide{Erhaltung des Baumcodes}
 \proofstepwidestar[]{e(\mathsf l(c))=(\mathsf l(c),L(c))}
  {\FormulaRefAuto{B48FVOperators}[def],
   \FormulaRefAuto{FullPlanarBinaryTreeRecursion}[thm]}
 \proofstepwidestar[]{H(\mathsf l(c))}{\rEI{1}}
 \proofstepwidestar[3]{H(p)\land H(q)}{\rA}
 \proofstepwidestar[4]{e(p)=(p,U)\land e(q)=(q,V)}{\rA}
 \proofstepwidestar[4]{\begin{aligned}
 e(\mathsf n(p,q))={}&
 (\mathsf n(p,q),U\cup(V\setminus I(p)))
 \end{aligned}}
  {\FormulaRefAuto{B48FVOperators}[def]{4},
   \FormulaRefAuto{FullPlanarBinaryTreeRecursion}[thm]}
 \proofstepwidestar[4]{H(\mathsf n(p,q))}{\rEI{5}}
 \proofstepwidestar[3]{H(\mathsf n(p,q))}{\rEEStar{3,4:6}}
 \proofstepwidestar[]{H(p)\land H(q)\to H(\mathsf n(p,q))}
  {\rRI{3,7}}
 \proofstepwidestar[]{\forall p\in\mathcal T\;H(p)}
  {\FormulaRefAuto{FullPlanarBinaryTreeStructuralInduction}[thm]
   {\rUI{2},\rUI{8}}}
 \closeproofpartwide
\end{tabproofsplitwide}
In den Zeugenzeilen sind \(U,V\subseteq\mathbb N\) die durch \(H\)
beschraenkten Zeugen. Zwei solche Werte sind gleich, denn aus
\((p,U)=(p,V)\) folgt \(U=V\) durch
\FormulaRefAuto{(a,b)=(c,d)\eqvdash a=c\land b=d}[thm].
Daher ist
\[
 V=\{(p,U)\in\mathcal T\times\powerset(\mathbb N)\mid e(p)=(p,U)\}
\]
nach \FormulaRefAuto{UniqueValuedGraphFunction}[thm] eine Funktion.
Einsetzen von \(e(p)=(p,V(p))\) und \(e(q)=(q,V(q))\) in die beiden
Rekursionsgleichungen liefert die behaupteten Gleichungen.
Eindeutigkeit folgt ebenfalls durch Bauminduktion: Zwei Kandidaten
stimmen an jedem Blatt durch \(L(c)\), und am Knoten nach Einsetzen
der beiden Induktionsgleichheiten ueberein. Funktionsextensionalitaet
\FormulaRefAuto{FunctionExtensionality}[thm] liefert ihre Gleichheit.
Damit ist die folgende Kennzeichnung gerechtfertigt.

\FormulaDefDeltaK[Freie Variablen eines Formelcodes]
{\operatorname{FV}(p)\coloneq V(p)}
{B48FreeVariables}{
 \DeltaRow{Funktion}{V\text{ aus }\FormulaRefAuto{B48FVExists}[thm]}
 \DeltaRow{Baumcode}{p\in\mathcal T}
}
Die Einschraenkung auf \(\mathcal F\) ist die Freivariablenfunktion
der Objektformeln. Die Auswertung auf Hilfsbaeumen bleibt lediglich
ein Hilfsmittel der Konstruktion.

\FormulaThmDeltaK[Rekursionsklauseln fuer freie Variablen]
{\begin{aligned}
 \operatorname{FV}(\mathsf e(i,j))&=\{i,j\},&
 \operatorname{FV}(\mathsf m(i,j))&=\{i,j\},\\
 \operatorname{FV}(\mathsf b)&=\varnothing,&
 \operatorname{FV}(\mathsf{neg}(p))&=\operatorname{FV}(p),\\
 \operatorname{FV}(\mathsf{and}(p,q))
  &=\operatorname{FV}(p)\cup\operatorname{FV}(q),\\
 \operatorname{FV}(\mathsf{ex}_i(p))
  &=\operatorname{FV}(p)\setminus\{i\}.
\end{aligned}}
{B48FVClauses}{
 \DeltaRow{Formelcodes}{p,q\in\mathcal F}
 \DeltaRow{Indizes}{i,j\in\mathbb N}
}
\begin{tabproof}
 \proofstep{}{\operatorname{FV}(\mathsf e(i,j))=\{i,j\}}
  {\FormulaRefAuto{B48FVExists}[thm],
   \FormulaRefAuto{B48FVOperators}[def]}
 \proofstep{}{\operatorname{FV}(\mathsf m(i,j))=\{i,j\}}
  {\FormulaRefAuto{B48FVExists}[thm],
   \FormulaRefAuto{B48FVOperators}[def]}
 \proofstep{}{\operatorname{FV}(\mathsf b)=\varnothing}
  {\FormulaRefAuto{B48FVExists}[thm],
   \FormulaRefAuto{B48TagNe05}[thm],
   \FormulaRefAuto{B48TagNe15}[thm]}
 \proofstep{}{\begin{gathered}
  \operatorname{FV}(\mathsf n_0)=
  \operatorname{FV}(\mathsf c_0)=
  \operatorname{FV}(\mathsf q_i)=\varnothing
 \end{gathered}}
  {\FormulaRefAuto{B48FVOperators}[def],
   \FormulaRefAuto{B48FVExists}[thm]}
 \proofstep{}{I(\mathsf n_0)=I(\mathsf c_0)=\varnothing}
  {\FormulaRefAuto{B48QNeNeg}[thm],
   \FormulaRefAuto{B48QNeAnd}[thm]}
 \proofstep{}{I(\mathsf n(\mathsf c_0,p))=\varnothing}
  {\FormulaRefAuto{B48LeafNotNode}[thm]}
 \proofstep{}{I(\mathsf q_i)=\{i\}}
  {\FormulaRefAuto{B48QuantifierRootInjective}[thm],\rII}
 \proofstep{}{\operatorname{FV}(\mathsf{neg}(p))=\operatorname{FV}(p)}
  {\FormulaRefAuto{B48FVExists}[thm]{4,5}}
 \proofstep{}{\operatorname{FV}(\mathsf n(\mathsf c_0,p))
                  =\operatorname{FV}(p)}
  {\FormulaRefAuto{B48FVExists}[thm]{4,5}}
 \proofstep{}{\operatorname{FV}(\mathsf{and}(p,q))
                  =\operatorname{FV}(p)\cup\operatorname{FV}(q)}
  {\FormulaRefAuto{B48FVExists}[thm]{9,6}}
 \proofstep{}{\operatorname{FV}(\mathsf{ex}_i(p))
                  =\operatorname{FV}(p)\setminus\{i\}}
  {\FormulaRefAuto{B48FVExists}[thm]{4,7}}
\end{tabproof}


\FormulaThmDeltaK[Gleichheit der Quantorenmarken]
{\mathsf q_i=\mathsf q_j\leftrightarrow i=j}
{B48QuantifierRootEquality}{\DeltaRow{Indizes}{i,j\in\mathbb N}}
\begin{tabproof}
 \proofstep{1}{\mathsf q_i=\mathsf q_j}{\rA}
 \proofstep{1}{i=j}{\FormulaRefAuto{B48QuantifierRootInjective}[thm]{1}}
 \proofstep{}{(\mathsf q_i=\mathsf q_j)\to i=j}{\rRI{1,2}}
 \proofstep{4}{i=j}{\rA}
 \proofstep{}{\mathsf q_i=\mathsf q_i}{\rII}
 \proofstep{4}{\mathsf q_i=\mathsf q_j}{\rIE{4,5}}
 \proofstep{}{i=j\to\mathsf q_i=\mathsf q_j}{\rRI{4,6}}
 \proofstep{}{\mathsf q_i=\mathsf q_j\leftrightarrow i=j}{\rLRI{3,7}}
\end{tabproof}

\FormulaDefDeltaK[Saetze der Mengensprache]
{\operatorname{Sent}(\mathcal L_\in)\coloneq
 \{p\in\mathcal F\mid\operatorname{FV}(p)=\varnothing\}}
{B48SentenceCodes}{\DeltaRow{Formelmenge}{\mathcal F}}

\FormulaThmDeltaK[Mitgliedschaft nach Entfernen eines anderen Elements]
{j\in U,\quad j\neq i\ \vdash\ j\in U\setminus\{i\}}
{B48MemberDifferenceSingleton}{\DeltaRow{Mengen}{U,i,j}}
\begin{tabproof}
 \proofstep{1}{j\in U}{\rA}
 \proofstep{2}{j\neq i}{\rA}
 \proofstep{}{j\in\{i\}\leftrightarrow j=i}
  {\FormulaRefAuto{x\in\{a\}\eqvdash x=a}[thm]}
 \proofstep{2}{j\notin\{i\}}{\rLRS{3,2}}
 \proofstep{1,2}{j\in U\land j\notin\{i\}}{\rAI{1,4}}
 \proofstep{1,2}{j\in U\setminus\{i\}}
  {\FormulaRefAuto{x\in A\setminus B\eqvdash x\in A\land x\notin B}[thm]{5}}
\end{tabproof}

\FormulaThmDeltaK[Fehlen ausserhalb des entfernten Elements]
{x\notin U\setminus\{i\},\quad x\neq i\ \vdash\ x\notin U}
{B48MissingFromDifference}{\DeltaRow{Mengen}{U,i,x}}
\begin{tabproof}
 \proofstep{1}{x\notin U\setminus\{i\}}{\rA}
 \proofstep{2}{x\neq i}{\rA}
 \proofstep{3}{x\in U}{\rA}
 \proofstep{2,3}{x\in U\setminus\{i\}}
  {\FormulaRefAuto{B48MemberDifferenceSingleton}[thm]{3,2}}
 \proofstep{1,2,3}{\bot}{\rBI{4,1}}
 \proofstep{1,2}{x\notin U}{\rCI{3,5}}
\end{tabproof}

\FormulaDefDeltaK[Uebereinstimmung auf einem Variablenbereich]
{s\equiv_U t\coloneq\forall j\in U\;(s(j)=t(j))}
{B48Agreement}{
 \DeltaRow{Belegungen}{s,t\in A_D}
 \DeltaRow{Variablenbereich}{U\subseteq\mathbb N}
}

\FormulaThmDeltaK[Punktweise Werte eines Belegungswechsels]
{\begin{aligned}
 s[i\mapsto a](i)&=a,\\
 j\neq i&\ \vdash\ s[i\mapsto a](j)=s(j)
\end{aligned}}
{B48UpdateValues}{
 \DeltaRow{Belegung}{s\in A_D}
 \DeltaRow{Indizes und Wert}{i,j\in\mathbb N,\quad a\in D}
}
\begin{tabproof}
 \proofstep{}{\forall k\in\mathbb N\;
   \BXLVIIIIf{k=i}{a}{s(k)}\in D}
  {\FormulaRefAuto{B48IfTyped}[thm],
   \FormulaRefAuto{F\colon A\to B,\ x\in A\vdash F(x)\in B}[thm]}
 \proofstep{}{\forall k\in\mathbb N\;
   s[i\mapsto a](k)=\BXLVIIIIf{k=i}{a}{s(k)}}
  {\FormulaRefAuto{TypedTermGraphValues}[thm]{1},
   \FormulaRefAuto{B48Assignments}[def]}
 \proofstep{}{i=i}{\rII}
 \proofstep{}{\BXLVIIIIf{i=i}{a}{s(i)}=a}
  {\FormulaRefAuto{B48IfTrue}[thm]{3}}
 \proofstep{}{s[i\mapsto a](i)=a}{\rIE{4,\rUE{2}}}
 \proofstep{6}{j\neq i}{\rA}
 \proofstep{6}{\BXLVIIIIf{j=i}{a}{s(j)}=s(j)}
  {\FormulaRefAuto{B48IfFalse}[thm]{6}}
 \proofstep{6}{s[i\mapsto a](j)=s(j)}{\rIE{7,\rUE{2}}}
\end{tabproof}

\FormulaThmDeltaK[Uebereinstimmung nach gleichem Belegungswechsel]
{s\equiv_{U\setminus\{i\}}t\ \vdash\
 s[i\mapsto a]\equiv_U t[i\mapsto a]}
{B48AgreementUpdate}{
 \DeltaRow{Belegungen}{s,t\in A_D}
 \DeltaRow{Bereich, Index und Wert}{U\subseteq\mathbb N,\quad i\in\mathbb N,\quad a\in D}
}
\begin{tabproof}
 \proofstep{1}{s\equiv_{U\setminus\{i\}}t}{\rA}
 \proofstep{2}{j\in U}{\rA}
 \proofstep{}{j=i\lor j\neq i}{\FormulaRefAuto{P\lor\neg P}[thm]}
 \proofstep{4}{j=i}{\rA}
 \proofstep{4}{s[i\mapsto a](j)=a}
  {\FormulaRefAuto{B48UpdateValues}[thm]{4}}
 \proofstep{4}{t[i\mapsto a](j)=a}
  {\FormulaRefAuto{B48UpdateValues}[thm]{4}}
 \proofstep{4}{s[i\mapsto a](j)=t[i\mapsto a](j)}
  {\FormulaRefAuto{a=b\dsep c=b\vdash a=c}[thm]{5,6}}
 \proofstep{8}{j\neq i}{\rA}
 \proofstep{2,8}{j\in U\setminus\{i\}}
  {\FormulaRefAuto{B48MemberDifferenceSingleton}[thm]{2,8}}
 \proofstep{1,2,8}{s(j)=t(j)}
  {\FormulaRefAuto{B48Agreement}[def]{1,9}}
 \proofstep{8}{s[i\mapsto a](j)=s(j)}
  {\FormulaRefAuto{B48UpdateValues}[thm]{8}}
 \proofstep{8}{t[i\mapsto a](j)=t(j)}
  {\FormulaRefAuto{B48UpdateValues}[thm]{8}}
 \proofstep{1,2,8}{s[i\mapsto a](j)=t[i\mapsto a](j)}
  {\FormulaRefAuto{a=b\dsep c=b\vdash a=c}[thm]{\rIE{10,11},12}}
 \proofstep{1,2}{s[i\mapsto a](j)=t[i\mapsto a](j)}
  {\rOE{3,4:7,8:13}}
 \proofstep{1}{\forall j\in U\;
                 s[i\mapsto a](j)=t[i\mapsto a](j)}
  {\rUI{\rRI{2,14}}}
 \proofstep{1}{s[i\mapsto a]\equiv_U t[i\mapsto a]}
  {\FormulaRefAuto{B48Agreement}[def]{15}}
\end{tabproof}

\FormulaDefDeltaK[Induktionsbehauptung des Koinzidenzlemmas]
{\begin{aligned}
 C(p)\coloneq\forall s,t\in A_D\;\bigl(
 s\equiv_{\operatorname{FV}(p)}t\ \to{}\\
 ((M,s\models p)\leftrightarrow(M,t\models p))\bigr)
\end{aligned}}
{B48CoincidenceProperty}{
 \DeltaRow{Struktur und Formelcode}{\mathsf{Str}(M),\quad p\in\mathcal F}
}

\FormulaThmDeltaK[Koinzidenzlemma]
{s\equiv_{\operatorname{FV}(p)}t\ \vdash\
 ((M,s\models p)\leftrightarrow(M,t\models p))}
{B48Coincidence}{
 \DeltaRow{Struktur}{\mathsf{Str}(M)}
 \DeltaRow{Formel und Belegungen}{p\in\mathcal F,\quad s,t\in A_D}
}
\begin{tabproofsplitwide}
 \proofpartwideindR[Atomare Formeln]{C(\mathsf e(i,j))\land
   C(\mathsf m(i,j))\land C(\mathsf b)}
  {C(\mathsf e(i,j))\land C(\mathsf m(i,j))\land C(\mathsf b)}
  [B48CoincidenceAtoms]
 \proofstepwidestar[1]{s\equiv_{\{i,j\}}t}{\rA}
 \proofstepwidestar[1]{s(i)=t(i)}
  {\FormulaRefAuto{B48Agreement}[def]{1}}
 \proofstepwidestar[1]{s(j)=t(j)}
  {\FormulaRefAuto{B48Agreement}[def]{1}}
 \proofstepwidestar[]{(s(i)=s(j))\leftrightarrow(s(i)=s(j))}
  {\FormulaRefAuto{P\leftrightarrow P}[thm]}
 \proofstepwidestar[1]{(s(i)=s(j))\leftrightarrow(t(i)=t(j))}
  {\rIE{2,3,4}}
 \proofstepwidestar[1]{\begin{aligned}
 (M,s\models\mathsf e(i,j))\leftrightarrow
 (M,t\models\mathsf e(i,j))
 \end{aligned}}
  {\FormulaRefAuto{B48Satisfaction}[def]{5}}
 \proofstepwidestar[]{((s(i),s(j))\in E)\leftrightarrow
                        ((s(i),s(j))\in E)}
  {\FormulaRefAuto{P\leftrightarrow P}[thm]}
 \proofstepwidestar[1]{((s(i),s(j))\in E)\leftrightarrow
                        ((t(i),t(j))\in E)}
  {\rIE{2,3,7}}
 \proofstepwidestar[1]{\begin{aligned}
 (M,s\models\mathsf m(i,j))\leftrightarrow
 (M,t\models\mathsf m(i,j))
 \end{aligned}}
  {\FormulaRefAuto{B48Satisfaction}[def]{8}}
 \proofstepwidestar[]{C(\mathsf e(i,j))}
  {\FormulaRefAuto{B48CoincidenceProperty}[def]
   {\rUI{\rRI{1,6}}},\FormulaRefAuto{B48FVClauses}[thm]}
 \proofstepwidestar[]{C(\mathsf m(i,j))}
  {\FormulaRefAuto{B48CoincidenceProperty}[def]
   {\rUI{\rRI{1,9}}},\FormulaRefAuto{B48FVClauses}[thm]}
 \proofstepwidestar[]{(M,s\models\mathsf b)\leftrightarrow
                        (M,t\models\mathsf b)}
  {\FormulaRefAuto{B48Satisfaction}[def],
   \FormulaRefAuto{P\leftrightarrow P}[thm]}
 \proofstepwidestar[]{C(\mathsf b)}
  {\FormulaRefAuto{B48CoincidenceProperty}[def]{\rUI{12}}}
 \proofstepwidestar[]{C(\mathsf e(i,j))\land
                      C(\mathsf m(i,j))\land C(\mathsf b)}
  {\rAI{10,\rAI{11,13}}}
 \closeproofpartwideindR

 \proofpartwideindR[Negation]{C(p)\vdash C(\mathsf{neg}(p))}
  {C(p)\vdash C(\mathsf{neg}(p))}[B48CoincidenceNeg]
 \proofstepwidestar[1]{C(p)}{\rA}
 \proofstepwidestar[2]{s\equiv_{\operatorname{FV}(\mathsf{neg}(p))}t}{\rA}
 \proofstepwidestar[2]{s\equiv_{\operatorname{FV}(p)}t}
  {\FormulaRefAuto{B48FVClauses}[thm]{2}}
 \proofstepwidestar[1,2]{(M,s\models p)\leftrightarrow(M,t\models p)}
  {\FormulaRefAuto{B48CoincidenceProperty}[def]{1,3}}
 \proofstepwidestar[1,2]{\neg(M,s\models p)\leftrightarrow
                       \neg(M,t\models p)}
  {\rLRS{4,\FormulaRefAuto{P\leftrightarrow P}[thm]}}
 \proofstepwidestar[1,2]{(M,s\models\neg p)\leftrightarrow
                       (M,t\models\neg p)}
  {\FormulaRefAuto{B48Satisfaction}[def]{5}}
 \proofstepwidestar[1]{C(\mathsf{neg}(p))}
  {\FormulaRefAuto{B48CoincidenceProperty}[def]{\rUI{\rRI{2,6}}}}
 \closeproofpartwideindR

 \proofpartwideindR[Konjunktion]{C(p),C(q)\vdash C(\mathsf{and}(p,q))}
  {C(p),C(q)\vdash C(\mathsf{and}(p,q))}[B48CoincidenceAnd]
 \proofstepwidestar[1]{C(p)}{\rA}
 \proofstepwidestar[2]{C(q)}{\rA}
 \proofstepwidestar[3]{s\equiv_{\operatorname{FV}(\mathsf{and}(p,q))}t}{\rA}
 \proofstepwidestar[3]{s\equiv_{\operatorname{FV}(p)}t}
  {\FormulaRefAuto{B48Agreement}[def]{3},
   \FormulaRefAuto{B48FVClauses}[thm]}
 \proofstepwidestar[3]{s\equiv_{\operatorname{FV}(q)}t}
  {\FormulaRefAuto{B48Agreement}[def]{3},
   \FormulaRefAuto{B48FVClauses}[thm]}
 \proofstepwidestar[1,3]{(M,s\models p)\leftrightarrow(M,t\models p)}
  {\FormulaRefAuto{B48CoincidenceProperty}[def]{1,4}}
 \proofstepwidestar[2,3]{(M,s\models q)\leftrightarrow(M,t\models q)}
  {\FormulaRefAuto{B48CoincidenceProperty}[def]{2,5}}
 \proofstepwidestar[1,2,3]{\begin{aligned}
 ((M,s\models p)\land(M,s\models q))\leftrightarrow{}\\
 ((M,t\models p)\land(M,t\models q))
 \end{aligned}}
  {\rLRS{6,\rLRS{7,\FormulaRefAuto{P\leftrightarrow P}[thm]}}}
 \proofstepwidestar[1,2,3]{(M,s\models p\land q)\leftrightarrow
                          (M,t\models p\land q)}
  {\FormulaRefAuto{B48Satisfaction}[def]{8}}
 \proofstepwidestar[1,2]{C(\mathsf{and}(p,q))}
  {\FormulaRefAuto{B48CoincidenceProperty}[def]{\rUI{\rRI{3,9}}}}
 \closeproofpartwideindR

 \proofpartwideindR[Existenzquantor]{C(p)\vdash C(\mathsf{ex}_i(p))}
  {C(p)\vdash C(\mathsf{ex}_i(p))}[B48CoincidenceExists]
 \proofstepwidestar[1]{C(p)}{\rA}
 \proofstepwidestar[2]{s\equiv_{\operatorname{FV}(\mathsf{ex}_i(p))}t}{\rA}
 \proofstepwidestar[2]{s\equiv_{\operatorname{FV}(p)\setminus\{i\}}t}
  {\FormulaRefAuto{B48FVClauses}[thm]{2}}
 \proofstepwidestar[4]{a\in D}{\rA}
 \proofstepwidestar[4]{s[i\mapsto a]\in A_D}
  {\FormulaRefAuto{B48AssignmentUpdate}[thm]{4}}
 \proofstepwidestar[4]{t[i\mapsto a]\in A_D}
  {\FormulaRefAuto{B48AssignmentUpdate}[thm]{4}}
 \proofstepwidestar[2,4]{s[i\mapsto a]\equiv_{\operatorname{FV}(p)}
                                  t[i\mapsto a]}
  {\FormulaRefAuto{B48AgreementUpdate}[thm]{3,4}}
 \proofstepwidestar[1,2,4]{\begin{aligned}
 (M,s[i\mapsto a]\models p)\leftrightarrow
 (M,t[i\mapsto a]\models p)
 \end{aligned}}
  {\FormulaRefAuto{B48CoincidenceProperty}[def]{1,5,6,7}}
 \proofstepwidestar[1,2]{\begin{aligned}
 \forall a\in D\;((M,s[i\mapsto a]\models p)\leftrightarrow{}\\
 (M,t[i\mapsto a]\models p))
 \end{aligned}}
  {\rUI{\rRI{4,8}}}
 \proofstepwidestar[1,2]{\begin{aligned}
 (\exists a\in D\;(M,s[i\mapsto a]\models p))\leftrightarrow{}\\
 (\exists a\in D\;(M,t[i\mapsto a]\models p))
 \end{aligned}}
  {\FormulaRefAuto{B48ExistsIff}[thm]{9}}
 \proofstepwidestar[1,2]{(M,s\models\exists v_i\,p)\leftrightarrow
                        (M,t\models\exists v_i\,p)}
  {\FormulaRefAuto{B48Satisfaction}[def]{10}}
 \proofstepwidestar[1]{C(\mathsf{ex}_i(p))}
  {\FormulaRefAuto{B48CoincidenceProperty}[def]{\rUI{\rRI{2,11}}}}
 \closeproofpartwideindR

 \proofpartwide{Abschluss der Formelinduktion}
 \proofstepwidestar[]{\{p\in\mathcal F\mid C(p)\}\subseteq\mathcal T}
  {\FormulaRefAuto{B48FormulaInTree}[thm]}
 \proofstepwidestar[]{\forall c\in\Sigma_{\rm at}\;
    \mathsf l(c)\in\{p\in\mathcal F\mid C(p)\}}
  {\FormulaRefAuto{B48CoincidenceAtoms}[thm],
   \FormulaRefAuto{B48FormulaAtom}[thm]}
 \proofstepwidestar[]{\forall p\in\mathcal F\;
     (C(p)\to C(\mathsf{neg}(p)))}
  {\FormulaRefAuto{B48CoincidenceNeg}[thm]}
 \proofstepwidestar[]{\forall p,q\in\mathcal F\;
     (C(p)\land C(q)\to C(\mathsf{and}(p,q)))}
  {\FormulaRefAuto{B48CoincidenceAnd}[thm]}
 \proofstepwidestar[]{\forall i\in\mathbb N\,\forall p\in\mathcal F\;
     (C(p)\to C(\mathsf{ex}_i(p)))}
  {\FormulaRefAuto{B48CoincidenceExists}[thm]}
 \proofstepwidestar[]{\mathsf{Cl}(\{p\in\mathcal F\mid C(p)\})}
  {\FormulaRefAuto{B48FormulaClosure}[def]{1,2,3,4,5},
   \FormulaRefAuto{B48FormulaNeg}[thm],
   \FormulaRefAuto{B48FormulaAnd}[thm],
   \FormulaRefAuto{B48FormulaExists}[thm]}
 \proofstepwidestar[]{\forall p\in\mathcal F\;C(p)}
  {\FormulaRefAuto{B48FormulaInduction}[thm]{6}}
 \proofstepwidestar[8]{s\equiv_{\operatorname{FV}(p)}t}{\rA}
 \proofstepwidestar[8]{(M,s\models p)\leftrightarrow(M,t\models p)}
  {\FormulaRefAuto{B48CoincidenceProperty}[def]{7,8}}
 \closeproofpartwide
\end{tabproofsplitwide}
Die Induktionsbehauptung quantifiziert ueber \emph{alle} Belegungen.
Daher darf sie im Quantorenschritt auf die beiden geaenderten
Belegungen angewendet werden.

\FormulaThmDeltaK[Eine nicht freie Variable darf neu belegt werden]
{i\notin\operatorname{FV}(p)\ \vdash\
 ((M,s\models p)\leftrightarrow(M,s[i\mapsto a]\models p))}
{B48FreshUpdate}{
 \DeltaRow{Struktur und Belegung}{\mathsf{Str}(M),\quad s\in A_D}
 \DeltaRow{Formel, Index und Wert}{p\in\mathcal F,\quad i\in\mathbb N,\quad a\in D}
}
\begin{tabproof}
 \proofstep{1}{i\notin\operatorname{FV}(p)}{\rA}
 \proofstep{2}{j\in\operatorname{FV}(p)}{\rA}
 \proofstep{3}{j=i}{\rA}
 \proofstep{2,3}{i\in\operatorname{FV}(p)}{\rIE{3,2}}
 \proofstep{1,2,3}{\bot}{\rBI{4,1}}
 \proofstep{1,2}{j\neq i}{\rCI{3,5}}
 \proofstep{1,2}{s[i\mapsto a](j)=s(j)}
  {\FormulaRefAuto{B48UpdateValues}[thm]{6}}
 \proofstep{1,2}{s(j)=s[i\mapsto a](j)}
  {\FormulaRefAuto{a=b\vdash b=a}[thm]{7}}
 \proofstep{1}{s\equiv_{\operatorname{FV}(p)}s[i\mapsto a]}
  {\FormulaRefAuto{B48Agreement}[def]{\rUI{\rRI{2,8}}}}
 \proofstep{}{s[i\mapsto a]\in A_D}
  {\FormulaRefAuto{B48AssignmentUpdate}[thm]}
 \proofstep{1}{(M,s\models p)\leftrightarrow
              (M,s[i\mapsto a]\models p)}
  {\FormulaRefAuto{B48Coincidence}[thm]{9,10}}
\end{tabproof}

\FormulaThmDeltaK[Saetze sind unabhaengig von der Belegung]
{\operatorname{FV}(p)=\varnothing\ \vdash\
 ((M,s\models p)\leftrightarrow(M,t\models p))}
{B48SentenceCoincidence}{
 \DeltaRow{Struktur}{\mathsf{Str}(M)}
 \DeltaRow{Formel und Belegungen}{p\in\mathcal F,\quad s,t\in A_D}
}
\begin{tabproof}
 \proofstep{1}{\operatorname{FV}(p)=\varnothing}{\rA}
 \proofstep{}{s\equiv_\varnothing t}{\FormulaRefAuto{B48Agreement}[def]}
 \proofstep{1}{s\equiv_{\operatorname{FV}(p)}t}{\rIE{1,2}}
 \proofstep{1}{(M,s\models p)\leftrightarrow(M,t\models p)}
  {\FormulaRefAuto{B48Coincidence}[thm]{3}}
\end{tabproof}


% Syntaktische Substitution in der konstantenfreien Sprache {=, in}.
\subsection{Kollisionsfreie Substitution}

In \(\mathcal L_\in\) sind die einzigen Terme Variablen. Deshalb ist
die hier benoetigte Einsetzung eine Variablenersetzung.
Konstanten der spaeteren Henkin-Sprache sind in diesem Abschnitt
noch nicht zugelassen. Eine Ersetzung wird nur dann in einer
Quantoren- oder Gleichheitsregel verwendet, wenn sie keine freie
Variable bindet.

\FormulaDefDeltaK[Operatoren der syntaktischen Variablenersetzung]
{\begin{aligned}
 \rho_{xy}(j)&\coloneq\BXLVIIIIf{j=x}{y}{j},\\
 r_{xy}((k,(i,j)))&\coloneq
 \BXLVIIIIf{k=t_{\rm e}\lor k=t_{\rm m}}
 {(k,(\rho_{xy}(i),\rho_{xy}(j)))}{(k,(i,j))},\\
 f_{xy}(c)&\coloneq(\mathsf l(c),\mathsf l(r_{xy}(c))),\\
 g_{xy}((p,u),(q,v))&\coloneq
 \bigl(\mathsf n(p,q),
 \BXLVIIIIf{p=\mathsf q_x}{\mathsf n(p,q)}{\mathsf n(u,v)}\bigr).
\end{aligned}}
{B48SubstitutionOperators}{
 \DeltaRow{Variablen}{x,y,i,j\in\mathbb N}
 \DeltaRow{Bereiche}{c\in\Sigma,\quad p,q,u,v\in\mathcal T}
}
Die Bindermarken bleiben unveraendert. Am Binder \(x\) wird der
urspruengliche gesamte Teilbaum zurueckgegeben. Das verhindert die
Ersetzung der dort gebundenen Vorkommen von \(x\).

\FormulaThmDeltaK[Existenz der syntaktischen Ersetzung]
{\begin{gathered}
 \exists!\sigma_{xy}\colon\mathcal T\to\mathcal T\quad
 \sigma_{xy}(\mathsf l(c))=\mathsf l(r_{xy}(c)),\\
 \sigma_{xy}(\mathsf n(p,q))=
 \BXLVIIIIf{p=\mathsf q_x}{\mathsf n(p,q)}
            {\mathsf n(\sigma_{xy}(p),\sigma_{xy}(q))}
\end{gathered}}
{B48SubstitutionExists}{
 \DeltaRow{Variablen}{x,y\in\mathbb N}
 \DeltaRow{Universell in den Gleichungen}{c\in\Sigma,\quad p,q\in\mathcal T}
}
Setze \(Z=\mathcal T\times\mathcal T\).
Die Operatoren bezeichnen wieder ihre beschraenkten Funktionsgraphen.
\begin{tabproof}
 \proofstep{}{\rho_{xy}(j)\in\mathbb N}
  {\FormulaRefAuto{B48IfTyped}[thm]}
 \proofstep{}{r_{xy}(c)\in\Sigma}
  {\FormulaRefAuto{B48SubstitutionOperators}[def]{1},
   \FormulaRefAuto{B48LabelType}[thm]}
 \proofstep{}{f_{xy}\colon\Sigma\to Z}
  {\FormulaRefAuto{TypedTermGraphFunction}[thm]{2},
   \FormulaRefAuto{B48LeafType}[thm]}
 \proofstep{}{g_{xy}\colon Z\times Z\to Z}
  {\FormulaRefAuto{TypedTermGraphFunction}[thm],
   \FormulaRefAuto{B48NodeType}[thm],
   \FormulaRefAuto{B48IfTyped}[thm]}
 \proofstep{}{\begin{gathered}
  \exists!e\colon\mathcal T\to Z\quad
  e(\mathsf l(c))=f_{xy}(c),\\
  e(\mathsf n(p,q))=g_{xy}(e(p),e(q))
 \end{gathered}}
  {\FormulaRefAuto{FullPlanarBinaryTreeRecursion}[thm]{3,4}}
\end{tabproof}
Fuer diese Auswertung sei \(H(p)\) die Aussage
\(\exists u\in\mathcal T\;(e(p)=(p,u))\).
\begin{tabproof}
 \proofstep{}{e(\mathsf l(c))=
  (\mathsf l(c),\mathsf l(r_{xy}(c)))}
  {\FormulaRefAuto{B48SubstitutionOperators}[def]}
 \proofstep{}{H(\mathsf l(c))}{\rEI{1}}
 \proofstep{3}{H(p)\land H(q)}{\rA}
 \proofstep{4}{e(p)=(p,u)\land e(q)=(q,v)}{\rA}
 \proofstep{4}{\begin{aligned}
 e(\mathsf n(p,q))={}&
 (\mathsf n(p,q),\\
 &\BXLVIIIIf{p=\mathsf q_x}{\mathsf n(p,q)}{\mathsf n(u,v)})
 \end{aligned}}
  {\FormulaRefAuto{B48SubstitutionOperators}[def]{4}}
 \proofstep{4}{H(\mathsf n(p,q))}{\rEI{5}}
 \proofstep{3}{H(\mathsf n(p,q))}{\rEEStar{3,4:6}}
 \proofstep{}{\forall p\in\mathcal T\;H(p)}
  {\FormulaRefAuto{FullPlanarBinaryTreeStructuralInduction}[thm]
   {\rUI{2},\rUI{\rRI{3,7}}}}
\end{tabproof}
Die beschraenkten Zeugen \(u,v\) sind frisch.
Die zweite Komponente ist eindeutig durch
\FormulaRefAuto{(a,b)=(c,d)\eqvdash a=c\land b=d}[thm].
Somit definiert der Graph
\(\{(p,u)\in\mathcal T^2\mid e(p)=(p,u)\}\) nach
\FormulaRefAuto{UniqueValuedGraphFunction}[thm] die gesuchte Funktion.
Ihre Gleichungen folgen durch Einsetzen der Komponenten in die
Rekursionsgleichungen von \(e\).
Fuer zwei Kandidaten sind die Blattwerte gleich. Am Knoten sind im
Fall \(p=\mathsf q_x\) beide Werte der urspruengliche Knoten,
andernfalls macht die Induktionsannahme beide Argumente von
\(\mathsf n\) gleich. Bauminduktion und
\FormulaRefAuto{FunctionExtensionality}[thm] beweisen Eindeutigkeit.

\FormulaDefDeltaK[Ersetzung freier Vorkommen]
{p[y/x]\coloneq\sigma_{xy}(p)}
{B48Substitution}{
 \DeltaRow{Formelcode}{p\in\mathcal F}
 \DeltaRow{Variablen}{x,y\in\mathbb N}
 \DeltaRow{Konstruierte Funktion}{\sigma_{xy}\text{ aus }
          \FormulaRefAuto{B48SubstitutionExists}[thm]}
}
Die syntaktische Operation ist auch dann definiert, wenn sie eine
Variable einfangen wuerde. Die nachstehende Bedingung schliesst
genau diese Faelle von den logischen Regeln aus.

\FormulaThmDeltaK[Rekursionsklauseln der Ersetzung]
{\begin{aligned}
 \mathsf e(i,j)[y/x]&=\mathsf e(\rho_{xy}(i),\rho_{xy}(j)),\\
 \mathsf m(i,j)[y/x]&=\mathsf m(\rho_{xy}(i),\rho_{xy}(j)),\\
 \mathsf b[y/x]&=\mathsf b,\\
 (\neg p)[y/x]&=\neg(p[y/x]),\\
 (p\land q)[y/x]&=p[y/x]\land q[y/x],\\
 (\exists v_i\,p)[y/x]&=
 \begin{cases}\exists v_i\,p,&i=x,\\
 \exists v_i\,(p[y/x]),&i\neq x.
 \end{cases}
\end{aligned}}
{B48SubstitutionClauses}{
 \DeltaRow{Formeln}{p,q\in\mathcal F}
 \DeltaRow{Variablen}{i,j,x,y\in\mathbb N}
}
\begin{tabproof}
 \proofstep{}{\sigma_{xy}(\mathsf n_0)=\mathsf n_0,\quad
 \sigma_{xy}(\mathsf c_0)=\mathsf c_0,\quad
 \sigma_{xy}(\mathsf q_i)=\mathsf q_i}
  {\FormulaRefAuto{B48SubstitutionOperators}[def],
   \FormulaRefAuto{B48SubstitutionExists}[thm]}
 \proofstep{}{\mathsf q_i=\mathsf q_x\leftrightarrow i=x}
  {\FormulaRefAuto{B48QuantifierRootEquality}[thm]}
 \proofstep{}{\sigma_{xy}(\mathsf n(\mathsf c_0,p))=
                         \mathsf n(\mathsf c_0,\sigma_{xy}(p))}
  {\FormulaRefAuto{B48SubstitutionExists}[thm]{1},
   \FormulaRefAuto{B48QNeAnd}[thm]}
 \proofstep{}{\begin{gathered}
 (\neg p)[y/x]=\neg(p[y/x]),\\
 (p\land q)[y/x]=p[y/x]\land q[y/x]
 \end{gathered}}
  {\FormulaRefAuto{B48SubstitutionExists}[thm]{1,3},
   \FormulaRefAuto{B48QNeNeg}[thm],
   \FormulaRefAuto{B48LeafNotNode}[thm]}
 \proofstep{}{\begin{aligned}
 (\exists v_i\,p)[y/x]=
 \BXLVIIIIf{i=x}{\exists v_i\,p}{\exists v_i\,(p[y/x])}
 \end{aligned}}
  {\FormulaRefAuto{B48SubstitutionExists}[thm]{1,2}}
\end{tabproof}
Die drei Atomgleichungen sind unmittelbar die Blattgleichung mit
\(r_{xy}\). Formelinduktion
\FormulaRefAuto{B48FormulaInduction}[thm] liefert aus den sechs
Klauseln zudem \(p[y/x]\in\mathcal F\): Atome bleiben Atome;
Negation und Konjunktion verwenden den jeweiligen Formelabschluss;
am Quantor bleibt entweder die urspruengliche Formel erhalten oder
der Existenzabschluss wird auf die ersetzte Teilformel angewendet.

\FormulaDefDeltaK[Frei fuer eine Variable]
{\begin{aligned}
 \mathsf{Fr}(\mathsf e(i,j),x,y)&\coloneq\top,\\
 \mathsf{Fr}(\mathsf m(i,j),x,y)&\coloneq\top,\qquad
 \mathsf{Fr}(\mathsf b,x,y)\coloneq\top,\\
 \mathsf{Fr}(\neg p,x,y)&\coloneq\mathsf{Fr}(p,x,y),\\
 \mathsf{Fr}(p\land q,x,y)&\coloneq
                  \mathsf{Fr}(p,x,y)\land\mathsf{Fr}(q,x,y),\\
 \mathsf{Fr}(\exists v_i\,p,x,y)&\coloneq
 i=x\lor\bigl(\mathsf{Fr}(p,x,y)\land
              (i\neq y\lor x\notin\operatorname{FV}(p))\bigr).
\end{aligned}}
{B48FreeFor}{
 \DeltaRow{Bedeutung}{y\text{ ist frei fuer }x\text{ in }p}
 \DeltaRow{Formeln und Variablen}{p,q\in\mathcal F,\quad i,j,x,y\in\mathbb N}
}
Hier ist \(\top\) die Abkuerzung \(\neg\bot\).
Auch diese rekursive Definition bezeichnet eine konstruierbare
Menge: Man fuehrt an jedem Baum gleichzeitig die Teilmenge
\(W(p)\subseteq\mathbb N^2\) aller zulaessigen Paare \((x,y)\) mit.
Die Blattwerte sind \(\mathbb N^2\); am Negations- und
Konjunktionshilfsknoten \(\mathsf c_0\) wird der rechte Wert
weitergereicht; am fertigen Konjunktionsknoten werden die beiden
Werte geschnitten. Am Quantorenknoten mit linkem Kind
\(\mathsf q_i\) ist der Wert
\[
 \{(x,y)\in\mathbb N^2\mid
    i=x\lor((x,y)\in W(p)\land
                 (i\neq y\lor x\notin\operatorname{FV}(p)))\}.
\]
Alle sonstigen Hilfsknoten erhalten den Wert \(\mathbb N^2\).
Diese Faelle sind durch
\FormulaRefAuto{B48LeafNotNode}[thm],
\FormulaRefAuto{B48AndNeNeg}[thm],
\FormulaRefAuto{B48QNeNeg}[thm],
\FormulaRefAuto{B48QNeAnd}[thm] und
\FormulaRefAuto{B48QuantifierRootInjective}[thm] eindeutig.
Die beiden totalen Operatoren haben damit Werte in
\(\mathcal T\times\powerset(\mathbb N^2)\).
\FormulaRefAuto{TypedTermGraphFunction}[thm] und
\FormulaRefAuto{FullPlanarBinaryTreeRecursion}[thm] liefern ihre
eindeutige Auswertung. Wie bei
\FormulaRefAuto{B48FVExists}[thm] bleibt die erste Komponente
durch Blatt- und Knoteninduktion der Baumcode.
Die Aussage \(\mathsf{Fr}(p,x,y)\) bedeutet
\((x,y)\in W(p)\). Die angezeigten Klauseln sind die
ausgerechneten Werte dieser Konstruktion.

\FormulaThmDeltaK[Eine nicht freie Variable wird nicht ersetzt]
{x\notin\operatorname{FV}(p)\ \vdash\ p[y/x]=p}
{B48SubstitutionVacuous}{
 \DeltaRow{Formel und Variablen}{p\in\mathcal F,\quad x,y\in\mathbb N}
}
Die Induktionsbehauptung wird fuer alle \(x,y\) gleichzeitig gefuehrt.
\begin{tabproofsplitwide}
 \proofpartwide{Atom- und Junktorenfaelle}
 \proofstepwidestar[1]{x\notin\{i,j\}}{\rA}
 \proofstepwidestar[1]{\rho_{xy}(i)=i,\quad\rho_{xy}(j)=j}
  {\FormulaRefAuto{B48SubstitutionOperators}[def]{1},
   \FormulaRefAuto{B48IfFalse}[thm]}
 \proofstepwidestar[1]{\mathsf e(i,j)[y/x]=\mathsf e(i,j),\quad
                      \mathsf m(i,j)[y/x]=\mathsf m(i,j)}
  {\FormulaRefAuto{B48SubstitutionClauses}[thm]{2}}
 \proofstepwidestar[]{\mathsf b[y/x]=\mathsf b}
  {\FormulaRefAuto{B48SubstitutionClauses}[thm]}
 \proofstepwidestar[5]{p[y/x]=p}{\rA}
 \proofstepwidestar[5]{(\neg p)[y/x]=\neg p}
  {\FormulaRefAuto{B48SubstitutionClauses}[thm]{5}}
 \proofstepwidestar[7]{q[y/x]=q}{\rA}
 \proofstepwidestar[5,7]{(p\land q)[y/x]=p\land q}
  {\FormulaRefAuto{B48SubstitutionClauses}[thm]{5,7}}
 \closeproofpartwide
 \proofpartwide{Quantorenfall und Abschluss}
 \proofstepwidestar[1]{x\notin\operatorname{FV}(p)\setminus\{i\}}{\rA}
 \proofstepwidestar[]{i=x\lor i\neq x}{\FormulaRefAuto{P\lor\neg P}[thm]}
 \proofstepwidestar[3]{i=x}{\rA}
 \proofstepwidestar[3]{(\exists v_i\,p)[y/x]=\exists v_i\,p}
  {\FormulaRefAuto{B48SubstitutionClauses}[thm]{3}}
 \proofstepwidestar[5]{i\neq x}{\rA}
 \proofstepwidestar[1,5]{x\notin\operatorname{FV}(p)}
  {\FormulaRefAuto{B48MissingFromDifference}[thm]{1,
    \FormulaRefAuto{a\neq b\vdash b\neq a}[thm]{5}}}
 \proofstepwidestar[7]{\forall x,y\in\mathbb N\;
   (x\notin\operatorname{FV}(p)\to p[y/x]=p)}{\rA}
 \proofstepwidestar[1,5,7]{p[y/x]=p}{\rRE{\rUE{7},6}}
 \proofstepwidestar[1,5,7]{(\exists v_i\,p)[y/x]=\exists v_i\,p}
  {\FormulaRefAuto{B48SubstitutionClauses}[thm]{5,8}}
 \proofstepwidestar[1,7]{(\exists v_i\,p)[y/x]=\exists v_i\,p}
  {\rOE{2,3:4,5:9}}
 \closeproofpartwide
\end{tabproofsplitwide}
Bei Negation ist die Freivariablenmenge unveraendert; bei
Konjunktion impliziert die Nichtzugehoerigkeit zu ihrer Vereinigung
die Nichtzugehoerigkeit zu beiden Teilmengen.
Damit liefern die Tabellen alle Abschlussfaelle fuer
\FormulaRefAuto{B48FormulaInduction}[thm], zusammen mit
\FormulaRefAuto{B48FVClauses}[thm] und den drei Formelabschluessen.

\FormulaThmDeltaK[Vertauschen verschiedener Belegungswechsel]
{i\neq x,\ i\neq y\ \vdash\
 s[i\mapsto a][x\mapsto s[i\mapsto a](y)]
     =s[x\mapsto s(y)][i\mapsto a]}
{B48UpdateCommute}{
 \DeltaRow{Belegung}{s\in A_D}
 \DeltaRow{Variablen und Wert}{i,x,y\in\mathbb N,\quad a\in D}
}
Setze \(u=s[i\mapsto a][x\mapsto s[i\mapsto a](y)]\) und
\(v=s[x\mapsto s(y)][i\mapsto a]\).
\begin{tabproofsplitwide}
 \proofpartwide{Beweis}
 \proofstepwidestar[1]{i\neq x}
  {\rA}
 \proofstepwidestar[2]{i\neq y}
  {\rA}
 \proofstepwidestar[1]{x\neq i}
  {\FormulaRefAuto{a\neq b\vdash b\neq a}[thm]{1}}
 \proofstepwidestar[2]{y\neq i}
  {\FormulaRefAuto{a\neq b\vdash b\neq a}[thm]{2}}
 \proofstepwidestar[2]{s[i\mapsto a](y)=s(y)}
  {\FormulaRefAuto{B48UpdateValues}[thm]{4}}
 \proofstepwidestar[]{u\colon\mathbb N\to D}
  {\FormulaRefAuto{B48AssignmentUpdate}[thm], \FormulaRefAuto{B48AssignmentCriterion}[thm]}
 \proofstepwidestar[]{v\colon\mathbb N\to D}
  {\FormulaRefAuto{B48AssignmentUpdate}[thm], \FormulaRefAuto{B48AssignmentCriterion}[thm]}
 \proofstepwidestar[]{j=x\lor j\neq x}
  {\FormulaRefAuto{P\lor\neg P}[thm]}
 \proofstepwidestar[9]{j=x}
  {\rA}
 \proofstepwidestar[2,9]{u(j)=s(y)}
  {\FormulaRefAuto{B48UpdateValues}[thm]{9,5}}
 \proofstepwidestar[1,9]{v(j)=s(y)}
  {\FormulaRefAuto{B48UpdateValues}[thm]{9,3}}
 \proofstepwidestar[1,2,9]{u(j)=v(j)}
  {\FormulaRefAuto{a=b\dsep c=b\vdash a=c}[thm]{10,11}}
 \proofstepwidestar[13]{j\neq x}
  {\rA}
 \proofstepwidestar[]{j=i\lor j\neq i}
  {\FormulaRefAuto{P\lor\neg P}[thm]}
 \proofstepwidestar[15]{j=i}
  {\rA}
 \proofstepwidestar[13,15]{u(j)=a}
  {\FormulaRefAuto{B48UpdateValues}[thm]{13,15}}
 \proofstepwidestar[15]{v(j)=a}
  {\FormulaRefAuto{B48UpdateValues}[thm]{15}}
 \proofstepwidestar[13,15]{u(j)=v(j)}
  {\FormulaRefAuto{a=b\dsep c=b\vdash a=c}[thm]{16,17}}
 \proofstepwidestar[19]{j\neq i}
  {\rA}
 \proofstepwidestar[13,19]{u(j)=s(j)}
  {\FormulaRefAuto{B48UpdateValues}[thm]{13,19}}
 \proofstepwidestar[13,19]{v(j)=s(j)}
  {\FormulaRefAuto{B48UpdateValues}[thm]{13,19}}
 \proofstepwidestar[13,19]{u(j)=v(j)}
  {\FormulaRefAuto{a=b\dsep c=b\vdash a=c}[thm]{20,21}}
 \proofstepwidestar[13]{u(j)=v(j)}
  {\rOE{14,15:18,19:22}}
 \proofstepwidestar[1,2]{u(j)=v(j)}
  {\rOE{8,9:12,13:23}}
 \proofstepwidestar[1,2]{\forall j\in\mathbb N\;(u(j)=v(j))}
  {\rUI{24}}
 \proofstepwidestar[1,2]{u=v}
  {\FormulaRefAuto{FunctionExtensionality}[thm]{6,7,25}}
 \closeproofpartwide
\end{tabproofsplitwide}

Die beiden Fallunterscheidungen pruefen zuerst \(j=x\), dann \(j=i\).
Es wird keine Vertauschung bei \(i=x\) behauptet.

\FormulaDefDeltaK[Induktionsbehauptung des Substitutionslemmas]
{\begin{aligned}
 S(p)\coloneq\forall x,y\in\mathbb N\,\forall s\in A_D\;
 \bigl(\mathsf{Fr}(p,x,y)\to{}\\
 ((M,s\models p[y/x])\leftrightarrow
  (M,s[x\mapsto s(y)]\models p))\bigr)
\end{aligned}}
{B48SubstitutionProperty}{
 \DeltaRow{Struktur und Formel}{\mathsf{Str}(M),\quad p\in\mathcal F}
}

\FormulaThmDeltaK[Substitutionslemma]
{\mathsf{Fr}(p,x,y)\ \vdash\
 ((M,s\models p[y/x])\leftrightarrow
  (M,s[x\mapsto s(y)]\models p))}
{B48SubstitutionLemma}{
 \DeltaRow{Struktur und Belegung}{\mathsf{Str}(M),\quad s\in A_D}
 \DeltaRow{Formel und Variablen}{p\in\mathcal F,\quad x,y\in\mathbb N}
}
\begin{tabproofsplitwide}
 \proofpartwideindR[Variablenwerte]{s(\rho_{xy}(j))=s[x\mapsto s(y)](j)}
  {s(\rho_{xy}(j))=s[x\mapsto s(y)](j)}[B48SubstitutionVariableValue]
 \proofstepwidestar[]{j=x\lor j\neq x}
  {\FormulaRefAuto{P\lor\neg P}[thm]}
 \proofstepwidestar[2]{j=x}
  {\rA}
 \proofstepwidestar[2]{\rho_{xy}(j)=y}
  {\FormulaRefAuto{B48SubstitutionOperators}[def]{2}, \FormulaRefAuto{B48IfTrue}[thm]{2}}
 \proofstepwidestar[]{s(\rho_{xy}(j))=s(\rho_{xy}(j))}
  {\rII}
 \proofstepwidestar[2]{s(\rho_{xy}(j))=s(y)}
  {\rIE{3,4}}
 \proofstepwidestar[2]{s[x\mapsto s(y)](j)=s(y)}
  {\FormulaRefAuto{B48UpdateValues}[thm]{2}}
 \proofstepwidestar[2]{s(\rho_{xy}(j))=s[x\mapsto s(y)](j)}
  {\FormulaRefAuto{a=b\dsep c=b\vdash a=c}[thm]{5,6}}
 \proofstepwidestar[8]{j\neq x}
  {\rA}
 \proofstepwidestar[8]{\rho_{xy}(j)=j}
  {\FormulaRefAuto{B48SubstitutionOperators}[def]{8}, \FormulaRefAuto{B48IfFalse}[thm]{8}}
 \proofstepwidestar[]{s(\rho_{xy}(j))=s(\rho_{xy}(j))}
  {\rII}
 \proofstepwidestar[8]{s(\rho_{xy}(j))=s(j)}
  {\rIE{9,10}}
 \proofstepwidestar[8]{s[x\mapsto s(y)](j)=s(j)}
  {\FormulaRefAuto{B48UpdateValues}[thm]{8}}
 \proofstepwidestar[8]{s(\rho_{xy}(j))=s[x\mapsto s(y)](j)}
  {\FormulaRefAuto{a=b\dsep c=b\vdash a=c}[thm]{11,12}}
 \proofstepwidestar[]{s(\rho_{xy}(j))=s[x\mapsto s(y)](j)}
  {\rOE{1,2:7,8:13}}
\closeproofpartwideindR

 \proofpartwideindR[Atome]{S(\mathsf e(i,j))\land S(\mathsf m(i,j))
                         \land S(\mathsf b)}
  {S(\mathsf e(i,j))\land S(\mathsf m(i,j))\land S(\mathsf b)}
  [B48SubstitutionAtoms]
 \proofstepwidestar[]{s(\rho_{xy}(i))=s[x\mapsto s(y)](i)}
  {\FormulaRefAuto{B48SubstitutionVariableValue}[thm]}
 \proofstepwidestar[]{s(\rho_{xy}(j))=s[x\mapsto s(y)](j)}
  {\FormulaRefAuto{B48SubstitutionVariableValue}[thm]}
 \proofstepwidestar[]{\begin{aligned}
 (M,s\models\mathsf e(i,j)[y/x])\leftrightarrow{}\\
 (M,s[x\mapsto s(y)]\models\mathsf e(i,j))
 \end{aligned}}
  {\FormulaRefAuto{B48SubstitutionClauses}[thm],
   \FormulaRefAuto{B48Satisfaction}[def],
   \rIE{1,2,\FormulaRefAuto{P\leftrightarrow P}[thm]}}
 \proofstepwidestar[]{\begin{aligned}
 (M,s\models\mathsf m(i,j)[y/x])\leftrightarrow{}\\
 (M,s[x\mapsto s(y)]\models\mathsf m(i,j))
 \end{aligned}}
  {\FormulaRefAuto{B48SubstitutionClauses}[thm],
   \FormulaRefAuto{B48Satisfaction}[def],
   \rIE{1,2,\FormulaRefAuto{P\leftrightarrow P}[thm]}}
 \proofstepwidestar[]{\begin{aligned}
 (M,s\models\mathsf b[y/x])\leftrightarrow{}\\
 (M,s[x\mapsto s(y)]\models\mathsf b)
 \end{aligned}}
  {\FormulaRefAuto{B48SubstitutionClauses}[thm],
   \FormulaRefAuto{B48Satisfaction}[def],
   \FormulaRefAuto{P\leftrightarrow P}[thm]}
 \proofstepwidestar[]{S(\mathsf e(i,j))\land S(\mathsf m(i,j))
                    \land S(\mathsf b)}
  {\FormulaRefAuto{B48SubstitutionProperty}[def]
   {\rAI{\rUI{3},\rAI{\rUI{4},\rUI{5}}}}}
 \closeproofpartwideindR

 \proofpartwideindR[Negation]{S(p)\vdash S(\neg p)}
  {S(p)\vdash S(\neg p)}[B48SubstitutionNeg]
 \proofstepwidestar[1]{S(p)}{\rA}
 \proofstepwidestar[2]{\mathsf{Fr}(\neg p,x,y)}{\rA}
 \proofstepwidestar[2]{\mathsf{Fr}(p,x,y)}
  {\FormulaRefAuto{B48FreeFor}[def]{2}}
 \proofstepwidestar[1,2]{(M,s\models p[y/x])\leftrightarrow
                   (M,s[x\mapsto s(y)]\models p)}
  {\FormulaRefAuto{B48SubstitutionProperty}[def]{1,3}}
 \proofstepwidestar[1,2]{(M,s\models(\neg p)[y/x])\leftrightarrow
                   (M,s[x\mapsto s(y)]\models\neg p)}
  {\rLRS{4,\FormulaRefAuto{P\leftrightarrow P}[thm]},
   \FormulaRefAuto{B48Satisfaction}[def],
   \FormulaRefAuto{B48SubstitutionClauses}[thm]}
 \proofstepwidestar[1]{S(\neg p)}
  {\FormulaRefAuto{B48SubstitutionProperty}[def]{\rUI{\rRI{2,5}}}}
 \closeproofpartwideindR

 \proofpartwideindR[Konjunktion]{S(p),S(q)\vdash S(p\land q)}
  {S(p),S(q)\vdash S(p\land q)}[B48SubstitutionAnd]
 \proofstepwidestar[1]{S(p)}{\rA}
 \proofstepwidestar[2]{S(q)}{\rA}
 \proofstepwidestar[3]{\mathsf{Fr}(p\land q,x,y)}{\rA}
 \proofstepwidestar[3]{\mathsf{Fr}(p,x,y)\land\mathsf{Fr}(q,x,y)}
  {\FormulaRefAuto{B48FreeFor}[def]{3}}
 \proofstepwidestar[1,3]{(M,s\models p[y/x])\leftrightarrow
                   (M,s[x\mapsto s(y)]\models p)}
  {\FormulaRefAuto{B48SubstitutionProperty}[def]{1,\rAEa{4}}}
 \proofstepwidestar[2,3]{(M,s\models q[y/x])\leftrightarrow
                   (M,s[x\mapsto s(y)]\models q)}
  {\FormulaRefAuto{B48SubstitutionProperty}[def]{2,\rAEb{4}}}
 \proofstepwidestar[1,2,3]{(M,s\models(p\land q)[y/x])\leftrightarrow
                   (M,s[x\mapsto s(y)]\models p\land q)}
  {\rLRS{5,\rLRS{6,\FormulaRefAuto{P\leftrightarrow P}[thm]}},
   \FormulaRefAuto{B48Satisfaction}[def],
   \FormulaRefAuto{B48SubstitutionClauses}[thm]}
 \proofstepwidestar[1,2]{S(p\land q)}
  {\FormulaRefAuto{B48SubstitutionProperty}[def]{\rUI{\rRI{3,7}}}}
 \closeproofpartwideindR

 \proofpartwideindR[Quantor mit demselben Binder]{i=x\vdash
 \begin{aligned}
 (M,s\models(\exists v_i\,p)[y/x])\leftrightarrow{}\\
 (M,s[x\mapsto s(y)]\models\exists v_i\,p)
 \end{aligned}}
 {i=x\vdash
 ((M,s\models(\exists v_i\,p)[y/x])\leftrightarrow
 (M,s[x\mapsto s(y)]\models\exists v_i\,p))}
 [B48SubstitutionBound]
 \proofstepwidestar[1]{i=x}{\rA}
 \proofstepwidestar[1]{(\exists v_i\,p)[y/x]=\exists v_i\,p}
  {\FormulaRefAuto{B48SubstitutionClauses}[thm]{1}}
 \proofstepwidestar[1]{x\notin\operatorname{FV}(\exists v_i\,p)}
  {\FormulaRefAuto{B48FVClauses}[thm]{1}}
 \proofstepwidestar[1]{\begin{aligned}
 (M,s\models\exists v_i\,p)\leftrightarrow{}\\
 (M,s[x\mapsto s(y)]\models\exists v_i\,p)
 \end{aligned}}
  {\FormulaRefAuto{B48FreshUpdate}[thm]{3}}
 \proofstepwidestar[1]{\begin{aligned}
 (M,s\models(\exists v_i\,p)[y/x])\leftrightarrow{}\\
 (M,s[x\mapsto s(y)]\models\exists v_i\,p)
 \end{aligned}}
  {\rIE{\FormulaRefAuto{a=b\vdash b=a}[thm]{2},4}}
 \closeproofpartwideindR

 \proofpartwideindR[Quantor ohne moegliche Kollision]{S(p),\
 i\neq x,\ i\neq y,\ \mathsf{Fr}(p,x,y)\vdash
 \begin{aligned}
 (M,s\models(\exists v_i\,p)[y/x])\leftrightarrow{}\\
 (M,s[x\mapsto s(y)]\models\exists v_i\,p)
 \end{aligned}}
 {S(p),\ i\neq x,\ i\neq y,\ \mathsf{Fr}(p,x,y)\vdash
 ((M,s\models(\exists v_i\,p)[y/x])\leftrightarrow
 (M,s[x\mapsto s(y)]\models\exists v_i\,p))}
 [B48SubstitutionQuantifierFresh]
 \proofstepwidestar[1]{S(p)}{\rA}
 \proofstepwidestar[2]{i\neq x}{\rA}
 \proofstepwidestar[3]{i\neq y}{\rA}
 \proofstepwidestar[4]{\mathsf{Fr}(p,x,y)}{\rA}
 \proofstepwidestar[5]{a\in D}{\rA}
 \proofstepwidestar[5]{s[i\mapsto a]\in A_D}
  {\FormulaRefAuto{B48AssignmentUpdate}[thm]{5}}
 \proofstepwidestar[1,4,5]{\begin{aligned}
 (M,s[i\mapsto a]\models p[y/x])\leftrightarrow{}\\
 (M,s[i\mapsto a][x\mapsto s[i\mapsto a](y)]\models p)
 \end{aligned}}
  {\FormulaRefAuto{B48SubstitutionProperty}[def]{1,4,6}}
 \proofstepwidestar[2,3,5]{s[i\mapsto a][x\mapsto s[i\mapsto a](y)]
                         =s[x\mapsto s(y)][i\mapsto a]}
  {\FormulaRefAuto{B48UpdateCommute}[thm]{2,3,5}}
 \proofstepwidestar[1,2,3,4,5]{\begin{aligned}
 (M,s[i\mapsto a]\models p[y/x])\leftrightarrow{}\\
 (M,s[x\mapsto s(y)][i\mapsto a]\models p)
 \end{aligned}}
  {\rIE{8,7}}
 \proofstepwidestar[1,2,3,4]{\begin{aligned}
 (\exists a\in D\;(M,s[i\mapsto a]\models p[y/x]))\leftrightarrow{}\\
 (\exists a\in D\;(M,s[x\mapsto s(y)][i\mapsto a]\models p))
 \end{aligned}}
  {\FormulaRefAuto{B48ExistsIff}[thm]{\rUI{\rRI{5,9}}}}
 \proofstepwidestar[1,2,3,4]{\begin{aligned}
 (M,s\models(\exists v_i\,p)[y/x])\leftrightarrow{}\\
 (M,s[x\mapsto s(y)]\models\exists v_i\,p)
 \end{aligned}}
  {\FormulaRefAuto{B48Satisfaction}[def]{10},
   \FormulaRefAuto{B48SubstitutionClauses}[thm]{2}}
 \closeproofpartwideindR

 \proofpartwideindR[Quantor bei leerer Ersetzungsstelle]{
 x\notin\operatorname{FV}(p)\vdash
 \begin{aligned}
 (M,s\models(\exists v_i\,p)[y/x])\leftrightarrow{}\\
 (M,s[x\mapsto s(y)]\models\exists v_i\,p)
 \end{aligned}}
 {x\notin\operatorname{FV}(p)\vdash
 ((M,s\models(\exists v_i\,p)[y/x])\leftrightarrow
 (M,s[x\mapsto s(y)]\models\exists v_i\,p))}
 [B48SubstitutionQuantifierVacuous]
 \proofstepwidestar[1]{x\notin\operatorname{FV}(p)}{\rA}
 \proofstepwidestar[1]{x\notin\operatorname{FV}(\exists v_i\,p)}
  {\FormulaRefAuto{B48FVClauses}[thm]{1}}
 \proofstepwidestar[1]{(\exists v_i\,p)[y/x]=\exists v_i\,p}
  {\FormulaRefAuto{B48SubstitutionVacuous}[thm]{2}}
 \proofstepwidestar[1]{\begin{aligned}
 (M,s\models\exists v_i\,p)\leftrightarrow{}\\
 (M,s[x\mapsto s(y)]\models\exists v_i\,p)
 \end{aligned}}
  {\FormulaRefAuto{B48FreshUpdate}[thm]{2}}
 \proofstepwidestar[1]{\begin{aligned}
 (M,s\models(\exists v_i\,p)[y/x])\leftrightarrow{}\\
 (M,s[x\mapsto s(y)]\models\exists v_i\,p)
 \end{aligned}}
  {\rIE{\FormulaRefAuto{a=b\vdash b=a}[thm]{3},4}}
 \closeproofpartwideindR

 \proofpartwide{Vollstaendiger Quantorenschritt}
 \proofstepwidestar[1]{S(p)}{\rA}
 \proofstepwidestar[2]{\mathsf{Fr}(\exists v_i\,p,x,y)}{\rA}
 \proofstepwidestar[2]{i=x\lor
  (\mathsf{Fr}(p,x,y)\land(i\neq y\lor x\notin\operatorname{FV}(p)))}
  {\FormulaRefAuto{B48FreeFor}[def]{2}}
 \proofstepwidestar[]{i=x\lor i\neq x}{\FormulaRefAuto{P\lor\neg P}[thm]}
 \proofstepwidestar[5]{i=x}{\rA}
 \proofstepwidestar[5]{\begin{aligned}
 (M,s\models(\exists v_i\,p)[y/x])\leftrightarrow{}\\
 (M,s[x\mapsto s(y)]\models\exists v_i\,p)
 \end{aligned}}
  {\FormulaRefAuto{B48SubstitutionBound}[thm]{5}}
 \proofstepwidestar[7]{i\neq x}{\rA}
 \proofstepwidestar[2,7]{\mathsf{Fr}(p,x,y)\land
             (i\neq y\lor x\notin\operatorname{FV}(p))}
  {\FormulaRefAuto{P\lor Q,\neg P\vdash Q}[thm]{3,7}}
 \proofstepwidestar[9]{i\neq y}{\rA}
 \proofstepwidestar[1,2,7,9]{\begin{aligned}
 (M,s\models(\exists v_i\,p)[y/x])\leftrightarrow{}\\
 (M,s[x\mapsto s(y)]\models\exists v_i\,p)
 \end{aligned}}
  {\FormulaRefAuto{B48SubstitutionQuantifierFresh}[thm]{1,7,9,\rAEa{8}}}
 \proofstepwidestar[11]{x\notin\operatorname{FV}(p)}{\rA}
 \proofstepwidestar[11]{\begin{aligned}
 (M,s\models(\exists v_i\,p)[y/x])\leftrightarrow{}\\
 (M,s[x\mapsto s(y)]\models\exists v_i\,p)
 \end{aligned}}
  {\FormulaRefAuto{B48SubstitutionQuantifierVacuous}[thm]{11}}
 \proofstepwidestar[1,2,7]{\begin{aligned}
 (M,s\models(\exists v_i\,p)[y/x])\leftrightarrow{}\\
 (M,s[x\mapsto s(y)]\models\exists v_i\,p)
 \end{aligned}}
  {\rOE{\rAEb{8},9:10,11:12}}
 \proofstepwidestar[1,2]{\begin{aligned}
 (M,s\models(\exists v_i\,p)[y/x])\leftrightarrow{}\\
 (M,s[x\mapsto s(y)]\models\exists v_i\,p)
 \end{aligned}}
  {\rOE{4,5:6,7:13}}
 \proofstepwidestar[1]{S(\exists v_i\,p)}
  {\FormulaRefAuto{B48SubstitutionProperty}[def]{\rUI{\rRI{2,14}}}}
 \closeproofpartwide
\end{tabproofsplitwide}
Die Atom-, Negations-, Konjunktions- und Quantorentabellen zeigen,
dass \(\{p\in\mathcal F\mid S(p)\}\) formelabgeschlossen ist.
Aus \FormulaRefAuto{B48FormulaInduction}[thm] folgt
\(\forall p\in\mathcal F\;S(p)\); Einsetzen von \(p,x,y,s\)
und Anwenden auf \(\mathsf{Fr}(p,x,y)\) ergibt den Satz.

\noindent
Beispielsweise ist \(y\) fuer \(x\) in \(x=y\) frei, aber fuer
\(x\) in \(\exists y\,\neg(x=y)\) nicht frei. Der zweite Fall darf
nicht durch stillschweigende Anwendung des Substitutionslemmas
behandelt werden.


\subsection{Belegungsfreie Erfuellung und Widerspruch}

Die folgende Abkuerzung gilt fuer Saetze, also Formeln ohne freie
Variablen. Die Nichtleerheit des Bereichs wird ausdruecklich
mitgefuehrt. Dadurch braucht die Existenz einer Belegung in den
folgenden Beweisen nicht stillschweigend vorausgesetzt zu werden.
Die verwendete Erfuellungsrelation ist durch
\FormulaRefAuto{B48SatisfactionExists}[thm] konstruiert; ihre
Schreibweise und ihre gerechtfertigten definierenden Klauseln stehen in
\FormulaRefAuto{B48Satisfaction}[def].

\FormulaDefDeltaK[Erfuellung von Saetzen und Theorien]
{\begin{aligned}
 M\models\varphi&\coloneqq
 \bigl(\exists d\,(d\in D)\bigr)\land
 \forall s\,\bigl((s\colon\operatorname{Var}\to D)
          \rightarrow(M,s\models\varphi)\bigr),\\
 M\models T&\coloneqq
 \forall\theta\,\bigl(\theta\in T\rightarrow(M\models\theta)\bigr)
\end{aligned}}
{B48SentenceSatisfaction}{
 \DeltaRow{Struktur}{M=(D,E)}
 \DeltaRow{Satz}{\varphi}
 \DeltaRow{Theorie}{T\subseteq\operatorname{Sent}(\mathcal L_\in)}
}

\FormulaThmDeltaK[Eine erfuellte Aussage liefert eine Belegung]
{M\models\varphi\vdash
 \exists s\,(s\colon\operatorname{Var}\to D)}
{B48SentenceValuation}{
 \DeltaRow{Struktur}{M=(D,E)}
 \DeltaRow{Satz}{\varphi}
}
Der im Schluss verwendete Zeuge \(a\) ist frisch fuer \(M,D,E,\varphi\)
und die ausserhalb des Zeugenbeweises offenen Voraussetzungen.
\begin{tabproofsplitwide}
 \proofpartwideindR[Konstante Belegung]{%
   d\in D\vdash\exists s\,(s\colon\operatorname{Var}\to D)
 }{%
   d\in D\vdash\exists s\,(s\colon\operatorname{Var}\to D)
 }[B48ConstantValuationWitness]
  \proofstepwidestar[1]{d\in D}{\rA}
  \proofstepwidestar[1]{%
   \ConstFun{\operatorname{Var}}{D}{d}
       \colon\operatorname{Var}\to D
  }{\FormulaRefAuto{b\in B\vdash \ConstFun{A}{B}{b}\colon A\to B}{1}}
  \proofstepwidestar[1]{%
   \exists s\,(s\colon\operatorname{Var}\to D)
  }{\rEI{2}}
 \closeproofpartwideindR

 \proofpartwide{Schluss}
  \proofstepwidestar[1]{M\models\varphi}{\rA}
  \proofstepwidestar[1]{%
   \bigl(\exists d\,(d\in D)\bigr)\land
   \forall s\,\bigl((s\colon\operatorname{Var}\to D)
           \rightarrow(M,s\models\varphi)\bigr)
  }{\FormulaRefAuto{B48SentenceSatisfaction}{1}}
  \proofstepwidestar[1]{\exists d\,(d\in D)}{\rAEa{2}}
  \proofstepwidestar[4]{a\in D}{\rA}
  \proofstepwidestar[4]{%
   \exists s\,(s\colon\operatorname{Var}\to D)
  }{\FormulaRefAuto{B48ConstantValuationWitness}{4}}
  \proofstepwidestar[1]{%
   \exists s\,(s\colon\operatorname{Var}\to D)
  }{\rEE{3,4:5}}
 \closeproofpartwide
\end{tabproofsplitwide}

\FormulaThmDeltaK[Satzwahrheit bei einer Belegung]
{M\models\varphi\dsep s\colon\operatorname{Var}\to D
 \vdash M,s\models\varphi}
{B48SentenceInstance}{
 \DeltaRow{Struktur}{M=(D,E)}
 \DeltaRow{Satz}{\varphi}
 \DeltaRow{Menge}{s}
}
\begin{tabproof}
 \proofstep{1}{M\models\varphi}{\rA}
 \proofstep{2}{s\colon\operatorname{Var}\to D}{\rA}
 \proofstep{1}{%
  \bigl(\exists d\,(d\in D)\bigr)\land
  \forall t\,\bigl((t\colon\operatorname{Var}\to D)
          \rightarrow(M,t\models\varphi)\bigr)
 }{\FormulaRefAuto{B48SentenceSatisfaction}{1}}
 \proofstep{1}{%
  \forall t\,\bigl((t\colon\operatorname{Var}\to D)
          \rightarrow(M,t\models\varphi)\bigr)
 }{\rAEb{3}}
 \proofstep{1}{%
  (s\colon\operatorname{Var}\to D)\rightarrow(M,s\models\varphi)
 }{\rUE{4}}
 \proofstep{1,2}{M,s\models\varphi}{\rRE{5,2}}
\end{tabproof}

\FormulaThmDeltaK[Keine Struktur erfuellt den Widerspruch]
{\neg(M\models\bot)}
{B48SentenceNotBottom}{
 \DeltaRow{Struktur}{M=(D,E)}
}
Im folgenden Beweis ist \(t\) frisch fuer \(M,D,E\) und die
ausserhalb des Beweises offenen Annahmen.
\begin{tabproof}
 \proofstep{1}{M\models\bot}{\rA}
 \proofstep{1}{\exists s\,(s\colon\operatorname{Var}\to D)}
  {\FormulaRefAuto{B48SentenceValuation}{1}}
 \proofstep{3}{t\colon\operatorname{Var}\to D}{\rA}
 \proofstep{1,3}{M,t\models\bot}
  {\FormulaRefAuto{B48SentenceInstance}{1,3}}
 \proofstep{1,3}{\bot}{\FormulaRefAuto{B48Satisfaction}{4}}
 \proofstep{1}{\bot}{\rEE{2,3:5}}
 \proofstep{}{\neg(M\models\bot)}{\rCI{1,6}}
\end{tabproof}

\FormulaThmDeltaK[Ein Satz und seine Negation sind nicht gemeinsam wahr]
{M\models\varphi\dsep M\models\neg\varphi\vdash\bot}
{B48SentenceContradiction}{
 \DeltaRow{Struktur}{M=(D,E)}
 \DeltaRow{Satz}{\varphi}
}
Hier ist \(t\) frisch fuer \(M,D,E,\varphi\) und die ausserhalb
des Beweises offenen Annahmen.
\begin{tabproof}
 \proofstep{1}{M\models\varphi}{\rA}
 \proofstep{2}{M\models\neg\varphi}{\rA}
 \proofstep{1}{\exists s\,(s\colon\operatorname{Var}\to D)}
  {\FormulaRefAuto{B48SentenceValuation}{1}}
 \proofstep{4}{t\colon\operatorname{Var}\to D}{\rA}
 \proofstep{1,4}{M,t\models\varphi}
  {\FormulaRefAuto{B48SentenceInstance}{1,4}}
 \proofstep{2,4}{M,t\models\neg\varphi}
  {\FormulaRefAuto{B48SentenceInstance}{2,4}}
 \proofstep{1,2,4}{\bot}
  {\FormulaRefAuto{B48SatContradiction}{5,6}}
 \proofstep{1,2}{\bot}{\rEE{3,4:7}}
\end{tabproof}


% Korrektheit fuer die konstantenfreie Sprache von Band 48.
\section{Korrektheit}

Die Objektsprache bleibt \(\mathcal L_\in\).
In der Metatheorie duerfen die Regeln aus Band~1 benutzt werden,
waehrend hier ihre Wirkung auf \emph{codierte} Herleitungen untersucht
wird. Bei der All-Einfuehrung ist die Variable in keiner offenen
Annahme frei; beim klassischen Widerspruchsschluss wird die
angenommene Negation entladen. Die in Band~1 erklaerten
Abhaengigkeitsmengen werden unten ausdruecklich mitgefuehrt.

\FormulaDefDeltaK[Abkuerzungen in der codierten Objektsprache]
{\begin{aligned}
 p\lor q&\coloneq\neg(\neg p\land\neg q),\\
 p\to q&\coloneq\neg(p\land\neg q),\\
 p\leftrightarrow q&\coloneq(p\to q)\land(q\to p),\\
 \forall x\,p&\coloneq\neg\exists x\,\neg p.
\end{aligned}}
{B48LogicalAbbreviations}{
 \DeltaRow{Formeln und Variable}{p,q\in\mathcal F,\quad x\in\mathbb N}
}
Dies praezisiert die zuvor verwendeten ueblichen Abkuerzungen.
Ihre semantischen Klauseln muessen aus der bereits konstruierten
Erfuellungsrelation folgen.

\FormulaThmDeltaK[Semantik der abgeleiteten logischen Zeichen]
{\begin{aligned}
 (M,s\models p\lor q)&\leftrightarrow
       ((M,s\models p)\lor(M,s\models q)),\\
 (M,s\models p\to q)&\leftrightarrow
       ((M,s\models p)\to(M,s\models q)),\\
 (M,s\models p\leftrightarrow q)&\leftrightarrow
       ((M,s\models p)\leftrightarrow(M,s\models q)),\\
 (M,s\models\forall x\,p)&\leftrightarrow
       \forall a\in D\;(M,s[x\mapsto a]\models p).
\end{aligned}}
{B48DerivedSemanticClauses}{
 \DeltaRow{Struktur und Belegung}{\mathsf{Str}(M),\quad s\in A_D}
 \DeltaRow{Formeln und Variable}{p,q\in\mathcal F,\quad x\in\mathbb N}
}
In den folgenden Tabellen kuerzen \(P\) und \(Q\) ausschliesslich
die metasprachlichen Aussagen \(M,s\models p\) und \(M,s\models q\)
ab.
\begin{tabproofsplitwide}
 \proofpartwideindR[Disjunktion]{(M,s\models p\lor q)\leftrightarrow(P\lor Q)}
  {(M,s\models p\lor q)\leftrightarrow(P\lor Q)}[B48SatOrClause]
 \proofstepwidestar[]{(M,s\models p\lor q)\leftrightarrow
                          \neg(\neg P\land\neg Q)}
  {\FormulaRefAuto{B48LogicalAbbreviations}[def],
   \FormulaRefAuto{B48Satisfaction}[def]}
 \proofstepwidestar[]{\neg(\neg P\land\neg Q)
                          \leftrightarrow(\neg\neg P\lor\neg\neg Q)}
  {\FormulaRefAuto{\neg(P\land Q)\eqvdash\neg P\lor\neg Q}[thm]}
 \proofstepwidestar[]{(\neg\neg P\lor\neg\neg Q)\leftrightarrow(P\lor Q)}
  {\FormulaRefAuto{P\lor Q\eqvdash\neg\neg P\lor\neg\neg Q}[thm],
   \FormulaRefAuto{B48IffSymmetry}[thm]}
 \proofstepwidestar[]{(M,s\models p\lor q)\leftrightarrow(P\lor Q)}
  {\rLRS{3,\rLRS{2,1}}}
 \closeproofpartwideindR

 \proofpartwideindR[Implikation]{(M,s\models p\to q)\leftrightarrow(P\to Q)}
  {(M,s\models p\to q)\leftrightarrow(P\to Q)}[B48SatImpClause]
 \proofstepwidestar[]{(M,s\models p\to q)\leftrightarrow\neg(P\land\neg Q)}
  {\FormulaRefAuto{B48LogicalAbbreviations}[def],
   \FormulaRefAuto{B48Satisfaction}[def]}
 \proofstepwidestar[]{\neg(P\land\neg Q)\leftrightarrow(\neg P\lor\neg\neg Q)}
  {\FormulaRefAuto{\neg(P\land Q)\eqvdash\neg P\lor\neg Q}[thm]}
 \proofstepwidestar[]{Q\leftrightarrow\neg\neg Q}
  {\FormulaRefAuto{P\eqvdash\neg\neg P}[thm]}
 \proofstepwidestar[]{\neg(P\land\neg Q)\leftrightarrow(\neg P\lor Q)}
  {\rLRS{3,2}}
 \proofstepwidestar[]{(P\to Q)\leftrightarrow(\neg P\lor Q)}
  {\FormulaRefAuto{P\to Q\eqvdash\neg P\lor Q}[thm]}
 \proofstepwidestar[]{(M,s\models p\to q)\leftrightarrow(P\to Q)}
  {\rLRS{5,\rLRS{4,1}}}
 \closeproofpartwideindR

 \proofpartwideindR[Bikonditional]{
 (M,s\models p\leftrightarrow q)\leftrightarrow(P\leftrightarrow Q)}
 {(M,s\models p\leftrightarrow q)\leftrightarrow(P\leftrightarrow Q)}
 [B48SatIffClause]
 \proofstepwidestar[]{\begin{aligned}
 (M,s\models p\leftrightarrow q)\leftrightarrow{}\\
 ((M,s\models p\to q)\land(M,s\models q\to p))
 \end{aligned}}
  {\FormulaRefAuto{B48LogicalAbbreviations}[def],
   \FormulaRefAuto{B48Satisfaction}[def]}
 \proofstepwidestar[]{(M,s\models p\to q)\leftrightarrow(P\to Q)}
  {\FormulaRefAuto{B48SatImpClause}[thm]}
 \proofstepwidestar[]{(M,s\models q\to p)\leftrightarrow(Q\to P)}
  {\FormulaRefAuto{B48SatImpClause}[thm]}
 \proofstepwidestar[]{(M,s\models p\leftrightarrow q)
                     \leftrightarrow((P\to Q)\land(Q\to P))}
  {\rLRS{2,\rLRS{3,1}}}
 \proofstepwidestar[]{(P\leftrightarrow Q)\leftrightarrow
                       ((P\to Q)\land(Q\to P))}
  {\FormulaRefAuto{P\leftrightarrow Q\eqvdash
                  (P\to Q)\land(Q\to P)}[thm]}
 \proofstepwidestar[]{(M,s\models p\leftrightarrow q)\leftrightarrow
                       (P\leftrightarrow Q)}
  {\rLRS{5,4}}
 \closeproofpartwideindR

 \proofpartwideindR[Allquantor]{(M,s\models\forall x\,p)\leftrightarrow
                \forall a\in D\;(M,s[x\mapsto a]\models p)}
 {(M,s\models\forall x\,p)\leftrightarrow
                \forall a\in D\;(M,s[x\mapsto a]\models p)}
 [B48SatAllClause]
 \proofstepwidestar[]{\begin{aligned}
 (M,s\models\forall x\,p)\leftrightarrow{}\\
 \neg\exists a\in D\;\neg(M,s[x\mapsto a]\models p)
 \end{aligned}}
  {\FormulaRefAuto{B48LogicalAbbreviations}[def],
   \FormulaRefAuto{B48Satisfaction}[def]}
 \proofstepwidestar[]{\begin{aligned}
 \forall a\;(a\in D\to(M,s[x\mapsto a]\models p))\leftrightarrow{}\\
 \neg\exists a\;(a\in D\land\neg(M,s[x\mapsto a]\models p))
 \end{aligned}}
  {\FormulaRefAuto{\forall x(P(x)\to Q(x))\eqvdash
                  \neg\exists x(P(x)\land\neg Q(x))}[thm]}
 \proofstepwidestar[]{(M,s\models\forall x\,p)\leftrightarrow
                \forall a\in D\;(M,s[x\mapsto a]\models p)}
  {\rLRS{2,1}}
 \closeproofpartwideindR
\end{tabproofsplitwide}

\FormulaDefDeltaK[Erfuellte Annahmen und Folgerung in einer Struktur]
{\begin{aligned}
 M,s\models\Gamma&\coloneq
       \forall p\in\Gamma\;(M,s\models p),\\
 \Gamma\models_M p&\coloneq
       \forall s\in A_D\;((M,s\models\Gamma)\to(M,s\models p)),\\
 x\mathrel{\#}\Gamma&\coloneq
       \forall q\in\Gamma\;(x\notin\operatorname{FV}(q)).
\end{aligned}}
{B48SemanticConsequence}{
 \DeltaRow{Struktur}{\mathsf{Str}(M)}
 \DeltaRow{Formeln und Annahmen}{p\in\mathcal F,\quad\Gamma\subseteq\mathcal F}
 \DeltaRow{Variable}{x\in\mathbb N}
}
Hier duerfen Annahmen freie Variablen haben. Die Quantifikation ueber
alle Belegungen ist Bestandteil der Folgerungsrelation.

\FormulaThmDeltaK[Anwendung einer semantischen Folgerung]
{\Gamma\models_M p,\ M,s\models\Gamma\ \vdash\ M,s\models p}
{B48ConsequenceInstance}{
 \DeltaRow{Struktur und Belegung}{\mathsf{Str}(M),\quad s\in A_D}
 \DeltaRow{Annahmen und Formel}{\Gamma\subseteq\mathcal F,\quad p\in\mathcal F}
}
\begin{tabproof}
 \proofstep{1}{\Gamma\models_M p}{\rA}
 \proofstep{2}{M,s\models\Gamma}{\rA}
 \proofstep{1}{\forall t\in A_D\;
    ((M,t\models\Gamma)\to(M,t\models p))}
  {\FormulaRefAuto{B48SemanticConsequence}[def]{1}}
 \proofstep{1}{(M,s\models\Gamma)\to(M,s\models p)}{\rUE{3}}
 \proofstep{1,2}{M,s\models p}{\rRE{4,2}}
\end{tabproof}

\FormulaThmDeltaK[Vereinigung erfuellter Annahmen]
{(M,s\models\Gamma\cup\Delta)\leftrightarrow
 ((M,s\models\Gamma)\land(M,s\models\Delta))}
{B48ContextUnion}{
 \DeltaRow{Struktur und Belegung}{\mathsf{Str}(M),\quad s\in A_D}
 \DeltaRow{Annahmenmengen}{\Gamma,\Delta\subseteq\mathcal F}
}
\begin{tabproof}
 \proofstep{1}{M,s\models\Gamma\cup\Delta}{\rA}
 \proofstep{2}{p\in\Gamma}{\rA}
 \proofstep{2}{p\in\Gamma\cup\Delta}
  {\FormulaRefAuto{x\in A\cup B\eqvdash x\in A\lor x\in B}[thm]{\rOIa{2}}}
 \proofstep{1,2}{M,s\models p}
  {\FormulaRefAuto{B48SemanticConsequence}[def]{1,3}}
 \proofstep{1}{M,s\models\Gamma}
  {\FormulaRefAuto{B48SemanticConsequence}[def]{\rUI{\rRI{2,4}}}}
 \proofstep{6}{p\in\Delta}{\rA}
 \proofstep{6}{p\in\Gamma\cup\Delta}
  {\FormulaRefAuto{x\in A\cup B\eqvdash x\in A\lor x\in B}[thm]{\rOIb{6}}}
 \proofstep{1,6}{M,s\models p}
  {\FormulaRefAuto{B48SemanticConsequence}[def]{1,7}}
 \proofstep{1}{M,s\models\Delta}
  {\FormulaRefAuto{B48SemanticConsequence}[def]{\rUI{\rRI{6,8}}}}
 \proofstep{1}{(M,s\models\Gamma)\land(M,s\models\Delta)}
  {\rAI{5,9}}
 \proofstep{}{(M,s\models\Gamma\cup\Delta)\to
        ((M,s\models\Gamma)\land(M,s\models\Delta))}
  {\rRI{1,10}}
 \proofstep{12}{(M,s\models\Gamma)\land(M,s\models\Delta)}{\rA}
 \proofstep{13}{p\in\Gamma\cup\Delta}{\rA}
 \proofstep{13}{p\in\Gamma\lor p\in\Delta}
  {\FormulaRefAuto{x\in A\cup B\eqvdash x\in A\lor x\in B}[thm]{13}}
 \proofstep{15}{p\in\Gamma}{\rA}
 \proofstep{12,15}{M,s\models p}
  {\FormulaRefAuto{B48SemanticConsequence}[def]{\rAEa{12},15}}
 \proofstep{17}{p\in\Delta}{\rA}
 \proofstep{12,17}{M,s\models p}
  {\FormulaRefAuto{B48SemanticConsequence}[def]{\rAEb{12},17}}
 \proofstep{12,13}{M,s\models p}{\rOE{14,15:16,17:18}}
 \proofstep{12}{M,s\models\Gamma\cup\Delta}
  {\FormulaRefAuto{B48SemanticConsequence}[def]{\rUI{\rRI{13,19}}}}
 \proofstep{}{((M,s\models\Gamma)\land(M,s\models\Delta))\to
                 (M,s\models\Gamma\cup\Delta)}
  {\rRI{12,20}}
 \proofstep{}{(M,s\models\Gamma\cup\Delta)\leftrightarrow
       ((M,s\models\Gamma)\land(M,s\models\Delta))}
  {\rLRI{11,21}}
\end{tabproof}

\FormulaThmDeltaK[Hinzunahme einer erfuellten Annahme]
{(M,s\models\Gamma\cup\{p\})\leftrightarrow
 ((M,s\models\Gamma)\land(M,s\models p))}
{B48ContextAssumption}{
 \DeltaRow{Struktur und Belegung}{\mathsf{Str}(M),\quad s\in A_D}
 \DeltaRow{Annahmen und Formel}{\Gamma\subseteq\mathcal F,\quad p\in\mathcal F}
}
\begin{tabproof}
 \proofstep{1}{M,s\models\{p\}}{\rA}
 \proofstep{1}{M,s\models p}
  {\FormulaRefAuto{B48SemanticConsequence}[def]{1}}
 \proofstep{}{(M,s\models\{p\})\to(M,s\models p)}{\rRI{1,2}}
 \proofstep{4}{M,s\models p}{\rA}
 \proofstep{5}{q\in\{p\}}{\rA}
 \proofstep{5}{q=p}{\FormulaRefAuto{x\in\{a\}\eqvdash x=a}[thm]{5}}
 \proofstep{4,5}{M,s\models q}
  {\rIE{\FormulaRefAuto{a=b\vdash b=a}[thm]{6},4}}
 \proofstep{4}{M,s\models\{p\}}
  {\FormulaRefAuto{B48SemanticConsequence}[def]{\rUI{\rRI{5,7}}}}
 \proofstep{}{(M,s\models p)\to(M,s\models\{p\})}{\rRI{4,8}}
 \proofstep{}{(M,s\models\{p\})\leftrightarrow(M,s\models p)}
  {\rLRI{3,9}}
 \proofstep{}{(M,s\models\Gamma\cup\{p\})\leftrightarrow
    ((M,s\models\Gamma)\land(M,s\models p))}
  {\rLRS{10,\FormulaRefAuto{B48ContextUnion}[thm]}}
\end{tabproof}

\FormulaThmDeltaK[Belegungswechsel in variablenfreien Annahmen]
{x\mathrel{\#}\Gamma,\ M,s\models\Gamma\ \vdash\
                         M,s[x\mapsto a]\models\Gamma}
{B48ContextFresh}{
 \DeltaRow{Struktur und Belegung}{\mathsf{Str}(M),\quad s\in A_D}
 \DeltaRow{Annahmen, Variable und Wert}{\Gamma\subseteq\mathcal F,\quad
                             x\in\mathbb N,\quad a\in D}
}
\begin{tabproof}
 \proofstep{1}{x\mathrel{\#}\Gamma}{\rA}
 \proofstep{2}{M,s\models\Gamma}{\rA}
 \proofstep{3}{p\in\Gamma}{\rA}
 \proofstep{1,3}{x\notin\operatorname{FV}(p)}
  {\FormulaRefAuto{B48SemanticConsequence}[def]{1,3}}
 \proofstep{2,3}{M,s\models p}
  {\FormulaRefAuto{B48SemanticConsequence}[def]{2,3}}
 \proofstep{1,3}{(M,s\models p)\leftrightarrow
                   (M,s[x\mapsto a]\models p)}
  {\FormulaRefAuto{B48FreshUpdate}[thm]{4}}
 \proofstep{1,2,3}{M,s[x\mapsto a]\models p}{\rLREa{6,5}}
 \proofstep{1,2}{M,s[x\mapsto a]\models\Gamma}
  {\FormulaRefAuto{B48SemanticConsequence}[def]{\rUI{\rRI{3,7}}}}
\end{tabproof}

\subsection{Einzeln nachgewiesene Regelfaelle}

In den folgenden Beweisen wird \(s\in A_D\) beliebig gewaehlt.
Die abschliessende All-Einfuehrung bindet diese Wahl. Die Formeln
\(p,q,r\) und die Annahmenmengen sind stets im zuvor angegebenen
Formelbereich.

\FormulaThmDeltaK[Semantischer Regelfall der ersten Konjunktionsbeseitigung]
{\Gamma\models_M p\land q\ \vdash\ \Gamma\models_M p}
{B48RuleAndLeft}{
 \DeltaRow{Struktur}{\mathsf{Str}(M)}
 \DeltaRow{Annahmen}{\Gamma,\Delta\subseteq\mathcal F}
 \DeltaRow{Formeln}{p,q,r\in\mathcal F}
}
\begin{tabproofsplitwide}
 \proofpartwide{Beweis}
 \proofstepwidestar[1]{\Gamma\models_M p\land q}
  {\rA}
 \proofstepwidestar[2]{M,s\models \Gamma}
  {\rA}
 \proofstepwidestar[1,2]{M,s\models p\land q}
  {\FormulaRefAuto{B48ConsequenceInstance}[thm]{1,2}}
 \proofstepwidestar[1,2]{M,s\models p}
  {\FormulaRefAuto{B48SatAndLeft}[thm]{3}}
 \proofstepwidestar[1]{(M,s\models \Gamma)\to(M,s\models p)}
  {\rRI{2,4}}
 \proofstepwidestar[1]{\Gamma\models_M p}
  {\FormulaRefAuto{B48SemanticConsequence}[def]{\rUI{5}}}
 \closeproofpartwide
\end{tabproofsplitwide}

\FormulaThmDeltaK[Semantischer Regelfall der zweiten Konjunktionsbeseitigung]
{\Gamma\models_M p\land q\ \vdash\ \Gamma\models_M q}
{B48RuleAndRight}{
 \DeltaRow{Struktur}{\mathsf{Str}(M)}
 \DeltaRow{Annahmen}{\Gamma,\Delta\subseteq\mathcal F}
 \DeltaRow{Formeln}{p,q,r\in\mathcal F}
}
\begin{tabproofsplitwide}
 \proofpartwide{Beweis}
 \proofstepwidestar[1]{\Gamma\models_M p\land q}
  {\rA}
 \proofstepwidestar[2]{M,s\models \Gamma}
  {\rA}
 \proofstepwidestar[1,2]{M,s\models p\land q}
  {\FormulaRefAuto{B48ConsequenceInstance}[thm]{1,2}}
 \proofstepwidestar[1,2]{M,s\models q}
  {\FormulaRefAuto{B48SatAndRight}[thm]{3}}
 \proofstepwidestar[1]{(M,s\models \Gamma)\to(M,s\models q)}
  {\rRI{2,4}}
 \proofstepwidestar[1]{\Gamma\models_M q}
  {\FormulaRefAuto{B48SemanticConsequence}[def]{\rUI{5}}}
 \closeproofpartwide
\end{tabproofsplitwide}

\FormulaThmDeltaK[Semantische Disjunktionseinfuehrung links]
{\Gamma\models_M p\ \vdash\ \Gamma\models_M p\lor q}
{B48RuleOrLeft}{
 \DeltaRow{Struktur}{\mathsf{Str}(M)}
 \DeltaRow{Annahmen}{\Gamma,\Delta\subseteq\mathcal F}
 \DeltaRow{Formeln}{p,q,r\in\mathcal F}
}
\begin{tabproofsplitwide}
 \proofpartwide{Beweis}
 \proofstepwidestar[1]{\Gamma\models_M p}
  {\rA}
 \proofstepwidestar[2]{M,s\models \Gamma}
  {\rA}
 \proofstepwidestar[1,2]{M,s\models p}
  {\FormulaRefAuto{B48ConsequenceInstance}[thm]{1,2}}
 \proofstepwidestar[1,2]{(M,s\models p)\lor(M,s\models q)}
  {\rOIa{3}}
 \proofstepwidestar[1,2]{M,s\models p\lor q}
  {\rLREb{\FormulaRefAuto{B48SatOrClause}[thm],4}}
 \proofstepwidestar[1]{(M,s\models \Gamma)\to(M,s\models p\lor q)}
  {\rRI{2,5}}
 \proofstepwidestar[1]{\Gamma\models_M p\lor q}
  {\FormulaRefAuto{B48SemanticConsequence}[def]{\rUI{6}}}
 \closeproofpartwide
\end{tabproofsplitwide}

\FormulaThmDeltaK[Semantische Disjunktionseinfuehrung rechts]
{\Gamma\models_M q\ \vdash\ \Gamma\models_M p\lor q}
{B48RuleOrRight}{
 \DeltaRow{Struktur}{\mathsf{Str}(M)}
 \DeltaRow{Annahmen}{\Gamma,\Delta\subseteq\mathcal F}
 \DeltaRow{Formeln}{p,q,r\in\mathcal F}
}
\begin{tabproofsplitwide}
 \proofpartwide{Beweis}
 \proofstepwidestar[1]{\Gamma\models_M q}
  {\rA}
 \proofstepwidestar[2]{M,s\models \Gamma}
  {\rA}
 \proofstepwidestar[1,2]{M,s\models q}
  {\FormulaRefAuto{B48ConsequenceInstance}[thm]{1,2}}
 \proofstepwidestar[1,2]{(M,s\models p)\lor(M,s\models q)}
  {\rOIb{3}}
 \proofstepwidestar[1,2]{M,s\models p\lor q}
  {\rLREb{\FormulaRefAuto{B48SatOrClause}[thm],4}}
 \proofstepwidestar[1]{(M,s\models \Gamma)\to(M,s\models p\lor q)}
  {\rRI{2,5}}
 \proofstepwidestar[1]{\Gamma\models_M p\lor q}
  {\FormulaRefAuto{B48SemanticConsequence}[def]{\rUI{6}}}
 \closeproofpartwide
\end{tabproofsplitwide}

\FormulaThmDeltaK[Semantische Bikonditionalbeseitigung vorwaerts]
{\Gamma\models_M p\leftrightarrow q\ \vdash\ \Gamma\models_M p\to q}
{B48RuleIffLeft}{
 \DeltaRow{Struktur}{\mathsf{Str}(M)}
 \DeltaRow{Annahmen}{\Gamma,\Delta\subseteq\mathcal F}
 \DeltaRow{Formeln}{p,q,r\in\mathcal F}
}
\begin{tabproofsplitwide}
 \proofpartwide{Beweis}
 \proofstepwidestar[1]{\Gamma\models_M p\leftrightarrow q}
  {\rA}
 \proofstepwidestar[2]{M,s\models \Gamma}
  {\rA}
 \proofstepwidestar[1,2]{M,s\models p\leftrightarrow q}
  {\FormulaRefAuto{B48ConsequenceInstance}[thm]{1,2}}
 \proofstepwidestar[1,2]{(M,s\models p)\leftrightarrow(M,s\models q)}
  {\rLREa{\FormulaRefAuto{B48SatIffClause}[thm],3}}
 \proofstepwidestar[1,2]{(M,s\models p)\to(M,s\models q)}
  {\rLREa{4}}
 \proofstepwidestar[1,2]{M,s\models p\to q}
  {\rLREb{\FormulaRefAuto{B48SatImpClause}[thm],5}}
 \proofstepwidestar[1]{(M,s\models \Gamma)\to(M,s\models p\to q)}
  {\rRI{2,6}}
 \proofstepwidestar[1]{\Gamma\models_M p\to q}
  {\FormulaRefAuto{B48SemanticConsequence}[def]{\rUI{7}}}
 \closeproofpartwide
\end{tabproofsplitwide}

\FormulaThmDeltaK[Semantische Bikonditionalbeseitigung rueckwaerts]
{\Gamma\models_M p\leftrightarrow q\ \vdash\ \Gamma\models_M q\to p}
{B48RuleIffRight}{
 \DeltaRow{Struktur}{\mathsf{Str}(M)}
 \DeltaRow{Annahmen}{\Gamma,\Delta\subseteq\mathcal F}
 \DeltaRow{Formeln}{p,q,r\in\mathcal F}
}
\begin{tabproofsplitwide}
 \proofpartwide{Beweis}
 \proofstepwidestar[1]{\Gamma\models_M p\leftrightarrow q}
  {\rA}
 \proofstepwidestar[2]{M,s\models \Gamma}
  {\rA}
 \proofstepwidestar[1,2]{M,s\models p\leftrightarrow q}
  {\FormulaRefAuto{B48ConsequenceInstance}[thm]{1,2}}
 \proofstepwidestar[1,2]{(M,s\models p)\leftrightarrow(M,s\models q)}
  {\rLREa{\FormulaRefAuto{B48SatIffClause}[thm],3}}
 \proofstepwidestar[1,2]{(M,s\models q)\to(M,s\models p)}
  {\rLREb{4}}
 \proofstepwidestar[1,2]{M,s\models q\to p}
  {\rLREb{\FormulaRefAuto{B48SatImpClause}[thm],5}}
 \proofstepwidestar[1]{(M,s\models \Gamma)\to(M,s\models q\to p)}
  {\rRI{2,6}}
 \proofstepwidestar[1]{\Gamma\models_M q\to p}
  {\FormulaRefAuto{B48SemanticConsequence}[def]{\rUI{7}}}
 \closeproofpartwide
\end{tabproofsplitwide}

\FormulaThmDeltaK[Semantische Konjunktionseinfuehrung fuer Annahmenmengen]
{\begin{gathered}\Gamma\models_M p,\quad\Delta\models_M q\\
 \vdash\ \Gamma\cup\Delta\models_M p\land q\end{gathered}}
{B48RuleAndIntro}{
 \DeltaRow{Struktur}{\mathsf{Str}(M)}
 \DeltaRow{Annahmen}{\Gamma,\Delta\subseteq\mathcal F}
 \DeltaRow{Formeln}{p,q,r\in\mathcal F}
}
\begin{tabproofsplitwide}
 \proofpartwide{Beweis}
 \proofstepwidestar[1]{\Gamma\models_M p}
  {\rA}
 \proofstepwidestar[2]{\Delta\models_M q}
  {\rA}
 \proofstepwidestar[3]{M,s\models \Gamma\cup\Delta}
  {\rA}
 \proofstepwidestar[3]{(M,s\models \Gamma)\land(M,s\models \Delta)}
  {\rLREa{\FormulaRefAuto{B48ContextUnion}[thm],3}}
 \proofstepwidestar[3]{M,s\models \Gamma}
  {\rAEa{4}}
 \proofstepwidestar[3]{M,s\models \Delta}
  {\rAEb{4}}
 \proofstepwidestar[1,3]{M,s\models p}
  {\FormulaRefAuto{B48ConsequenceInstance}[thm]{1,5}}
 \proofstepwidestar[2,3]{M,s\models q}
  {\FormulaRefAuto{B48ConsequenceInstance}[thm]{2,6}}
 \proofstepwidestar[1,2,3]{M,s\models p\land q}
  {\FormulaRefAuto{B48SatAndIntro}[thm]{7,8}}
 \proofstepwidestar[1,2]{(M,s\models \Gamma\cup\Delta)\to(M,s\models p\land q)}
  {\rRI{3,9}}
 \proofstepwidestar[1,2]{\Gamma\cup\Delta\models_M p\land q}
  {\FormulaRefAuto{B48SemanticConsequence}[def]{\rUI{10}}}
 \closeproofpartwide
\end{tabproofsplitwide}

\FormulaThmDeltaK[Semantische Implikationsbeseitigung]
{\begin{gathered}\Gamma\models_M p\to q,\quad\Delta\models_M p\\
 \vdash\ \Gamma\cup\Delta\models_M q\end{gathered}}
{B48RuleImpElim}{
 \DeltaRow{Struktur}{\mathsf{Str}(M)}
 \DeltaRow{Annahmen}{\Gamma,\Delta\subseteq\mathcal F}
 \DeltaRow{Formeln}{p,q,r\in\mathcal F}
}
\begin{tabproofsplitwide}
 \proofpartwide{Beweis}
 \proofstepwidestar[1]{\Gamma\models_M p\to q}
  {\rA}
 \proofstepwidestar[2]{\Delta\models_M p}
  {\rA}
 \proofstepwidestar[3]{M,s\models \Gamma\cup\Delta}
  {\rA}
 \proofstepwidestar[3]{(M,s\models \Gamma)\land(M,s\models \Delta)}
  {\rLREa{\FormulaRefAuto{B48ContextUnion}[thm],3}}
 \proofstepwidestar[3]{M,s\models \Gamma}
  {\rAEa{4}}
 \proofstepwidestar[3]{M,s\models \Delta}
  {\rAEb{4}}
 \proofstepwidestar[1,3]{M,s\models p\to q}
  {\FormulaRefAuto{B48ConsequenceInstance}[thm]{1,5}}
 \proofstepwidestar[2,3]{M,s\models p}
  {\FormulaRefAuto{B48ConsequenceInstance}[thm]{2,6}}
 \proofstepwidestar[1,3]{(M,s\models p)\to(M,s\models q)}
  {\rLREa{\FormulaRefAuto{B48SatImpClause}[thm],7}}
 \proofstepwidestar[1,2,3]{M,s\models q}
  {\rRE{9,8}}
 \proofstepwidestar[1,2]{(M,s\models \Gamma\cup\Delta)\to(M,s\models q)}
  {\rRI{3,10}}
 \proofstepwidestar[1,2]{\Gamma\cup\Delta\models_M q}
  {\FormulaRefAuto{B48SemanticConsequence}[def]{\rUI{11}}}
 \closeproofpartwide
\end{tabproofsplitwide}

\FormulaThmDeltaK[Semantische Bikonditionaleinfuehrung]
{\begin{gathered}\Gamma\models_M p\to q,\quad\Delta\models_M q\to p\\
 \vdash\ \Gamma\cup\Delta\models_M p\leftrightarrow q\end{gathered}}
{B48RuleIffIntro}{
 \DeltaRow{Struktur}{\mathsf{Str}(M)}
 \DeltaRow{Annahmen}{\Gamma,\Delta\subseteq\mathcal F}
 \DeltaRow{Formeln}{p,q,r\in\mathcal F}
}
\begin{tabproofsplitwide}
 \proofpartwide{Beweis}
 \proofstepwidestar[1]{\Gamma\models_M p\to q}
  {\rA}
 \proofstepwidestar[2]{\Delta\models_M q\to p}
  {\rA}
 \proofstepwidestar[3]{M,s\models \Gamma\cup\Delta}
  {\rA}
 \proofstepwidestar[3]{(M,s\models \Gamma)\land(M,s\models \Delta)}
  {\rLREa{\FormulaRefAuto{B48ContextUnion}[thm],3}}
 \proofstepwidestar[3]{M,s\models \Gamma}
  {\rAEa{4}}
 \proofstepwidestar[3]{M,s\models \Delta}
  {\rAEb{4}}
 \proofstepwidestar[1,3]{M,s\models p\to q}
  {\FormulaRefAuto{B48ConsequenceInstance}[thm]{1,5}}
 \proofstepwidestar[2,3]{M,s\models q\to p}
  {\FormulaRefAuto{B48ConsequenceInstance}[thm]{2,6}}
 \proofstepwidestar[1,3]{(M,s\models p)\to(M,s\models q)}
  {\rLREa{\FormulaRefAuto{B48SatImpClause}[thm],7}}
 \proofstepwidestar[2,3]{(M,s\models q)\to(M,s\models p)}
  {\rLREa{\FormulaRefAuto{B48SatImpClause}[thm],8}}
 \proofstepwidestar[1,2,3]{(M,s\models p)\leftrightarrow(M,s\models q)}
  {\rLRI{9,10}}
 \proofstepwidestar[1,2,3]{M,s\models p\leftrightarrow q}
  {\rLREb{\FormulaRefAuto{B48SatIffClause}[thm],11}}
 \proofstepwidestar[1,2]{(M,s\models \Gamma\cup\Delta)\to(M,s\models p\leftrightarrow q)}
  {\rRI{3,12}}
 \proofstepwidestar[1,2]{\Gamma\cup\Delta\models_M p\leftrightarrow q}
  {\FormulaRefAuto{B48SemanticConsequence}[def]{\rUI{13}}}
 \closeproofpartwide
\end{tabproofsplitwide}

\FormulaThmDeltaK[Semantische Widerspruchseinfuehrung]
{\begin{gathered}\Gamma\models_M p,\quad\Delta\models_M \neg p\\
 \vdash\ \Gamma\cup\Delta\models_M \bot\end{gathered}}
{B48RuleBottom}{
 \DeltaRow{Struktur}{\mathsf{Str}(M)}
 \DeltaRow{Annahmen}{\Gamma,\Delta\subseteq\mathcal F}
 \DeltaRow{Formeln}{p,q,r\in\mathcal F}
}
\begin{tabproofsplitwide}
 \proofpartwide{Beweis}
 \proofstepwidestar[1]{\Gamma\models_M p}
  {\rA}
 \proofstepwidestar[2]{\Delta\models_M \neg p}
  {\rA}
 \proofstepwidestar[3]{M,s\models \Gamma\cup\Delta}
  {\rA}
 \proofstepwidestar[3]{(M,s\models \Gamma)\land(M,s\models \Delta)}
  {\rLREa{\FormulaRefAuto{B48ContextUnion}[thm],3}}
 \proofstepwidestar[3]{M,s\models \Gamma}
  {\rAEa{4}}
 \proofstepwidestar[3]{M,s\models \Delta}
  {\rAEb{4}}
 \proofstepwidestar[1,3]{M,s\models p}
  {\FormulaRefAuto{B48ConsequenceInstance}[thm]{1,5}}
 \proofstepwidestar[2,3]{M,s\models \neg p}
  {\FormulaRefAuto{B48ConsequenceInstance}[thm]{2,6}}
 \proofstepwidestar[1,2,3]{\bot}
  {\FormulaRefAuto{B48SatContradiction}[thm]{7,8}}
 \proofstepwidestar[1,2,3]{M,s\models \bot}
  {\FormulaRefAuto{B48Satisfaction}[def]{9}}
 \proofstepwidestar[1,2]{(M,s\models \Gamma\cup\Delta)\to(M,s\models \bot)}
  {\rRI{3,10}}
 \proofstepwidestar[1,2]{\Gamma\cup\Delta\models_M \bot}
  {\FormulaRefAuto{B48SemanticConsequence}[def]{\rUI{11}}}
 \closeproofpartwide
\end{tabproofsplitwide}

\FormulaThmDeltaK[Semantische Implikationseinfuehrung]
{\Gamma\cup\{p\}\models_M q\ \vdash\ \Gamma\models_M p\to q}
{B48RuleImpIntro}{
 \DeltaRow{Struktur}{\mathsf{Str}(M)}
 \DeltaRow{Annahmen}{\Gamma,\Delta,\Theta\subseteq\mathcal F}
 \DeltaRow{Formeln und Variablen}{p,q,r\in\mathcal F,\quad x,y,z\in\mathbb N}
}
\begin{tabproofsplitwide}
 \proofpartwide{Beweis}
 \proofstepwidestar[1]{\Gamma\cup\{p\}\models_M q}
  {\rA}
 \proofstepwidestar[2]{M,s\models \Gamma}
  {\rA}
 \proofstepwidestar[3]{M,s\models p}
  {\rA}
 \proofstepwidestar[2,3]{(M,s\models \Gamma)\land(M,s\models p)}
  {\rAI{2,3}}
 \proofstepwidestar[2,3]{M,s\models \Gamma\cup\{p\}}
  {\rLREb{\FormulaRefAuto{B48ContextAssumption}[thm],4}}
 \proofstepwidestar[1,2,3]{M,s\models q}
  {\FormulaRefAuto{B48ConsequenceInstance}[thm]{1,5}}
 \proofstepwidestar[1,2]{(M,s\models p)\to(M,s\models q)}
  {\rRI{3,6}}
 \proofstepwidestar[1,2]{M,s\models p\to q}
  {\rLREb{\FormulaRefAuto{B48SatImpClause}[thm],7}}
 \proofstepwidestar[1]{(M,s\models \Gamma)\to(M,s\models p\to q)}
  {\rRI{2,8}}
 \proofstepwidestar[1]{\Gamma\models_M p\to q}
  {\FormulaRefAuto{B48SemanticConsequence}[def]{\rUI{9}}}
 \closeproofpartwide
\end{tabproofsplitwide}

\FormulaThmDeltaK[Semantische Negationseinfuehrung mit Entladung]
{\Gamma\cup\{p\}\models_M \bot\ \vdash\ \Gamma\models_M \neg p}
{B48RuleNegIntro}{
 \DeltaRow{Struktur}{\mathsf{Str}(M)}
 \DeltaRow{Annahmen}{\Gamma,\Delta,\Theta\subseteq\mathcal F}
 \DeltaRow{Formeln und Variablen}{p,q,r\in\mathcal F,\quad x,y,z\in\mathbb N}
}
\begin{tabproofsplitwide}
 \proofpartwide{Beweis}
 \proofstepwidestar[1]{\Gamma\cup\{p\}\models_M \bot}
  {\rA}
 \proofstepwidestar[2]{M,s\models \Gamma}
  {\rA}
 \proofstepwidestar[3]{M,s\models p}
  {\rA}
 \proofstepwidestar[2,3]{(M,s\models \Gamma)\land(M,s\models p)}
  {\rAI{2,3}}
 \proofstepwidestar[2,3]{M,s\models \Gamma\cup\{p\}}
  {\rLREb{\FormulaRefAuto{B48ContextAssumption}[thm],4}}
 \proofstepwidestar[1,2,3]{M,s\models \bot}
  {\FormulaRefAuto{B48ConsequenceInstance}[thm]{1,5}}
 \proofstepwidestar[1,2,3]{\bot}
  {\FormulaRefAuto{B48Satisfaction}[def]{6}}
 \proofstepwidestar[1,2]{\neg(M,s\models p)}
  {\rCI{3,7}}
 \proofstepwidestar[1,2]{M,s\models \neg p}
  {\FormulaRefAuto{B48Satisfaction}[def]{8}}
 \proofstepwidestar[1]{(M,s\models \Gamma)\to(M,s\models \neg p)}
  {\rRI{2,9}}
 \proofstepwidestar[1]{\Gamma\models_M \neg p}
  {\FormulaRefAuto{B48SemanticConsequence}[def]{\rUI{10}}}
 \closeproofpartwide
\end{tabproofsplitwide}

\FormulaThmDeltaK[Semantischer klassischer Regelfall mit Entladung]
{\Gamma\cup\{\neg p\}\models_M \bot\ \vdash\ \Gamma\models_M p}
{B48RuleClassical}{
 \DeltaRow{Struktur}{\mathsf{Str}(M)}
 \DeltaRow{Annahmen}{\Gamma,\Delta,\Theta\subseteq\mathcal F}
 \DeltaRow{Formeln und Variablen}{p,q,r\in\mathcal F,\quad x,y,z\in\mathbb N}
}
\begin{tabproofsplitwide}
 \proofpartwide{Beweis}
 \proofstepwidestar[1]{\Gamma\cup\{\neg p\}\models_M \bot}
  {\rA}
 \proofstepwidestar[2]{M,s\models \Gamma}
  {\rA}
 \proofstepwidestar[3]{\neg(M,s\models p)}
  {\rA}
 \proofstepwidestar[3]{M,s\models \neg p}
  {\FormulaRefAuto{B48Satisfaction}[def]{3}}
 \proofstepwidestar[2,3]{(M,s\models \Gamma)\land(M,s\models \neg p)}
  {\rAI{2,4}}
 \proofstepwidestar[2,3]{M,s\models \Gamma\cup\{\neg p\}}
  {\rLREb{\FormulaRefAuto{B48ContextAssumption}[thm],5}}
 \proofstepwidestar[1,2,3]{M,s\models \bot}
  {\FormulaRefAuto{B48ConsequenceInstance}[thm]{1,6}}
 \proofstepwidestar[1,2,3]{\bot}
  {\FormulaRefAuto{B48Satisfaction}[def]{7}}
 \proofstepwidestar[1,2]{M,s\models p}
  {\rCE{3,8}}
 \proofstepwidestar[1]{(M,s\models \Gamma)\to(M,s\models p)}
  {\rRI{2,9}}
 \proofstepwidestar[1]{\Gamma\models_M p}
  {\FormulaRefAuto{B48SemanticConsequence}[def]{\rUI{10}}}
 \closeproofpartwide
\end{tabproofsplitwide}

\FormulaThmDeltaK[Semantische Disjunktionsbeseitigung mit beiden Entladungen]
{\begin{gathered}
 \Gamma\models_M p\lor q,\quad
 \Delta\cup\{p\}\models_M r,\quad\Theta\cup\{q\}\models_M r\\
 \vdash\ (\Gamma\cup\Delta)\cup\Theta\models_M r
 \end{gathered}}
{B48RuleOrElim}{
 \DeltaRow{Struktur}{\mathsf{Str}(M)}
 \DeltaRow{Annahmen}{\Gamma,\Delta,\Theta\subseteq\mathcal F}
 \DeltaRow{Formeln und Variablen}{p,q,r\in\mathcal F,\quad x,y,z\in\mathbb N}
}
\begin{tabproofsplitwide}
 \proofpartwide{Beweis}
 \proofstepwidestar[1]{\Gamma\models_M p\lor q}
  {\rA}
 \proofstepwidestar[2]{\Delta\cup\{p\}\models_M r}
  {\rA}
 \proofstepwidestar[3]{\Theta\cup\{q\}\models_M r}
  {\rA}
 \proofstepwidestar[4]{M,s\models (\Gamma\cup\Delta)\cup\Theta}
  {\rA}
 \proofstepwidestar[4]{(M,s\models \Gamma\cup\Delta)\land(M,s\models \Theta)}
  {\rLREa{\FormulaRefAuto{B48ContextUnion}[thm],4}}
 \proofstepwidestar[4]{(M,s\models \Gamma)\land(M,s\models \Delta)}
  {\rLREa{\FormulaRefAuto{B48ContextUnion}[thm],\rAEa{5}}}
 \proofstepwidestar[1,4]{M,s\models p\lor q}
  {\FormulaRefAuto{B48ConsequenceInstance}[thm]{1,\rAEa{6}}}
 \proofstepwidestar[1,4]{(M,s\models p)\lor(M,s\models q)}
  {\rLREa{\FormulaRefAuto{B48SatOrClause}[thm],7}}
 \proofstepwidestar[9]{M,s\models p}
  {\rA}
 \proofstepwidestar[4,9]{(M,s\models \Delta)\land(M,s\models p)}
  {\rAI{\rAEb{6},9}}
 \proofstepwidestar[4,9]{M,s\models \Delta\cup\{p\}}
  {\rLREb{\FormulaRefAuto{B48ContextAssumption}[thm],10}}
 \proofstepwidestar[2,4,9]{M,s\models r}
  {\FormulaRefAuto{B48ConsequenceInstance}[thm]{2,11}}
 \proofstepwidestar[13]{M,s\models q}
  {\rA}
 \proofstepwidestar[4,13]{(M,s\models \Theta)\land(M,s\models q)}
  {\rAI{\rAEb{5},13}}
 \proofstepwidestar[4,13]{M,s\models \Theta\cup\{q\}}
  {\rLREb{\FormulaRefAuto{B48ContextAssumption}[thm],14}}
 \proofstepwidestar[3,4,13]{M,s\models r}
  {\FormulaRefAuto{B48ConsequenceInstance}[thm]{3,15}}
 \proofstepwidestar[1,2,3,4]{M,s\models r}
  {\rOE{8,9:12,13:16}}
 \proofstepwidestar[1,2,3]{(M,s\models (\Gamma\cup\Delta)\cup\Theta)\to(M,s\models r)}
  {\rRI{4,17}}
 \proofstepwidestar[1,2,3]{(\Gamma\cup\Delta)\cup\Theta\models_M r}
  {\FormulaRefAuto{B48SemanticConsequence}[def]{\rUI{18}}}
 \closeproofpartwide
\end{tabproofsplitwide}

\FormulaThmDeltaK[Semantische All-Einfuehrung]
{\Gamma\models_M p,\quad x\mathrel{\#}\Gamma\
   \vdash\ \Gamma\models_M\forall x\,p}
{B48RuleAllIntro}{
 \DeltaRow{Struktur}{\mathsf{Str}(M)}
 \DeltaRow{Annahmen}{\Gamma,\Delta,\Theta\subseteq\mathcal F}
 \DeltaRow{Formeln und Variablen}{p,q,r\in\mathcal F,\quad x,y,z\in\mathbb N}
}
\begin{tabproofsplitwide}
 \proofpartwide{Beweis}
 \proofstepwidestar[1]{\Gamma\models_M p}
  {\rA}
 \proofstepwidestar[2]{x\mathrel{\#}\Gamma}
  {\rA}
 \proofstepwidestar[3]{M,s\models \Gamma}
  {\rA}
 \proofstepwidestar[4]{a\in D}
  {\rA}
 \proofstepwidestar[4]{s[x\mapsto a]\in A_D}
  {\FormulaRefAuto{B48AssignmentUpdate}[thm]{4}}
 \proofstepwidestar[2,3,4]{M,s[x\mapsto a]\models\Gamma}
  {\FormulaRefAuto{B48ContextFresh}[thm]{2,3,4}}
 \proofstepwidestar[1,2,3,4]{M,s[x\mapsto a]\models p}
  {\FormulaRefAuto{B48ConsequenceInstance}[thm]{1,6,5}}
 \proofstepwidestar[1,2,3]{\forall a\in D\;(M,s[x\mapsto a]\models p)}
  {\rUI{\rRI{4,7}}}
 \proofstepwidestar[1,2,3]{M,s\models \forall x\,p}
  {\rLREb{\FormulaRefAuto{B48SatAllClause}[thm],8}}
 \proofstepwidestar[1,2]{(M,s\models \Gamma)\to(M,s\models \forall x\,p)}
  {\rRI{3,9}}
 \proofstepwidestar[1,2]{\Gamma\models_M \forall x\,p}
  {\FormulaRefAuto{B48SemanticConsequence}[def]{\rUI{10}}}
 \closeproofpartwide
\end{tabproofsplitwide}

\FormulaThmDeltaK[Semantische All-Beseitigung mit Substitution]
{\Gamma\models_M \forall x\,p,\quad\mathsf{Fr}(p,x,y)\
 \vdash\ \Gamma\models_M p[y/x]}
{B48RuleAllElim}{
 \DeltaRow{Struktur}{\mathsf{Str}(M)}
 \DeltaRow{Annahmen}{\Gamma,\Delta,\Theta\subseteq\mathcal F}
 \DeltaRow{Formeln und Variablen}{p,q,r\in\mathcal F,\quad x,y,z\in\mathbb N}
}
\begin{tabproofsplitwide}
 \proofpartwide{Beweis}
 \proofstepwidestar[1]{\Gamma\models_M \forall x\,p}
  {\rA}
 \proofstepwidestar[2]{\mathsf{Fr}(p,x,y)}
  {\rA}
 \proofstepwidestar[3]{M,s\models \Gamma}
  {\rA}
 \proofstepwidestar[1,3]{M,s\models \forall x\,p}
  {\FormulaRefAuto{B48ConsequenceInstance}[thm]{1,3}}
 \proofstepwidestar[]{s\colon\mathbb N\to D}
  {\FormulaRefAuto{B48AssignmentCriterion}[thm]}
 \proofstepwidestar[]{s(y)\in D}
  {\FormulaRefAuto{F\colon A\to B,\ x\in A\vdash F(x)\in B}[thm]{5}}
 \proofstepwidestar[2]{(M,s\models p[y/x])\leftrightarrow(M,s[x\mapsto s(y)]\models p)}
  {\FormulaRefAuto{B48SubstitutionLemma}[thm]{2}}
 \proofstepwidestar[1,3]{\forall a\in D\;(M,s[x\mapsto a]\models p)}
  {\rLREa{\FormulaRefAuto{B48SatAllClause}[thm],4}}
 \proofstepwidestar[1,3]{M,s[x\mapsto s(y)]\models p}
  {\rRE{\rUE{8},6}}
 \proofstepwidestar[1,2,3]{M,s\models p[y/x]}
  {\rLREb{7,9}}
 \proofstepwidestar[1,2]{(M,s\models \Gamma)\to(M,s\models p[y/x])}
  {\rRI{3,10}}
 \proofstepwidestar[1,2]{\Gamma\models_M p[y/x]}
  {\FormulaRefAuto{B48SemanticConsequence}[def]{\rUI{11}}}
 \closeproofpartwide
\end{tabproofsplitwide}

\FormulaThmDeltaK[Semantische Existenz-Einfuehrung mit Substitution]
{\Gamma\models_M p[y/x],\quad\mathsf{Fr}(p,x,y)\
 \vdash\ \Gamma\models_M \exists x\,p}
{B48RuleExistsIntro}{
 \DeltaRow{Struktur}{\mathsf{Str}(M)}
 \DeltaRow{Annahmen}{\Gamma,\Delta,\Theta\subseteq\mathcal F}
 \DeltaRow{Formeln und Variablen}{p,q,r\in\mathcal F,\quad x,y,z\in\mathbb N}
}
\begin{tabproofsplitwide}
 \proofpartwide{Beweis}
 \proofstepwidestar[1]{\Gamma\models_M p[y/x]}
  {\rA}
 \proofstepwidestar[2]{\mathsf{Fr}(p,x,y)}
  {\rA}
 \proofstepwidestar[3]{M,s\models \Gamma}
  {\rA}
 \proofstepwidestar[1,3]{M,s\models p[y/x]}
  {\FormulaRefAuto{B48ConsequenceInstance}[thm]{1,3}}
 \proofstepwidestar[]{s\colon\mathbb N\to D}
  {\FormulaRefAuto{B48AssignmentCriterion}[thm]}
 \proofstepwidestar[]{s(y)\in D}
  {\FormulaRefAuto{F\colon A\to B,\ x\in A\vdash F(x)\in B}[thm]{5}}
 \proofstepwidestar[2]{(M,s\models p[y/x])\leftrightarrow(M,s[x\mapsto s(y)]\models p)}
  {\FormulaRefAuto{B48SubstitutionLemma}[thm]{2}}
 \proofstepwidestar[1,2,3]{M,s[x\mapsto s(y)]\models p}
  {\rLREa{7,4}}
 \proofstepwidestar[1,2,3]{M,s\models \exists x\,p}
  {\FormulaRefAuto{B48SatExistsIntro}[thm]{6,8}}
 \proofstepwidestar[1,2]{(M,s\models \Gamma)\to(M,s\models \exists x\,p)}
  {\rRI{3,9}}
 \proofstepwidestar[1,2]{\Gamma\models_M \exists x\,p}
  {\FormulaRefAuto{B48SemanticConsequence}[def]{\rUI{10}}}
 \closeproofpartwide
\end{tabproofsplitwide}

\FormulaThmDeltaK[Semantische Existenz-Beseitigung mit Eigenvariable]
{\begin{gathered}
 \Gamma\models_M\exists x\,p,\quad\Delta\cup\{p\}\models_M q,\\
 x\mathrel{\#}\Delta,\quad x\notin\operatorname{FV}(q)\\
 \vdash\ \Gamma\cup\Delta\models_M q
 \end{gathered}}
{B48RuleExistsElim}{
 \DeltaRow{Struktur}{\mathsf{Str}(M)}
 \DeltaRow{Annahmen}{\Gamma,\Delta,\Theta\subseteq\mathcal F}
 \DeltaRow{Formeln und Variablen}{p,q,r\in\mathcal F,\quad x,y,z\in\mathbb N}
}
Der Zeuge \(a\) ist frisch fuer die ausserhalb seines Teilbeweises offenen Annahmen.
\begin{tabproofsplitwide}
 \proofpartwide{Beweis}
 \proofstepwidestar[1]{\Gamma\models_M\exists x\,p}
  {\rA}
 \proofstepwidestar[2]{\Delta\cup\{p\}\models_M q}
  {\rA}
 \proofstepwidestar[3]{x\mathrel{\#}\Delta}
  {\rA}
 \proofstepwidestar[4]{x\notin\operatorname{FV}(q)}
  {\rA}
 \proofstepwidestar[5]{M,s\models \Gamma\cup\Delta}
  {\rA}
 \proofstepwidestar[5]{(M,s\models \Gamma)\land(M,s\models \Delta)}
  {\rLREa{\FormulaRefAuto{B48ContextUnion}[thm],5}}
 \proofstepwidestar[1,5]{M,s\models \exists x\,p}
  {\FormulaRefAuto{B48ConsequenceInstance}[thm]{1,\rAEa{6}}}
 \proofstepwidestar[1,5]{\exists a\;(a\in D\land(M,s[x\mapsto a]\models p))}
  {\FormulaRefAuto{B48Satisfaction}[def]{7}}
 \proofstepwidestar[9]{a\in D\land(M,s[x\mapsto a]\models p)}
  {\rA}
 \proofstepwidestar[9]{a\in D}
  {\rAEa{9}}
 \proofstepwidestar[9]{M,s[x\mapsto a]\models p}
  {\rAEb{9}}
 \proofstepwidestar[9]{s[x\mapsto a]\in A_D}
  {\FormulaRefAuto{B48AssignmentUpdate}[thm]{10}}
 \proofstepwidestar[3,5,9]{M,s[x\mapsto a]\models\Delta}
  {\FormulaRefAuto{B48ContextFresh}[thm]{3,\rAEb{6},10}}
 \proofstepwidestar[3,5,9]{M,s[x\mapsto a]\models\Delta\cup\{p\}}
  {\rLREb{\FormulaRefAuto{B48ContextAssumption}[thm],\rAI{13,11}}}
 \proofstepwidestar[2,3,5,9]{M,s[x\mapsto a]\models q}
  {\FormulaRefAuto{B48ConsequenceInstance}[thm]{2,14,12}}
 \proofstepwidestar[4,9]{(M,s\models q)\leftrightarrow(M,s[x\mapsto a]\models q)}
  {\FormulaRefAuto{B48FreshUpdate}[thm]{4,10}}
 \proofstepwidestar[2,3,4,5,9]{M,s\models q}
  {\rLREb{16,15}}
 \proofstepwidestar[1,2,3,4,5]{M,s\models q}
  {\rEE{8,9:17}}
 \proofstepwidestar[1,2,3,4]{(M,s\models \Gamma\cup\Delta)\to(M,s\models q)}
  {\rRI{5,18}}
 \proofstepwidestar[1,2,3,4]{\Gamma\cup\Delta\models_M q}
  {\FormulaRefAuto{B48SemanticConsequence}[def]{\rUI{19}}}
 \closeproofpartwide
\end{tabproofsplitwide}

\FormulaThmDeltaK[Semantische Gleichheitseinfuehrung]
{\Gamma\models_M v_x=v_x}
{B48RuleEqualityIntro}{
 \DeltaRow{Struktur}{\mathsf{Str}(M)}
 \DeltaRow{Annahmen}{\Gamma,\Delta,\Theta\subseteq\mathcal F}
 \DeltaRow{Formeln und Variablen}{p,q,r\in\mathcal F,\quad x,y,z\in\mathbb N}
}
\begin{tabproofsplitwide}
 \proofpartwide{Beweis}
 \proofstepwidestar[1]{M,s\models \Gamma}
  {\rA}
 \proofstepwidestar[]{s(x)=s(x)}
  {\rII}
 \proofstepwidestar[]{M,s\models v_x=v_x}
  {\FormulaRefAuto{B48Satisfaction}[def]{2}}
 \proofstepwidestar[]{(M,s\models \Gamma)\to(M,s\models v_x=v_x)}
  {\rRI{1,3}}
 \proofstepwidestar[]{\Gamma\models_M v_x=v_x}
  {\FormulaRefAuto{B48SemanticConsequence}[def]{\rUI{4}}}
 \closeproofpartwide
\end{tabproofsplitwide}

\FormulaThmDeltaK[Semantische Gleichheitsbeseitigung]
{\begin{gathered}
 \Gamma\models_M v_x=v_y,\quad\Delta\models_M p[x/z],\\
 \mathsf{Fr}(p,z,x),\quad\mathsf{Fr}(p,z,y)\\
 \vdash\ \Gamma\cup\Delta\models_M p[y/z]
 \end{gathered}}
{B48RuleEqualityElim}{
 \DeltaRow{Struktur}{\mathsf{Str}(M)}
 \DeltaRow{Annahmen}{\Gamma,\Delta,\Theta\subseteq\mathcal F}
 \DeltaRow{Formeln und Variablen}{p,q,r\in\mathcal F,\quad x,y,z\in\mathbb N}
}
\begin{tabproofsplitwide}
 \proofpartwide{Beweis}
 \proofstepwidestar[1]{\Gamma\models_M v_x=v_y}
  {\rA}
 \proofstepwidestar[2]{\Delta\models_M p[x/z]}
  {\rA}
 \proofstepwidestar[3]{\mathsf{Fr}(p,z,x)}
  {\rA}
 \proofstepwidestar[4]{\mathsf{Fr}(p,z,y)}
  {\rA}
 \proofstepwidestar[5]{M,s\models \Gamma\cup\Delta}
  {\rA}
 \proofstepwidestar[5]{(M,s\models \Gamma)\land(M,s\models \Delta)}
  {\rLREa{\FormulaRefAuto{B48ContextUnion}[thm],5}}
 \proofstepwidestar[1,5]{M,s\models v_x=v_y}
  {\FormulaRefAuto{B48ConsequenceInstance}[thm]{1,\rAEa{6}}}
 \proofstepwidestar[1,5]{s(x)=s(y)}
  {\FormulaRefAuto{B48Satisfaction}[def]{7}}
 \proofstepwidestar[2,5]{M,s\models p[x/z]}
  {\FormulaRefAuto{B48ConsequenceInstance}[thm]{2,\rAEb{6}}}
 \proofstepwidestar[3]{(M,s\models p[x/z])\leftrightarrow(M,s[z\mapsto s(x)]\models p)}
  {\FormulaRefAuto{B48SubstitutionLemma}[thm]{3}}
 \proofstepwidestar[2,3,5]{M,s[z\mapsto s(x)]\models p}
  {\rLREa{10,9}}
 \proofstepwidestar[1,2,3,5]{M,s[z\mapsto s(y)]\models p}
  {\rIE{8,11}}
 \proofstepwidestar[4]{(M,s\models p[y/z])\leftrightarrow(M,s[z\mapsto s(y)]\models p)}
  {\FormulaRefAuto{B48SubstitutionLemma}[thm]{4}}
 \proofstepwidestar[1,2,3,4,5]{M,s\models p[y/z]}
  {\rLREb{13,12}}
 \proofstepwidestar[1,2,3,4]{(M,s\models \Gamma\cup\Delta)\to(M,s\models p[y/z])}
  {\rRI{5,14}}
 \proofstepwidestar[1,2,3,4]{\Gamma\cup\Delta\models_M p[y/z]}
  {\FormulaRefAuto{B48SemanticConsequence}[def]{\rUI{15}}}
 \closeproofpartwide
\end{tabproofsplitwide}

\FormulaThmDeltaK[Semantischer Annahmenfall]
{p\in\Gamma\ \vdash\ \Gamma\models_M p}
{B48RuleAssumption}{
 \DeltaRow{Struktur}{\mathsf{Str}(M)}
 \DeltaRow{Annahmen}{\Gamma,\Delta,\Theta\subseteq\mathcal F}
 \DeltaRow{Formeln und Variablen}{p,q,r\in\mathcal F,\quad x,y,z\in\mathbb N}
}
\begin{tabproofsplitwide}
 \proofpartwide{Beweis}
 \proofstepwidestar[1]{p\in\Gamma}
  {\rA}
 \proofstepwidestar[2]{M,s\models \Gamma}
  {\rA}
 \proofstepwidestar[1,2]{M,s\models p}
  {\FormulaRefAuto{B48SemanticConsequence}[def]{2,1}}
 \proofstepwidestar[1]{(M,s\models \Gamma)\to(M,s\models p)}
  {\rRI{2,3}}
 \proofstepwidestar[1]{\Gamma\models_M p}
  {\FormulaRefAuto{B48SemanticConsequence}[def]{\rUI{4}}}
 \closeproofpartwide
\end{tabproofsplitwide}

\subsection{Induktion ueber die endliche Herleitung}

\FormulaDefDeltaK[Daten einer elementaren Beweistabelle]
{\begin{aligned}
 d&=((I_j,\theta_j,r_j,J_j))_{1\leq j\leq N},\\
 \Gamma_j&\coloneq\{\theta_k\mid k\in I_j\},\\
 H_d(n)&\coloneq
  \forall j\in\mathbb N\;
  ((1\leq j\leq N\land j\leq n)\to\Gamma_j\models_M\theta_j).
\end{aligned}}
{B48DerivationData}{
 \DeltaRow{Laenge}{N\in\mathbb N,\quad N\geq1}
 \DeltaRow{Formeln}{\theta_j\in\mathcal F}
 \DeltaRow{Annahmenindizes}{I_j\subseteq\{1,\ldots,j\}}
 \DeltaRow{Regel und fruehere Verweise}{r_j,\quad J_j}
}
Der Datensatz \(d\) ist eine \emph{gueltige} elementare Beweistabelle
nach Band~1: Eine Annahmezeile \(j\) hat \(I_j=\{j\}\).
Jeder sonstige Verweis in \(J_j\) zeigt auf eine Zeile mit Nummer
kleiner als \(j\). Die Formel und die Abhaengigkeitsmenge jeder
Schlusszeile entstehen nach ihrer angegebenen Regel.
Bei einer Entladung wird der Index der betreffenden Annahme aus
der Abhaengigkeitsmenge des Subbeweises entfernt; danach werden die
verbleibenden Mengen vereinigt.
Bei \(\forall I\) kommt die Variable in keiner offenen Annahme frei
vor. Bei \(\exists E\) kommt sie weder in der Schlussformel noch in
den nach der Entladung verbleibenden Annahmen frei vor.
Alle Einsetzungen erfuellen \(\mathsf{Fr}\).

Die \(\Gamma_j\) sind Mengen von Formeln, aber die Entladung wird
\emph{vorher an den Indizes} ausgefuehrt. Zwei verschiedene
Annahmezeilen mit derselben Formel duerfen deshalb nicht
miteinander verwechselt werden. Fuer eine entladene Annahme mit
Formel \(p\) gilt stets
\(\Gamma_{\rm Subbeweis}=\Gamma_{\rm Rest}\cup\{p\}\), auch wenn
\(p\) wegen einer anderen Annahme weiterhin in \(\Gamma_{\rm Rest}\)
liegt.

\FormulaThmDeltaK[Korrektheit jeder Zeile einer Beweistabelle]
{\forall j\in\{1,\ldots,N\}\;(\Gamma_j\models_M\theta_j)}
{B48DerivationSoundness}{
 \DeltaRow{Struktur}{\mathsf{Str}(M)}
 \DeltaRow{Gueltige elementare Beweistabelle}{d\text{ nach }
                          \FormulaRefAuto{B48DerivationData}[def]}
}
Wir beweisen \(H_d(n)\) durch Induktion ueber \(n\).
Fuer \(n=0\) ist der Bereich der zu pruefenden Zeilen leer.
Sei \(H_d(n)\) gegeben. Fuer \(n+1>N\) kommt keine Zeile hinzu.
Andernfalls betrachten wir die Regel der Zeile \(n+1\).
Alle unten als semantische Folgerungen aufgefuehrten Premissen
folgen aus \(H_d(n)\), weil die referenzierten Zeilen frueher
liegen. Die Einsetzungs- und Eigenvariablenbedingungen folgen
aus der Gueltigkeit von \(d\). Die Mengen
\(\Gamma,\Delta,\Theta\) sind jeweils die verbleibenden
Annahmenmengen der betreffenden Praemissen beziehungsweise
Subbeweise.

\begin{longtable}{@{}p{0.13\linewidth}p{0.56\linewidth}p{0.24\linewidth}@{}}
\toprule
Regel & Angewendeter semantischer Schluss & Nachweis\\
\midrule\endhead
\(\RuleRef{A}{A}\) &
 \(p\in\Gamma\ \vdash\ \Gamma\models_M p\) &
 \FormulaRefAuto{B48RuleAssumption}[thm]\\[5pt]
\(\RuleRef{AI}{\land I}\) &
 \(\begin{gathered}\Gamma\models_M p,\ \Delta\models_M q\\
 \vdash\Gamma\cup\Delta\models_M p\land q\end{gathered}\) &
 \FormulaRefAuto{B48RuleAndIntro}[thm]\\[5pt]
\(\RuleRef{AE1}{\land E_1}\) &
 \(\Gamma\models_M p\land q\ \vdash\Gamma\models_M p\) &
 \FormulaRefAuto{B48RuleAndLeft}[thm]\\[5pt]
\(\RuleRef{AE2}{\land E_2}\) &
 \(\Gamma\models_M p\land q\ \vdash\Gamma\models_M q\) &
 \FormulaRefAuto{B48RuleAndRight}[thm]\\[5pt]
\(\RuleRef{OI1}{\lor I_1}\) &
 \(\Gamma\models_M p\ \vdash\Gamma\models_M p\lor q\) &
 \FormulaRefAuto{B48RuleOrLeft}[thm]\\[5pt]
\(\RuleRef{OI2}{\lor I_2}\) &
 \(\Gamma\models_M q\ \vdash\Gamma\models_M p\lor q\) &
 \FormulaRefAuto{B48RuleOrRight}[thm]\\[5pt]
\(\RuleRef{OE}{\lor E}\) &
 \(\begin{gathered}\Gamma\models_M p\lor q,\\
 \Delta\cup\{p\}\models_M r,\ \Theta\cup\{q\}\models_M r\\
 \vdash(\Gamma\cup\Delta)\cup\Theta\models_M r\end{gathered}\) &
 \FormulaRefAuto{B48RuleOrElim}[thm]\\[5pt]
\(\RuleRef{RI}{\to I}\) &
 \(\Gamma\cup\{p\}\models_M q\ \vdash\Gamma\models_M p\to q\) &
 \FormulaRefAuto{B48RuleImpIntro}[thm]\\[5pt]
\(\RuleRef{RE}{\to E}\) &
 \(\begin{gathered}\Gamma\models_M p\to q,\ \Delta\models_M p\\
 \vdash\Gamma\cup\Delta\models_M q\end{gathered}\) &
 \FormulaRefAuto{B48RuleImpElim}[thm]\\[5pt]
\(\RuleRef{LRI}{\leftrightarrow I}\) &
 \(\begin{gathered}\Gamma\models_M p\to q,\ \Delta\models_M q\to p\\
 \vdash\Gamma\cup\Delta\models_M p\leftrightarrow q\end{gathered}\) &
 \FormulaRefAuto{B48RuleIffIntro}[thm]\\[5pt]
\(\RuleRef{LRE1}{\leftrightarrow E_1}\) &
 \(\Gamma\models_M p\leftrightarrow q\ \vdash\Gamma\models_M p\to q\) &
 \FormulaRefAuto{B48RuleIffLeft}[thm]\\[5pt]
\(\RuleRef{LRE2}{\leftrightarrow E_2}\) &
 \(\Gamma\models_M p\leftrightarrow q\ \vdash\Gamma\models_M q\to p\) &
 \FormulaRefAuto{B48RuleIffRight}[thm]\\[5pt]
\(\RuleRef{BI}{\bot I}\) &
 \(\begin{gathered}\Gamma\models_M p,\ \Delta\models_M\neg p\\
 \vdash\Gamma\cup\Delta\models_M\bot\end{gathered}\) &
 \FormulaRefAuto{B48RuleBottom}[thm]\\[5pt]
\(\RuleRef{CI}{\neg I}\) &
 \(\Gamma\cup\{p\}\models_M\bot\ \vdash\Gamma\models_M\neg p\) &
 \FormulaRefAuto{B48RuleNegIntro}[thm]\\[5pt]
\(\RuleRef{CE}{\neg E}\) &
 \(\Gamma\cup\{\neg p\}\models_M\bot\ \vdash\Gamma\models_M p\) &
 \FormulaRefAuto{B48RuleClassical}[thm]\\[5pt]
\(\RuleRef{UI}{\forall I}\) &
 \(\begin{gathered}\Gamma\models_M p,\ x\mathrel{\#}\Gamma\\
 \vdash\Gamma\models_M\forall x\,p\end{gathered}\) &
 \FormulaRefAuto{B48RuleAllIntro}[thm]\\[5pt]
\(\RuleRef{UE}{\forall E}\) &
 \(\begin{gathered}\Gamma\models_M\forall x\,p,\ \mathsf{Fr}(p,x,y)\\
 \vdash\Gamma\models_M p[y/x]\end{gathered}\) &
 \FormulaRefAuto{B48RuleAllElim}[thm]\\[5pt]
\(\RuleRef{EI}{\exists I}\) &
 \(\begin{gathered}\Gamma\models_M p[y/x],\ \mathsf{Fr}(p,x,y)\\
 \vdash\Gamma\models_M\exists x\,p\end{gathered}\) &
 \FormulaRefAuto{B48RuleExistsIntro}[thm]\\[5pt]
\(\RuleRef{EE}{\exists E}\) &
 \(\begin{gathered}\Gamma\models_M\exists x\,p,\ \Delta\cup\{p\}\models_M q,\\
 x\mathrel{\#}\Delta,\ x\notin\operatorname{FV}(q)\\
 \vdash\Gamma\cup\Delta\models_M q\end{gathered}\) &
 \FormulaRefAuto{B48RuleExistsElim}[thm]\\[5pt]
\(\RuleRef{II}{=I}\) &
 \(\Gamma\models_M v_x=v_x\) &
 \FormulaRefAuto{B48RuleEqualityIntro}[thm]\\[5pt]
\(\RuleRef{IE}{=E}\) &
 \(\begin{gathered}\Gamma\models_M v_x=v_y,\ \Delta\models_M p[x/z],\\
 \mathsf{Fr}(p,z,x),\ \mathsf{Fr}(p,z,y)\\
 \vdash\Gamma\cup\Delta\models_M p[y/z]\end{gathered}\) &
 \FormulaRefAuto{B48RuleEqualityElim}[thm]\\
\bottomrule
\end{longtable}

Damit ist in jedem moeglichen Regelfall
\(\Gamma_{n+1}\models_M\theta_{n+1}\) nachgewiesen.
Fuer \(j\leq n\) gilt die Behauptung weiter nach \(H_d(n)\);
fuer \(j=n+1\) gilt der gerade bewiesene Fall. Also gilt
\(H_d(n+1)\). Die Induktion
\FormulaRefAuto{PeanoInductionRuleAddOne}[thm] liefert
\(\forall n\in\mathbb N\;H_d(n)\). Einsetzen von \(n=N\) ist
genau die behauptete Korrektheit aller Zeilen.

\noindent
Die aufgefuehrten 21 Faelle decken alle elementaren Regeln aus
Band~1 ab. Die dortigen iterierten Gleichheits-, Konjunktions-
und Existenznotationen bezeichnen endliche Folgen solcher
Regelschritte. Der Satz setzt keine Korrektheit ungepruefter
Bibliotheksverweise voraus: \(\operatorname{Prf}_1\) bezeichnet hier
die elementare Herleitung, in der verwendete Abkuerzungen bereits
aufgeloest sind. Die fuer Henkin benoetigten allgemeinen
Umformungssaetze ueber Beweisobjekte bleiben eine eigene Aufgabe.

\FormulaThmDeltaK[Ein Modell erfuellt jede Teilmenge seiner Theorie]
{M\models T,\quad\Gamma\subseteq T\ \vdash\ M,s\models\Gamma}
{B48ModelContext}{
 \DeltaRow{Struktur und Belegung}{\mathsf{Str}(M),\quad s\in A_D}
 \DeltaRow{Theorie und Annahmen}{T\subseteq\operatorname{Sent}(\mathcal L_\in),
                               \quad\Gamma\subseteq\mathcal F}
}
\begin{tabproof}
 \proofstep{1}{M\models T}{\rA}
 \proofstep{2}{\Gamma\subseteq T}{\rA}
 \proofstep{3}{p\in\Gamma}{\rA}
 \proofstep{2,3}{p\in T}
  {\FormulaRefAuto{A\subseteq B,\ x\in A\vdash x\in B}[thm]{2,3}}
 \proofstep{1,2,3}{M\models p}
  {\FormulaRefAuto{B48SentenceSatisfaction}[def]{1,4}}
 \proofstep{}{s\colon\operatorname{Var}\to D}
  {\FormulaRefAuto{B48AssignmentCriterion}[thm]}
 \proofstep{1,2,3}{M,s\models p}
  {\FormulaRefAuto{B48SentenceInstance}[thm]{5,6}}
 \proofstep{1,2}{M,s\models\Gamma}
  {\FormulaRefAuto{B48SemanticConsequence}[def]{\rUI{\rRI{3,7}}}}
\end{tabproof}

\FormulaThmDeltaK[Korrektheit der Herleitung]
{T\vdash\varphi,\quad M\models T\ \vdash\ M\models\varphi}
{B48Soundness}{
 \DeltaRow{Struktur}{\mathsf{Str}(M)}
 \DeltaRow{Theorie}{T\subseteq\operatorname{Sent}(\mathcal L_\in)}
 \DeltaRow{Satz}{\varphi\in\operatorname{Sent}(\mathcal L_\in)}
}
Im Zeugenbeweis ist \(d\) frisch. Seine Laenge ist \(N\), seine
letzte Formel ist \(\varphi\), und seine offenen Annahmen
\(\Gamma_N\) liegen in \(T\). Die Belegung \(s\) wird beliebig
gewaehlt und vor der Entladung von \(d\) wieder gebunden.
\begin{tabproofsplitwide}
 \proofpartwide{Schluss aus der letzten Beweiszeile}
 \proofstepwidestar[1]{T\vdash\varphi}{\rA}
 \proofstepwidestar[2]{M\models T}{\rA}
 \proofstepwidestar[1]{\exists d\;\operatorname{Prf}_1(d,T,\varphi)}
  {\FormulaRefAuto{B48Syntax}[def]{1}}
 \proofstepwidestar[4]{\operatorname{Prf}_1(d,T,\varphi)}{\rA}
 \proofstepwidestar[4]{\Gamma_N\models_M\varphi}
  {\FormulaRefAuto{B48DerivationSoundness}[thm]{4},
   \FormulaRefAuto{B48DerivationData}[def]}
 \proofstepwidestar[4]{\Gamma_N\subseteq T}
  {\FormulaRefAuto{B48Syntax}[def]{4}}
 \proofstepwidestar[7]{s\in A_D}{\rA}
 \proofstepwidestar[2,4,7]{M,s\models\Gamma_N}
  {\FormulaRefAuto{B48ModelContext}[thm]{2,6,7}}
 \proofstepwidestar[2,4,7]{M,s\models\varphi}
  {\FormulaRefAuto{B48ConsequenceInstance}[thm]{5,8}}
 \proofstepwidestar[2,4]{\forall s\in A_D\;(M,s\models\varphi)}
  {\rUI{\rRI{7,9}}}
 \proofstepwidestar[]{\exists a\;(a\in D)}
  {\FormulaRefAuto{B48StrDomainAxiom}[ax]}
 \proofstepwidestar[2,4]{M\models\varphi}
  {\FormulaRefAuto{B48SentenceSatisfaction}[def]{11,10},
   \FormulaRefAuto{B48AssignmentCriterion}[thm]}
 \proofstepwidestar[1,2]{M\models\varphi}{\rEE{3,4:12}}
 \closeproofpartwide
\end{tabproofsplitwide}



\FormulaThmDeltaK[Ein Modell sichert Konsistenz]
{M\models T\ \vdash\ \BXLVIIICon(T)}
{B48ModelConsistency}{
  \DeltaRow{Theorie}{T}
  \DeltaRow{Struktur}{M}
}
\noindent
Der folgende lokale Schluss verwendet die Korrektheitsaussage
\FormulaRefAuto{B48Soundness}[thm]. Die Korrektheit elementarer
Herleitungen ist im vorstehenden Abschnitt aus Koinzidenz,
Substitution und den einzelnen Regelfaellen hergeleitet.
Die hier benoetigte semantische Unmoeglichkeit ist dagegen bereits in
\FormulaRefAuto{B48SentenceNotBottom}[thm] hergeleitet.
Im Schluss ist \(q\) frisch fuer \(M,T\) und die ausserhalb des
Beweises offenen Annahmen.
\begin{tabproofsplitwide}
 \proofpartwideindR[Ein Widerspruchsbeweis fuehrt zum Widerspruch]{%
   M\models T\dsep \operatorname{Prf}_{1}(p,T,\bot)\vdash\bot
 }{%
   M\models T\dsep \operatorname{Prf}_{1}(p,T,\bot)\vdash\bot
 }[B48ModelConsistencyProofWitness]
  \proofstepwidestar[1]{M\models T}{\rA}
  \proofstepwidestar[2]{\operatorname{Prf}_{1}(p,T,\bot)}{\rA}
  \proofstepwidestar[2]{%
   \exists r\,\operatorname{Prf}_{1}(r,T,\bot)
  }{\rEI{2}}
  \proofstepwidestar[2]{T\vdash\bot}
   {\FormulaRefAuto{B48Syntax}{3}}
  \proofstepwidestar[1,2]{M\models\bot}
   {\FormulaRefAuto{B48Soundness}{4,1}}
  \proofstepwidestar[]{\neg(M\models\bot)}
   {\FormulaRefAuto{B48SentenceNotBottom}}
  \proofstepwidestar[1,2]{\bot}{\rBI{5,6}}
 \closeproofpartwideindR

 \proofpartwide{Schluss}
  \proofstepwidestar[1]{M\models T}{\rA}
  \proofstepwidestar[2]{T\vdash\bot}{\rA}
  \proofstepwidestar[2]{%
   \exists p\,\operatorname{Prf}_{1}(p,T,\bot)
  }{\FormulaRefAuto{B48Syntax}{2}}
  \proofstepwidestar[4]{\operatorname{Prf}_{1}(q,T,\bot)}{\rA}
  \proofstepwidestar[1,4]{\bot}
   {\FormulaRefAuto{B48ModelConsistencyProofWitness}{1,4}}
  \proofstepwidestar[1,2]{\bot}{\rEE{3,4:5}}
  \proofstepwidestar[1]{\neg(T\vdash\bot)}{\rCI{2,6}}
  \proofstepwidestar[1]{\BXLVIIICon(T)}
   {\FormulaRefAuto{B48ConsistencyNotation}{7}}
 \closeproofpartwide
\end{tabproofsplitwide}

\section{Henkin-Konstruktion und Vollständigkeit}

Für die Konstruktion dürfen neue Konstantensymbole hinzugefügt werden.
Ein Satz hat keine freien Variablen; Konstanten sind erlaubt.

\FormulaThmDeltaK[Ein frischer Henkin-Zeuge erhält Konsistenz]
{\BXLVIIICon(S)\ \vdash\
 \BXLVIIICon\bigl(S\cup\{\exists x\,\psi(x)\to\psi(c)\}\bigr)}
{B48HenkinWitness}{
  \DeltaRow{Theorie}{S}
  \DeltaRow{Frische Konstante}{c\notin\operatorname{Symb}(S,\psi)}
}
\noindent\textbf{Beweisstand: syntaktische Hilfssaetze fehlen.}
Der folgende Entwurf setzt die Entladung auf Beweisobjekten und
die Ersetzung einer frischen Konstante durch eine Eigenvariable
voraus. Beides ist noch nicht durch eine vorherige Beweistabelle
gerechtfertigt. Die Schlussregel \(\forall I\) fuer Formeln darf
nicht ohne diesen Uebertragungssatz auf die Aussage
\(S\vdash\psi(c)\) angewendet werden.
\begin{tabproof}
 \proofstep{1}{\BXLVIIICon(S)}{\rA}
 \proofstep{2}{S+\bigl(\exists x\,\psi(x)\to\psi(c)\bigr)\vdash\bot}{\rA}
 \proofstep{2}{S\vdash\exists x\,\psi(x)\land\neg\psi(c)}
 {\BXLVIIIWhy{Deduktion und klassische Negation einer Implikation}}
 \proofstep{2}{S\vdash\forall x\,\neg\psi(x)}
 {\BXLVIIIWhy{\(c\) kommt in \(S,\psi\) nicht vor: durch eine frische
 Variable ersetzen und generalisieren}}
 \proofstep{1,2}{\bot}{\BXLVIIIWhy{Existenz und universelle Negation
 widersprechen der Konsistenz von \(S\)}}
 \proofstep{1}{\BXLVIIICon\bigl(S+(\exists x\,\psi(x)\to\psi(c))\bigr)}
 {\BXLVIIIWhy{2 entladen}}
\end{tabproof}

\FormulaThmDeltaK[Vollständige Henkin-Erweiterung]
{\BXLVIIICon(T)\ \vdash\ \exists T^*\supseteq T\
 \text{vollständig, konsistent und mit Henkin-Zeugen}}
{B48HenkinCompletion}{
  \DeltaRow{Sprache}{\mathcal L\text{ abzählbar}}
  \DeltaRow{Theorie}{T\subseteq\operatorname{Sent}(\mathcal L)}
}
\noindent\textbf{Beweisstand: Henkin-Konstruktion noch offen.}
Erforderlich sind getrennte Nachweise fuer Sprachkonservativitaet,
endliche Beweistraeger, konsistente aufsteigende Vereinigungen,
den Entscheidungsschritt und den deduktiven Abschluss.
Der nachstehende Entwurf ersetzt diese Nachweise nicht.
\begin{tabproof}
 \proofstep{1}{\BXLVIIICon(T)}{\rA}
 \proofstep{1}{\BXLVIIIText{Füge zunächst eine unbenutzte Konstante hinzu.
 Zähle alle Existenzsätze der Sprache auf und füge nacheinander die
 Zeugenaxiome mit jeweils frischen Konstanten hinzu.}}
 {\BXLVIIIWhy{Abzählbarkeit; vorige Zeugenlemma}}
 \proofstep{1}{\BXLVIIIText{Wiederhole diesen Vorgang in abzählbar vielen
 Runden, damit auch Formeln mit neu hinzugefügten Konstanten Zeugen
 erhalten. Jede endliche Menge der gesammelten Axiome liegt bereits in
 einer konsistenten endlichen Konstruktionsstufe.}}
 {\BXLVIIIWhy{Endlichkeit von Formeln und Beweisen}}
 \proofstep{1}{\BXLVIIIText{Die Vereinigung \(H\) ist konsistent.
 Zähle alle Sätze der endgültigen, weiterhin abzählbaren Sprache auf.
 Füge jeweils \(\theta\) hinzu, wenn dies konsistent ist, sonst
 \(\neg\theta\).}}
 {\BXLVIIIWhy{Rekursion auf \(\omega\)}}
 \proofstep{1}{\BXLVIIIText{Sind sowohl \(S+\theta\) als auch
 \(S+\neg\theta\) inkonsistent, so beweist \(S\) durch Deduktion
 sowohl \(\neg\theta\) als auch \(\theta\). Deshalb erhält jeder
 Entscheidungsschritt die Konsistenz.}}
 {\BXLVIIIWhy{Klassische Deduktion}}
 \proofstep{1}{\BXLVIIIText{Die deduktive Hülle der Vereinigung ist
 \(T^*\). Sie enthält \(T\), entscheidet jeden Satz und enthält alle
 Zeugenaxiome. Ein Widerspruch hätte schon endlich viele ihrer
 Voraussetzungen benutzt und ist daher ausgeschlossen.}}
 {\BXLVIIIWhy{Endlichkeit der Herleitungen}}
\end{tabproof}

\FormulaThmDeltaK[Vollständigkeit und abzählbares Modell]
{\BXLVIIICon(T)\ \Longleftrightarrow\
 \exists M\ (M\models T\land |M|\le\aleph_0)}
{B48Completeness}{
  \DeltaRow{Sprache}{\mathcal L\text{ abzählbar, relational}}
  \DeltaRow{Theorie}{T\subseteq\operatorname{Sent}(\mathcal L)}
}
\noindent\textbf{Beweisstand: kein abgeschlossener Vollstaendigkeitsbeweis.}
Neben der Henkin-Erweiterung sind die Gleichheitskongruenz, die
Unabhaengigkeit der Relationsinterpretation von Repraesentanten
und das Termmodell-Wahrheitslemma noch auszuarbeiten. Die folgende
Tabelle ist weiterhin die Konstruktionsskizze.
\begin{tabproof}
 \proofstep{1}{\BXLVIIICon(T)}{\rA}
 \proofstep{1}{\BXLVIIIText{Wähle eine vollständige konsistente
 Henkin-Theorie \(T^*\). Setze \(c\sim d\) genau dann, wenn
 \(T^*\vdash c=d\). Die Gleichheitsregeln machen dies zu einer
 Äquivalenzrelation.}}
 {\FormulaRefAuto{B48HenkinCompletion}[thm]}
 \proofstep{1}{\BXLVIIIText{Die Grundmenge \(D\) besteht aus den
 Äquivalenzklassen der Konstanten; jede Konstante \(c\) wird
 durch \([c]\) interpretiert. Für jedes \(n\)-stellige
 Relationssymbol \(R\) setze
 \(R^M([c_1],\ldots,[c_n])\) genau dann, wenn
 \(T^*\vdash R(c_1,\ldots,c_n)\).
 Gleichheitssubstitution sichert die Unabhängigkeit von Repräsentanten.}}
 {\BXLVIIIWhy{Quotientenbildung}}
 \proofstep{1}{\BXLVIIIText{Wahrheitslemma:
 \(M\models\theta([c_1],\ldots,[c_n])\) genau dann, wenn
 \(\theta(c_1,\ldots,c_n)\in T^*\).
 Für atomare Formeln gilt dies per Definition; für Negation und
 Konjunktion durch Vollständigkeit und Konsistenz von \(T^*\).}}
 {\BXLVIIIWhy{Strukturelle Induktion über \(\theta\)}}
 \proofstep{1}{\BXLVIIIText{Ein in \(T^*\) enthaltener Existenzsatz hat
 durch sein Zeugenaxiom eine Zeugeninstanz in \(T^*\).
 Umgekehrt folgt aus jeder Zeugeninstanz der Existenzsatz.
 Dies beweist den Existenzschritt des Wahrheitslemmas.}}
 {\BXLVIIIWhy{Henkin-Eigenschaft, Existenz-Einführung}}
 \proofstep{1}{M\models T,\qquad |M|\le\aleph_0}
 {\BXLVIIIWhy{Wahrheitslemma; abzählbar viele Konstanten}}
 \proofstep{7}{\exists M\ (M\models T\land |M|\le\aleph_0)}{\rA}
 \proofstep{7}{\BXLVIIICon(T)}
 {\FormulaRefAuto{B48ModelConsistency}[thm]}
 \proofstep{}{\BXLVIIICon(T)\ \Longleftrightarrow\
 \exists M\ (M\models T\land |M|\le\aleph_0)}
 {\BXLVIIIWhy{Beide Richtungen, 1 und 7 entladen}}
\end{tabproof}

\noindent\textbf{Wichtige Grenze.}
Dieser Satz liefert keine transitive und keine von außen fundierte
Mengenstruktur. Insbesondere wird im Folgenden niemals aus
\(\BXLVIIICon(\BXLVIIIZFC)\) die Existenz eines abzählbaren
\emph{transitiven} ZFC-Modells geschlossen.

\section{Referenzierte Schlussformen f\"ur Modellargumente}

In diesem Abschnitt sind die Aussagen \(M\models\varphi\) und
\(T\vdash\varphi\) Aussagen der zuvor festgelegten Metatheorie.
Auch auf dieser Ebene werden die Schlussregeln aus Band~1 angewendet.
Die letzte Spalte nennt jeweils die Regel oder den zuvor bewiesenen
Satz und die benutzten Zeilen. Eigenvariablen werden in jedem
Subbeweis frisch gew\"ahlt.

\FormulaDefDeltaK[Ausgeschriebene Modell- und Erweiterungsnotation]
{\begin{gathered}
 T+\varphi\coloneqq T\cup\{\varphi\},\\
 M\models T\ \coloneqq\
 \forall\theta\,(\theta\in T\longrightarrow M\models\theta),\\
 T\nvdash\varphi\ \coloneqq\ \neg(T\vdash\varphi)
 \end{gathered}}
{B48ModelReasonNotation}{
 \DeltaRow{Theorie}{T\subseteq\operatorname{Sent}(\mathcal L)}
 \DeltaRow{Satz}{\varphi\in\operatorname{Sent}(\mathcal L)}
 \DeltaRow{Struktur}{M}
}
Diese Definition schreibt die bereits benutzten Abk\"urzungen aus;
sie f\"uhrt kein zus\"atzliches Mengenaxiom ein. Insbesondere wird
die Erf\"ullung einer erweiterten Theorie nicht als neue
Schlussregel vorausgesetzt, sondern im folgenden Satz hergeleitet.

\FormulaThmDeltaK[Mitgliedschaft in einer Theorieerweiterung]
{\theta\in T+\varphi\ \longleftrightarrow\
 (\theta\in T\lor\theta=\varphi)}
{B48TheoryExtensionMember}{
 \DeltaRow{Theorie}{T\subseteq\operatorname{Sent}(\mathcal L)}
 \DeltaRow{S\"atze}{\theta,\varphi\in\operatorname{Sent}(\mathcal L)}
}
\begin{tabproof}
 \proofstep{1}{\theta\in T+\varphi}{\rA}
 \proofstep{1}{\theta\in T\cup\{\varphi\}}
 {\FormulaRefAuto{B48ModelReasonNotation}{1}}
 \proofstep{1}{\theta\in T\lor\theta\in\{\varphi\}}
 {\FormulaRefAuto{x\in A\cup B\eqvdash x\in A\lor x\in B}{2}}
 \proofstep{4}{\theta\in T}{\rA}
 \proofstep{4}{\theta\in T\lor\theta=\varphi}
 {\RuleRef{OI1}{\lor I_1}\,4}
 \proofstep{6}{\theta\in\{\varphi\}}{\rA}
 \proofstep{6}{\theta=\varphi}
 {\FormulaRefAuto{x\in\{a\}\eqvdash x=a}{6}}
 \proofstep{6}{\theta\in T\lor\theta=\varphi}
 {\RuleRef{OI2}{\lor I_2}\,7}
 \proofstep{1}{\theta\in T\lor\theta=\varphi}
 {\RuleRef{OE}{\lor E}\,3,4\!:\!5,6\!:\!8}
 \proofstep{}{\theta\in T+\varphi\rightarrow
   (\theta\in T\lor\theta=\varphi)}
 {\RuleRef{RI}{\rightarrow I}\,1\!:\!9}
 \proofstep{11}{\theta\in T\lor\theta=\varphi}{\rA}
 \proofstep{12}{\theta\in T}{\rA}
 \proofstep{12}{\theta\in T\lor\theta\in\{\varphi\}}
 {\RuleRef{OI1}{\lor I_1}\,12}
 \proofstep{14}{\theta=\varphi}{\rA}
 \proofstep{14}{\theta\in\{\varphi\}}
 {\FormulaRefAuto{x\in\{a\}\eqvdash x=a}{14}}
 \proofstep{14}{\theta\in T\lor\theta\in\{\varphi\}}
 {\RuleRef{OI2}{\lor I_2}\,15}
 \proofstep{11}{\theta\in T\lor\theta\in\{\varphi\}}
 {\RuleRef{OE}{\lor E}\,11,12\!:\!13,14\!:\!16}
 \proofstep{11}{\theta\in T\cup\{\varphi\}}
 {\FormulaRefAuto{x\in A\cup B\eqvdash x\in A\lor x\in B}{17}}
 \proofstep{11}{\theta\in T+\varphi}
 {\FormulaRefAuto{B48ModelReasonNotation}{18}}
 \proofstep{}{(\theta\in T\lor\theta=\varphi)
   \rightarrow\theta\in T+\varphi}
 {\RuleRef{RI}{\rightarrow I}\,11\!:\!19}
 \proofstep{}{\theta\in T+\varphi\leftrightarrow
   (\theta\in T\lor\theta=\varphi)}
 {\RuleRef{LRI}{\leftrightarrow I}\,10,20}
\end{tabproof}

\FormulaThmDeltaK[Modell einer Theorieerweiterung]
{\begin{gathered}
 M\models T+\varphi\ \longleftrightarrow\\
 (M\models T)\land(M\models\varphi)
 \end{gathered}}
{B48ModelTheoryExtension}{
 \DeltaRow{Theorie}{T\subseteq\operatorname{Sent}(\mathcal L)}
 \DeltaRow{Satz}{\varphi\in\operatorname{Sent}(\mathcal L)}
 \DeltaRow{Struktur}{M}
}
\begin{tabproof}
 \proofstep{1}{M\models T+\varphi}{\rA}
 \proofstep{1}{\forall\theta\,(\theta\in T+\varphi
   \rightarrow M\models\theta)}
 {\FormulaRefAuto{B48ModelReasonNotation}{1}}
 \proofstep{3}{\theta\in T}{\rA}
 \proofstep{3}{\theta\in T\lor\theta=\varphi}
 {\RuleRef{OI1}{\lor I_1}\,3}
 \proofstep{3}{\theta\in T+\varphi}
 {\FormulaRefAuto{B48TheoryExtensionMember}{4}}
 \proofstep{1}{\theta\in T+\varphi\rightarrow M\models\theta}
 {\RuleRef{UE}{\forall E}\,2}
 \proofstep{1,3}{M\models\theta}
 {\RuleRef{RE}{\rightarrow E}\,6,5}
 \proofstep{1}{\theta\in T\rightarrow M\models\theta}
 {\RuleRef{RI}{\rightarrow I}\,3\!:\!7}
 \proofstep{1}{\forall\theta\,(\theta\in T\rightarrow M\models\theta)}
 {\RuleRef{UI}{\forall I}\,8}
 \proofstep{1}{M\models T}
 {\FormulaRefAuto{B48ModelReasonNotation}{9}}
 \proofstep{}{\varphi=\varphi}{\RuleRef{II}{=I}}
 \proofstep{}{\varphi\in T\lor\varphi=\varphi}
 {\RuleRef{OI2}{\lor I_2}\,11}
 \proofstep{}{\varphi\in T+\varphi}
 {\FormulaRefAuto{B48TheoryExtensionMember}{12}}
 \proofstep{1}{\varphi\in T+\varphi\rightarrow M\models\varphi}
 {\RuleRef{UE}{\forall E}\,2}
 \proofstep{1}{M\models\varphi}
 {\RuleRef{RE}{\rightarrow E}\,14,13}
 \proofstep{1}{(M\models T)\land(M\models\varphi)}
 {\RuleRef{AI}{\land I}\,10,15}
 \proofstep{}{\begin{gathered}M\models T+\varphi\rightarrow\\
   ((M\models T)\land(M\models\varphi))\end{gathered}}
 {\RuleRef{RI}{\rightarrow I}\,1\!:\!16}
 \proofstep{18}{(M\models T)\land(M\models\varphi)}{\rA}
 \proofstep{18}{M\models T}{\RuleRef{AE1}{\land E_1}\,18}
 \proofstep{18}{M\models\varphi}{\RuleRef{AE2}{\land E_2}\,18}
 \proofstep{18}{\forall\theta\,(\theta\in T\rightarrow M\models\theta)}
 {\FormulaRefAuto{B48ModelReasonNotation}{19}}
 \proofstep{22}{\theta\in T+\varphi}{\rA}
 \proofstep{22}{\theta\in T\lor\theta=\varphi}
 {\FormulaRefAuto{B48TheoryExtensionMember}{22}}
 \proofstep{24}{\theta\in T}{\rA}
 \proofstep{18}{\theta\in T\rightarrow M\models\theta}
 {\RuleRef{UE}{\forall E}\,21}
 \proofstep{18,24}{M\models\theta}
 {\RuleRef{RE}{\rightarrow E}\,25,24}
 \proofstep{27}{\theta=\varphi}{\rA}
 \proofstep{}{\theta=\theta}{\RuleRef{II}{=I}}
 \proofstep{27}{\varphi=\theta}{\RuleRef{IE}{=E}\,27,28}
 \proofstep{18,27}{M\models\theta}{\RuleRef{IE}{=E}\,29,20}
 \proofstep{18,22}{M\models\theta}
 {\RuleRef{OE}{\lor E}\,23,24\!:\!26,27\!:\!30}
 \proofstep{18}{\theta\in T+\varphi\rightarrow M\models\theta}
 {\RuleRef{RI}{\rightarrow I}\,22\!:\!31}
 \proofstep{18}{\forall\theta\,(\theta\in T+\varphi
   \rightarrow M\models\theta)}
 {\RuleRef{UI}{\forall I}\,32}
 \proofstep{18}{M\models T+\varphi}
 {\FormulaRefAuto{B48ModelReasonNotation}{33}}
 \proofstep{}{\begin{gathered}((M\models T)\land(M\models\varphi))\\
   \rightarrow M\models T+\varphi\end{gathered}}
 {\RuleRef{RI}{\rightarrow I}\,18\!:\!34}
 \proofstep{}{\begin{gathered}M\models T+\varphi\leftrightarrow\\
   ((M\models T)\land(M\models\varphi))\end{gathered}}
 {\RuleRef{LRI}{\leftrightarrow I}\,17,35}
\end{tabproof}

\FormulaThmDeltaK[Ein Modell verhindert eine Widerlegung]
{M\models T,\ M\models\varphi\ \vdash\ T\nvdash\neg\varphi}
{B48ModelBlocksRefutation}{
 \DeltaRow{Theorie}{T\subseteq\operatorname{Sent}(\mathcal L)}
 \DeltaRow{Satz}{\varphi\in\operatorname{Sent}(\mathcal L)}
 \DeltaRow{Struktur}{M}
}
\begin{tabproof}
 \proofstep{1}{M\models T}{\rA}
 \proofstep{2}{M\models\varphi}{\rA}
 \proofstep{3}{T\vdash\neg\varphi}{\rA}
 \proofstep{1,3}{M\models\neg\varphi}
  {\FormulaRefAuto{B48Soundness}{3,1}}
 \proofstep{1,2,3}{\bot}
  {\FormulaRefAuto{B48SentenceContradiction}{2,4}}
 \proofstep{1,2}{\neg(T\vdash\neg\varphi)}{\rCI{3,5}}
 \proofstep{1,2}{T\nvdash\neg\varphi}
  {\FormulaRefAuto{B48ModelReasonNotation}{6}}
\end{tabproof}

\FormulaThmDeltaK[Ein Gegenmodell verhindert einen Beweis]
{M\models T,\ M\models\neg\varphi\ \vdash\ T\nvdash\varphi}
{B48CountermodelBlocksProof}{
 \DeltaRow{Theorie}{T\subseteq\operatorname{Sent}(\mathcal L)}
 \DeltaRow{Satz}{\varphi\in\operatorname{Sent}(\mathcal L)}
 \DeltaRow{Struktur}{M}
}
\begin{tabproof}
 \proofstep{1}{M\models T}{\rA}
 \proofstep{2}{M\models\neg\varphi}{\rA}
 \proofstep{3}{T\vdash\varphi}{\rA}
 \proofstep{1,3}{M\models\varphi}
  {\FormulaRefAuto{B48Soundness}{3,1}}
 \proofstep{1,2,3}{\bot}
  {\FormulaRefAuto{B48SentenceContradiction}{4,2}}
 \proofstep{1,2}{\neg(T\vdash\varphi)}{\rCI{3,5}}
 \proofstep{1,2}{T\nvdash\varphi}
  {\FormulaRefAuto{B48ModelReasonNotation}{6}}
\end{tabproof}

\FormulaThmDeltaK[Relative Konsistenz verhindert Widerlegbarkeit]
{\BXLVIIICon(T+\varphi)\longrightarrow T\nvdash\neg\varphi}
{B48RelativeNonRefutation}{
 \DeltaRow{Sprache}{\mathcal L\text{ abz\"ahlbar, relational}}
 \DeltaRow{Theorie}{T\subseteq\operatorname{Sent}(\mathcal L)}
 \DeltaRow{Satz}{\varphi\in\operatorname{Sent}(\mathcal L)}
}
\begin{tabproof}
 \proofstep{1}{\BXLVIIICon(T+\varphi)}{\rA}
 \proofstep{1}{\exists M\,(M\models T+\varphi\land |M|\le\aleph_0)}
 {\FormulaRefAuto{B48Completeness}{1}}
 \proofstep{3}{M\models T+\varphi\land |M|\le\aleph_0}{\rA}
 \proofstep{3}{M\models T+\varphi}{\RuleRef{AE1}{\land E_1}\,3}
 \proofstep{3}{(M\models T)\land(M\models\varphi)}
 {\FormulaRefAuto{B48ModelTheoryExtension}{4}}
 \proofstep{3}{M\models T}{\RuleRef{AE1}{\land E_1}\,5}
 \proofstep{3}{M\models\varphi}{\RuleRef{AE2}{\land E_2}\,5}
 \proofstep{3}{T\nvdash\neg\varphi}
 {\FormulaRefAuto{B48ModelBlocksRefutation}{6,7}}
 \proofstep{1}{T\nvdash\neg\varphi}
 {\RuleRef{EE}{\exists E}\,2,3\!:\!8}
 \proofstep{}{\BXLVIIICon(T+\varphi)\rightarrow T\nvdash\neg\varphi}
 {\RuleRef{RI}{\rightarrow I}\,1\!:\!9}
\end{tabproof}

\FormulaThmDeltaK[Relative Konsistenz der Negation verhindert Beweisbarkeit]
{\BXLVIIICon(T+\neg\varphi)\longrightarrow T\nvdash\varphi}
{B48RelativeNonProof}{
 \DeltaRow{Sprache}{\mathcal L\text{ abz\"ahlbar, relational}}
 \DeltaRow{Theorie}{T\subseteq\operatorname{Sent}(\mathcal L)}
 \DeltaRow{Satz}{\varphi\in\operatorname{Sent}(\mathcal L)}
}
\begin{tabproof}
 \proofstep{1}{\BXLVIIICon(T+\neg\varphi)}{\rA}
 \proofstep{1}{\exists M\,(M\models T+\neg\varphi\land |M|\le\aleph_0)}
 {\FormulaRefAuto{B48Completeness}{1}}
 \proofstep{3}{M\models T+\neg\varphi\land |M|\le\aleph_0}{\rA}
 \proofstep{3}{M\models T+\neg\varphi}{\RuleRef{AE1}{\land E_1}\,3}
 \proofstep{3}{(M\models T)\land(M\models\neg\varphi)}
 {\FormulaRefAuto{B48ModelTheoryExtension}{4}}
 \proofstep{3}{M\models T}{\RuleRef{AE1}{\land E_1}\,5}
 \proofstep{3}{M\models\neg\varphi}{\RuleRef{AE2}{\land E_2}\,5}
 \proofstep{3}{T\nvdash\varphi}
 {\FormulaRefAuto{B48CountermodelBlocksProof}{6,7}}
 \proofstep{1}{T\nvdash\varphi}
 {\RuleRef{EE}{\exists E}\,2,3\!:\!8}
 \proofstep{}{\BXLVIIICon(T+\neg\varphi)\rightarrow T\nvdash\varphi}
 {\RuleRef{RI}{\rightarrow I}\,1\!:\!9}
\end{tabproof}


\section{Von relativer Konsistenz zu Unabhängigkeit}

\FormulaThmDeltaK[Deduktionskriterium]
{\begin{aligned}
 \BXLVIIICon(T+\varphi)&\Longleftrightarrow T\nvdash\neg\varphi,\\
 \BXLVIIICon(T+\neg\varphi)&\Longleftrightarrow T\nvdash\varphi
\end{aligned}}
{B48ConsistencyDeduction}{
  \DeltaRow{Satz}{\varphi}
  \DeltaRow{Theorie}{T}
}
\noindent\textbf{Beweisstand: Beweisumformungen noch auszuarbeiten.}
Die nachstehende Skizze ist noch kein referenzierter Beweis dieses
Kriteriums. Benoetigt werden zuvor bewiesene Saetze ueber endliche
Beweisobjekte, insbesondere Entladung und Zusammenfuegung.
Diese Eigenschaften duerfen nicht lediglich als Definitionen neuer
Beweisoperationen postuliert werden.
\begin{tabproof}
 \proofstep{}{\BXLVIIIText{Ein Beweis von \(\bot\) aus \(T+\varphi\)
 lässt sich durch Entladen der zusätzlichen Annahme in einen Beweis
 von \(\neg\varphi\) aus \(T\) umwandeln.}}
 {\BXLVIIIWhy{Deduktion beziehungsweise Negationseinführung}}
 \proofstep{}{\BXLVIIIText{Ein Beweis von \(\neg\varphi\) aus \(T\)
 ergibt zusammen mit der zusätzlichen Annahme \(\varphi\) einen
 Widerspruch.}}
 {\BXLVIIIWhy{Negationsbeseitigung}}
 \proofstep{}{\BXLVIIICon(T+\varphi)\Longleftrightarrow T\nvdash\neg\varphi}
 {\BXLVIIIWhy{Negation der beiden äquivalenten Beweisexistenzen}}
 \proofstep{}{\BXLVIIICon(T+\neg\varphi)\Longleftrightarrow T\nvdash\varphi}
 {\BXLVIIIWhy{Vorige Zeile mit \(\neg\varphi\); doppelte Negation}}
\end{tabproof}

\FormulaThmDeltaK[Unabhängigkeitskriterium]
{\begin{aligned}
 &\BXLVIIICon(T+\varphi)\land\BXLVIIICon(T+\neg\varphi)\\
 &\hspace{12mm}\vdash (T\nvdash\varphi)\land(T\nvdash\neg\varphi)
\end{aligned}}
{B48IndependenceCriterion}{
  \DeltaRow{Theorie}{T}
  \DeltaRow{Satz}{\varphi}
}
\begin{tabproof}
\proofstep{1}{\BXLVIIICon(T+\varphi)\land\BXLVIIICon(T+\neg\varphi)}{\rA}
 \proofstep{1}{\BXLVIIICon(T+\varphi)}{\RuleRef{AE1}{\land E_1}\,1}
 \proofstep{1}{\BXLVIIICon(T+\neg\varphi)}{\RuleRef{AE2}{\land E_2}\,1}
 \proofstep{1}{T\nvdash\varphi}{\FormulaRefAuto{B48RelativeNonProof}{3}}
 \proofstep{1}{T\nvdash\neg\varphi}{\FormulaRefAuto{B48RelativeNonRefutation}{2}}
 \proofstep{1}{(T\nvdash\varphi)\land(T\nvdash\neg\varphi)}
 {\RuleRef{AI}{\land I}\,4,5}
\end{tabproof}
